%LaTeX
%\documentclass[twoside]{article}
\documentclass[twoside]{amsart}
%This makes the margins little smaller than the default
%\usepackage{fullpage}
%fullpage is not installed on andrew, so we'll just use these lines.
\oddsidemargin-0.75cm
\evensidemargin0cm
\topmargin-1.6cm     %I recommend adding these three lines to increase the 
\textwidth17.35cm   %amount of usable space on the page (and save trees)
\textheight23.85cm  


%if you need more complicated math stuff, you should use the next line
\usepackage{amsmath}
%This next line defines a variety of special math symbols which you
%may need
\usepackage{amssymb}
\usepackage{graphics}
\usepackage{mathtools}
\usepackage{hyperref}

%This next line (when uncommented) allow you to use encapsulated
%postscript files for figures in your document
%\usepackage{epsfig}

\newtheorem{definition}{Definition}

%plain makes sure that we have page numbers
\pagestyle{plain}


\title{
Notes and Solutions to \emph{Quantum Field Theory} by Mark Srednicki, 2007.  
}
\author{
  Ernest Yeung
       }
\date{zima 2009}

%This defines a new command \questionhead which takes one argument and
%prints out Question #. with some space.
\newcommand{\questionhead}[1]
  {\bigskip\bigskip
   \noindent{\large\bf Question #1.}
   \bigskip}

\newcommand{\problemhead}[1]
  {\smallskip
   \noindent{\large\bf Problem #1.}
   }

\newcommand{\exercisehead}[1]
  {\smallskip
   \noindent{\large\bf Exercise #1.}
   }

\newcommand{\solutionhead}[1]
  {\smallskip
   \noindent{\large\bf Solution #1.}
   }


%-----------------------------------
\begin{document}
%-----------------------------------

\maketitle

Fund Science! \& Help Ernest finish his Physics Research! : quantum super-A-polynomials - a thesis by Ernest Yeung \\

\url{http://igg.me/at/ernestyalumni2014} \\

\noindent Facebook  : ernestyalumni \\
gmail         : ernestyalumni \\
linkedin      : ernestyalumni \\
stackexchange : ernestyalumni2014  \\ 
tumblr        : ernestyalumni \\
twitter       : ernestyalumni \\
weibo         : ernestyalumni \\ 
youtube       : ernestyalumni \\
indiegogo     : ernestyalumni \\


%-----------------------------------%-----------------------------------%-----------------------------------%-----------------------------------%-----------------------------------
Solutions for \emph{Quantum Field Theory}.  Mark Srednicki.  Cambridge University Press.  2007.  ISBN-13 978-0-527-86449-7 hardback
%-----------------------------------%-----------------------------------%-----------------------------------%-----------------------------------%-----------------------------------

\section*{Part I Spin Zero}

\section{Attempts at relativistic quantum mechanics}

\subsection*{Problems}

\solutionhead{1.1} Recall
\[
\lbrace \alpha^j, \alpha^k \rbrace_{ab} = 2 \delta^{jk} \delta_{ab}, \quad \quad \, \lbrace \alpha^j, \beta \rbrace_{ab} = 0 \quad \quad \,  (\beta^2)_{ab} = \delta_{ab} \quad \quad \quad  \, (1.27) 
\]

\noindent Suppose $\beta | \psi \rangle = \lambda | \psi \rangle$ \\ 
\phantom{Suppose} $\beta^2 | \psi \rangle = \lambda^2 | \psi \rangle = 1 | \psi \rangle = | \psi \rangle$  $\lambda^2 = 1$ or $\lambda = \pm 1$  

Use hint.  $Tr{\alpha_1^2 \beta} = - Tr{ \alpha_1 \beta \alpha_1} = - Tr{ \alpha_1^2 \beta} \Longrightarrow Tr{ \alpha_1^2 \beta} = 0$

\[
\lbrace \alpha^1, \alpha^1 \rbrace_{ab} = 2 (\alpha^1 )^2_{ab} = 2 \delta_{ab}
\]
so $( \alpha^1)^2 = 1 \Longrightarrow Tr{ \beta } =0$

Since $Tr{\beta} = 0$ and $\lambda_{\beta} = \pm 1$ and $Tr{\beta} = \sum \lambda_{\beta}$, then $\beta$ must be even-dimensional for a set of $\pm 1$'s to sum up and cancel each other.

Likewise,

$\alpha_i | \psi \rangle = \lambda | \psi \rangle$

\[
\alpha_i^2 | \psi \rangle = \lambda \alpha_i | \psi \rangle = \lambda^2 | \psi \rangle = | \psi \rangle \quad  \quad \, \Longrightarrow \lambda = \pm 1 
\]
\[
Tr{ \alpha_1^2 \alpha_j} = - Tr{ \alpha_1 \alpha_j \alpha_1} = - Tr{ \alpha_1^2 \alpha_j } \Longrightarrow Tr{ \alpha_1^2 \alpha_j} =0 \text{ so } Tr{ (\alpha_j) } = 0
\]

So then $\alpha^j$ even dimensional as well.  


\solutionhead{1.2} Recall the position basis Schr\"{o}dinger equation for $n$ particles.  
\[
i \hbar \frac{ \partial }{ \partial t } \psi = \left[ \sum^n_{j=1} \left( \frac{ - \hbar^2}{ 2 m } \nabla_j^2 + U(x_j) \right) + \sum_{j=1}^n \sum_{k=1}^{j-1} V(x_j - x_k) \right] \psi \quad \quad \quad \, (1.30)
\]
Recall
\[
H = \int d^3x a^{\dag}(\mathbf{x}) \left( \frac{ - \hbar^2}{ 2m } \nabla^2 + U(\mathbf{x} ) \right) a(\mathbf{x}) + \frac{1}{2} \int d^3x d^3y V(\mathbf{x} - \mathbf{y})a^{\dag}\mathbf{x})a^{\dag}(\mathbf{y}) a(\mathbf{y})a(\mathbf{x}) \quad \quad \quad \, (1.32) 
\]
\[
|\psi, t \rangle = \int d^3x_1, \dots, d^3x_n \psi(\mathbf{x}_1, \dots \mathbf{x}_n,t) a^{\dag}(\mathbf{x}_1)\dots a^{\dag}(\mathbf{x}_n) |0 \rangle \quad \quad \quad \, (1.33) 
\]
and
\[
i \hbar \frac{ \partial}{ \partial t} | \psi, t \rangle = H | \psi, t \rangle \quad \quad \quad \, (1.1) 
\]
Also recall the commutation and anticommutation rules:
\[
\begin{aligned}
  \left[ a(\mathbf{x}), a(\mathbf{x}') \right] & = 0 \\ 
  [ a^{\dag}(\mathbf{x}), a^{\dag}(\mathbf{x}') ] & =0 \\ 
  [ a(\mathbf{x}), a^{\dag}(\mathbf{x}') ] & = \delta^3(\mathbf{x}- \mathbf{x}) 
\end{aligned} \quad \quad \, (1.31) \quad \quad \quad \quad \begin{aligned}
  \lbrace a(\mathbf{x}), a(\mathbf{x}') \rbrace & =0  \\
  \lbrace a^{\dag}(\mathbf{x}), a^{\dag}(\mathbf{x}') \rbrace & =0  \\
 \lbrace a(\mathbf{x}), a^{\dag}(\mathbf{x}') \rbrace & = \delta^3(\mathbf{x}- \mathbf{x}')
\end{aligned} \quad \quad \, (1.38)
\]

Using some shorthand notation,
\[
\begin{gathered}
  H | \psi,t \rangle = \left( \int d^3x a^{\dag}_{\mathbf{x}} \left( \frac{- \hbar^2}{2m } \nabla^2 + U_{\mathbf{x}} \right) a_{\mathbf{x}} + \frac{1}{2} \int d^3x d^3y V_{\mathbf{x} - \mathbf{y}} a^{\dag}_{ \mathbf{x} }a^{\dag}_{\mathbf{y}} a_{\mathbf{y}}a_{\mathbf{x}} \right) \cdot  \int d^3x_1 \dots d^3x_n \psi(\mathbf{x}_1, \dots , \mathbf{x}_n, t) a^{\dag}_{\mathbf{x}_1 }, \dots, a^{\dag}_{\mathbf{x}_n } | 0 \rangle
\end{gathered}
\]
Now 
\[
\begin{gathered}
  a_{\mathbf{x}}a^{\dag}_{\mathbf{x}_1} \dots a^{\dag}_{\mathbf{x}_n} = ((\pm 1) a^{\dag}_{\mathbf{x}_1} a_{\mathbf{x}} + \delta_{\mathbf{x}- \mathbf{x}_1} ) a^{\dag}_{\mathbf{x}_2} \dots a^{\dag}_{\mathbf{x}_n} =  (\pm 1) a^{\dag}_{\mathbf{x}_1} ( ( \pm 1) a^{\dag}_{\mathbf{x}_2} a_{\mathbf{x}} + \delta_{\mathbf{x} - \mathbf{x}_2} ) a^{\dag}_{\mathbf{x}_3} \dots a^{\dag}_{\mathbf{x}_n} + \delta_{\mathbf{x} - \mathbf{x}_1} a^{\dag}_{\mathbf{x}_2} \dots a^{\dag}_{\mathbf{x}_n }
\end{gathered}
\]

If this term is applied to $|0 \rangle$, note $a_x |0 \rangle =0$.  

\[
a_{\mathbf{x}} a^{\dag}_{\mathbf{x}_1} \dots a^{\dag}_{\mathbf{x}_n} | 0 \rangle = \sum_{j=1}^n (\pm 1)^{j-1} \delta_{\mathbf{x} - \mathbf{x}_j} \prod_{k\neq j}^n a_{\mathbf{x}_k}^{\dag} | 0 \rangle
\]

For the kinetic and self-potential terms of $H | \psi, t \rangle$, 
\[
\int d^3x a^{\dag}_{\mathbf{x}}  \left( \frac{ - \hbar^2}{2m } \nabla^2 + U_{\mathbf{x}} \right) \int d^3x_1 \dots d^3x_n \psi(\mathbf{x}_1, \dots, \mathbf{x}_n, t ) \sum_{j=1}^n (\pm 1)^{j-1} \delta_{\mathbf{x} - \mathbf{x}_j} \prod_{k\neq j}^n a^{\dag}_{\mathbf{x}_k} |  0 \rangle
\]

Evaluate $\int d^3x$, in particular, with $\delta_{\mathbf{x} - \mathbf{x}_j}$, 
\[
 \sum_{j=1}^n a^{\dag}_{\mathbf{x}_j} \left( \frac{ - \hbar^2 }{ 2m } \nabla_j^2 + U_{\mathbf{x}_j } \right) \int d^3x_1, \dots , d^3x_n \psi(\mathbf{x}_1, \dots, \mathbf{x}_n , t ) (\pm 1)^{j-1} \prod_{k\neq j }^n a^{\dag}_{x_k } | 0 \rangle
\]

To (anti)commmute $a^{\dag}_{\mathbf{x}_j}$ over to its ordered place in $\prod_{k\neq j}^n a^{\dag}_{\mathbf{x}_k}( \pm 1)^{j-1}$, $j-1$ (anti)commutations must take place.  
\[
\Longrightarrow \sum_{j=1}^n \left(  - \frac{ \hbar^2}{2m } \nabla_j^2 + U_{\mathbf{x}_j} \right) \int d^3x_1, \dots , d^3x_n \psi(\mathbf{x}_1, \dots, \mathbf{x}_n, t) a^{\dag}_{\mathbf{x}_1}, \dots , a^{\dag}_{\mathbf{x}_n} | 0 \rangle 
\]
Now, on $a_{\mathbf{x}} a^{\dag}_{\mathbf{x}_1} , \dots, a^{\dag}_{\mathbf{x}_n} | 0 \rangle $.  

\[
\begin{gathered}
  \int d^3x V_{\mathbf{x} - \mathbf{y}} a_{\mathbf{x}}^{\dag} a^{\dag}_{\mathbf{y}} a_{\mathbf{y}} a_{\mathbf{x}} a^{\dag}_{\mathbf{x}_1} , \dots , a^{\dag}_{\mathbf{x}_n} | 0 \rangle = \sum_{j=1}^n ( \pm 1)^{j-1} \int d^3x V_{\mathbf{x} - \mathbf{y}} a^{\dag}_{\mathbf{x}} a^{\dag}_{\mathbf{y}} a_{\mathbf{y}} \delta_{\mathbf{x} - \mathbf{x}_j} \prod_{ k \neq j }^n a_{\mathbf{x}_k}^{\dag} | 0 \rangle  = \\
  = \sum_{j=1}^n (\pm 1)^{j-1} V_{\mathbf{x}_j - \mathbf{y} } a^{\dag}_{\mathbf{x}_j } a^{\dag}_{\mathbf{y}} a_{\mathbf{y}} \prod_{ k \neq j}^n a_{\mathbf{x}_k}^{\dag} | 0 \rangle
\end{gathered}
\]
Now
\[
a_{ \mathbf{y}} \prod_{k \neq j }^n a_{\mathbf{x}_k}^{\dag} | 0 \rangle = \sum_{k \neq j }^n \delta_{\mathbf{x}_k - \mathbf{y}} (-1)^m  \prod_{l \neq j,k }^n a^{\dag}_{\mathbf{x}_l } | 0 \rangle 
\]
So then
\[
\begin{gathered}
  \int d^3x d^3y V_{\mathbf{x} - \mathbf{y} } a^{\dag}_{\mathbf{x}} a^{\dag}_{\mathbf{y}} a_{\mathbf{y}} a_{\mathbf{x}} a^{\dag}_{\mathbf{x}_1} , \dots , a^{\dag}_{\mathbf{x}_n } | 0 \rangle = \int d^3 y \sum_{j=1}^n (\pm 1)^{j-1} V_{\mathbf{x}_j - \mathbf{y}} a^{\dag}_{\mathbf{x}_j } a^{\dag}_{\mathbf{y}} \sum_{k \neq j }^n \delta_{\mathbf{x}_k - y} (-1)^m \prod_{l \neq j,k }^n a^{ \dag}_{\mathbf{x}_l} | 0 \rangle = \\
  = \sum_{j=1}^n \sum_{ k \neq j }^n V_{\mathbf{x}_j - \mathbf{x}_k } (\pm 1 )^{j-1 + m } a^{\dag}_{\mathbf{x}_j } a_{\mathbf{x}_k}^{\dag} \prod_{l \neq j,k}^n a^{\dag}_{\mathbf{x}_l} | 0 \rangle
\end{gathered}
\]

To bring $a_{\mathbf{x}_k}^{\dag}$ back in its ordered place with $\prod_{l \neq j,k}^n a^{\dag}_{\mathbf{x}_l}$, then $m$ (anti)commutations must occur (we really are repeating what we did before, exactly, again).  So there's another factor of $(-1)^m$.  Likewise for $a^{\dag}_{\mathbf{x}_j}$, to be back in its ordered place, another factor of $(-1)^{j-1}$ results.  

\[
\Longrightarrow \sum_{j=1}^n \sum_{k \neq j}^n V(\mathbf{x}_j - \mathbf{x}_k) a^{\dag}_{\mathbf{x}_1} , \dots, a_{\mathbf{x}_n}^{\dag} | 0 \rangle
\]

Now $\sum_{j=1}^n \sum_{k \neq j}^n$ counts each unique pair twice so 
\[
\sum_{j=1}^n \sum_{k \neq j}^n V(\mathbf{x}_j - \mathbf{x}_k ) a^{\dag}_{\mathbf{x}_1} , \dots , a^{\dag}_{\mathbf{x}_n} | 0 \rangle = 2 \sum_{j=1}^n \sum_{k=1}^{j-1} V(\mathbf{x}_j - \mathbf{x}_k  ) a^{\dag}_{\mathbf{x}_1}, \dots, a_{\mathbf{x}_n}^{\dag} | 0 \rangle
\]
\[
\Longrightarrow 
\begin{gathered} 
  \frac{1}{2} \int d^3x d^3y V(\mathbf{x}-  \mathbf{y}) a^{\dag}_{\mathbf{x} } a^{\dag}_{\mathbf{y}} a_{\mathbf{y}} a_{\mathbf{x}} \int d^3x_1, \dots , d^3x_n \psi(\mathbf{x}_1, \dots, \mathbf{x}_n, t) a^{\dag}_{\mathbf{x}_1}, \dots, a^{\dag}_{ \mathbf{x}_n } | 0 \rangle = \\
 = \sum_{j=1}^n \sum_{k=1}^{j-1} V(\mathbf{x}_j - \mathbf{x}_k) \int d^3x_1 , \dots, d^3x_n \psi(\mathbf{x}_1, \dots \mathbf{x}_n , t) a^{\dag}_{\mathbf{x}_1}, \dots , a^{\dag}_{\mathbf{x}_n} | 0 \rangle
\end{gathered}
\]

\[
\Longrightarrow i \hbar \frac{ \partial}{ \partial t} | \psi, t \rangle = H| \psi, t \rangle \text{ is satisfied.} 
\]


\problemhead{1.3} Show explicitly that $[N,H]=0$, where $H$ is given by eq. (1.32) and $N$ by eq. (1.35).  


\solutionhead{1.3} 
\[
 \begin{aligned} H & = \int d^3 x a^{\dag}(x) \left( \frac{-1}{2m} \nabla^2 + U(x) \right) a(x) + \frac{1}{2} \int d^3x d^3 y V(x-y) a^{\dag}(x) a^{\dag}(y) a(y) a(x) = \quad \, (1.32) \\
   & = \int d\vec{x} a^{\dag}_x \left( \frac{-1}{2m} \nabla_x + U_x \right) a_x + \frac{1}{2} \int d\vec{x} d\vec{y} V(x-y) a^{\dag}_x a^{\dag}_y a_y a_x  \quad \, \text{(notation)}
\end{aligned}
\]
\[
N = \int d^3x a^{\dag}(x) a(x) = \int d\vec{x} a^{\dag}_x a_x \quad \, \text{(notation)} \quad \quad \, (1.33)
\]

\[
\begin{aligned}
  NH & = \int d\vec{z} a_z^{\dag} a_z \int d\vec{x} a^{\dag}_x \left( \frac{-1}{2m} \nabla_x + U_x \right) a_x + \frac{1}{2} \int d\vec{z} a_z^{\dag} a_z \int d\vec{x} d\vec{y} V(x-y) a_x^{\dag} a_y^{\dag} a_y a_x \\ 
  HN & = \int d\vec{x} a^{\dag}_x \left( \frac{-1}{2m} \nabla_x + U_x \right) a_x \int d\vec{z} a_z^{\dag} a_z + \frac{1}{2} \int d\vec{x} d\vec{y} V(x-y) a_x^{\dag}a_y^{\dag} a_y a_x \int d\vec{z} a_z^{\dag} a_z
\end{aligned}
\]

Recall
\[
\begin{aligned}
  \left[a(x),a(x') \right]_{\pm} & = 0  \\
  [a^{\dag}(x), a^{\dag}(x') ]_{\pm} & = 0  \\
  [a(x), a^{\dag}(x')]_{\pm} & = \delta^3(x-x')
\end{aligned} \quad \quad \, \begin{aligned} & a_x a_z^+ = \pm a_z^{\dag} a_x + \delta^3(x-z) \\ & [a^{\dag}_x, a_z ]_{\pm} = - [a_z, a_x^{\dag}]_{\pm} = -\delta^3(z-x) \Longrightarrow a_x^{\dag} a_z = \pm a_z a_x^{\dag} - \delta^3(z-x) \\ & [ \frac{-1}{2m} \nabla_x + U_x, a_z^{\dag} ] = [\mathcal{H}_{0x}, a_z^{\dag} ] = 0 
\end{aligned}
\]
Now
\[
\left[ \mathcal{H}_{0x}, a_z \right] =0 
\]
Treat $x,y,z$ as independent coordinates.  

\[
\int d\vec{z} d\vec{x} a_x^{\dag} \mathcal{H}_{0x} a_x a_z^{\dag} a_z = \int d\vec{z} d\vec{x} \left( a_z^{\dag} a_z a_x^{\dag} \mathcal{H}_{0x} a_x \right) + \int d\vec{z} a_z^{\dag} \mathcal{H}_{0z} a_z - \int d\vec{z} a_z^{\dag} \mathcal{0z} a_z = \int d\vec{z} d\vec{x} \left( a_z^{\dag} a_z a_x^{\dag} \mathcal{H}_{0x} a_x \right) 
\]

Notice even if the anticommutator rules applies, $a_z^{\dag}$ switched with $a_x^{\dag}$ and $a_x$, for a factor of $(-1)^2 =1$ overall.  Likewise for $a_z$ with $a_x$ and $a_x^{\dag}$.  

Note that $\begin{aligned} \left[ V(x-y),a_z^{\dag} \right] & = 0 \\ [V(x-y), a_z ] & = 0 \end{aligned}$.  

\[
\begin{gathered}
  V(x-y) a_x^{\dag} a_y^{\dag} a_y a_x a_z^{\dag} a_z = V(x-y) a_x^{\dag}a_y^{\dag} a_y \left( (\pm 1) a_z^{\dag} a_x + \delta(x-z) \right)a_z = \\
  = V(x-y)a_x^{\dag}a_y^{\dag} (a_z^{\dag} a_y + \delta(z-y) ) a_x a_z + V(x-y) a_x^{\dag} a_y^{\dag} a_y a_z \delta(x-z)  \\
  \xrightarrow{ \int d\vec{z} } \int d\vec{z} a_z^{\dag} V(x-y) a_x^{\dag} a_y^{\dag} a_z a_y a_x + V(-xy) a_x^{\dag} a_y^{\dag} a_x a_y + V(x-y) a_x^{\dag}a_y^{\dag}a_y a_x = \\
  = \int d\vec{z} a_z^{\dag}V(x-y) a_x^{\dag} ((\pm 1) a_z a_y^{\dag} - \delta(z-y)) a_y a_x + V(x-y) a_x^{\dag}a_y^{\dag} \lbrace a_x ,a_y \rbrace = 
\end{gathered}
\]
\[
\begin{gathered}
= \int d\vec{z} a_z^{\dag} V(x-y) a_z a_x^{\dag} a_y^{\dag} a_y a_x + \int d\vec{z} a_z^{\dag} V(x-y) (- \delta(z-x) ) a_y^{\dag} a_y a_x + \int d\vec{z} a_z^{\dag} V(x-y) a_x^{\dag} (-\delta(z-y) )a_y a_x) + \\
+ V(x-y) a_x^{\dag} a_y^{\dag} \lbrace a_x , a_y \rbrace = \\ 
 = \int d\vec{z} a_z^{\dag}a_z V(x-y) a_x^{\dag} a_y^{\dag} a_y a_x + -V(x-y) \lbrace a_x^{\dag}, a_y^{\dag} \rbrace a_y a_x + V(x-y) a_x^{\dag} a_y^{\dag} \lbrace a_x, a_y \rbrace \\
=\int d\vec{z} a_z^{\dag} a_z V(-xy) a_x^{\dag} a_y^{\dag} a_y a_x + V(x-y) ( a_x^{\dag} a_y^{\dag} a_x a_y - a_y^{\dag} a_x^{\dag} a_y a_x) = \int d\vec{z} a_z^{\dag} a_z V(-xy) a_x^{\dag} a_y^{\dag} a_y a_x \quad \,  \\ \text{(for both commutation and anticommutator rules)}
\end{gathered}
\]

Then indeed $[N,H]=0$.  

\section{Lorentz invariance}

\subsection*{Problems}

\problemhead{2.1} Verify that eq. (2.8) follows from eq. (2.3).  

\solutionhead{2.1} Recall $\begin{aligned}
& g_{\mu \nu} \Lambda^{\mu}_{\rho} \Lambda^{\nu}_{\sigma} = g_{\rho \sigma} &    \quad \quad \, (2.3)  \\
\Lambda^{\mu}_{\nu} & = \delta^{\mu}_{\nu} + \delta \omega^{\mu}_{\nu}  & \quad \quad \, (2.8)  
\end{aligned}$.  \\

Then
\[
g_{\mu \nu} ( \delta^{\mu}_{\rho} + \delta^{\mu}_{\rho} )( \delta^{\nu}_{\sigma} + \delta \omega^{\nu}_{\sigma} ) = g_{\rho \sigma} = g_{\mu \nu} ( \delta^{\mu}_{\rho} \delta^{\nu}_{\sigma} + \delta^{\nu}_{\sigma} \delta \omega^{\mu}_{\rho} + \delta^{\mu}_{\rho} \delta \omega^{\nu}_{\sigma} + \delta \omega^{\mu}_{\rho} \delta \omega^{\nu}_{\sigma} )
\]
Neglect 2nd. order infinitesimals: $\delta \omega^{\mu}_{\rho }\delta \omega^{\nu}_{\sigma} \to 0$
\[
g_{\rho\sigma} = g_{\rho\sigma} + g{\mu \sigma} \delta \omega^{\mu}_{\rho} + g_{\rho \nu} \delta \omega^{\nu}_{\sigma}  \quad \, \Longrightarrow g_{\mu \sigma} \delta \omega^{\mu}_{\rho} = -g_{\rho \nu} \delta \omega^{\nu}_{\sigma}
\]

Recall $A^i_k = g^{il} A_{lk}$; $g_{ik} g^{kl} = \delta^l_i$, and $g_{\mu \sigma}$ symmetric.  

\[
\begin{gathered}
g_{\mu \sigma}g^{\mu l} \delta \omega_{l \rho} = -g_{\rho \nu} g^{\nu m} \delta \omega_{m \sigma} = - \delta^m_{\rho} \delta \omega_{m \sigma} = -\delta \omega_{\rho \sigma} \\
g_{\sigma \mu} g_{\mu l} \delta \omega_{l \rho} = \delta^l_{\rho} \delta_{l \rho} = \delta \omega_{\sigma \rho} \quad \, \Longrightarrow \boxed{ \delta \omega_{\sigma \rho} = - \delta \omega_{\rho \sigma} } \quad \, \text{(antisymmetric tensor)}
\end{gathered}
\]

\problemhead{2.2} Verify that eq. (2.14) follows from $U(\Lambda)^{-1}U(\Lambda')U(\Lambda) = U(\Lambda^{-1} \Lambda' \Lambda)$.  

\solutionhead{2.2} $\Lambda^{-1} \Lambda' \Lambda = \Lambda^{-1} (1 + \delta \omega') \Lambda = 1 + \Lambda^{-1} \delta \omega' \Lambda$ \\
$\begin{aligned}
U(\Lambda^{-1} \Lambda' \Lambda) & = 1 + \frac{i}{2} ( \Lambda^{-1} \delta \omega' \Lambda)_{\rho \sigma} M^{\rho \sigma} = 1 + \frac{i}{2} (\Lambda^{-1})_{\rho}^a (\delta \omega')_{ab} (\Lambda)^b_{\sigma} M^{\rho \sigma}
\end{aligned}$

\[
\begin{aligned}
  & U(\Lambda') = U(1 + \delta \omega') = 1 + \frac{i}{2} \delta \omega_{\mu \nu} M^{\mu \nu} \\ 
  & U(\Lambda)^{-1} U(\Lambda') U(\Lambda) = U(\Lambda)^{-1} U(\Lambda) + \frac{i}{2} U(\Lambda)^{-1} \delta \omega_{\mu \nu} M^{\mu \nu} U(\Lambda)
\end{aligned}
\]

Using $U(\Lambda)^{-1}U(\Lambda')U(\Lambda) = U(\Lambda^{-1} \Lambda' \Lambda)$,
\[
\Longrightarrow U(\Lambda^{-1} )\delta \omega_{\mu \nu} M^{\mu \nu} U(\Lambda) = \delta \omega_{\mu \nu} U(\Lambda^{-1} )M^{\mu \nu} U(\Lambda) = \delta \omega_{\mu \nu} (\Lambda^{-1})_{\rho}^{\mu} (\Lambda)^{\nu}_{\sigma} M^{\rho \sigma}
\]
Note that $\delta \omega_{\mu \nu}$ is just an entry in a rank-2 tensor, a number, and for arbitrary Lorentz transformation, each a uniquelly determined quantity $\delta \omega_{\mu \nu}$.  (6 of them for antisymmetric $\delta \omega_{\mu \nu}$).  
\[
\Longrightarrow U(\Lambda^{-1}) M^{\mu \nu} U(\Lambda) = (\Lambda^{-1})_{\rho}^{\mu} (\Lambda)^{\nu}_{\sigma} M^{\rho \sigma}
\]

\problemhead{2.3} Verify that eq. (2.16) follows from eq. (2.14).  

\solutionhead{2.3} Recall that $U(\Lambda)^{-1} M^{\mu \nu}U(\Lambda) = \Lambda^{\mu}_{\rho}\Lambda^{\nu}_{\sigma} M^{\rho \sigma} \quad \quad \, (2.14)$.  

\[
\begin{gathered}
  U(\Lambda)^{-1} M^{\mu \nu} U(\Lambda) = (1 - \frac{i}{2} \delta \omega_{ab} M^{ab} ) M^{\mu \nu} ( 1 + \frac{i}{2} \delta \omega_{cd} M^{cd}) = (M^{\mu \nu} - \frac{i}{2} \delta \omega_{ab} M^{ab} M^{\mu \nu} )( 1 + \frac{i}{2} \delta \omega_{cd} M^{cd} ) = \\
  = M^{\mu \nu} - \frac{i}{2} \delta \omega_{ab} M^{ab} M^{\mu \nu} + \frac{i}{2} M^{\mu \nu} \delta \omega_{cd} M^{cd} = M^{\mu \nu } + \frac{i}{2} ( \delta \omega_{cd} M^{\mu \nu} M^{cd} - \delta \omega_{ab} M^{ab} M^{\mu \nu } )
\end{gathered}
\]
We're only concerned with terms in first order of $\delta \omega$.  

Now
\[
[ M^{\mu \nu}, M^{\rho \sigma} ] = M^{\mu \nu } M^{\rho \sigma} - M^{\rho \sigma} M^{\mu \nu } 
\]
So then running through all the indices and picking out indices we're interested in (particularly $\rho, \sigma$),
\[
\begin{gathered}
\delta \omega_{cd} M^{\mu \nu } M^{cd} - \delta \omega_{ab} M^{ab} M^{\mu \nu } \Longrightarrow \delta \omega_{\rho \sigma} M^{\mu \nu} M^{\rho \sigma} - \delta \omega_{\rho \sigma} M^{\rho \sigma} M^{\mu \nu } + \delta \omega_{\sigma \rho} M^{\mu \nu}M^{\sigma \rho} - \delta \omega_{\sigma \rho} M^{\sigma rho }M^{\mu \nu} = \\ = 2 \delta \omega_{\rho \sigma} [ M^{\mu \nu}, M^{\rho \sigma} ]
\end{gathered}
\]
where we used $\begin{aligned}  \delta \omega_{\rho \sigma} & = -\delta \omega_{\sigma \rho} \\ M^{\rho \sigma} & = - M^{\sigma \rho } \end{aligned}$ (antisymmetry).  

Now
\[
\begin{gathered}
\Lambda^{\mu}_{\rho} \Lambda^{\nu}_{\sigma} M^{\rho \sigma} = \Lambda^{\mu}_a \Lambda^{\nu}_b M^{ab} = (1 + \delta \omega^{\mu}_a )( 1 + \delta \omega^{\nu}_b) M^{ab} = (1 + \delta \omega^{\mu}_a)( M^{ab} + \delta \omega^{\nu}_b M^{ab} ) = M^{ab} + \delta \omega_{lb}g^{\nu l} M^{ab} + \delta \omega_{ma} g^{\mu m} M^{ab}
\end{gathered}
\]
Considering only the first order infinitesimal, (and note that for $M^{ax}, M^{yb}$, $a,b$ run from $0$ to $3$),
\[
\begin{aligned}
  \delta \omega_{lb} g^{\nu l } M^{ab} + \delta \omega_{ma} g^{\mu m} M^{ab} & = \delta \omega_{\rho \sigma} g^{\nu \rho } M^{a \sigma} + \delta \omega_{\sigma \rho} g^{\nu \sigma }M^{a\rho} + \delta \omega_{\rho \sigma} g^{\mu \rho }M^{\sigma b} + \delta \omega_{\sigma \rho}g^{\mu \sigma} M^{\rho b} = \\
&  = \delta \omega_{\rho \sigma} \lbrace g^{\mu \rho} M^{\sigma b} - g^{\nu \rho}M^{a\sigma} + - ( g^{\mu \sigma} M^{\rho b} + g^{\nu \sigma }M^{a\rho} ) \rbrace
\end{aligned}
\]

For $a \to \mu$ and $b \to \nu$, 
\[
\Longrightarrow \begin{aligned} \left[ M^{\mu \nu }, M^{\rho \sigma} \right] & = i \lbrace g^{\mu \rho} M^{b\sigma } - g^{\nu \rho}M^{a\sigma}  - ( g^{\mu \sigma} M^{b\rho } - g^{\nu \sigma }M^{a\rho} ) \rbrace = \\
  & =\boxed{ i \lbrace g^{\mu \rho} M^{\nu \sigma } - g^{\nu \rho}M^{\mu\sigma}  - ( g^{\mu \sigma} M^{\nu \rho } - g^{\nu \sigma }M^{\mu\rho} ) \rbrace }
\end{aligned}
\]
Note that if $a=b$, we get $0$.  If $b= \mu$, $a = \nu$, then 
\[
g^{\mu \rho} M^{\mu \sigma} - g^{\nu \rho} M^{\nu \sigma} - (g^{\mu \sigma} M^{\mu \rho} - g^{\nu \sigma}M^{\nu \rho} ) = M^{\rho \sigma} - M^{\rho \sigma} - M^{\sigma \rho} + M^{\sigma \rho} = 0 \quad \, \text{(modulo a $(-1)$ factor)}
\]
If $b = \rho$, $a = \sigma$, 
\[
g^{\mu \rho} M^{\rho \sigma} - g^{\nu \rho} M^{\sigma \sigma} - (g^{\mu \sigma} M^{\rho \rho} - g^{\nu \sigma} M^{\sigma \rho }) = g^{\mu \rho} M^{\rho \sigma} - g^{\nu \sigma} M^{\rho \sigma} = (g^{\mu \rho } - g^{\nu \sigma}) M^{\rho \sigma}
\]
If vice-versa,
\[
g^{\mu \rho} M^{\sigma \sigma} - g^{\nu \rho}M^{\rho \sigma} - (g^{\mu \sigma} M^{\sigma \rho} - g^{\nu \sigma} M^{\rho \rho} ) = -g^{\nu \rho }M^{\rho \sigma} - g^{\mu \sigma} M^{\sigma \rho} = (g^{\mu \sigma} - g^{\nu \rho }) M^{\rho \sigma}
\]




\problemhead{2.4} Verify that eq. (2.17) follows from eq. (2.16).

\solutionhead{2.4} Recall \\
$ \left[ M_{\mu \nu}, M^{\rho \sigma} \right] = i \lbrace g^{\mu \rho} M^{\nu \sigma } - g^{\nu \rho}M^{\mu\sigma}  - ( g^{\mu \sigma} M^{\nu \rho } - g^{\nu \sigma }M^{\mu\rho} ) \rbrace  \quad \, (2.16)$  \\

$\begin{aligned} \left[ J_i,J_j \right] & = i \epsilon_{ijk} J_k \\ [J_i,K_j] & = i \epsilon_{ijk} K_k \\ [K_i,K_j] & = -i \epsilon_{ijk} J_k \quad \, (2.17)  \end{aligned}$ \\
\quad \, \\

Recall $J_i \equiv \frac{1}{2} \epsilon_{ijk} M^{jk}$.  Then
\[
i\epsilon_{ijk} J_k = i \epsilon_{kij} \frac{1}{2} \epsilon_{klm} M^{lm} = \frac{i}{2} (\delta_{il} \delta_{jm} - \delta_{im} \delta_{jl}) M^{lm} = \frac{i}{2} (M^{ij} - M^{ji}) = iM^{ij}
\]

Recall $K_i = M^{i0}$.  

\[
\begin{aligned}
\left[ K_i, K_j \right] & = [M^{i0}, M^{j0} ] = i ( g^{ij} M^{00} - g^{0j} M^{i0} - g^{i0} M^{0j} + g^{00} M^{ij} ) = i M^{ij}
\end{aligned}
\] 

\quad \, \\ 
\quad \, \\ 
$i\epsilon_{ijk} K_k = i \epsilon_{ijk} M^{k0}$.  
\[
\begin{aligned}
  \left[ J_i, K_j \right] & = [ \frac{1}{2} \epsilon_{ikl} M^{kl}, M^{j0} ] = \frac{\epsilon_{ikl}}{2} [ M^{kl}, M^{j0} ] = \frac{ i \epsilon_{ikl}}{2} \left( g^{kj} M^{l0} - g^{lj} M^{k0} - g^{k0}M^{lj} + g^{l0} M^{kj} \right) = \\
 & = \frac{ i \epsilon_{ikl}}{2} (M^{k0} - M^{k0} ) = i\epsilon_{ijk} M^{k0} = \boxed{i\epsilon_{ijk} K_k}
\end{aligned}
\]
$g^{k0} = g^{l0} =0$ for $k,l \neq 0$.   \\





\problemhead{2.5} Verify that eq. (2.18) follows from eq. (2.15).

\solutionhead{2.5} Recall $U(\Lambda)^{-1} P^{\mu} U(\Lambda) = \Lambda^{\mu}_{\nu} P^{\nu} \quad \, (2.15)$.\\
Let $\Lambda = 1 + \delta \omega$.  \\

Keep first order infinitesimals.  

\[
U(\Lambda)^{-1} P^{\mu} U(\Lambda) = (1 - \frac{i}{2} \delta \omega_{ab} M^{ab} )P^{\mu}(1 - \frac{i}{2} \delta \omega_{cd} M^{cd} ) = P^{\mu} - \frac{i}{2} \delta \omega_{ab} M^{ab} P^{\mu} + \frac{i}{2} P^{\mu} \delta \omega_{cd} M^{cd}
\]

$a,b,c,d$ run from $0$ to $3$.  We sought obtaining term $\delta \omega_{\rho \sigma}$

\[
\begin{gathered}
  \delta \omega_{\rho \sigma} \frac{i}{2} P^{\mu} M^{\rho \sigma} + \delta \omega_{\sigma \rho} \frac{i}{2} P^{\mu} M^{\sigma \rho} - \delta \omega_{\rho \sigma} \frac{i}{2} M^{\rho \sigma} P^{\mu} - \delta \omega_{\sigma \rho} \frac{i}{2} M^{\sigma \rho} P^{\mu} = \delta \omega_{\rho \sigma} i [ P^{\mu}, M^{\rho \sigma} ]
\end{gathered}
\]  

For the right hand side of (2.15), 
\[
\begin{gathered}
\Lambda^{\mu}_{\nu} P^{\nu} = (1 + \delta \omega^{\mu}_{\nu} )P^{\nu} = P^{\nu} + g^{\mu k} \delta \omega_{k\nu} P^{\nu}  \\
\Longrightarrow \delta \omega_{k \nu} g^{\mu k} P^{\nu} \end{gathered}
\]

$\nu, k$ run from $0$ to $3$.  Extract terms involving $\rho, \sigma$ from the sum.  

\[
\Longrightarrow  g^{\mu \rho} \delta \omega_{\rho \sigma} P^{\sigma} + g^{\mu \sigma} \delta \omega_{\sigma \rho} P^{\rho} = \delta \omega_{\rho \sigma} ( g^{\mu \rho} P^{\sigma} - g^{\mu \sigma} P^{\rho} )
\]

So then
\[
\boxed{ \left[ P^{\mu}, M^{\mu \nu} \right] = i(g^{\mu \sigma} P^{\rho} - g^{\mu \rho}P^{\sigma} )  }
\]


\problemhead{2.6} Verify that eq. (2.19) follows from eq. (2.18).

\solutionhead{2.6} Using (2.18): $\left[ P^{\mu}, M^{\rho \sigma} \right] = i ( g^{\mu \sigma} P^{\rho} - g^{\mu \rho} P^{\sigma} )$,
\[
\left[ J_i, H \right] = \frac{1}{2} \epsilon_{ijk} \left[ M^{jk}, H \right] = \frac{-1}{2} \epsilon_{ijk} i (g^{0k} P^j - g^{0j} P^k ) = 0 \quad \, j,k \neq 0
\]
\[
\left[ J_i, P_j \right] = \frac{-1}{2} \epsilon_{ijk} g_{jl} \left[ P^l, M^{jk} \right] = \text{ (assume $l=1,2,3$)} \frac{-1}{2} \epsilon_{ijk} g_{jl} (g^{lk} P^j - g^{lj} P^k) = -\frac{1}{2} \epsilon_{ijk} (-P^k - P^k) = i \epsilon_{ijk} P^k 
\]
\[
\left[ K_i, H \right] = \left[ M^{i0}, H \right] = - \left[ H, M^{i0} \right] = - i (g^{00} P^i - G^{0i} P^0 ) = -i P^i 
\]
\[
\left[ K_i, P_j \right] = \left[ M^{i0}, P_j \right] = g_{jk} \left[ M^{i0}, P^k \right] = -g_{jk} [ P^k , M^{i0} ] = -g_{jk} i ( g^{k0} P^i - g^{ki} P^0 ) = i \delta_{ij} H
\]
Note, I use the metric, $g^{\mu \nu} = \left( \begin{matrix} 1 & & & \\ & -1 & & \\ &  & -1 & \\ & & & -1 \end{matrix} \right)$.  





\problemhead{2.8} \begin{itemize}
\item[(a)] Let $\Lambda = 1 + \delta \omega$ in eq. (2.26), and show that 
\[
\left[ \varphi(x), M^{\mu \nu} \right] = \mathcal{L}^{\mu \nu} \varphi(x), \quad \quad \, (2.29)
\]
where 
\[
\mathcal{L}^{\mu \nu } \equiv \frac{\hbar }{ i } (x^{\mu \nu} - x^{\nu \mu}), \quad \quad \, (2.30)
\]
\item[(b)] Show that $\left[ \left[ \varphi(x), M^{\mu \nu} \right], M^{\rho \sigma} \right] = \mathcal{L}^{\mu \nu} \mathcal{L}^{\rho \sigma} \varphi(x)$.  
\item[(c)] Prove the \emph{Jacobi identity}.  $\left[ \left[ A, B \right], C \right] + [[B,C],A] + [[C,A],B] =0$.  Hint: write out all the commutators.  
\item[(d)] Use your results from parts (b) and (c) to show that 
\[
\left[ \varphi(x), [ M^{\mu \nu}, M^{\rho \sigma} ] \right] = ( \mathcal{L}^{\mu \nu} \mathcal{\rho \sigma} - \mathcal{L}^{\rho \sigma} \mathcal{L}^{\mu \nu } ) \varphi(x). \quad \quad \, (2.31)
\]
\item[(e)] Simplify the right-hand side of eq. (2.31) as much as possible.  
\item[(f)] Use your results from part (e) to verify eq. (2.16), up to the possibility of a term on the right-hand side that commutes with $\varphi(x)$ and its derivatives.  (Such a term, called a \emph{central charge}, in fact does not arise for the Lorentz algebra.)
\end{itemize}

\solutionhead{2.8} 
\begin{itemize}
\item[(a)] Given $U(\Lambda)^{-1} \varphi(x) U(\Lambda) = \varphi(\Lambda^{-1} x)$, \\
$\Lambda = 1 + \delta \omega$,
\[
\Longrightarrow \begin{aligned} U(\Lambda)^{-1} \varphi(x) U(\Lambda) & = ( I - \frac{i}{2\hbar} \delta \omega_{\mu \nu} M^{\mu \nu} ) \left( \varphi(x) + \frac{i \varphi(x)}{2\hbar} \delta_{ab} M^{ab} \right) = \\ 
  & = \varphi(x)- \frac{i}{2\hbar} \delta \omega_{\mu \nu} M^{\mu \nu} \varphi(x) + \frac{ i \delta \omega_{\mu \nu} }{2 \hbar} \varphi(x) M^{\mu \nu} + \frac{1}{ (2\hbar)^2 } \delta \omega_{\mu \nu } \delta \omega_{ab} M^{\mu \nu} \varphi(x) M^{ab} \end{aligned}
\]
Keep up to 1st. order terms:
\[
= \varphi(x) - \frac{i}{\hbar} \delta \omega_{\mu \nu} [ M^{\mu \nu} , \varphi(x) ]
\]
($\mu, \nu$ run through $0,\dots 3$ in the summation.  There'll be a term $\delta \omega_{\nu \mu} M^{\nu \mu } = - \delta \omega_{\mu \nu} (-M^{\mu \nu } ) = \delta \omega_{\mu \nu} M^{\mu \nu}$ to be included.  Hence, a factor of 2).  

Now, 
\[
\begin{gathered}
  \Lambda^{-1} x \to (\Lambda^{-1} )^{\rho}_{\nu} x^{\nu} = \Lambda_{\nu}^{\rho}x^{\nu} = (\delta_{\nu}^{\rho} + \delta \omega_{\nu}^{\rho} ) x^{\nu} 
\end{gathered}
\]
\[
\begin{aligned}
  \varphi(\Lambda^{-1}x ) & = \varphi(x^{\rho} + \delta \omega_{\nu}^{\rho} x^{\nu} \simeq \varphi(x^{\rho} ) + \delta \omega_{\nu}^{\rho} x^{ \nu} \partial_{\rho} \varphi = \\
  & = \varphi(x) + g^{\rho \alpha} \delta \omega_{\nu \alpha} x^{\nu} g_{\rho \beta} \partial^{\beta} \varphi = \varphi(x) + g_{\rho \beta} g^{\rho \alpha} \delta \omega_{\nu \alpha}x^{\nu} \partial^{\beta} \varphi = \\
  & = \varphi(x) + \delta \omega_{\nu \alpha} x^{\nu } \partial^{\alpha} \varphi
\end{aligned}
\]
Let $\nu, \alpha$ run through $0\dots 3$, picking out terms involving $\mu, \nu$
\[
\Longrightarrow \delta \omega_{\mu \nu} x^{\mu} \partial^{\nu} + \delta \omega_{\nu \mu} x^{\nu} \partial^{\mu} = \delta \omega_{\mu \nu} (x^{\mu}\partial_{\nu} - x^{\nu} \partial^{\mu} ) 
\]
\[
\Longrightarrow \boxed{ \left[ M^{\mu \nu}, \varphi(x) \right] = i \hbar (x^{\mu} \partial^{\nu} - x^{\nu} \partial^{\mu} ) \varphi = - \mathcal{L}^{\mu \nu} \varphi(x) }
\]
\item[(b)] \[
\begin{gathered}
  \left[ \left[ \varphi(x), M^{\mu \nu} \right], M^{\rho \sigma} \right] = \left[ \mathcal{L}^{\mu \nu} \varphi(x), M^{\rho \sigma} \right] = \mathcal{L}^{\mu \nu } \mathcal{L}^{\rho \sigma} \varphi(x) + \left[ \mathcal{L}^{\mu \nu }, M^{\rho \sigma} \right] \varphi(x) = \mathcal{L}^{\mu \nu} \mathcal{L}^{\rho \sigma} \varphi(x)
\end{gathered}
\] 
Since $M^{\rho \sigma}$ doesn't depend on $x$.  (Recall $M^{\mu \nu}$ is a hermitian operator that's a generator of the Lorentz group.  
\item[(c)] \[
\begin{gathered}
  \left[ \left[ A, B \right], C \right] + \left[ \left[ B, C \right] , A \right] + \left[ \left[ C,A \right], B \right] = \left[ A, B \right] C - C [A,B] + [B,C] A - A [B,C] + [C,A] B- B[C,A ] = \\ 
  = (AB-BA) C - C(AB- BA) + (BC - CB) A - A(BC - CB ) + (CA - AC) B - B(CA - AC) = \\
ABC - BAC - CAB + CBA + BCA - CBA - ABC + ACB + CAB - ACB - BCA + BAC = 0 
\end{gathered}
\]
\item[(d)] \[
\begin{gathered}
  [ \varphi(x) , [ M^{\mu \nu }, M^{\rho \sigma} ] = [[ M^{\rho \sigma}, \varphi(x) ], M^{\mu \nu } ] + [[ \varphi(x), M^{\mu \nu } ], M^{\rho \sigma} ] = - \mathcal{L}^{\rho \sigma} \mathcal{ \mu \nu } \varphi(x) + \mathcal{L}^{\mu \nu } \mathcal{L}^{\rho \sigma} \varphi(x)
\end{gathered}
\] 
\item[(e)]
\[
\begin{gathered}
  [ \varphi, [ M^{\mu \nu}, M^{\rho \sigma} ] ] = ( \mathcal{L}^{\mu \nu} \mathcal{L}^{\rho \sigma} -  \mathcal{L}^{\rho \sigma} \mathcal{L}^{\mu \nu} ) \varphi(x) = \\
  = (-1) ( x^{\mu} g^{\nu \rho } \partial^{\sigma} \varphi + x^{\mu } x^{\rho} \partial^{\nu \sigma} \varphi - x^{\mu} g^{\nu \sigma} \partial^{\rho} \varphi - x^{\mu} x^{\sigma} \partial^{\nu \rho} \varphi - x^{\nu} g^{\mu \rho} \partial^{\sigma} \varphi - x^{\nu} x^{ \rho} \partial^{\mu \sigma} \varphi + x^{\nu} g^{\mu \sigma} \partial^{\rho} \varphi + x^{\nu} x^{\sigma} \partial^{\mu \rho} \varphi ) - \\
  -(-1) ( x^{\rho} g^{\sigma \mu} \partial^{\nu} \varphi + x^{\rho} x^{\mu} \partial^{\sigma \nu } \varphi - x^{\rho } g^{\sigma \nu } \partial^{\mu} \varphi - x^{\rho} x^{\nu} \partial^{\sigma \mu} \varphi - x^{\sigma} g^{\rho \mu} \partial^{\nu} \varphi - x^{\sigma} x^{ \mu} \partial^{\rho \nu} \varphi + x^{\sigma} g^{\rho \nu} \partial^{\mu} \varphi + x^{\sigma} x^{\nu} \partial^{\rho \mu} \varphi ) = \\
  = g^{\mu \rho} ( x^{\nu} \partial^{\sigma} \varphi - x^{\sigma} \partial^{\nu} \varphi ) - g^{\nu \rho} ( x^{\mu} \partial^{\sigma} \varphi - x^{\sigma} \partial^{\mu} \varphi) - (g^{\mu \sigma} x^{\nu} \partial^{\rho} \varphi - g^{\mu \sigma} x^{\rho} \partial^{\nu} \varphi) + g^{\nu \sigma} ( x^{\mu} \partial^{\rho} \varphi - x^{\rho} \partial^{\mu} \varphi ) = \\
  = i g^{\mu \rho} \mathcal{L}^{\nu \sigma} \varphi - ig^{\nu \rho} \mathcal{L}^{\mu \sigma} \varphi - ig^{\mu \sigma} \mathcal{L}^{\nu \rho} \varphi + i g^{\nu \sigma} \mathcal{L}^{\mu \rho } \varphi
\end{gathered}
\]
\item[(f)] Recall
\[
[ M^{\mu \nu}, M^{\rho \sigma} ] = i ( g^{\mu \rho } M^{\nu \sigma} - g^{\nu \rho } M^{\mu \sigma } ) - i ( g^{\mu \sigma} M^{\nu \rho } - g^{ \nu \sigma} M^{\mu \rho } ) \quad \quad \quad \, (2.16)
\]
If we take (2.16) to be true,
\[
\begin{gathered}
  [ \varphi, [ M^{\mu \nu}, M^{\rho \sigma} ] ] = [ \varphi, i g^{\mu \rho} M^{\nu \sigma} - ig^{\nu \rho } M^{\mu \sigma} - i g^{\mu \sigma} M^{\nu \rho } + i g^{\nu \sigma}M^{\mu \rho } ] = \\
  = ig^{\mu \rho } \mathcal{L}^{\nu \sigma} \varphi - ig^{\nu \rho } \mathcal{L}^{\mu \sigma} \varphi - i g^{\mu \sigma} \mathcal{L}^{\nu \rho } \varphi + i g^{\nu \sigma} \mathcal{L}^{\mu \rho } \varphi
\end{gathered}
\]
From the previous results from part (e), we see that both agree.  Indeed, (2.16) verified up to terms that commute with $\varphi$ and its derivatives, i.e. add terms that commute with $\varphi$ and derivatives of $\varphi$ to the RHS of (2.16).  Put the RHS of (2.16) into $[[M^{\mu \nu}, M^{\rho \sigma} ], \varphi ]$.  Those terms will commute with $\varphi$.
\end{itemize}


\section{Canonical quantization of scalar fields}

\[
\begin{gathered}
L(q_i, \dot{q}_i) \\
p_i = \frac{ \partial L}{ \partial \dot{q}_i} \\
H = \sum_i p_i \dot{q}_i - L 
\end{gathered}
\]
\begin{align}
  & \boxed{ \Pi(x) = \frac{ \partial \mathcal{L}}{ \partial \dot{\varphi}(x)} }   \quad \quad \quad (3.22) \\ 
  & \boxed{ \mathcal{H} = \pi \dot{\varphi} - \mathcal{L}  }\quad \quad \quad (3.23)
\end{align}




\subsection*{Problems}

\problemhead{3.1} Derive eq. (3.29) from eqs. (3.21), (3.24). 

\solutionhead{3.1} Using 
\[
\begin{aligned}
  a(k) & = \int d^3 x e^{-ikx} [ i \partial_0 \varphi(x) + \omega \varphi(x) ] = \\ 
  & = i \int d^3 x e^{-ikx} \overleftrightarrow{\partial_0} \varphi(x) \quad \quad \, (3.21) 
\end{aligned}
\]
\[
\Pi(x) = \dot{\varphi}(x) \quad \quad \, (3.24)
\]
Then we want 
\[
\begin{aligned}
  & \left[ a(k), a(k') \right] = 0 \\ 
  & [a^{\dag}(k), a^{\dag}(k') ] = 0 \\ 
  & [a(k),a^{\dag}(k') ] = (2\pi)^3 2 \omega \delta^3(k-k') \quad \quad \, (3.29)
\end{aligned}
\]

I will also use
\[
\begin{aligned}
  & \left[ \varphi(x,t), \varphi(x',t) \right] =0 \\ 
  & [ \Pi(x,t), \Pi(x',t) ] = 0 \\ 
  & [\varphi(x,t), \Pi(x',t) ] = i \delta^3(x-x')
\end{aligned}
\]
\[
\begin{gathered}
\left[ a(k), a(k') \right] = \int d^3x d^3x' e^{-i(kx + k'x') } [ i \partial_0 \varphi(x) + \omega \varphi(x), i \partial_0 \varphi(x') + \omega \varphi(x') ] \\ 
\Longrightarrow [i\partial_0 \varphi(x) + \omega \varphi(x), i \partial_0 \varphi(x') + \omega \varphi(x') ] = [i\partial_0 \varphi_x, i \partial_0 \varphi_{x'} ] + [ \omega \varphi_x, i \partial_0 \varphi_{x'}] + [i\partial_0 \varphi_x, \omega \varphi_{x'} ] + \omega^2 [\varphi_x, \varphi_{x'} ]
\end{gathered}
\]
Now $\left[ \varphi_x, \varphi_{x'} \right] =0$, and \\
 $\partial_0 \varphi  = \partial \varphi/\partial t = \dot{\varphi} = \Pi$, for this particular Lagrangian.  We also have

\[
\begin{gathered}
  \left[ \partial_0 \varphi_x, \partial_0 \varphi_{x'} \right] = [ \dot{\varphi}_x, \dot{\varphi}_{x'} ] =0 \\ 
  [\varphi_x, \dot{\varphi}_{x'} ] = i \delta^3(x-x') \\ 
  [ \varphi_0 \varphi_x, \varphi_{x'} ] = [ \dot{\varphi}_x, \varphi_{x'} ] = i - \delta^3(x'-x) = -i \delta^3(x-x')
\end{gathered}
\]
\[
\Longrightarrow \boxed{ [a(k),a(k')] =0  }
\]

$\left[ a^{\dag}(k), a^{\dag}(k') \right] = $?  $a^{\dag}(k) = \int d^3 x e^{ikx} [-i\partial_0 \varphi(x) + \omega \varphi(x)]$  (recall we had imposed the condition, $\varphi = \varphi^*$).  

\[
\begin{gathered}
  \left[ a^{\dag}(k), a^{\dag}(k') \right] = [a^{\dag}_k, a^{\dag}_{k'} ] = \int d^3x d^3 x' e^{i(kx + k'x')} [ -i \partial_0 \varphi(x) + \omega \varphi(x), -i \partial_0 \varphi(x') + \omega \varphi(x') ] = \\
  = \int d^3x d^3 x' e^{i (kx + k'x') } \left( [ -i \partial_0 \varphi_x, -i \partial_0 \varphi_{x_1} ] + [\omega \varphi_x, -i \partial_0 \varphi_{x'} ] + [ -i \partial_0 \varphi_x, \omega \varphi_{x'} ] + \omega^2 [\varphi_x, \varphi_{x'} ] \right)
\end{gathered}
\]
Again $\left[ \varphi_x, \varphi_{x'} \right] = [\dot{\varphi}_x, \dot{\varphi}_{x'} ] = 0$.  \\
$\begin{aligned} \left[ \varphi_x, \dot{\varphi}_{x'} \right] & = i \delta^3(x-x') \\ \left[ \dot{\varphi}_x , \varphi_{x'} \right] & = -i \delta^3(x'-x) = -i\delta^3(x-x') \end{aligned} $

\[
\Longrightarrow \boxed{ \left[ a^{\dag}(k), a^{\dag}(k') \right] =0 }
\]

\[
\begin{gathered}
\left[ a(k), a^{\dag}(k') \right] = \int d^3x d^3x' e^{-i (kx-k'x') } \left( i \partial_0 \varphi(x) + \omega \varphi(x), -i \partial_0 \varphi(x') + \omega \varphi(x') \right)
\end{gathered}
\]
Now
\[
\begin{gathered}
  \left[ i \partial_0 \varphi_x + \omega \varphi_x, -i \partial_0 \varphi_{x'} + \omega \varphi_{x'} \right] = [ \partial_0 \varphi_x , \partial_0 \varphi_{x'} ] - i \omega [ \varphi_x, \partial_0 \varphi_{x'} ] + i\omega [ \partial_0 \varphi_x, \varphi_{x'} ] + \omega^2 [\varphi_x, \varphi_{x'} ] = \\
  = - i \omega (i \delta^3(x - x') ) + i \omega (-i\delta^3(x' - x) ) = 2 \omega \delta^3(x-x')
\end{gathered}
\]

\[
\Longrightarrow \left[ a(k), a^{\dag}(k') \right] = \int d^3x d^3x' e^{-i(kx-k'x') }(2\omega \delta^3(x-x') ) = \int d^3x (2 \omega) e^{-i(k-k')x } = \boxed{ (2\pi)^3 2\omega \delta^3(k-k') }
\]

\problemhead{3.2} Use the commutation relations, eq. (3.29), to show explicitly that a state of the form 
\[
\left| k_1 \dots k_n \rangle \right. \equiv a^{\dag}(\mathbf{k}_1) \dots a^{\dag}(\mathbf{k}_n) \left| 0 \rangle \right.
\]
is an eigenstate of the hamiltonian, eq. (3.30), with eigenvalue $\omega_1 + \dots + \omega_n$.  The vacuum $\left| 0\rangle \right.$ is annihilated by $a(\mathbf{k})$, $a(\mathbf{k})\left| 0 \rangle \right.$, and we take $\Omega_0 = \mathcal{E}_0$ in eq. (3.30).  

\solutionhead{3.2} Recall that 
\[
\begin{aligned}
  & \left[ a(k), a(k') \right] = 0 \\ 
  & [a^{\dag}(k), a^{\dag}(k') ] = 0 \\ 
  & [a(k), a^{\dag}(k') ] = (2\pi)^3 2 \omega \delta^3(k-k')  \quad \quad (3.29)
\end{aligned}
\]
\[
H = \int \tilde{dk} \omega a^{\dag}(k) a(k) + (\mathcal{E}_0 - \Omega_0) V \xrightarrow{\mathcal{E}_0 = \Omega_0 } \int \tilde{dk} \omega a^{\dag}(k) a(k)
\]
Letting $A_0 = (2\pi)^3 2\omega$, $\delta^3(k-k_1) = \delta_{k1}$, 
\[
\begin{aligned}
  H | k_1 \dots k_n \rangle  = \int \tilde{dk} \omega a^{\dag}(k)a(k) a^{\dag}(k_1) \dots a^{\dag}(k_n) | 0 
\end{aligned}
\]
Consider
\[
\begin{aligned}
  a^{\dag}(k) a(k) a^{\dag}(k_1) \dots a^{\dag}(k_n) & = a^{\dag}_k a_k a^{\dag}_1 \dots a_n^{\dag} = a^{\dag}_k (a_1^{\dag} a_k + A_0 \delta^3_{k1} ) a_2^{\dag} \dots a_n^{\dag} = A_0 \delta_{k1} a_k^{\dag}a_2^{\dag}\dots a_n^{\dag} + a^{\dag}_k a_1^{\dag} (a_2^{\dag} a_k + A_0 \delta_{k2}) a_3^{\dag} \dots a_n^{\dag} = \\ 
  & = A_0 \delta_{k1} a_k^{\dag} a_2^{\dag} \dots a_n^{\dag} + A_0 \delta_{k2} a_k^{ \dag} a_1^{\dag} a_3^{\dag} \dots a_n^{\dag} + a_k^{\dag} a_1^{ \dag} a_2^{\dag} (a_3^{\dag} a_k + A_0 \delta_{k3}) a_4^{\dag} \dots a_n^{\dag} = \dots
\end{aligned}
\]

By induction, 
\[
\Longrightarrow \left(A_0 \sum_{j=1}^n \delta_{kj} a_k^{\dag} \prod_{l\neq k} a_k^{\dag} + a_k^{\dag}a_1^{\dag} \dots a_n^{\dag} a_k \right) | 0 \rangle
\]

Note that $a_k | 0 \rangle$.  

So note $\omega = \omega(k) = (k^2- m^2)^{1/2}$ and letting $\omega(k_j) = \omega_j$, 
\[
H |k_1 \dots k_n \rangle = \sum_{j=1}^n \omega_j a_1^{\dag} \dots a_n^{\dag}  | 0 \rangle = \left( \sum_{j=1}^n \omega_j \right) | k_1 \dots k_n \rangle
\]
Note that $A_0$ factor, $A_0 = (2\pi)^3 2\omega$ canceled by our Lorentz invariant differential, $\tilde{ dk} = \frac{d^3k}{ (2\pi)^3 2 \omega}$ in the integration.  

\problemhead{3.3} Use $U(\Lambda)^{-1} \varphi(x) U(\Lambda) = \varphi(\Lambda^{-1} x)$ to show that 
\[
\begin{aligned}
  U(\Lambda)^{-1} a(\mathbf{k})U(\Lambda) & = a(\Lambda^{-1} \mathbf{k}), \\ 
  U(\Lambda)^{-1} a^{\dag}(\mathbf{k})U(\Lambda) & = a^{\dag}(\Lambda^{-1} \mathbf{k}) \quad \quad \, (3.34)
\end{aligned}
\]
and hence that 
\[
U(\Lambda)|k_1 \dots k_n \rangle = | \Lambda k_1 \dots \Lambda k_2 \rangle, \quad \quad \, (3.35)
\]
where $|k_1 \dots k_n \rangle = a^{\dag}(\mathbf{k}_1 \dots a^{\dag}(\mathbf{k}_n) |0 \rangle$ is a state of $n$ particles with momenta $k_1, \dots, k_n$.  

\solutionhead{3.3} (Incomplete, 20090205) Given $U(\Lambda)^{-1} \varphi(x) U(\Lambda) = \varphi(\Lambda^{-1} x)$, 
\[
\begin{aligned}
  U(\Lambda)^{-1} a(k) U(\Lambda) & = U(\Lambda)^{-1} \int d^3 x e^{-ikx} [ i \partial_0 \varphi(x) + \omega \varphi(x) ] U(\Lambda) = \int d\vec{x} e^{-ikx} \left( i U(\Lambda)^{-1} \partial_0 \varphi(x) U(\Lambda) + \omega U(\Lambda)^{-1} \varphi(x) U(\Lambda) \right) = \\ 
  & = \int d\vec{x} e^{-ikx} \left( i U(\Lambda)^{-1} \partial_0 \varphi(x) U(\Lambda) + \omega \varphi(\Lambda^{-1} x) \right)
\end{aligned}
\]
Now  
\[
U(\Lambda)^{-1} \partial_0 \varphi(x) U(\Lambda) = \partial_0 U(\Lambda)^{-1} \varphi(x) U(\Lambda) = \partial_0 \varphi(\Lambda^{-1}x) 
\]
So then for the expression before,
\[
\Longrightarrow = \int d\vec{x} e^{-ikx} \left( i \partial_0 \varphi(\Lambda^{-1} x) + \omega \varphi(\Lambda^{-1} x) \right) = \int d\vec{x} e^{-i(\Lambda^{-1}k )( \Lambda^{-1} x ) } \left( i \partial_0 \varphi(\Lambda^{-1} x ) + \omega \varphi(\Lambda^{-1} x) \right) = a(\Lambda^{-1} k)
\]

\problemhead{3.4} Recall that $T(a)^{-1}\varphi(x) T(a) = \varphi(x-a)$, where $T(a) \equiv \exp{(-i P^{\mu}a_{\mu})}$ is the space-time translation operator, and $P^0$ is identified as the hamiltonian $H$.  
\begin{itemize}
\item[(a)] Let $a^{\mu}$ be infinitesimal, and derive an expression for $[ P^{\mu}, \varphi(x)]$.  
\item[(b)] Show that the time component of your result is equivalent to the Heisenberg equation of motion $\dot{\varphi} = i[H,\varphi ]$.  
\item[(c)] For a free field, use the Heisenberg equation to derive the Klein-Gordon equation. 
\item[(d)] define a spatial momentum operator
\[
\mathbf{P} \equiv - \int d^3x \Pi(x) \nabla \varphi(x)
\]
Use the canoncial commutation relations to show that $\mathbf{P}$ obeys the relation you derived in part (a). 
\item[(e)]Express $\mathbf{P}$ in terms of $a(\mathbf{k})$ and $a^{\dag}(\mathbf{k})$.  
\end{itemize}

\solutionhead{3.4} INCOMPLETE (20090218).  
\begin{itemize}
\item[(a)] Now $T(a)^{-1} \varphi(x) T(a) = \varphi(x-a)$.  $T(a) \equiv \exp{ (-iP^{\mu} a_{\mu})}$.  

Let $a^{\mu}$ be infinitesimal.  Then
\[
\begin{gathered} T(a)^{-1} \varphi(x) T(a) = ( 1 + i P^{\mu} a_{\mu})\varphi(x) (1 - i P^{\mu} a_{\mu}) = \varphi(x) + i P^{\mu} a_{\mu} \varphi(x) - i \varphi(x) P^{\mu} a_{\mu} = \varphi(x) + (\partial^{\mu} \varphi)a_{\mu} \\ 
\quad \quad \text{(keep up to $O(a_{\mu})$ terms)}
\end{gathered}
\]
\[
\Longrightarrow i[P^{\mu}, \varphi(x) ] = \partial^{\mu} \varphi
\]
\item[(b)]
\item[(c)]
\item[(d)]
\item[(e)]
\end{itemize}

\problemhead{3.5} Consider a complex (that is, nonhermitian) scalar field $\varphi$ with lagrangian density 
\[
\mathcal{L} = - \partial^{\mu} \varphi^{\dag} \partial_{\mu} \varphi - m^2 \varphi^{\dag} \varphi + \Omega_0.  \quad \quad \, (3.37)
\]
\begin{itemize}
\item[(a)] Show that $\varphi$ obeys the Klein-Gordon equation.  
\item[(b)] Treat $\varphi$ and $\varphi^{\dag}$ as independent fields, and find the conjugate momentum for each.  Compute the hamiltonian density in terms of these conjugate momenta and the fields themselves (but not their time derivatives).  
\item[(c)] Write the mode expansion of $\varphi$ as 
\[
\varphi(x) = \int \tilde{dk} [ a(\mathbf{k}) e^{ikx} + b^{\dag}(\mathbf{k}) e^{-ikx} ]. \quad \quad \,  (3.38)
\]
\item[(d)] Assuming canonical commutation rleations for the fields and their conjugate momenta, find the commutation relations obeyed by $a(\mathbf{k})$ and $b(\mathbf{k})$ and their hermitian conjugates.  
\item[(e)] Express the hamiltonian in terms of $a(\mathbf{k})$ and $b(\mathbf{k})$ and their hermitian conjugates.  What value must $\Omega_0$ have in order for the ground state to have zero energy?  
\end{itemize}

\solutionhead{3.5} \begin{itemize}
\item[(a)] The action is 
\[
S = \int d^4 \mathcal{L} = \int d^4x (-\partial^{\mu} \varphi^{\dag} \partial_{\mu} \varphi - m^2 \varphi^{\dag} \varphi + \Omega_0 ) 
\]
\[
\Longrightarrow \delta S = \int d^4 x ( - \partial^{\mu} \delta \varphi^{\dag} \partial_{\mu} \varphi - \partial^{\mu} \varphi^{\dag} \partial_{\mu} \delta \varphi - m^2 \delta \varphi^{\dag} \varphi - m^2 \varphi^{\dag} \delta \varphi ) 
\]
From integration by parts, 
\[
\begin{aligned}
  \int d^4 x \partial^{\mu} \delta \varphi^{\dag} \partial_{\mu} \varphi & = \int_{\partial \Omega}(dS \cdot \delta \varphi^{\dag} ) \partial_{\mu} \varphi - \int d^4 x \delta \varphi^{\dag} \partial^{\mu} \partial_{\mu} \varphi \\ 
  \int d^4 x \partial^{\mu}  \varphi^{\dag} \partial_{\mu}\delta  \varphi & = \int_{\partial \Omega}(dS \cdot \delta \varphi ) \partial^{\mu} \varphi^{\dag} - \int d^4 x \delta \varphi \partial_{\mu} \partial^{\mu} \varphi^{\dag} \\ 
\end{aligned}
\]
\[
\Longrightarrow \delta S = \int d^4 x \lbrace (\partial^{\mu} \partial_{\mu} \varphi - m^2 \varphi) \delta \varphi^{\dag} + (\partial_{\mu} \partial^{\mu} \varphi^{\dag} - m^2 \varphi^{\dag} ) \delta \varphi \rbrace = 0 \quad \, \text{(extremize action)} 
\]
We obtain Klein-Gordon:
\[
\Longrightarrow (\partial^{\mu} \partial_{\mu} - m^2) \varphi = (\partial_{\mu} \partial^{\mu} - m^2 ) \varphi^{\dag} =0 
\]
If you wanted to make this explicit, note that $\delta \varphi + \delta \varphi^{\dag} = 2 \Re{(\delta \varphi)}$, or assume treat $\varphi$, $\varphi^{\dag}$ as independent (which we assume in the next part).  
\item[(b)] $\partial^{\mu} \varphi^{\dag} \partial_{\mu} \varphi$ includes $\dot{\varphi}^{\dag} \dot{\varphi}$ term in the sum.  

In Srednicki's notation, $\partial^{\mu} \varphi^{\dag} \partial_{\mu} \varphi = - \dot{\varphi}^{\dag} \dot{\varphi} + \mathbf{\nabla} \varphi^{\dag} \cdot \mathbf{\nabla} \varphi$.  

The conjugate momenta are
\[
\begin{aligned}
  \Pi_{\varphi} & = \frac{ \partial \mathcal{L} }{\partial \dot{\varphi }} =  \dot{\varphi}^{\dag} \\  
  \Pi_{\varphi^{\dag}} & =  \dot{\varphi}
\end{aligned}
\]
Recall $\mathcal{H} = \Pi \dot{\varphi} - \mathcal{L}$, definition of Hamiltonian.  
\[
\begin{aligned}
  \mathcal{H} & = \Pi_{\varphi} \dot{\varphi} + \Pi_{\varphi^{\dag}} \dot{\varphi^{\dag}} - \left( - \partial^{\mu} \varphi^{\dag} \partial_{\mu} \varphi - m^2 \varphi^{\dag} \varphi + \Omega_0 \right) = + \Pi_{\varphi} \Pi_{\varphi^{\dag}} + \Pi_{\varphi^{\dag}} \Pi_{\varphi} + \partial^{\mu} \varphi^{\dag} \partial_{\mu} \varphi + m^2 \varphi^{\dag} \varphi - \Omega_0 = \\ 
  & = \boxed{  \Pi_{\varphi} \Pi_{\varphi^{\dag}} + \mathbf{\nabla} \varphi^{\dag} \cdot \mathbf{\nabla} \varphi + m^2 \varphi^{\dag} \varphi - \Omega_0 }
\end{aligned}
\]
\item[(c)]Writing the mode expansion of $\varphi$ as 
\[
\varphi(x) = \int \widetilde{dk} [ a(\mathbf{k}) e^{ikx} + b^{\dag}(\mathbf{k}) e^{-ikx} ] = \int \widetilde{dk} (a_k e^{ikx} + b^{\dag}_k e^{-ikx} )
\]
Then 
\[
\begin{aligned}
  \varphi^{\dag} & = \int \tilde{dk} (a_k^{\dag} e^{-ikx} + b_k e^{ikx} ) \\ 
  \dot{\varphi} & = \int \tilde{dk} ( a_k(-i\omega) e^{ikx} + b^{\dag}_k(i\omega) e^{-ikx} ) \\ 
  \dot{\varphi}^{\dag} & = \int \tilde{dk} (a_k^{\dag} (i\omega) e^{-ikx} + b_k (-i\omega) e^{ikx} 
\end{aligned}
\]

\[
\begin{aligned}
  \int d^3x e^{-ik'x} \varphi^{\dag}(x) & = \int d^3x e^{-ik'x} \int \tilde{k} (a_k^{\dag}e^{-ikx} + b_k e^{ikx} ) = \int \tilde{k} a_k^{\dag} e^{i(\omega + \omega') t} (2\pi)^3 \delta(k+k') + \int \tilde{dk} b_k(2\pi)^3 \delta(k-k')  = \\ 
  & = \frac{ a^{\dag}(-\mathbf{k}) }{2\pi } e^{2i\omega t} + \frac{b(\mathbf{k}) }{2\omega} 
\end{aligned}
\]

\[
\begin{aligned}
  \int d^3x e^{-ik'x} \partial_0 \varphi^{\dag}(x) & = \int d^3x e^{-ik'x} \int \tilde{k} (a_k^{\dag}(i\omega) e^{-ikx} + b_k (-i\omega) e^{ikx} ) = i\int \tilde{k} \omega \int d^3 x \left( e^{-ix (k'+k) } a_k^{\dag} - e^{ix(k -k')} b_k \right) = \\
  & = \frac{i}{2} e^{2 i \omega t} a^{\dag}(-\mathbf{k}) - \frac{i}{2} b(\mathbf{k}) 
\end{aligned}
\]
Likewise,
\[
\begin{aligned}
  \int d^3 x e^{-ikx } \varphi(x) & = \frac{ a(\mathbf{k})}{2 \omega} + \frac{ b^{\dag}(-\mathbf{k}) e^{2i \omega t}}{ 2\omega} \\ 
  \int d^3 x e^{-ikx } \partial_0 \varphi(x) & = \frac{-i a(\mathbf{k})}{2} + \frac{i b^{\dag}(-\mathbf{k}) e^{2 i \omega t}}{ 2} \\ 
  \int d^3 x e^{-ikx } \varphi^{\dag}(x) & = \frac{ a^{\dag}(-\mathbf{k}) }{2 \omega}e^{2i\omega t} + \frac{ b(\mathbf{k})}{ 2\omega} \\ 
  \int d^3 x e^{-ikx } \partial_0 \varphi^{\dag}(x) & = \frac{ ia^{\dag}(-\mathbf{k})}{2 } e^{2i\omega t}+ \frac{ -i b(\mathbf{k})}{ 2} \\ 
\end{aligned} \quad \quad \, \Longrightarrow \boxed{ \begin{aligned} a(\mathbf{k}) & = \int d^3x e^{-ikx} (\omega \varphi(x) + i \partial_0 \varphi )  \\ b(\mathbf{k}) & = \int d^3x e^{-ikx} (\omega \varphi^{\dag}(x) + i \partial_0 \varphi^{\dag}(x) ) \end{aligned} }
\]
\item[(d)] Assume \\
\[
\begin{aligned}
   \left[ \varphi(x,t), \varphi(x',t) \right] & =0 \\ 
   [ \Pi(x,t), \Pi(x',t) ] & = 0 \\ 
   [\varphi(x,t), \Pi(x',t) ] & = i \delta^3(x-x')
\end{aligned}
\]
In our case $\varphi$, $\varphi^{\dag}$ treated as independent fields.  Recall $\Pi_{\varphi} = - \dot{\varphi}^{\dag}$, $\Pi_{\varphi^{\dag}} = - \dot{\varphi}$.  

Assume further, treating $\varphi, \varphi^{\dag}$ as independent, 
\[
\begin{aligned}
& \left[ \varphi(x,t), \varphi^{\dag}(x',t) \right] =0   \\ 
& \left[ \Pi_{ \varphi}(x,t), \Pi_{\varphi}(x',t) \right] = \left[ \Pi_{ \varphi^{\dag}}(x,t), \Pi_{\varphi^{\dag}}(x',t) \right] =  \left[ \Pi_{ \varphi}(x,t), \Pi_{\varphi^{\dag}}(x',t) \right] = 0 \\ 
&  [\varphi(x,t), \Pi_{\varphi}(x',t) ]  = i \delta^3(x-x'), \quad \,  [\varphi^{\dag}(x,t), \Pi_{\varphi^{\dag}}(x',t) ]  = i \delta^3(x-x')
\end{aligned}
\]
\[
\begin{aligned}
  \left[ a(k), a^{\dag}(k')  \right] & = \int d^3x \int d^3x' e^{-ikx} e^{ik'x'} [ \omega \varphi(x) + i \partial_0 \varphi, \omega_{k'} \varphi^{\dag}(x') - i \partial_0 \varphi^{\dag} ] = \\ 
  & = \int d^3x d^3x' e^{-i (kx-k'x') } \left( i \omega_{k'} [-\Pi_{\varphi^{\dag}}, \varphi^{\dag} ] + i \omega [ \varphi(x), \Pi_{\varphi} ] \right) = \\ 
  & = \int d^3x d^3x' e^{-i (kx-k'x') } \left( i \omega_{k'} i \delta^3(x'-x) + - \omega \delta^3(x-x') \right) = - \int d^3x e^{-i(kx - k'x) } 2\omega_k = \boxed{ -(2\pi)^3 (2\omega) \delta^3(k-k') }
\end{aligned}
\]

\[
\begin{aligned}
  \left[ b(k), b^{\dag}(k')  \right] & = \int d^3x \int d^3x' e^{-ikx} e^{ik'x'} [ \omega \varphi^{dag}(x) + i \partial_0 \varphi^{\dag}, \omega_{k'} \varphi(x') - i \partial_0 \varphi(x') ] = \\ 
  & = \int d^3x d^3x' e^{-i (kx-k'x') } \left( i \omega [-\Pi_{\varphi}, \varphi ] + i \omega [ \varphi^{\dag}(x), \Pi_{\varphi^{\dag}} ] \right) = \\ 
  & = \int d^3x d^3x' e^{-i (kx-k'x') }  i \omega \left(i \delta^3(x'-x) +  i\omega \delta^3(x-x') \right) = - \int d^3x e^{-i(kx - k'x) } 2\omega = \boxed{ -(2\pi)^3 (2\omega) \delta^3(k-k') }
\end{aligned}
\]

\[
\left[ a(k), a(k') \right] = \int d^3x d^3x' e^{-ikx} e^{-ik'x'} [ \omega \varphi(x) + i \partial_0 \varphi, \omega_{k'} \varphi(x') + i \partial_0 \varphi(x') ] = 0 \text{ since } \begin{aligned} \left[ \varphi(x), \varphi(x') \right] & = 0 \\ \left[ \partial_0 \varphi(x), \partial_0 \varphi^{\dag}(x') \right] & = 0 \\ \left[ \varphi(x), \Pi_{\varphi^{\dag} }(x') \right] & = 0 \end{aligned}  
\]

\[
\left[ a^{\dag}(k), a^{\dag}(k') \right] = \int d^3x d^3x' e^{ikx} e^{ik'x'} [ \omega \varphi^{\dag}(x)  i \partial_0 \varphi^{\dag}, \omega_{k'} \varphi^{\dag}(x') - i \partial_0 \varphi^{\dag}(x') ] = 0 \text{ since } \begin{aligned} \left[ \varphi^{\dag}(x), \varphi^{\dag}(x') \right] & = 0 \\ \left[ \partial_0 \varphi^{\dag}(x), \partial_0 \varphi^{\dag}(x') \right] & = 0 \\ \left[ \varphi^{\dag}(x), \Pi_{\varphi }(x') \right] & = 0 \end{aligned}  
\]

Looking at $b(k)$ comprised of $\varphi^{\dag}(x)$, $\partial_0 \varphi^{\dag}(x)$ and $b^{\dag}(k)$ comprised of $\varphi(x)$, $\partial_0 \varphi(x)$.  \\
\phantom{Looking} Then, likewise with $a,a^{\dag}$, 
\[
\begin{aligned}
  \left[ b(k), b(k') \right] & =0  \\ 
  \left[ b^{\dag}(k), b^{\dag}(k') \right] & = 0 
\end{aligned}
\]

\[
\left[ a(k), b(k') \right] = \int d^3x d^3x' e^{-ikx} e^{-ik'x'} \left[ \omega \varphi(x) + i \partial_0 \varphi(x), \omega_{k'} \varphi(x') + i \partial_0 \varphi(x') \right] = 0 
\]
since 
\[
\begin{aligned}
  &  \left[ \varphi(x), \varphi(x') \right] = 0 \\ 
  & \left[ \Pi_{\varphi}(x), \Pi_{\varphi}(x') \right] = 0 
\end{aligned} 
\quad \quad \, \begin{aligned} \left[ \varphi(x), \partial_0 \varphi(x') \right] = \left[ \varphi(x) , - \Pi_{\varphi^{\dag} }(x') \right] = 0 \\ [ \partial_0 \varphi(x), \varphi(x') ] =0 \end{aligned}
\]

Likewise for $\left[ a^{\dag}(k), b^{\dag}(k') \right] = 0$.  

It can be shown that 
\[
\begin{aligned}
  \left[ a(k), b^{\dag}(k') \right] & = 0
\end{aligned}
\]

\[
\begin{aligned}
  \left[ a^{\dag}(k), b(k') \right] & = 0
\end{aligned}
\]

because, recall
\[
\begin{aligned}
  & a(k) = \int d^3x e^{-ikx} [ \omega \phi(x) + i \partial_0 \phi ] = -i \int d^3x e^{-ikx} \overleftrightarrow{\partial_0} \phi = \int d^3x e^{-ikx} [ \omega \phi(x) + i (- \Pi_{\phi^{\dag}} ) ] \\
  & b^{\dag}(k') = \int d^3x' e^{ik'x'} [ \omega \phi(x') - i \partial_0 \phi(x') ] = i \int d^3x' e^{ik'x'} \overleftrightarrow{\partial_{0'}} \phi(x') = \int d^3x e^{ik'x} [ \omega \phi + i \Pi_{\phi^{\dag}} ]
\end{aligned}
\]

\[
\begin{aligned}
  \Pi_{\phi} & = - \dot{\phi}^{\dag} \\ 
  \Pi_{\phi^{\dag}} & = - \dot{\phi}
\end{aligned}
\]

\[
\begin{aligned}
  \Longrightarrow \left[ a(k), b^{\dag}(k') \right] & = \int d^3x \int d^3x' e^{ikx} e^{ik'x'} [ \omega \phi_x - i \Pi_{\phi^{\dag}(x)} , \omega \phi_{x'} + i \Pi_{\phi^{\dag}(x') } ] = \\
  & = \int d^3x d^3x' e^{ikx} e^{ik'x'} ( i\omega [\phi(x), \Pi_{\phi^{\dag}}(x') ] - i \omega [ \Pi_{\phi^{\dag}(x) }, \phi(x') ] ) = 0 
\end{aligned}
\]
(must assume treating $\phi$, $\phi^{\dag}$ independent: then $[\phi(x), \Pi_{\phi^{\dag} }(x')] = [\phi^{\dag}(x),\Pi_{\phi}(x)] = 0 $)

\item[(e)] 
\[
\begin{gathered}
\mathcal{H}  =  \Pi_{\varphi }\Pi_{\varphi^{\dag}} + \vec{\nabla} \varphi^{\dag} \cdot \vec{\nabla} \varphi + m^2 \varphi^{\dag} \varphi - \Omega_0 =   \int \widetilde{dk}' ((i\omega')a_{k'}^{\dag}  e^{-ik'x} + (-i\omega')b_{k'} e^{ik'x} )\int \widetilde{dk}  ( (-i\omega)a(k)  e^{ikx} + (i\omega)b^{\dag}(k)  e^{-ikx} ) + \\
+ \int \widetilde{dk} \lbrace -i \mathbf{k} a_k^{\dag} e^{-ikx} + i\mathbf{k} b_k e^{ikx} \rbrace \cdot \int \widetilde{dk'} \lbrace i\mathbf{k}' a_{k'} e^{ik'x} - i \mathbf{k}' b_{k'} e^{-ik'x} \rbrace + \\ + m^2 \int \widetilde{dk} \int \widetilde{dk}' (a_{k'}^{\dag} e^{-ik'x} + b_{k'} e^{ik'x} )( a_k e^{ikx} + b_k^{\dag}e^{-ikx} ) - \Omega_0 
\end{gathered}
\]
where $\nabla = \partial_x$, \\
$ikx = i (-\omega t + k_j x_k )$, \\
$\partial_{x_j} e^{ikx} = ik_j e^{ikx}$.  

\[
\begin{gathered}
  \mathcal{H} = \int \widetilde{dk}' \widetilde{dk} \lbrace \omega' \omega a_{k'}^{\dag} a_k e^{-i (k' -k)x } - \omega' \omega b_{k'} a_k e^{i (k' + k)x} - \omega' \omega a_{k'}^{\dag} b_k^{\dag} e^{-i(k'+k) x} + \omega' \omega b_{k'} b_k^{\dag} e^{i(k'-k) x} \rbrace + \\ 
  + \int \widetilde{dk} \widetilde{dk}' \lbrace \mathbf{k} \cdot \mathbf{k}' a_k^{\dag} a_k e^{-i(k -k')x} - \mathbf{k}\cdot \mathbf{k}' b_k a_{k'} e^{i(k+k')x} - \mathbf{k}\cdot \mathbf{k}' a_k^{\dag} b_{k'}^{\dag} e^{-i(k+k')x} + \mathbf{k}\cdot \mathbf{k}' b_k b_{k'}^{\dag} e^{i(k-k')x} \rbrace  + \\
  + m^2 \int \widetilde{dk} \widetilde{dk}' \lbrace a_{k'}^{\dag} a_k e^{-i(k'-k)x} + b_{k'} a_k e^{i (k'+k) x} + a_{k'}^{\dag} b_k^{\dag} e^{-i(k'+k)x} + b_{k'}b_k^{\dag} e^{i(k'-k)x} \rbrace
\end{gathered}
\]

Consider doing $\int \widetilde{dk}' \int d^3 x = \int \frac{d^3k'}{2\omega'(2\pi)^3 } \int d^3x$.  \\

$\int d^3x e^{-i(k'-k)x} = e^{i(\omega' - \omega)t} (2\pi)^3 \delta^3(\mathbf{k}' - \mathbf{k}) = (2\pi)^3 \delta^3(\mathbf{k'}-\mathbf{k})$.  \\
Note how $\omega_{\mathbf{k}} = \sqrt{ \mathbf{k}^2 + m^2}$, for the last step.  

Then also
\[
\int d^3x e^{i(k'+k)x} = e^{-i(\omega' + \omega) t} (2\pi)^3 \delta^3(\mathbf{k}' + \mathbf{k}) = e^{-2i\omega t}(2\pi)^3 \delta^3(\mathbf{k}'+\mathbf{k})
\]
\emph{Remember}, $\delta^3(\mathbf{k}'+\mathbf{k})$ makes $\mathbf{k}' = - \mathbf{k}$.  

Likewise,
\[
\int d^3x e^{i(k'-k)x} = e^{i(\omega - \omega')t} (2\pi)^3 \delta^3(\mathbf{k}' - \mathbf{k}) = (2\pi)^3 \delta^3(\mathbf{k}'-\mathbf{k})
\]
\[
\int d^3x e^{-i(k'+k)x} = e^{i(\omega' + \omega) t} (2\pi)^3 \delta^3(\mathbf{k}' + \mathbf{k}) = e^{2i\omega t}(2\pi)^3 \delta^3(\mathbf{k}'+\mathbf{k})
\]

\[
\begin{gathered}
  H = \int d^3x \mathcal{H} = \int \widetilde{dk} \lbrace \frac{ \omega}{2} a_k^{\dag} a_k - \frac{ \omega}{2} b_k a_k e^{-2i\omega t} - \frac{ \omega}{2} a_k^{\dag} b_k^{\dag} e^{2i\omega t} + \frac{ \omega }{2} b_k b_k^{\dag} \rbrace + \\ 
  + \int \widetilde{dk} \lbrace \frac{ \mathbf{k}^2 }{2\omega } a_k^{\dag} a_k + \frac{ \mathbf{k}^2 b_k a_{k'} e^{-2i\omega t} }{ 2\omega } + \frac{ \mathbf{k}^2 a_k^{\dag} b_{k'}^{\dag} e^{2i\omega t} }{2\omega} + \frac{ \mathbf{k}^2 b_k b_k^{\dag} }{ 2\omega } \rbrace + \\ 
  + m^2 \int \widetilde{dk} \lbrace \frac{ a_k^{\dag} a_k}{ 2\omega} + \frac{ b_k a_k}{2\omega} e^{-2i\omega t} + \frac{ a_k^{\dag} b_k^{\dag}}{2\omega } e^{2i\omega t} + \frac{ b_k b_k^{\dag}}{2\omega} \rbrace + - \Omega_0 V
\end{gathered}
\]

Note that \\
$\begin{aligned}
& m^2 = \omega^2 - \mathbf{k}^2 \\ 
& \omega^2 + \mathbf{k}^2 + m^2 = \omega^2 + \mathbf{k}^2 + \omega^2 - \mathbf{k}^2 = 2\omega^2
\end{aligned}$

\[
\Longrightarrow H  =\int \widetilde{dk} \omega \lbrace a_k^{\dag} a_k + b_k b_k^{\dag} \rbrace + - \Omega_0 V
\]

Now $b_k b_k^{\dag} = b_k^{\dag} b_k + (2\pi)^3 2\omega \delta^3(\mathbf{k} - \mathbf{k}') = b_k^{\dag} b_k + (2\pi)^3 2\omega$.  

\[
H = \int \widetilde{dk} \omega \lbrace a^{\dag}_k a_k + b_k^{\dag} b_k \rbrace + (\mathcal{\epsilon}_0 - \Omega_0) V 
\]
where 
\[
\mathcal{\epsilon}_0 V = \int \widetilde{dk} (2\pi)^3 2\omega^2 = \int \frac{ d^3k}{(2\pi)^3 2\omega } (2\pi)^3 2\omega^2 = \int d^3k \omega_{\mathbf{k}}
\]

For $\langle 0 | H | 0 \rangle = 0$, then $\mathcal{\epsilon}_0 V = \Omega_0 V$ or $\int d^3k \omega_{\mathbf{k}} = \Omega_0 V$.  \\

$V = (2\pi)^3 \delta^3(0)$ interpreted as volume of space $V$.  
\[
\mathcal{\epsilon}_0 = \frac{1}{(2\pi)^3} \int d^3k \omega_{\mathbf{k}} = 2\left( \frac{1}{2} (2\pi)^{-3} \int d^3k \omega \right)
\]
Twice the zero-point energy was found, due to 2 kinds of particles associated with the $\varphi$, $\varphi^{\dag}$ fields.  
\end{itemize}

\section{The spin-statistics theorem}

\problemhead{4.1} Verify eq. (4.12).  Verify its limit as $m \to 0$.  

\solutionhead{4.1} Recall that we defined a nonhermitian field
\[
\begin{aligned}
	\varphi^{\dag}(x,0) & \equiv \int \widetilde{dk} e^{ik\cdot x} a(\mathbf{k}) & \quad \quad \, (4.3) \\ 
	\text{ its Hermitian conjugate } \varphi^-(x,0) & \equiv \int \widetilde{dk} e^{-ik\cdot x} a^{\dag}(\mathbf{k}) & \quad \quad \, (4.4)
\end{aligned}
\]
Time-evolved with $H_0$: 
\[
\begin{aligned}
	\varphi^{+}(\mathbf{x},t) & = e^{iH_0 t} \varphi^{+}(\mathbf{x},0) e^{-i H_0 t} & = \int \widetilde{dk} e^{ikx} a(\mathbf{k}) \\ 
	\varphi^-(\mathbf{x},t) & = e^{iH_0 t} \varphi^{-}(\mathbf{x},0)e^{-i H_0 t} & = \int \widetilde{dk} e^{-ikx} a^{\dag}(\mathbf{k})  \quad \quad \, (4.5)
\end{aligned}
\]

\[
\begin{gathered}
		[\varphi^+(x), \varphi^-(x')]_{\mp} = \int \widetilde{dk} \int \widetilde{dk'} e^{i (kx -k'x')} [a(\mathbf{k}), a^{\dag}(\mathbf{k}')]_{\mp} = \\
			= \int \widetilde{dk} \widetilde{dk'} e^{i(kx-k'x')} (2\pi)^3 2\omega \delta^3(\mathbf{k}-\mathbf{k}') \quad \quad \, \text{(by fundamental canoncial commutation (or anticommutation) relations)} = \\	
	= \int \widetilde{dk} e^{ik(x-x')}
\end{gathered}
\]

INCOMPLETE (20090218) Look up how to do modified Bessel functions for functions with symmetry about k.  





\section{The LSZ reduction formula}

\problemhead{5.1} Work out the LSZ reduction formula for the complex scalar field that was introduced in problem 3.5.  Note that we must specify the type ($a$ or $b$) of each incoming and outgoing particle.

\solutionhead{5.1} 

Recall,
\[
\begin{aligned}
  a^{\dag}(\mathbf{k}) & = \int d^3 x e^{ikx} [ \omega \phi^{\dag}(x) - i \partial_0 \phi^{\dag} ]  = - i \int d^3x e^{ikx} \overleftrightarrow{\partial_0} \phi^{\dag} \\ 
  b^{\dag}(\mathbf{k}) & = \int d^3 x e^{ikx} [ \omega \phi(x) - i \partial_0 \phi(x) ]  = - i \int d^3x e^{ikx} \overleftrightarrow{\partial_0} \phi(x)  \\
  a(\mathbf{k}) & = \int d^3x e^{-ikx} [ \omega \phi(x) + i \partial_0 \phi(x)] = i \int d^3x e^{-ikx} \overleftrightarrow{\partial_0} \phi
\end{aligned}
\]

Note there are 2 different types of particles.  

Define $f_1(\mathbf{k}) \propto \exp{ [-(\mathbf{k} - \mathbf{k}_1)^2/4\sigma^2 ]} $
\[
\begin{aligned}
  a_1^{\dag} & = \int d^3k f_1(\mathbf{k}) a^{\dag}(k) \\ 
  b_2^{\dag} & = \int d^3k f_2(\mathbf{k}) b^{\dag}(k) 
\end{aligned}
\]

Guess, for interacting theory,
\[
\begin{aligned}
  |i\rangle & = \lim_{t \to -\infty} a_1^{\dag}(t) a_2^{\dag}(t) |0 \rangle \text{ or } \\ 
  |i\rangle & = \lim_{t \to -\infty} b_1^{\dag}(t) b_2^{\dag}(t) |0 \rangle \text{ or } \\ 
  |i\rangle & = \lim_{t \to -\infty} a_1^{\dag}(t) b_2^{\dag}(t) |0 \rangle \text{ or in general} \\ 
  |i \rangle & = \lim_{t \to -\infty} a_1^{\dag}(t) \dots a_m^{\dag}(t) b_{m+1}^{\dag}(t) \dots b_n^{\dag}(t) |0 \rangle
\end{aligned}
\]
(note $[a_1^{\dag}, b_2^{\dag}] = \int d^3k f_1(\mathbf{k}) \int d^3k' f_2(\mathbf{k}') [a^{\dag}(k), b^{\dag}(k')] = 0$)

Recall
\[
\begin{gathered}
  a^{\dag}_1(+\infty) - a_1^{\dag}(-\infty) = \int_{-\infty}^{\infty} dt \partial_0 a_1^{\dag}(t) = - i \int d^3k f_1(\mathbf{k}) \int d^4x \partial_0 (e^{ikx} (i \omega \phi^{\dag}(x) + \partial_0 \phi^{\dag}(x) ) ) = \\
  = -i \int d^3k f_1(\mathbf{k}) \int d^4x e^{ikx} (\partial_0^2 + \omega^2) \phi^{\dag}(x) = -i \int d^3k f_1(\mathbf{k}) \int d^4x e^{ikx} (\partial_0^2 + \mathbf{k}^2 + m^2) \phi^{\dag}(x) =\\ 
   =-i \int d^3k f_1(\mathbf{k}) \int d^4x e^{ikx} (-\partial_x^2 + m^2) \phi^{\dag}(x)
\end{gathered}
\]
So, taking the limit for $f(\mathbf{k})$, 
\[
a_1^{\dag}(+\infty) - a_1^{\dag}(-\infty) = -i \int d^4 x e^{ik_1 x} (-\partial_1^2 + m^2 ) \phi^{\dag}(x)
\]

For the Hermitian conjugate, 
\[
\begin{gathered}
  a_1(+\infty) - a_1(-\infty) = \int_{-\infty}^{\infty} dt \partial_0 a_1(t) = -i \int d^3k f_1(\mathbf{k}) \int d^4x \partial_0 (e^{-ikx} [ i\omega \phi(x) - \partial_0 \phi(x)] ) = \\ 
  = -i \int d^3k f_1(\mathbf{k}) \int d^4x e^{-ikx} (-\omega^2  - \partial_0^2 ) \phi = + i \int d^3k f_1(k) \int d^4x e^{-ikx} (\mathbf{k}^2 + m^2 + \partial_0^2) \phi(x) = \\ 
= i \int d^4x e^{-ik_1x} (-\partial_1^2 + m^2 )\phi(x)
\end{gathered}
\]

Likewise,
\[
\begin{aligned}
  & b_2^{\dag}(+\infty) - b_2^{\dag}(-\infty) = -i \int d^4x e^{ik_2 x} (-\partial_2^2 + m^2)\phi(x) \\ 
  & b_2(+\infty) - b_2(-\infty) = i \int d^4x e^{-ik_2 x} (-\partial_2^2 + m^2) \phi^{\dag}(x)
\end{aligned}
\]


Now
\[
\begin{aligned}
  & \langle f | = \lim_{t \to +\infty} \langle 0 | a_{1'}(t) a_{2'}(t) \text{ or } \\ 
  & \langle f | = \lim_{t \to +\infty} \langle 0 | b_{1'}(t) b_{2'}(t)  \\ 
  & \langle f | = \lim_{t \to +\infty} \langle 0 | a_{1'}(t) b_{2'}(t) \\ 
  & \langle f | = \lim_{t \to +\infty} \langle 0 | a_{1'}(t) \dots a_{m'}(t) b_{1'}(t) \dots b_{n'}(t)
\end{aligned}
\]

Note that the only nontrivial commutation relations are $[ a(k), a^{\dag}(k') ] = [b(k),b^{\dag}(k')] = (2\pi)^3 2\omega \delta^3(k-k')$.  $a$'s, $b$'s commute, essentially.  

\[
\begin{gathered}
\Longrightarrow  \\ 
\begin{gathered}
  \langle f| i \rangle  = \langle 0 | T a_{1'}(+ \infty) \dots a_{m'}(+\infty) b_{m+1'}(+\infty) \dots b_{n'}(+\infty) a_1^{\dag}(-\infty)\dots a_m^{\dag}(-\infty) b_{m+1}^{\dag}(-\infty)\dots b_n^{\dag}(-\infty)| 0 \rangle = \label{Eq:5-1_LSZ_complex_ab} \\ 
   = i^{n'+n} \int d^4x_{1'} e^{-ik_{1'} x_{1'}} (-\partial_{1'}^2 + m_a^2) \dots \int d^4 x_{m'} e^{-ik_{m'} x_{m'} }( - \partial_{m'}^2 + m_a^2) * \\
* \int d^4 x_{m+1'} e^{-ik_{m+1'} x_{m+1'} } (-\partial_{m+1'}^2 + m_b^2) \dots \int d^4 x_{n'} e^{-ik_{n'} x_{n'}} (-\partial_{n'}^2 + m_b^2) * \\ 
     \phantom{=} *\int d^4x_1 e^{ik_1 x_1} (-\partial_1^2 + m_a^2) \dots \int d^4x_m e^{ik_m x_m} (-\partial_m^2 + m_a^2) * \\
\int d^4 x_{m+1} e^{ik_{m+1} x_{m+1} } (-\partial_{m+1}^2 + m_b^2) \dots \int d^4x_n e^{ik_n x_n} (-\partial_n^2 + m_b^2) * \\ 
     \phantom{=} * \langle 0 | T \phi(x_{1'})\dots \phi(x_{m'}) \phi^{\dag}(x_{m+1'}) \dots \phi^{\dag}(x_{n'}) \phi^{\dag}(x_1) \dots \phi^{\dag}(x_m) \phi(x_{m+1}) \dots \phi(x_n) | 0  \rangle
\end{gathered}
\end{gathered}
\]

e.g. A pair of  $a$ and $b$ particles come in; a pair of $a$ and $b$ particles come out.
\[
\begin{gathered}
  \langle f | i \rangle  = i^4 \int d^4 x_{1'} e^{-ik_{1'} x_{1'} } (-\partial_{1'}^2 + m_a^2) \int d^4x_{2'} e^{-ik_{2'} x_{2'} }(-\partial_{2'}^2 + m_b^2) * \\ 
   \phantom{ = * } * \int d^4x_1 e^{ik_1 x_1} (- \partial_1^2 + m_a^2) \int d^4 x_2 e^{ik_2 x_2} (-\partial_2^2 + m_b^2)  \langle 0 | T\phi(x_{1'}) \phi^{\dag}(x_{2'}) \phi^{\dag}(x_1) \phi(x_2)|0\rangle
\end{gathered}
\]

%Notice how it wasn't possible to have a pair of $a$ particles come in and a pair of $b$ particles come out because then you'd have the factor

%\[
%\langle 0 | T\phi^{ \dag}(x_{1'}) \phi^{\dag}(x_{2'}) \phi^{\dag}(x_1) \phi^{\dag}(x_2)|0\rangle
%\]
%and because of the time ordering operator, the creation operators will act on $\langle  0|$ and destroy it.  And of course having the same pair of particles come in and come out is boring - there's no interaction!  

%$a$ and $b$ commute, so treat the type $a$ and type $b$ particles as different.  Looking at Eq. (\ref{Eq:5-1_LSZ_complex_ab}), 
%\[
%\langle 0 | T a_{1'}(+ \infty) \dots a_{m'}(+\infty) b_{m+1'}(+\infty) \dots b_{n'}(+\infty) a_1^{\dag}(-\infty)\dots a_m^{\dag}(-\infty) b_{m+1}^{\dag}(-\infty)\dots b_n^{\dag}(-\infty)| 0 \rangle = \label{Eq:5-1_LSZ_complex_ab} 
%\]

% \hrulefill 

% However, in the light of Problem 10.2, with the need to find an approach that will take into consideration both momentum arrows and ``charge'' arrows, that would distinguish between the 2 sources, $J$, $J^{\dag}$, or, essentially, between $\varphi$ and $\varphi^{\dag}$, so that only 1 kind of arrow is required, then I claim that without loss of generality, the following modified LSZ formula for complex scalar theory is valid.  

% \[
% \begin{gathered}
%  \langle f | i \rangle = i^{n' +n } \int d^4x_{1'} e^{- i k_{1'} x_{1'} } ( - \partial_{1'}^2 + m_a^2 )  \dots  \int d^4x_{m'} e^{-ik_m x_{m'} } (- \partial^2_{m'} + m_a^2 ) \cdot \\ 
%  \cdot \int d^4x_{m+1'} e^{ + i k_{m+1'} x_{m+1'} } ( - \partial^2_{m+1'} + m_b^2 )  \dots  \int d^4x_{n'} e^{ + i k_{n'} x_{n'} } (- \partial_{n'}^2 + m_b^2 ) \cdot \\ 
%  \cdot \int d^4x_1 e^{ik_1 x_1} (- \partial_1^2 + m_a^2) \dots  \int d^4x_m e^{+ i k_m x_m } (- \partial^2_m + m_a^2 ) \cdot \\
%  \cdot \int d^4x_{m+1} e^{-i k_{m+1} x_{m+1} } ( - \partial_{m+1}^2 + m_b^2 ) \dots \int d^4x_n e^{-ik_n x_n } (- \partial_n^2 + m_b^2 ) \cdot \\  
%  \cdot \langle 0 | T\varphi(x_1') \dots \varphi(x_{m'} ) \varphi^{\dag}(x_{m+1}' ) \dots \varphi^{\dag}(x_{n'}) \varphi^{\dag}(x_1) \dots \varphi^{\dag}(x_m) \varphi(x_{m+1}) \dots \varphi(x_n) | 0 \rangle 
% \end{gathered}
% \]

% Note that the time-ordering operator helps already with distinguishing which are incident and which are final particle states.  



\section{Path integrals in quantum mechanics}

\problemhead{6.1} \begin{itemize}
\item[(a)] Find an explicit formula for $\mathcal{D}q$ in eq. (6.9).  Your formula should be of the form $\mathcal{D}q = C \Pi_{j=1}^N dq_j$, where $C$ is a constant that you should compute.  
\item[(b)] For the case of a free particle, $V(Q) =0$, evaluate the path integral of eq. (6.9) explicitly.  Hint: integrate over $q_1$, then $q_2$, etc., and look for a pattern.  Express your final answer in terms of $q',t',q'',t''$, and $m$.  Restore $\hbar$ by dimensional analysis.  
\item[(c)] Compute $\langle q'', t''|q', t' \rangle = \langle q''|e^{-iH(t''-t')} | q'\rangle$ by inserting a complete set of momentum eigenstates, and performing the integral over the momentum.  Compare with your result in part (b).  
\end{itemize}

\solutionhead{6.1} 
\begin{itemize}
\item[(a)] Recall, for $\delta t \to 0$, use Weyl ordering.  
\[
\langle q'', t''| q',t' \rangle = \int \prod^N_{k=1} dq_k \prod^N_{j=0} \frac{dp_j}{2\pi} e^{ip_j(q_{j+1} - q_j) } e^{-i H(p_j, \overline{q}_j) \delta t} \quad \quad \, \text{ where } \begin{aligned} \overline{q}_j & = \frac{1}{2} ( q_j + q_{j+1} ) \\ q_0 & = q' \\ q_{N+1} & = q'' \end{aligned} 
\]

Define $\dot{q}_j \equiv \frac{ q_{j+1} - q_j}{\delta t}$

Now $H(p_j,\overline{q}_j)$ only quadratic in $p_j$, $H(p_j,\overline{q}_j) = \frac{p_j^2}{2m} + V(\overline{q}_j)$.  

Consider
\[
\int \frac{dp_j}{2\pi} e^{ip_j(q_{j+1}-q_j)} e^{-i\delta t \frac{p_j^2}{2m} } = *
\]
using $\int_{-\infty}^{\infty} dx e^{-Ax^2 + Bx} = \sqrt{ \frac{\pi}{A} } \exp{ \left( \frac{B^2}{2A} \right)}$

\[
\begin{gathered}
  * = \left( \frac{1}{2\pi} \right) \left( \frac{ \pi}{ \frac{ i \delta t}{2m } } \right)^{1/2} \exp{ \left( \frac{ - (q_{j+1} - q_j)^2}{ \frac{ 4 i \delta t}{2m} } \right) } = \left( \frac{ m }{ 2\pi i \delta t} \right)^{1/2} \exp{ \left( \frac{ - m (q_{j+1} - q_j)^2 \delta t }{ 2i (\delta t)^2 } \right) }
\end{gathered}
\]
\[
\begin{gathered}
  \Longrightarrow \langle q'', t''|q',t' \rangle = \int \prod^N_{k=1} dq_k \left( \frac{ m}{2\pi i \delta t} \right)^{\frac{N+1}{2} } \prod^N_{j=0} \exp{ \left( \frac{ i m}{2} \left( \frac{ q_{j+1} - q_j}{ \delta t } \right)^2 \delta t - i \delta t V(\overline{q}_j) \right) } = \\
  \begin{aligned} & = \int \prod_{k=1}^N dq_k \left( \frac{ m}{ 2 \pi i \delta t} \right)^{ \frac{N+1}{2} } \exp{ \left( i \lbrace \sum_{j=0}^N \left( \frac{m}{2} \dot{q}_j^2 \right) - V(\overline{q}_j) \rbrace \delta t \right) } = \\ 
    & = \int \left( \frac{ m }{2\pi i \delta t } \right)^{\frac{N+1}{2} } \prod_{k=1}^N dq_k \exp{ \left[ i \int_{t'}^{t''} dt L(\dot{q}(t), q(t)) \right] }
\end{aligned}
\end{gathered}
\]
So $\boxed{ \mathcal{D}q = \lim_{N\to \infty} \left( \frac{m}{2\pi i \delta t} \right)^{\frac{N+1}{2}} \prod_{k=1}^N dq_k }$; $\delta t = \frac{T}{N+1}$.  
\item[(b)] For $V=0$,
\[
\begin{aligned}
\langle q'',t''|q',t' \rangle & = \int \left( \frac{ m }{2\pi i \delta t} \right)^{\frac{N+1}{2} } \prod^N_{k=1} dq_k \exp{ \left( i \sum_{j=0}^N \frac{m}{2} \left( \frac{ q_{j+1} - q_j }{ \delta t} \right)^2 \delta t \right) } = \\
& = \int \sqrt{ \left( \frac{ - \mu }{\pi} \right)^{N+1} } \prod_{k=1}^N dq_k \exp{ \left( \sum_{j=0}^N \mu (q_{j+1} - q_j)^2 \right) }
\end{aligned}
\]
where $\mu = \frac{ i m}{ 2 \delta t}$.  

Now this will be useful: $\boxed{ \int_{-\infty}^{\infty} dy e^{\alpha (x-y)^2 + \beta (z-y)^2 } = \left( \frac{ -\pi}{ \alpha + \beta} \right)^{1/2} \exp{ \left[ \frac{ \alpha \beta}{ \alpha + \beta} (x-z)^2\right] } }$  \\
since $\begin{gathered} \exp{ (\alpha(x-y)^2 + \beta(z-y)^2) } = \exp{ ( \alpha(x^2 - 2xy +y^2) + \beta( z^2  -2yz + y^2) )} = \\ \exp{ ( \alpha x^2 + \beta z^2)} \exp{ ( ( \alpha + \beta)y^2) } \end{gathered}$
\[
\begin{gathered}
  \xrightarrow{\int_{-\infty}^{\infty} dy } \sqrt{  \frac{ -\pi}{ \alpha + \beta } } \exp{ \left[ \frac{ -4 (x\alpha + \beta z)^2}{ 4( \alpha + \beta)} \right] } \exp{ ( \alpha x^2 + \beta z^2) } = \\ 
  \begin{aligned} & = \sqrt{ \frac{ -\pi }{\alpha + \beta}} \exp{ \left[ \frac{ - (x^2 \alpha^2 + \beta^2 z^2 + 2x\alpha \beta z) + \alpha^2 + x^2 + \beta \alpha x^2 + \alpha \beta z^2 + \beta^2 z^2 }{ \alpha + \beta} \right] } = \\
      & = \sqrt{  \frac{ - \pi}{\alpha + \beta} } \exp{ \left[ \frac{ \beta \alpha x^2 - 2x\alpha \beta z + \alpha \beta z^2}{ \alpha + \beta} \right] } = \sqrt{ \frac{ - \pi}{\alpha + \beta} } \exp{ \left[ \frac{ \alpha \beta}{\alpha + \beta} (x-z)^2 \right] } \end{aligned}
\end{gathered}
\]
where my formula, $\int_{-\infty}^{\infty} dx e^{-Ax^2 + Bx} = \sqrt{ \frac{ \pi}{A} } \exp{ \left( \frac{B^2}{4A} \right)}$ was used.  

Consider the following:
\[
\begin{aligned}
  & \int dq_1 \exp{ \left( \mu (q_2 - q_1)^2 + \mu(q_1 -q_0)^2 \right) } = \left( \frac{ - \pi}{2\mu } \right)^{1/2} \exp{ \left[ \frac{\mu}{2} (q_2-q_0)^2 \right] } \\ 
  & \begin{gathered} \int dq_2 \exp{ (\mu (q_3 - q_2)^2 )} \left( \frac{-\pi}{2\mu } \right)^{1/2} \exp{ \left( \frac{ \mu}{2} (q_2-q_0)^2 \right)} = \left( \frac{-\pi}{2\mu} \right)^{1/2} \left( \frac{ -\pi }{ - \frac{ 3\mu}{2} }\right)^{1/2} \exp{ \left[ \frac{ \mu^2/2}{ 3\mu /2} (q_3-q_0)^2 \right]} = \\ 
      = \left( \left( \frac{ -\pi}{\mu} \right)^2 \frac{1}{3} \right)^{1/2} \exp{ \left[ \frac{ \mu}{3} (q_3- q_0)^2 \right] }
\end{gathered}
\end{aligned}
\]

Suppose $\int \prod_{j=1}^{n-1} dq_j \exp{ \left( \sum_{j=0}^{n-1} \mu (q_{n+1} -q_j )^2 \right)}  = \sqrt{ \left( \frac{-\pi}{\mu} \right)^{n-1} \frac{1}{n} } \exp{ \left[ \frac{ \mu}{n} (q_n-q_0)^2  \right]}$.
\[
\begin{gathered}
  \int \prod_{j=1}^n dq_j \exp{ \left( \sum_{j=0}^n \mu (q_{j+1} - q_j)^2 \right)} = \int dq_n e^{\mu (q_{n+1} - q_n)^2} \sqrt{ \left( \frac{-\pi}{\mu} \right)^{n+1} \frac{1}{n} } \exp{ \left[ \frac{\mu}{n} (q_n-q_0)^2 \right] } = \\
  = \sqrt{ \left( \frac{-\pi}{\mu} \right)^{n-1} \frac{1}{n} } \sqrt{ \frac{ -\pi}{ \mu \left( \frac{n+1}{n} \right) } } \exp{ \left[ \frac{ \mu 1/n }{ \frac{n+1}{n} } (q_{n+1} - q_0)^2 \right] } = \sqrt{ \left( \frac{-\pi}{ \mu} \right)^n \frac{1}{n+1} } \exp{ \left[ \frac{ \mu}{ n+1} (q_{n+1} - q_0)^2 \right] }
\end{gathered}
\]
Then by induction,
\[
\begin{gathered}
  \sqrt{ \left( \frac{ -\mu}{\pi} \right)^{N+1} } \int \prod_{k=1}^N dq_k \exp{ \left( \sum_{j=0}^N \mu (q_{j+1} - q_j)^2 \right)} = \sqrt{ \left( \frac{-\mu}{\pi} \right)^{N+1} } \sqrt{ \left( \frac{- \pi}{\mu} \right)^N \frac{1}{N+1} } \exp{ \left[ \frac{ \mu}{N +1} (q_{N+1} - q_0)^2 \right] } = \\ 
  = \sqrt{ \frac{-\pi}{\pi} \left( \frac{1}{N+1} \right)} \exp{ \left[ \frac{ \mu}{N+1} ( q_{N+1} - q_0)^2 \right]} = \sqrt{ \frac{- i m }{ 2\pi \delta t} \left( \frac{1}{N+1} \right) } \exp{ \left[ \frac{ - i m}{2\delta t \pi} \left( \frac{1}{N+1} \right) (q_{N+1} - q_0)^2 \right]} 
\end{gathered}
\]
\[
\boxed{ \langle q'', t''| q',t' \rangle = \sqrt{  \frac{ - i m }{2T} } \exp{ \left[ \frac{-i m }{2T} (q_{N+1} -q_0)^2 \right] } }
\]
where $\mu = \frac{ i m}{2 \delta t}$, $\delta t = \frac{T}{N+1}$.  

$T = t'' -t'$; $q_{N+1} = q''$, $q_0 = q'$.  $\langle q'',t''|q',t' \rangle = \sqrt{ \frac{ - i m}{ 2(t''-t')\pi} } \exp{ \left[ \frac{ - i m}{ 2(t''-t') }(q''-q')^2 \right] }$

By dimensional analysis,
\[
\boxed{ \langle q'',t''|q',t' \rangle = \sqrt{ \frac{ -i m }{ 2\hbar(t''-t')\pi} } \exp{ \left[ \frac{ -im}{2\hbar} \frac{ (q''-q')^2}{ t''-t'} \right] } }
\]
\item[(c)] \[
\begin{gathered}
  \langle q'',t''|q',t' \rangle = \langle q'' | e^{ -i H (t''-t') } |q' \rangle = \langle q'' | \int dp | p\rangle \langle p | e^{-i H(t''-t')} \int dp' |p'\rangle \langle p' | q' \rangle = \\
  = \int dp dp' \frac{ e^{ipq''} }{\sqrt{ 2\pi}} \langle p | e^{ -iH (t''-t')} |p' \rangle \frac{ e^{-ip'q'}}{\sqrt{ 2\pi}} \text{ where } \langle q|p \rangle = \frac{1}{\sqrt{2\pi}} e^{ipq}
\end{gathered}
\]
\[
\Longrightarrow = \int dp dp' \frac{ e^{ipq''} e^{-ip'q'}}{ 2\pi } \langle p | e^{-iH (t''-t')} | p'\rangle = \frac{ \int dp dp'}{2\pi} e^{ipq''}e^{-ip'q'} e^{-i \frac{ (p^{'2} (t''-t') )}{2m } } \delta(p-p')
\]
with $H =\frac{p^2}{2m}$ for the free particle Hamiltonian.  

\[
\Longrightarrow \begin{aligned} & = \frac{1}{2\pi} \int dp e^{ipq''} e^{-ipq'} e^{-i(p^2(t''-t') )/2m} = \frac{1}{2\pi} \int dp e^{ip(q''-q')} e^{-ip^2(t''-t')/2m} = \\
  & = \frac{1}{2\pi} \sqrt{ \frac{ \pi}{ \frac{i}{2m} (t''-t')} } \exp{ \left( \frac{m (q''-q')^2}{ 2i (t''-t')} \right)} = \sqrt{ \frac{- i m}{ 2(t''-t')\pi} } \exp{ \left( \frac{ - im}{2} \frac{ (q''-q')^2}{t''-t'} \right)} \end{aligned}
\]
\end{itemize}

\section{The path integral for the harmonic oscillator}

\problemhead{7.1} Starting with eq. (7.12), do the contour integral to verify eq. (7.14).  

\solutionhead{7.1} Recall
\[
\int_{-\infty}^{\infty} \frac{dE}{2\pi} \frac{e^{-i E(t-t)} }{ -E^2 + \omega^2 - i \epsilon}  \quad \quad \quad \, (7.12)
\]
Now 
\[
\omega^2 - i\epsilon  - E^2 = ( \sqrt{ \omega^2 - i \epsilon } - E)(\sqrt{ \omega^2 - i\epsilon } + E) 
\]

On the semicircle $C$, $E = Re^{i \theta}$.  $-iE(t-t') = -i R(\cos{\theta} + i \sin{\theta})\Delta t = - iRc{\theta} \Delta t + Rs(\theta) \delta t$.  \\

We see for $\Delta t >0$, we sought the lower half semicircle, \\
\phantom{We see for } $\Delta t < 0$, we sought the upper half semicircle.  \\

For $R \to \infty$, $\int_C \to 0$, due to $e^{-R s(\theta) \Delta t}$ factor in integrand.  

\[
\begin{gathered}
	(7.12) = \begin{cases} (i) \frac{ e^{- i ( \sqrt{ \omega^2 - i \epsilon }) \Delta t} }{ 2 \sqrt{ \omega^2  - i\epsilon } } & \text{ for } \Delta t > 0 \\ i \frac{ e^{i \sqrt{ \omega^2 - i \epsilon }  \Delta t} }{ 2 \sqrt{ \omega^2 - i \epsilon } } & \text{ for } \Delta t < 0 \end{cases} \xrightarrow{ \epsilon \to 0} \begin{cases} \frac{ i e^{-i\omega \Delta t} }{ 2\omega }  & \text{ for } \Delta t > 0 \\ \frac{ i e^{i \omega \Delta t} }{2\omega} & \text{ for } \Delta t < 0 \end{cases} \\
= \boxed{ \frac{i}{2\omega} e^{-i \omega |t- t'| } }
\end{gathered}
\]

\problemhead{7.2} Starting with eq. (7.14), verify eq. (7.13).  

\solutionhead{7.2} 20090225 INCOMPLETE 

Recall
\[
G(t-t') = \frac{i}{2\omega} \exp{ (-i \omega |t-t'| ) }  \quad \quad \, (7.14)
\]
If $t\neq t'$, 
\[
\partial_{tt}^2 G(t-t') = \begin{cases} \frac{-i\omega }{2} e^{-i \omega (t-t') }  & \text{ if } t > t' \\ \frac{-i\omega }{2} e^{i \omega (t-t') }  & \text{ if } t < t' \end{cases} = - \omega^2 G
\]





\problemhead{7.3} \begin{itemize}
\item[(a)] Use the Heisenberg equation of motion, $\dot{A} = i [H,A]$, to find explicit expressions for $\dot{Q}$ and $\dot{P}$.  Solve these to get the Heisenberg-picture operators $Q(t)$ and $P(t)$ in terms of the Schr\"{o}dinger-picture operators $Q$ and $P$.  
\item[(b)] Write the Schr\"{o}dinger-picture operators $Q$ and $P$ in terms of the creation and annihilation operators $a$ and $a^{\dag}$, where $H = \hbar \omega (a^{\dag} a + \frac{1}{2} )$.  Then, using your result from part (a), write the Heisenberg-picture operators $Q(t)$ and $P(t)$ in terms of $a$ and $a^{\dag}$.  
\item[(c)] Using your result from part (b), and $a|0\rangle = \langle 0 | a^{\dag} =0$, verify eqs. (7.16) and (7.17).  
\end{itemize}

\solutionhead{7.3} \begin{itemize}
\item[(a)] Recall $H(P,Q) = \frac{1}{2m} P^2 + \frac{1}{2} m\omega^2 Q^2$ (the Hamiltonian for the harmonic oscilator).  

Recall $[Q,P] = i$ (canonical commutation relation).  

\[
\begin{aligned}
  \dot{Q} & = i [H,Q] = i \lbrace \frac{1}{2m} [P^2, Q] + \frac{1}{2} m \omega^2 [Q^2,Q] \rbrace = \frac{i}{2m} P[P,Q] + [P,Q]P = \\ 
  & = \frac{i}{2m} (-i P - i P) = \frac{P}{m} 
\end{aligned}
\]

Likewise, 
\[
\dot{P} = i[H,P] = i \lbrace \frac{1}{2} m \omega^2 ( Q[Q,P] + [Q,P] Q) \rbrace = \frac{i}{2} m \omega^2 (2i Q) = - m \omega^2 Q
\]

So then 
\[
\begin{aligned}
\ddot{Q} & = - \omega^2 Q \\ 
\ddot{P} & = - \omega^2 P 
\end{aligned}
\]

So the form of $Q(t). P(t)$ must be
\[
\begin{aligned}
  Q(t) & = (a_q e^{i\omega t} + b_q e^{-i\omega t} ) \\ 
  P(t) & = (a_p e^{i\omega t} + b_p e^{-i\omega t} ) 
\end{aligned}
\]

Note that $\dot{Q} = i\omega (a_q e^{i\omega t} - b_q e^{-i\omega t})$

I \textbf{assume} the following initial conditions for the Heisenberg representation of $Q,P$:
\[
\begin{aligned}
  Q(0) & = a_q + b_q = Q \\ 
  \dot{Q}(0) & = i \omega (a_q - b_q) = \frac{P}{m} \text{ or } a_q - b_q = \frac{P}{im \omega}
\end{aligned}
\]
since, for the last statement, $Q$ ought to obey the Heisenberg equation of motion.  

Then, it's easy to solve for $a_q, b_q$:
\[
\left[ \begin{matrix} 1 & 1 \\ 1 & -1 \end{matrix} \right] \left[ \begin{matrix} a_q \\ b_q \end{matrix} \right] = \left[ \begin{matrix} Q \\ P/im\omega \end{matrix} \right] \quad \quad \Longrightarrow \frac{-1}{2} \left[ \begin{matrix}  -1 & -1  \\ -1 & 1 \end{matrix} \right] \left[ \begin{matrix} Q  \\ P/im \omega \end{matrix} \right] = \left[ \begin{matrix} \frac{1}{2} \left( Q + \frac{P}{ i m \omega} \right) \\ \frac{1}{2} \left( Q - \frac{P}{im\omega } \right) \end{matrix} \right] = \left[ \begin{matrix} a_q \\ b_q \end{matrix} \right]
\]

So
\[
\boxed{ Q(t) = \frac{1}{2} \left( Q + \frac{P}{im \omega} \right) e^{i\omega t} + \frac{1}{2} \left( Q - \frac{P}{im \omega } \right) e^{-i \omega t} }
\]

Similarly for $P(t)$,

\[
\left[ \begin{matrix} 1 & 1 \\ 1 & -1 \end{matrix} \right] \left[ \begin{matrix} a_p \\ b_p \end{matrix} \right] = \left[ \begin{matrix} P \\ -m\omega^2 Q \end{matrix} \right] \quad \quad \Longrightarrow \frac{-1}{2} \left[ \begin{matrix}  -1 & -1  \\ -1 & 1 \end{matrix} \right] \left[ \begin{matrix} P  \\ -\frac{m\omega Q}{i} \end{matrix} \right] = \left[ \begin{matrix} \frac{1}{2} \left( P - \frac{m \omega Q}{ i } \right) \\ \frac{1}{2} \left( P + \frac{m \omega Q}{i } \right) \end{matrix} \right] = \left[ \begin{matrix} a_p \\ b_p \end{matrix} \right]
\]

\[
\boxed{ P(t) = \frac{1}{2} \left( P - \frac{ m \omega Q}{i} \right) e^{i \omega t} + \frac{1}{2} \left( P + \frac{ m \omega Q}{i} \right) e^{-i \omega t} }
\]
\item[(b)] Recall
\[
\begin{aligned} a & = \sqrt{ \frac{ m \omega_0 }{2} } \left( x + \frac{ i p }{m\omega_0} \right) \\ 
  a^{\dag} & = \sqrt{ \frac{ m \omega_0}{2}} \left( x - \frac{ i p }{m \omega_0} \right) \end{aligned} \quad \quad \quad \, \begin{aligned} Q & = \frac{ a + a^{\dag}}{ \sqrt{ 2m\omega_0} } \\ P & = \frac{ \sqrt{ m \omega_0}}{ i \sqrt{2}} (a- a^{\dag}) \end{aligned}
\]

\[
\begin{aligned}
  Q(t) & = \frac{1}{2} \left( \frac{ a + a^{\dag}}{ \sqrt{ 2 m \omega_0 } } + \frac{ - (a-a^{\dag} )}{ \sqrt{ 2 m \omega_0} } \right) e^{i \omega t} + \frac{1}{2} \left( \frac{ a + a^{\dag}}{ \sqrt{ 2 m \omega_0} } + \frac{ a- a^{ \dag}}{ \sqrt{ 2 m \omega_0} } \right) e^{-i \omega t} = \\
  & = \left( \frac{ 1}{ \sqrt{ 2 m \omega_0 } } \right) a^{\dag} e^{i \omega t} + \frac{1}{ \sqrt{ 2 m \omega_0} } ae^{- i \omega t} = \left( \frac{ 1}{ \sqrt{ 2 m \omega_0} } \right) (a^{\dag} e^{i \omega t} + ae^{-i\omega t} )
\end{aligned}
\]

\[
\begin{aligned}
  P(t) & = \frac{1}{2} \left( \frac{ \sqrt{ m \omega_0}}{ i \sqrt{2}} (a- a^{\dag}) - \frac{ \sqrt{ m \omega_0}}{ i \sqrt{2}} (a+ a^{\dag}) \right) e^{i \omega t} + \frac{1}{2} \left( \frac{1}{i} \sqrt{ \frac{m \omega_0}{2} } (a-a^{\dag} ) + \frac{1}{i} \sqrt{ \frac{ m\omega_0 }{2} } (a+a^{\dag}) \right) e^{-i \omega t} = \\
  & = \frac{1}{i} \sqrt{ \frac{ m \omega_0}{2} } \left( -a^{\dag} e^{i \omega t} +  ae^{-i\omega t} \right)
\end{aligned}
\]

\item[(c)] \[
\begin{aligned}
  \langle 0 | TQ(t_1) Q(t_2) |0 \rangle & = \frac{1}{2m \omega_0} \langle 0 | (a^{\dag} e^{i \omega t_1} + ae^{-i \omega t_1} )(a^{\dag} e^{i \omega t_2} + ae^{-i \omega t_2} ) | 0 \rangle = \\
  & = \frac{1}{2 m \omega_0} e^{- i \omega (t_1 - t_2) }  = \frac{1}{i} G(t_2- t_1)
\end{aligned}
\]
since $G(t-t') = \frac{i}{2\omega} \exp{ ( - i \omega | t-t'| ) }$  (7.14), and 
\[
G(t_2 - t_1) = \frac{i}{2 \omega} \exp{ ( -i \omega |t_2 -t_1| ) } = \frac{i}{2\omega} \exp{ (-i \omega (t_1-t_2) )}
\]

\[
\begin{gathered}
  \langle 0 | TQ(t_1) Q(t_2) Q(t_3) Q(t_4) | 0 \rangle = \langle 0 | Q(t_1) Q(t_2) \frac{1}{ 2 m \omega_0} (a^{\dag} e^{ i \omega t_3} + ae^{- i \omega t_3 } ) e^{ i \omega t_4} |1 \rangle = \\
  = \frac{ \langle 0 | Q(t_1) Q(t_2) }{ (2m \omega_0 )} (e^{ i \omega (t_3 + t_4)} | 2\rangle + e^{- i \omega (t_3 - t_4)} | 0 \rangle = \frac{ \langle 0 | Q(t_1) }{ (2m \omega_0)^{3/2} } (a^{\dag} e^{ i \omega t_2} + a e^{- i \omega t_2} )  ( e^{i \omega (t_3+t_4)} |2\rangle + e^{-i\omega(t_3 - t_4)} |0\rangle ) = \\ 
    =\frac{ \langle 0 | Q(t_1)}{ (2m \omega_0)^{3/2} } (e^{ i \omega (t_3+t_4 -t_2)} | 1 \rangle + e^{-i \omega (t_3-t_4-t_2) } | 1\rangle ) = \frac{1}{ (2m \omega_0)^2} (e^{i \omega (t_3 + t_4 - t_2 -t_1)} + e^{- i \omega (t_3 - t_4 -t_2 + t_1) } )
\end{gathered}
\]

\end{itemize}

\section{The path integral for free-field theory}

introduce 4-dim. Fourier transformations
\begin{align}
  & \widetilde{\varphi}(k) = \int d^4x e^{-ikx} \varphi(x) \\ 
  & \varphi(x) = \int \frac{ d^4k}{ (2\pi)^4} e^{ikx} \widetilde{\varphi}(k) \quad \quad \quad (8.6)
\end{align}

\problemhead{8.1} Starting with eq. (8.11), verify eq. (8.12).

\solutionhead{8.1} Recall $\begin{aligned} \partial^{\mu} & = (-\partial_t, \vec{\nabla} ) \\ \partial_{\mu} & = ( \partial_t, \vec{\nabla} ) \end{aligned}$ \\ 
$\begin{aligned} \partial^{\mu} e^{ikx} & = (ik^0, i\mathbf{k} ) e^{ikx} \\ \partial_{\mu} e^{ikx} & = (-ik^0, i \mathbf{k}) e^{ikx} \end{aligned}$
\[
\Longrightarrow \partial_x^2 e^{ikx} = \partial^{\mu} \partial_{\mu} e^{ikx} = (k^{0 \, 2} - \mathbf{k}^2 ) e^{ikx} = -k^2 e^{ikx} 
\] 
for $k^2 = \mathbf{k}^2 - (k^0)^2$.  

\[
(-\partial_x^2 + m^2) \nabla(x-x') = \int \frac{d^4 k}{ (2\pi)^4} \frac{ (k^2 + m^2) e^{ik(x-x') } }{ k^2 + m^2 - i\epsilon }  = \int \frac{ d^4k}{ (2\pi)^4 } \frac{ e^{ik(x-x') } }{ 1 - \frac{ i \epsilon }{ k^2 + m^2 } } \xrightarrow{ \epsilon \to 0 } \delta^4(x-x')
\]
since $\int \frac{d^4k }{(2\pi)^4 } e^{ik(x-x') }$ is one Dirac delta function, representation.  

\problemhead{8.2} Starting with eq. (8.11), verify eq. (8.13).  

\solutionhead{8.2}  
\[
\Delta(x-x') = \int \frac{d^4 k}{ (2\pi)^4} \frac{e^{ik(x-x')} }{k^2 + m^2 - i \epsilon } \quad \quad \quad \, (8.11)
\]
Now
\[
\begin{aligned}
  & k^2 = \mathbf{k}^2 - (k^0)^2 \\ 
  & k^2 + m^2 - i \epsilon = \overline{k}^2 + m^2 - i \epsilon -(k^0)^2 = (\sqrt{ \mathbf{k}^2 + m^2 - i\epsilon } - k^0)( \sqrt{ \mathbf{k}^2 + m^2 - i\epsilon } + k^0 ) \\ 
  & ik(x-x') = i( -k^0 (t-t') + \mathbf{k}\cdot (\mathbf{x} - \mathbf{x}')) = -i k^0 \tau + i \mathbf{k} \cdot (\mathbf{x} - \mathbf{x}')
\end{aligned}
\]

So from $e^{ik(x-x')}$, consider $e^{-ik^0 \tau}$.   \\

On semicircle $C$, let $k^0 = Re^{i\theta}$, 
\[
\exp{ (-iRe^{i\theta} \tau) } = \exp{ (-iR(\cos{\theta} + i \sin{\theta}) \tau)} = \exp{ (-iRc{(\theta)} )}\exp{(R\sin{\theta} \tau)}
\]

If $\tau >0$, then for $\lim_{R\to \infty}\int \to 0$, $C$ must be in the lower half plane.  \\
If $\tau <0$, then for $\lim_{R\to \infty}\int \to 0$, $C$ must be in the upper half plane.   \\

By residue thm.,
\[
\int \frac{dk^0}{2\pi} \frac{e^{-i k^0 \tau} }{ k^2 + m^2 - i \epsilon } = \begin{cases} \frac{ i e^{-i \sqrt{ \mathbf{k}^2 + m^2 - i\epsilon}  \tau} }{ 2 \sqrt{ \mathbf{k}^2  + m^2 - i\epsilon }} & \text{ if } \tau > 0 \\ \frac{ i e^{i \sqrt{ \mathbf{k}^2 + m^2 - i\epsilon}  \tau} }{ 2 \sqrt{ \mathbf{k}^2  + m^2 - i\epsilon }} & \text{ if } \tau < 0 \end{cases} \xrightarrow{ \epsilon \to 0} \frac{ i e^{-i \sqrt{ \mathbf{k}^2 + m^2} |t-t'| }}{ 2 \sqrt{ \mathbf{k}^2 + m^2 } }
\]
Now $\omega^2 - \mathbf{k}^2 = m^2$ with $\omega = \omega(\mathbf{k})$  (by special relativity).  

\[
\begin{aligned}
  \Longrightarrow \Delta (x-x') & = i \int \frac{d^3\mathbf{k}}{(2\pi)^3} \frac{e^{i \mathbf{k}\cdot (\mathbf{x} - \mathbf{x}')}  e^{-i\omega |t-t'| } }{ 2\omega } = i \int \widetilde{dk} e^{i\mathbf{k} \cdot (\mathbf{x}-\mathbf{x}') - i \omega |t-t' | } = \\
& = i \theta(t-t') \int \widetilde{dk} e^{ik(x-x')} + i\theta(t'-t) \int \widetilde{dk} e^{-ik(x-x') }
\end{aligned}
\]

where for the 2nd. term of the very last expression, consider $\tau <0$.  Then $\Delta = i \int \frac{d^3k}{(2\pi)^3} \left( \frac{1}{2\omega} \right) e^{i\mathbf{k}\cdot(\mathbf{x} - \mathbf{x}' )+ i \omega (t-t')}$.  \\

By symmetry about $(\mathbf{x} - \mathbf{x}')$, $\mathbf{k} \to -\mathbf{k}'$ leaves $\int d^3\mathbf{k}$ unchanged, +1 overall factor.  


\problemhead{8.3} Starting with eq. (8.13), verify eq. (8.12).  Note that the time derivative in the Klein-Gordon wave operator can act on either the field (which obeys the Klein-Gordon equation) or the time-ordering step functions.

\solutionhead{8.3} Recall
\[
\Delta(x-x') = i \theta(t-t') \int \widetilde{dk} e^{ik(x-x')} + i \theta(t'-t) \int \widetilde{dk} e^{-ik(x-x')} \quad \quad \quad \, (8.13)
\]
Now
\[
\nabla^2 \Delta(x-x') = -\mathbf{k}^2 \Delta
\]
Recall
\[
\partial_x^2 = \nabla^2 - \partial_{tt}^2 \text{ and } k(x-x') = \mathbf{k}\cdot (\mathbf{x} - \mathbf{x}') - \omega_{\mathbf{k}}(t-t') \quad \quad \, (\text{Srednicki})
\]
\[
\begin{aligned}
  \partial_t \Delta & = i \delta(t-t') \int \widetilde{dk} e^{ik(x-x') } + i\theta(t-t') \int \widetilde{dk} (-i\omega) e^{ik(x-x') } + \\ 
  & + -i \delta(t'-t) \int \widetilde{dk} e^{ik(x-x')} + i\theta(t'-t) \int \widetilde{dk} e^{ik(x-x')}
\end{aligned}
\]

\[
\begin{aligned}
  \partial_{tt}^2 \Delta & = i\delta'(t-t') \int \widetilde{Dk} e^{ik(x-x')} + 2i \delta(t-t') \int \widetilde{dk}(-i\omega) e^{ik(x-x')} + i \theta(t-t') \int \widetilde{dk} (-\omega^2) e^{ik(x-x')} + \\ 
  & + - i \delta'(t'-t) \int \widetilde{dk} e^{ik(x-x')} - 2i \delta(t'-t) \int \widetilde{dk} (i\omega) e^{-ik(x-x')} + i\theta(t'-t) \int \widetilde{dk} e^{-ik(x-x')}
\end{aligned}
\]

Note that $\frac{d\theta(t'-t)}{dt} = - \frac{d\theta(t-t')}{dt} = - \delta(t-t') = - \delta(t'-t)$, \\
and by my notation, $\delta'(t'-t) = \frac{d}{dt} \delta(t'-t)$ in all cases for this particular problem.  

Observe a $-\omega^2 \Delta$ term.  \\
$\partial_x^2 \Delta$ yields the terms $(-\mathbf{k}^2 + \omega^2) \Delta = m^2 \Delta$.  So $-m^2 \Delta$ will cancel these terms.

$\partial_{tt}^2 \Delta$ has other terms.

Note that $\widetilde{dk} \equiv \frac{d^3k }{ (2\pi)^3 2 \omega}$.  
\[
2 \omega \delta(t-t') \int \frac{d^3k}{ (2\pi)^3 2\omega } e^{ik(x-x')} = \delta(t-t') \int \frac{d^3k}{(2\pi)^3} e^{i\mathbf{k} \cdot (\mathbf{x}- \mathbf{x}') - i \omega (t-t') }
\]

\[
2 \omega \delta(t'-t) \int \widetilde{dk} e^{ik(x-x')} = 2 \omega \delta(t'-t) \int \frac{d^3k}{(2\pi)^3 2\omega} e^{ik(x-x')} = \delta(t'-t) \int \frac{d^3k}{(2\pi)^3} e^{-ik(x-x')} = \delta(t'-t) \int \frac{d^3k}{(2\pi)^3} e^{-i\mathbf{k}\cdot (\mathbf{x}-\mathbf{x}') + i \omega(t-t')}
\]
\[
i \delta'(t-t') \int \widetilde{dk} e^{ik(x-x')} = i \delta(t-t') \int \frac{d^3\mathbf{k}}{(2\pi)^3 2\omega} e^{i\mathbf{k}\cdot (\mathbf{x}-\mathbf{x}') - i \omega (t-t')}
\]

\[
\Longrightarrow \quad \quad \, \begin{aligned}
  (-\partial_x^2 + m^2) \Delta & = \delta(t-t') \int \frac{d^3k}{(2\pi)^3} e^{ik(x-x')} + \delta(t'-t) \int \frac{d^3k}{(2\pi)^3} e^{-ik(x-x')} + \\ & + i \delta'(t-t') \int \widetilde{dk} e^{ik(x-x')} - i \delta'(t'-t) \int \widetilde{dk} e^{-ik(x-x')} = *
\end{aligned}
\]

If $t=t'$, or $x=x'$, then $* \to \infty $ due to $\delta(t-t')$, $\delta(t'-t)$. \\
If $x\neq x'$, then $* =0$ due to $\delta(t-t')$, $\delta(t'-t)$.  Surely $\delta'(t-t') = \delta'(t'-t) =0$ for $t\neq t'$ since slope is zero for $\delta$ for $t\neq t'$.  

Now 
\[
\int d^4x f(x) \delta(x-x') = f(x')
\]

$\int d^4 x f(x) * = $ ?

\[
\begin{gathered}
  \int dt f(x) \delta(t-t') \int \frac{d^3k}{(2\pi)^3} e^{ik(x-x')} = \int \frac{ d^3k}{(2\pi)^3} f(\mathbf{x}, t') e^{i \mathbf{k} \cdot (\mathbf{x}-\mathbf{x}')} = f(\mathbf{x},t') \delta^3(\mathbf{x}-\mathbf{x}') \\ 
  \xrightarrow{ \int d^3 x} f(x')
\end{gathered}
\]

Likewise
\[
\int d^4x f(x) \delta(t'-t) \int \frac{d^3k}{(2\pi)^3} e^{-ik(x-x')} = f(x')
\]

\[
\begin{gathered}
  \int dt f(x) \left( \frac{d}{dt} \delta(t-t')\right) \frac{ e^{-i\omega (t-t')} }{2\omega} = 0 - \int \delta(t-t') \lbrace \frac{df(x)}{dt} \frac{ e^{-i \omega (t-t')} }{2\omega} + f(x) \left( \frac{-i}{2} \right) e^{-i \omega(t-t')} \rbrace = \\
  = \left. - \left( \frac{df(x)}{dt} \right) \right|_{t=t'} \frac{1}{2\omega} + \frac{i}{2} f(\mathbf{x},t)
\end{gathered}
\]
$\int \frac{d^3k}{(2\pi)^3} e^{i \mathbf{k}\cdot (\mathbf{x} - \mathbf{x}') } = \delta^3(\mathbf{x}-\mathbf{x}')$ so 
\[
\xrightarrow{ + i \int d^3 x \delta^3(\mathbf{x} - \mathbf{x}') } \left. \frac{-i }{2\omega} \left( \frac{df(x)}{dt} \right) \right|_{x=x'} - \frac{f(x')}{2}
\]

\[
\begin{gathered}
  \int dt f(x) \left( \frac{d}{dt} \delta(t'-t) \right) \frac{e^{i \omega (t-t')} }{2\omega} = 0 - \int \delta(t'-t) \lbrace \frac{df(x)}{dt}  \frac{e^{i\omega(t-t')}}{2\omega} + \frac{ f(x) (i\omega) }{2\omega} e^{i \omega (t-t')} \rbrace = \\ 
  = - \left. \frac{df(x)}{dt} \right|_{t=t'} \frac{1}{2\omega} - \frac{f(\mathbf{x},t') i}{2}
\end{gathered}
\]
$\int \frac{d^3k}{(2\pi)^3} e^{-i\mathbf{k}\cdot (\mathbf{x}-\mathbf{x}') } = \delta^3(\mathbf{x} - \mathbf{x}')$ so
\[
\xrightarrow{ -i \int d^3x \delta^3(\mathbf{x}-\mathbf{x}')} + \left. \frac{df(x)}{dt} \right|_{x=x'} \frac{i}{2\omega} - \frac{f(x')}{2}
\]

\[
\Longrightarrow \int d^3x f(x) (-\partial_x^2 + m^2) \Delta(x-x') = f(x') + f(x') + 0 - \frac{f(x')}{2} - \frac{f(x')}{2} = f(x') = \int d^4x f(x) \delta(x-x')
\]

So since $* = (-\partial^2_x+m^2)\Delta(x-x') $ obeys the properties of a Dirac delta function $\delta^4(x-x')$, it is equal to the Dirac delta function. 



\problemhead{8.4} Use eqs. (3.19), (3.29), and (5.3) (and its hermitian conjugate) to verify the last line of eq. (8.15).  

\solutionhead{8.4} Recall that 
\[
\varphi(x) = \int \widetilde{dk} [a(\mathbf{k}) e^{ikx} + a^*(\mathbf{k}) e^{-ikx} ]  \quad \quad \quad (3.19)
\]
\[
\begin{aligned}
  \left[ a(\mathbf{k}), a(\mathbf{k}') \right] & = 0 \\ 
  [a^{\dag}(\mathbf{k}), a^{\dag}(\mathbf{k}') ] & = 0 \\ 
  [a(\mathbf{k}), a^{\dag}(\mathbf{k}') ] & = (2\pi)^3 2\omega \delta^3(\mathbf{k} - \mathbf{k}')  \quad \quad (3.29)
\end{aligned} \quad \quad \quad  a(\mathbf{k})|0 \rangle = 0 \quad \quad (5.3)
\]

Recall 
\[
Z_0(J) = \exp{ \left[ \frac{i}{2} \int \frac{d^4k}{(2\pi)^4} \frac{\tilde{J}(k) \tilde{J}(-k)}{ k^2 + m^2 - i\epsilon } \right] } = \exp{ \left[ \frac{i}{2} \int d^4x d^3x' J(x) \Delta(x-x')J(x') \right] } \quad \quad \, (8.10)
\]
\[
\Delta(x-x') = \int \frac{d^4k}{(2\pi)^4} \frac{e^{ik(x-x')} }{k^2 + m^2 -i\epsilon} \quad \quad \, (8.11)
\]

Then recall,
\[
\begin{aligned}
  \langle 0 | T\varphi(x_1) \varphi(x_2) |0\rangle & = \frac{1}{i} \frac{\delta}{\delta J(x_1)}  \frac{1}{i} \frac{\delta}{\delta J(x_2)}  \left. Z_0(J) \right|_{J=0}  = \frac{1}{i} \frac{\delta}{\delta J(x_1)} \left[ \int d^4x' \Delta(x_2-x')J(x') \right] \left. Z_0(J) \right|_{J=0} = \\ 
  & = \frac{1}{i} \left[ \Delta(x_2 - x_1) + \left[ \int d^3x' \Delta(x_2-x') J(x') \right]\frac{\delta}{\delta J(x_1) } \right] \left. Z_0(J) \right|_{J=0} \xrightarrow{J=0} \frac{1}{i} \Delta(x_2 - x_1)
\end{aligned}
\]

Now what we want is to show the same result but with the mode expansion of the field.  \\

\hrulefill

Consider $\lbrace 0 | T\varphi(x_1) | 0 \rangle = \left. \frac{1}{i} \frac{ \delta}{ \delta J(x_1) } Z_0(J) \right|_{J=0}$.  
\[
\langle 0 | T\varphi(x_1) | 0 \rangle = \langle 0 | \int \widetilde{dk} [ a(\mathbf{k}) e^{ikx_1} + a^*(\mathbf{k}) e^{-ikx_1} ] |0 \rangle = 0 
\]
since $a(\mathbf{k}) | 0 \rangle$; $\langle 0 | a^{\dag}(\mathbf{k}) = 0$.  

Then we should get the correct result for the quadratic term:

\[
\begin{gathered}
  \langle 0 | T\varphi(x_2) \varphi(x_1) | 0 \rangle  = \langle 0 | \int \widetilde{dk}' \widetilde{dk} ( a(k')e^{ik'x_2} + a^{\dag}(k') e^{-ik'x_2} ) (a(k) e^{ikx_1} + a^{\dag}(k) e^{-ikx_1} ) |0 \rangle = \\
  \begin{gathered} = \int \widetilde{dk}' \widetilde{dk} \langle 0  |  e^{i(k'x_2+ -kx_1)} a(k') a(k) + e^{-i(k'x_2 - kx_1) } a^{\dag}_{k'} a_k  + \\ 
 + e^{i(k'x_2- kx_1) } a_{k'} a_k^{\dag} + e^{-i(k' x_2+ kx_1) } a_{k'}^{\dag} a_k^{\dag} | 0 \rangle 
\end{gathered}
\end{gathered}
\]
Since $a_{k'} a_k^{\dag} = a_k^{\dag} a_{k'} + (2\pi)^3 2\omega \delta^3(k-k')$.  
\[
\Longrightarrow \int \widetilde{dk} e^{ik(x_2-x_1)}
\]

By Time ordering, we require
\[
\begin{aligned} \int \widetilde{dk} e^{ik(x_2-x_1)} & \text{ if } x_2^0 - x_1^0 > 0 \\ 
  \int \widetilde{dk} e^{ik(x_1-x_2)} & \text{ if } x_1^0 - x_2^0 > 0  \end{aligned}  
\]

So
\[
  \langle 0 | T\varphi(x_2) \varphi(x_1) | 0 \rangle = \theta(x_2^0 - x_1^0) \int \widetilde{dk} e^{ik(x_2-x_1)} + \theta(x_1^0 - x_2^0) \int \widetilde{dk} e^{ik(x_1-x_2)} = \frac{1}{i} \Delta(x_2-x_1)
\]

So we get back the same result from the mode expansion.  

\problemhead{8.5} The retarded and advanced Green's functions for the Klein-Gordon wave operator satisfy $\Delta_{ret}(x-y) =0$ for $x^0 \geq y^0$ and $\Delta_{adv}(x-y) =0$ for $x^0 \leq y^0$.  Find the pole prescriptions on the right-hand side of eq. (8.11) that yield these Green's functions.  

\solutionhead{8.5} Recall
\[
\Delta(x-x') = \int \frac{d^4 k}{(2\pi)^4} \frac{e^{ik(x-x')} }{ k^2 + m^2 -i \epsilon } \quad \quad \quad \, (8.11)
\]
The rationale for adding $-i\epsilon$ was to avoid the poles.  Then in doing the contour integration in complex $k^0$, we can take the integration across the real axis.  \\

Now $k^2 = \mathbf{k}^2 - (k^0)^2$ (Srednicki's convention)
\[
k^2 + m^2 - i \epsilon = \mathbf{k}^2 + m^2 - i \epsilon - (k^0)^2 = (\sqrt{ \omega_{\mathbf{k}}^2 - i\epsilon } - k^0)( \sqrt{ \omega_{\mathbf{k}}^2 - i\epsilon } + k^0 ) = (w - k^0)(w+k^0)
\]
$w$ in 4th quadrant for $\epsilon \ll \omega_k^2$.  $-w$ in 2nd. quadrant.  

\[
k (x-y ) = \mathbf{k}\cdot (\mathbf{x} - \mathbf{y}) - k^0(x^0 - y^0)
\]
Consider $e^{-ik^0 \tau}$; $\tau = x^0 - y^0$.  

If $\tau >0$, for $k^0 = Re^{i\theta} = R\cos{\theta} + i R\sin{\theta}$, $e^{-iR \cos{\theta} \tau} e^{R\sin{\theta} \tau}$.  Then close the integration in lower half-plane (i.e. $\sin{\theta} <0$).  
\[
\int_{\Gamma} \frac{dk^0 e^{-ik^0 \tau} }{ (w- k^0)(w+k^0) } = 2\pi i \left( \frac{ e^{-i w \tau}}{2 w} \right) = \frac{ \pi i }{w} e^{-i w \tau } = \int_{-\infty}^{\infty} dk^0 \frac{e^{-ik^0 \tau}}{ w^2 - (k^{02}) } + \int_{C_R}
\]

\[
\left| \int_{C_R} \frac{e^{-i k^0 \tau}}{ w^2 - k^{02} } dk^0 \right| \leq \int_{C_R} \left| \frac{ e^{R\sin{\theta} \tau} }{ R^2 - w^2 } \right| R d\theta \to 0 \text{ for } \sin{\theta} < 0 
\]

\[
\int_{-\infty}^{\infty} dk^0 \frac{ e^{-ik^0 \tau}}{ w^2 - k^{02} } = \frac{ \pi i e^{-i \sqrt{ \omega_k^2 - i \epsilon } (x_0 -y_0) } }{ \sqrt{ \omega_k^2 - i \epsilon } } \xrightarrow{ \epsilon \to 0} \frac{ \pi i e^{-i \sqrt{ \mathbf{k}^2 + m^2 } (x_0 - y_0) } }{ \sqrt{ \mathbf{k}^2 + m^2 } } \text{ for } x_0 > y_0 
\]

If $\tau <0$, then close the integration in upper half plane.  Then residue $k^0 = - w$ enclosed.  

\[
\int_{-\infty}^{\infty} \frac{ e^{-ik^0 \tau} }{ w^2 - k^{02}} = 2\pi i \left( \frac{ e^{iw \tau}}{ -2w} \right) = \frac{ - \pi i e^{ i \sqrt{ \mathbf{k}^2 + m^2  }(x_0 -y_0 )} }{ \sqrt{ \mathbf{k}^2 + m^2  }}
\]

So finally, \\
$\Delta_{ret}(x-y) = 0$ for $x^0 \geq y^0$, 
\[
\Delta_{ret}(x-y) = \theta(y^0 - x^0) \frac{ - \pi i e^{ i \sqrt{ \mathbf{k}^2 + m^2  }(x_0 -y_0 )} }{ \sqrt{ \mathbf{k}^2 + m^2  }} \text{ for integration closing the upper half plane }
\]

$\Delta_{adv}(x-y) = 0$ for $x^0 \leq y^0$, 
\[
\Delta_{adv}(x-y) = \theta(x^0 - y^0) \frac{  \pi i e^{ -i \sqrt{ \mathbf{k}^2 + m^2  }(x_0 -y_0 )} }{ \sqrt{ \mathbf{k}^2 + m^2  }} \text{ for integration closing the lower half plane }
\]


\hrulefill

Consider not appending $-i\epsilon$.  Then there are poles at $w = \sqrt{ \mathbf{k}^2 + m^2 }$, $-w$.  

There are a number of issues to consider:
\begin{itemize}
\item We must consider semicircle contour integrals about the poles of infinitesimally small radius.  
\item What branch are we integrating on, especially on those semicircle integrals around each of the poles?
\end{itemize}

If $\tau > 0$, for $C_{\delta_1}$, $C_{\delta_2}$ semicircle integrals around each of the poles enclosing the poles from the upper half, 
\[
\begin{gathered}
  \int_{\Gamma} = \int_{C_R} + \int_{-\infty}^{-\delta_1 - w} + \int_{C_{\delta_1}} + \int_{-w + \delta}^{w- \delta} + \int_{w+ \delta}^{\infty} = (2\pi i) \lbrace \frac{ e^{-iw\tau }  }{2w} + \frac{e^{i w \tau}}{ 2 w} \rbrace \\ 
  \int_{C_{\delta_1}} i d\theta \frac{ \delta e^{i \theta} e^{-i (-w + \delta e^{ i \theta} )\tau }}{ \delta e^{i \theta} (2w - \delta e^{i\theta} )} = \int_{\pi}^0 i d\theta \frac{ e^{i w \tau} e^{-i \delta e^{i\theta}  \tau } }{ 2w - \delta e^{i \theta} } \xrightarrow{ \delta \to 0 } \frac{ i e^{i w\tau} (- \pi )}{ 2w } = - \pi i \frac{ e^{i w \tau }}{2w } 
\end{gathered}
\]
for $k^0 = -w + \delta e^{i\theta}$.  

\[
\int_{C_{\delta_2 } } \frac{ i d\theta \delta e^{i \theta} e^{-i (w + \delta e^{i \theta}  ) \tau } }{ -\delta e^{i \theta} (2 w + \delta e^{i \theta}  )} \xrightarrow{ \delta \to 0} \frac{ (-i\theta) e^{-iw\tau }}{ -2w } = \frac{ i \pi e^{-iw\tau} }{2w }
\]
for $k^0 = w + \delta e^{i \theta}$.  
\[
\Longrightarrow \quad \, \begin{aligned} \int_{-\infty}^{\infty} & = (\pi i) \lbrace \frac{ e^{-iw \tau }}{w} \left( 1 - \frac{1}{2} \right) + \frac{ e^{iw \tau}}{ w} \left( 1 - \frac{1}{2} \right) \rbrace = \pi i \lbrace \frac{ e^{- iw\tau } }{2w} + \frac{ e^{i w \tau }}{ 2w } \rbrace = \\
  & = \frac{ \pi i }{w} \cos{(w \tau )}
\end{aligned}
\]
As you can see, there are issues with which branch to choose that semicircles $C_{\delta_1}$, $C_{\delta_2}$ lie on.  

If $\tau > 0$, for $C_{\delta_1}$, $C_{\delta_2}$ semicircle integrals around each of the poles \emph{not} enclosing the poles from the lower half, 

\[
\begin{aligned}
  & \int_{-\infty}^{\infty} \frac{ ie^{i \omega \tau} \pi}{ 2 \omega } + \frac{ - i \pi e^{-i \omega \tau} }{ 2 \omega} = 0 \\
  & \int_{-\infty}^{\infty} = i \pi \left( \frac{ e^{-i \omega \tau} }{2\omega } - \frac{ e^{i\omega \tau}}{2\omega } \right) = \frac{ \pi}{ \omega } \sin{(\omega \tau )}
\end{aligned}
\]

If $\tau < 0$, for $C_{\delta_1}$, $C_{\delta_2}$ semicircle integrals around each of the poles enclosing the poles from below the poles, semicircle in the upper half,
\[
\begin{aligned}
  & \int_{-\infty}^{\infty} + \frac{ \pi i e^{i \omega \tau} }{ 2 \omega } + \frac{ -i \pi e^{-i \omega \tau }}{ 2\omega } = (2\pi i ) \lbrace \frac{ e^{-i \omega \tau } }{2\omega} + \frac{ e^{ i \omega \tau}}{2\omega } \rbrace \\ 
  &  \int_{-\infty}^{\infty} = \frac{ \pi i }{\omega} \lbrace e^{-i \omega \tau} \left( 1 + \frac{1}{2} \right) + e^{i \omega \tau} \left( 1 - \frac{1}{2} \right) \rbrace = \frac{ \pi i }{\omega} ( e^{-i\omega \tau} \left( \frac{3}{2} \right) + e^{i \omega \tau} \frac{1}{2} )
\end{aligned}
\]

If $\tau < 0$, for $C_{\delta_1}$, $C_{\delta_2}$ semicircle integrals around each of the poles \emph{not} enclosing the poles from above the poles, semicircle in the upper half,
\[
\begin{aligned}
  \int_{-\infty}^{\infty} & = - \lbrace -\pi i \frac{ e^{i\omega \tau}}{2 \omega} + \frac{ i \pi e^{-i \omega \tau}}{2\omega} \rbrace = \pi i \lbrace \frac{ e^{i \omega \tau}}{2\omega} - \frac{ e^{-i\omega \tau}}{2\omega} \rbrace = \\
  & = \frac{-\pi}{\omega} \sin{(\omega \tau)}
\end{aligned}
\]


\solutionhead{8.6} Supposing $Z_0(J) = \exp{ i W_0(J)}$, recall, after some algebraic manipulations,
\[
Z_0(J) = \exp{ \left( \frac{i}{2} \int d^4x d^4x' J(x) \Delta(x-x') J(x') \right) }
\]
\[
\Longrightarrow W_0(J) = \frac{1}{2} \int d^4x d^4x' J(x) \Delta(x-x') J(x')
\]
Suppose $J(x)$, a classical source, is a real function.  

Now $\int \widetilde{dk} e^{i k(x-x')} = \int \widetilde{dk} ( \cos{ (k(x-x'))} + i \sin{ (k(x-x') )} )$.  

Likewise, 
\[
\begin{aligned}
  \Delta(x-x') & = - \theta(t-t') \int \widetilde{dk} \sin{ (k(x-x') )} + \theta(t'-t) \int \widetilde{dk} \sin{( k(x-x'))} + \\ 
  & \phantom{ = - } + i \theta(t-t') \int \widetilde{dk} \cos{(k(x-x'))} + i \theta(t'-t) \int \widetilde{dk} \cos{ (k(x-x'))} = \\
  & = ( - \theta(t-t') + \theta(t'-t)) \int \widetilde{dk} \sin{ (k(x-x'))} + i \int \widetilde{dk} \cos{(k(x-x'))}
\end{aligned}
\]
\[
\boxed{ \begin{aligned}
  & \Im{ (W_0(J))} = \frac{1}{2} \int d^4x d^4x' J(x) \int \widetilde{dk} \cos{ (k(x-x'))} J(x') \\
  & \Re{ (W_0(J))} = \frac{1}{2} \int d^4x d^4x' J(x) ( - \theta(t-t') + \theta(t'-t) ) \int \widetilde{dk} \sin{(k(x-x') )} J(x')
\end{aligned} }
\]




\problemhead{8.7} Repeat the analysis of this section for the complex scalar field that was introduced in problem 3.5, and further studied in problem 5.1.  Write your source term in the form $J^{\dag} \varphi + J \varphi^{\dag}$, and find an explicit formula, analogous to eq. (8.10), for $Z_0(J^{\dag}, J)$.  Write down the appropriate generalization of eq. (8.14), and use it to compute $\langle 0 | T\varphi(x_1) \varphi(x_2) | 0 \rangle$, $\langle 0 | T\varphi^{\dag}(x_1) \varphi(x_2) | 0 \rangle$, and $\langle 0 | T\varphi^{\dag}(x_1) \varphi^{\dag}(x_2) | 0 \rangle$.  Then verify your results by using the method of problem 8.4.  Finally, give the appropriate generalization of eq. (8.17).  

\solutionhead{8.7} 
\[
\mathcal{L} = - \partial^{\mu} \phi^{\dag} \partial_{\mu} \phi - m^2 \phi^{\dag} \phi + \Omega_0 + J^{\dag} \phi + J \phi^{\dag}
\]
\[
S = i \int d^4 x \mathcal{L} = i \int \frac{ d^4k}{(2\pi)^4} \left[ - \widetilde{\phi}^{\dag}(k)(k^2 + m^2) \widetilde{\phi}(k) + \Omega_0 \widetilde{\phi}^{\dag}(k) \widetilde{\phi}(k) + \widetilde{J}^{\dag}(k) \phi(k) + \widetilde{J}(k) \phi^{\dag}(k) \right]
\]
where $J(x) = \int \frac{ d^4k}{(2\pi)^4} e^{ikx} \widetilde{J}(k)$, \\
so that
\[
J(x) \phi^{\dag}(x) = \int \frac{d^4k}{(2\pi)^8} e^{ikx} \widetilde{J}(k) d^4k' e^{-ik'x} \widetilde{\phi}^{\dag}(k') \xrightarrow{ \int d^4 x} \int \frac{d^4k}{(2\pi)^4} \widetilde{J}(k) \widetilde{\phi}^{\dag}(k)
\]

For the other terms,

\[
\begin{aligned}
  \phi(x) & = \int \frac{ d^4k}{ (2\pi)^4} e^{ikx} \widetilde{\phi}(k) \\ 
  \phi^{\dag}(x) & = \int \frac{d^4k}{(2\pi)^4} e^{-ikx} \widetilde{\phi}^{\dag}(k) 
\end{aligned} \quad \quad \quad \, \begin{aligned}
  \widetilde{\phi}(k) & = \int d^4x e^{-ikx} \phi(x) \\ 
  \widetilde{\phi}^{\dag}(k) & = \int d^4x e^{ikx} \phi^{\dag}(x)
\end{aligned}
\]

\[
e^{-ikx} = e^{-i(\mathbf{k}\cdot \mathbf{x} - k_0 t) } = e^{-i (\mathbf{k}\cdot \mathbf{x} ) + i k^0 t}
\]
\[
\begin{gathered}
  \partial^{\mu} \phi^{\dag} \partial_{\mu} \phi = \int \frac{d^4k}{(2\pi)^4} (\partial^{\mu} e^{-ikx} ) \widetilde{\phi}^{\dag}(k) \int \frac{d^4 k'}{(2\pi)^4} (\partial_{\mu} e^{ik'x} ) \widetilde{\phi}(k') = \\ 
  = \frac{1}{(2\pi)^8} \int d^4k d^4k' (-i \omega, -i \mathbf{k}) (-i \omega', i \mathbf{k}) \widetilde{\phi}^{\dag}(k) \widetilde{\phi}(k') e^{-i(k-k') x} = \frac{1}{(2\pi)^8} \int d^4k d^4k' (-\omega \omega' + \mathbf{k}\cdot \mathbf{k}' ) \widetilde{\phi}^{\dag}(k) \widetilde{\phi}(k') e^{-i (k-k')} = \\
  = \frac{1}{(2\pi)^4} \int d^4k (-\omega^2 + \mathbf{k}^2) \widetilde{\phi}^{\dag}(k) \widetilde{\phi}(k) = \frac{1}{(2\pi)^4} \int d^4k (k^2 \widetilde{\phi}^{\dag}(k)\widetilde{\phi}(k) )
\end{gathered}
\]

Drop tilde notation now.  

Make this substitute for the integration variable:
\[
\begin{aligned}
  \chi(k) & = \phi(k) - \frac{ J(k) }{k^2 + m^2} \\ 
  \chi^{\dag} & = \phi^{\dag} - \frac{J^{\dag}}{k^2 +m^2}
\end{aligned}
\]

So
\[
\begin{gathered}
   - \left( \chi^{\dag} + \frac{J^{\dag} }{ k^2 + m^2} \right)(k^2 + m^2) \left( \chi + \frac{J}{k^2 +m^2} \right) + \Omega_0 \left( \chi^{\dag} + \frac{J^{\dag}}{ k^2 + m^2} \right) \left( \chi + \frac{J}{k^2 + m^2} \right) + \\
   + J^{\dag} \left( \chi + \frac{J}{k^2 + m^2} \right) + J \left( \chi^{\dag} + \frac{J^{\dag}}{ k^2 + m^2} \right) = \\
 =  - \left( \chi^{\dag} \chi + \frac{ \chi^{\dag} J + J^{\dag} \chi}{k^2 + m^2 } + \frac{J^{\dag} J}{(k^2 +m^2)} \right)(k^2 + m^2) + \Omega_0 \left( \chi^{\dag} \chi + \frac{ \chi^{\dag} J + J^{\dag} \chi}{k^2 + m^2} + \frac{J^{\dag} J}{(k^2 + m^2)^2} \right) + (J^{\dag} \chi + J \chi^{\dag} ) + \frac{J^{\dag} J + JJ^{\dag}}{ k^2 + m^2} = \\ 
  =  (-(k^2 + m^2) + \Omega_0) \chi^{\dag} \chi + \Omega_0 (\chi^{\dag} J + J^{\dag} \chi) + \frac{ \Omega_0}{(k^2 + m^2)^2 } J^{\dag} J + \frac{J J^{\dag}}{ k^2 + m^2 }
\end{gathered}
\]

Now 
$\mathcal{D} \phi \propto \prod_x d\phi(x)$ \\
$\phi$ simply shifted by a constant: $\mathcal{D}\phi = \mathcal{D} \chi$.  

Suppose the integration over $\chi$ simply yields $Z_0(0) = \langle 0 | 0 \rangle_{J,J^{\dag} =0} = 1$, and neglect terms containing $\Omega_0$ (let $\Omega_0 =0$)

\[
\boxed{ Z_0(J^{\dag},J) = \exp{ \left[ i \int \frac{d^4k}{(2\pi)^4} \frac{ \widetilde{J}(k) \widetilde{J}^{\dag}(+k) }{ k^2 + m^2 - i \epsilon} \right] }  = \exp{ \left[ i \int d^4x d^4x' J(x) \Delta(x-x') J^{\dag}(x') \right]} }
\]
where $\Delta(x-x') = \int \frac{d^4k}{(2\pi)^4}  \frac{ e^{-ik(x-x')}}{k^2  +m^2 - i \epsilon }$.  

Now
\[
\boxed{ \langle 0 | T \phi^{\dag}(x_1) \dots \phi(x_m) \dots | 0 \rangle = \left. \frac{1}{i} \delta_{J(x_1)} \dots \delta_{J^{\dag}(x_m)} Z_0(J^{\dag},J)  \right|_{J,J^{\dag}=0} =\left. \frac{1}{i} \frac{ \delta}{ \delta J(x_1)} \dots \frac{ \delta}{ \delta J^{\dag}(x_m)} Z_0(J^{\dag},J) \right|_{J,J^{\dag} =0} }
\]

\[
\begin{aligned}
\langle 0 | T\phi(x_1)\phi(x_2) | 0 \rangle & = \frac{1}{i} \frac{ \delta }{ \delta J^{\dag}(x_1)} \left[ \int d^4x J(x) \Delta(x-x_2) \right] \left. Z_0(J^{\dag},J) \right|_{J^{\dag},J =0} = \\
& = \left( \int d^4x J(x) \Delta(x-x_2) \right) (\int d^4x' J(x') \Delta(x'-x_1) ) \left. Z_0(J^{\dag},J) \right|_{J^{\dag},J =0} = 0 
\end{aligned}
\]

Likewise,
\[
\begin{aligned}
  & \langle 0 | T\phi^{\dag}(x_1) \phi^{\dag}(x_2) | 0 \rangle = 0 \\ 
  & \langle 0 | T\phi^{\dag}(x_1) \phi(x_2) | 0 \rangle = \frac{1}{i} \frac{ \delta}{\delta J(x_1)} \left[ \int d^4x J(x) \Delta(x-x_2) \right] \left. Z_0(J^{\dag},J) \right|_{J^{\dag},J =0} = \frac{1}{i} \Delta(x_1-x_2)
\end{aligned}
\]

Now $\phi(x) = \int \widetilde{dk} [ a(k) e^{ikx} + b^{\dag}(k) e^{-ikx} ]$.  

Immediately,
\[
\langle 0 | T \phi(x_1)\phi(x_2) | 0 \rangle = \langle 0 | T\phi^{\dag}(x_1) \phi^{\dag}(x_2) | 0 \rangle = 0 \text{ since }
\]
\[
a(k_2)| 0 \rangle = \langle 0 | b^{\dag}(k_1) b^{\dag}(k_2) | 0 \rangle = \langle 0 | a(k_1) b^{\dag}(k_2)|0\rangle = b(k_2)|0\rangle = \langle 0 |b(k_1)a^{\dag}(k_2)|0 \rangle = \langle 0 | a^{\dag}(k_1)a^{\dag}(k_2) | 0 \rangle = 0
\]

\[
\begin{gathered}
  \langle 0 | T\phi^{\dag}(x_1)\phi(x_2)| 0 \rangle = \langle 0 | \int \widetilde{dk}_1 \widetilde{dk}_2 (a^{\dag}(k_1) e^{-ik_1 x_1} + b(k_1) e^{ik_1 x_1})(a(k_2) e^{ik_2 x_2} + b^{\dag}(k_2) e^{-ik_2 x_2} ) |0\rangle = \\
  = \int \widetilde{dk}_1 \widetilde{dk}_2 e^{i(k_1 x_1 - k_2 x_2)} \langle 0 | b(k_1) b^{\dag}(k_2)|0 \rangle
\end{gathered}
\]

Recall that $i (k_1 x_1 - k_2 x_2) = i (\mathbf{k}_1 \cdot \mathbf{x}_1 - \mathbf{k}_2 \cdot \mathbf{x}_2 -(k_1^0 t_1 - k_2^0 t_2))$ and
\[
\begin{aligned}
  & [b(k), b^{\dag}(k')] = (2\pi)^3 2\omega \delta^3(k-k') \text{ so that } \\
  & \langle 0 | b(k_1) b^{\dag}(k_2) | 0 \rangle = (2\pi)^3 2\omega \delta^3(\mathbf{k}_1 - \mathbf{k}_2)
\end{aligned}
\]

\[
\Longrightarrow \int \widetilde{dk}_1 e^{i (\mathbf{k}_1 \cdot \mathbf{x}_1 - \mathbf{k}_1 \cdot \mathbf{x}_2 ) - i (k_1^0 t_1 - k_2^0 t_2) } = \int \widetilde{dk} e^{i( \mathbf{k}\cdot (\mathbf{x}_1 - \mathbf{x}_2) ) - i (k_1^0 t_1 - k_2^0 t_2) }
\]

Now
\[
\begin{aligned}
  \Delta(x-x') & = \int \frac{d^3k}{(2\pi)^3} e^{-i\mathbf{k}\cdot (\mathbf{x} - \mathbf{x}')} \int \frac{ dk^0}{(2\pi)} \frac{e^{i(k^0(t-t'))} }{ (\mathbf{k}^2 +m^2) - k^{02} } = \\
  & = \int \frac{d^3k}{(2\pi)^3} e^{-i\mathbf{k} \cdot(\mathbf{x} - \mathbf{x}')} \lbrace \frac{ e^{i \sqrt{ \mathbf{k}^2 + m^2} |t-t'| } }{2\omega} \rbrace = i \int \widetilde{dk} e^{-i \mathbf{k}\cdot (\mathbf{x} - \mathbf{x}')} e^{i \omega|t-t'| }
\end{aligned}
\]
where $\omega^2 - \mathbf{k}^2 = m^2$.  Note, from the form of the Lagrangian, $m_a = m_b$ (mass of particle of type $a$, is same as mass of particle of type $b$).  

\section{The path integral for interacting field theory}




\solutionhead{9.2}
\begin{itemize}
\item[(a)] 
\[
\begin{aligned}
  \mathcal{L} & = \mathcal{L}_0 + \mathcal{L}_1 \\
  \mathcal{L}_0 & = - \frac{1}{2} \partial^{\mu}\phi \partial_{\mu} \phi - \frac{1}{2} m^2 \phi^2 \\
 \mathcal{L}_1 & = - \frac{1}{24} Z_{\lambda} \lambda \phi^4 + \mathcal{L}_{ct} = -\frac{1}{4!} Z_{\lambda} \lambda \phi^4 + \mathcal{L}_{ct}
\end{aligned}
\]

Recall for the path integral, 
\[
\begin{aligned}
  Z(J) & \propto e^{ i \int d^4x \mathcal{L}_1\left( \frac{1}{i} \frac{ \delta }{ \delta J(x) } \right)} Z_0(J) \\
  & = e^{ i \int d^4x \frac{-1}{24} Z_{\lambda}\lambda \left( \frac{1}{i} \delta_{J(x)} \right)^4} \sum_{P=0}^{\infty} \frac{1}{P!} \left[ \frac{i}{2} \int d^4y d^4z J(y) \Delta(y-z) J(z) \right]^P
\end{aligned}
\]
For $V=1$, 1 vertex,
\[
\int d^4x \frac{-i}{24} Z_{\lambda} \lambda \left( \frac{1}{i} \delta_{J(x)} \right)^4 \frac{1}{P!} \left[ \frac{i}{2} \int d^4y d^4z J(y) \Delta(y-z) J(z) \right]^P
\]
Consider only diagrams with sources for all legs.  

Then $P=4$.  

We count the number of terms that result in a particular diagram.  

Rearrange $\left( \frac{1}{i} \delta_{J(x)} \right)^4$ by $4!$ ways.  

Rearrange $V!$ vertices.  $V=1$, $1! = 1$.  

Rearrange 2 sources at end of each propagator, for each $\delta_{J(x)}$.  $2!$, $2!^4$.  

Rearrange the propagator themselves.  $P! = 4!$.  

We can rearrange the $4$ propagators in $4!$ ways, all these arrangements duplicated by exchanging the derivatives at the vertices.  $S=4!$ symmetry factor.  $\frac{1}{S!} = \frac{1}{4!}$.  

However, for some given set of coordinates, $y_1$, $y_2$, $y_3$, $y_4$, $4!$ ways to choose which propagator is exactly $y_1$, $y_2$, $y_3$, or $y_4$.  $4! \cdot \frac{1}{4!} = 1$.  

\[
\frac{1}{24} \cdot 4! \cdot 1! \cdot \frac{1}{4!} \left( \frac{1}{2} \right)^4 \cdot (2!)^4 \cdot 4! \cdot 4! \cdot \frac{1}{4!} =1
\]

% Factor of $4!$ comes from how $\delta_{J(x)}^4$ could act on sources.  

% Note that there are 2 possible sources for $\delta_{J(x)}$ to act on, and so factor of $\frac{1}{2}$ from sources of $Z_0$ are cancelled.  

% Look at the diagram.  Note the symmetry factor of $4!$.  

\setlength{\unitlength}{0.5cm}

\begin{picture}(4,4)
\put(2,2){\line(1,1){2}}
\put(4,4){\circle*{0.5}}
%\put(3,3){\line(1,1){1}}
\put(2,2){\line(1,-1){2}}
\put(4,0){\circle*{0.5}}
%\put(3,1){\line(1,-1){1}}
\put(0,4){\line(1,-1){2}}
\put(0,4){\circle*{0.5}}
\put(0,0){\line(1,1){2}}
\put(0,0){\circle*{0.5}}
%\put(2,2){\circle*{0.5}} 
\end{picture}

\hrulefill 

Doing this problem more carefully, \\
\phantom{Doing} start from the given Lagrangian.

\[
\begin{aligned}
  \mathcal{L}_0 & = - \frac{1}{2} \partial^{\mu} \varphi \partial_{\mu} \varphi - \frac{1}{2} m^2 \varphi^2 \\ 
  \mathcal{L}_1 & = - \frac{1}{24} Z_{\lambda} \lambda \varphi^4 + \mathcal{L}_{ct} \\ 
  \mathcal{L}_{ct} & = -\frac{1}{2} ( Z_{\varphi} - 1 ) \partial^{\mu} \varphi \partial_{\mu} \varphi - \frac{1}{2} ( Z_m - 1 ) m^2 \varphi^2 = \frac{-1}{2} A \partial^{\mu} \varphi \partial_{\mu} \varphi - \frac{1}{2} B m^2 \varphi^2
\end{aligned}
\]

Assuming $Z_{\lambda} = 1 + \mathcal{O}(\lambda^2)$,

\[
\begin{aligned}
  Z_1 & = e^{ i \int d^4x \mathcal{L}_1\left( \frac{1}{i} \delta_{J_x} \right) } Z_0 \simeq i \int d^4x \frac{ -Z_{\lambda} \lambda }{ 4!} \left( \frac{1}{i } \delta_{J_x} \right)^4 \frac{1}{4!} \left( \frac{ i }{2} \int d^4y d^4z J_y \Delta_{y-z} J_z \right)^4 = \\ 
  & = \frac{ -i Z_{\lambda} \lambda }{ 4!} \int d^4x \left( \int d^4y J_y \Delta_{y-x} \right)^4
\end{aligned}
\]
with $4! \cdot 2^4$ counting factor matching $\delta_{J_x}$'s to propagators.  

For $\varphi \varphi \to \varphi \varphi$ scattering, use $\frac{1}{i} \delta_{J_{x_1}} \frac{1}{i} \delta_{J_{x_2}} \frac{1}{i} \delta_{J_{x_3}} \frac{1}{i} \delta_{J_{x_4}}$, 

\[
\langle 0 | T \varphi(x_1) \varphi(x_2) \varphi(x_3) \varphi(x_4) | 0 \rangle \simeq \frac{ - i Z_{\lambda} \lambda }{4!} 4! \int d^4x \Delta_{x_1 -x} \Delta_{x_2 -x} \Delta_{x_3 -x} \Delta_{x_4 -x} \frac{1}{i^4} = \frac{-i}{i^4 } Z_{\lambda} \lambda \int d^4x \Delta_{x_1 -x} \Delta_{x_2 -x} \Delta_{x_3 -x } \Delta_{x_4 -x}
\]

\[
\begin{gathered}
  \langle f | i \rangle = i^4 \int d^4x_1 e^{ i k_1 x_1} ( - \partial_1^2 + m^2 ) \int d^4x_2 e^{ i k_2 x_2} ( - \partial_2^2 + m^2 ) \int d^4x_3 e^{-i k_3 x_3} ( -\partial_3^2 + m^2 ) \int d^4x_4 e^{-ik_4 x_4} ( - \partial_4^2 + m^2 ) \cdot \\
  \cdot \langle 0 | T \varphi_1 \varphi_2 \varphi_3 \varphi_4 | 0 \rangle  = - i Z_{\lambda} \lambda (2\pi)^4 \delta^4(k_1 + k_2 - k_3 - k_4) 
\end{gathered}
\]

For $A,B$ counterterms, consider $P=2$ propagators.  
\[
\begin{aligned}
  Z(J) &  \simeq \frac{-i}{2} \int d^4x \left( \frac{1}{i} \delta_{J_x} \right) ( - A \partial_x^2 +  Bm^2 ) \left( \frac{1}{i} \delta_{J_x} \right) \cdot Z_0  \simeq \frac{i}{2} \int d^4x \delta_{J_x} ( - A \partial_x^2 + B m^2 ) \delta_{J_x} \frac{1}{2!} \left( \frac{i}{2} \int d^4y d^4z J_y \Delta_{y-z} J_z \right)^2 = \\
  & = \frac{-1}{2} \int d^4x \delta_{J_x} ( - A \partial_x^2 + Bm^2 ) \int d^4y J_y \Delta_{y-x} \left( \frac{i}{2} \int d^4z d^4w J_z \Delta_{z-w} J_w \right) = \\ 
  & = - \frac{i}{2} \int d^4x \int d^4z J_z \Delta_{z-x} ( - A \partial_x^2 + Bm^2) \int d^4y J_y \Delta_{y-x}
\end{aligned}
\]
Remember, only consider connected diagrams.  

Also, note that $ \partial_x^2 \Delta_{y-x} = \partial_x^2 \int \frac{d^4k}{(2\pi)^4} \frac{ e^{ ik (y-x)} }{ k^2 + m^2 - i \epsilon } = \int \frac{ d^4k}{(2\pi)^4} \frac{ ( -k^2) e^{ik(y-x)}}{ k^2 + m^2 - i \epsilon }$, so order is important.  

\[
\langle 0 | T \varphi(x_1) \varphi(x_2) | 0 \rangle = \frac{-i}{2} \int d^4x \Delta_{x_2 - x} ( - A \partial_x^2 + Bm^2) \Delta_{x_1- x} \frac{1}{i^2} \quad \quad \, \text{(for this particular case, this particular order)}
\]

\begin{gather}
  \langle f | i \rangle = i^2 \int d^4x e^{ik_1 x_1} ( - \partial_1^2 + m^2) \int d^4x_2 e^{-ik_2x_2} ( - \partial_2^2 + m^2) \frac{- i }{2}\int d^4x \Delta_{x_2 - x} (  - A \partial_x^2 + Bm^2 ) \Delta_{x_1 - x} \frac{1}{i^2} = \\ \notag 
  = \frac{-i }{2} \int d^4x e^{-ik_2 x} \int d^4x_1 e^{ik_1 x_1 } ( - A \partial_x^2 + Bm^2) \delta^4(x_1 - x) + \dots \label{Eq:phi4_counter}
\end{gather}


To deal with $\partial_x^2$, use integration by parts: \\
$ \int d^4x e^{-ik_2x} \partial_x^2 \delta^4(x_1 - x) = -k_2^2 \int d^4x e^{-ik_2x} \delta^4(x_1-x) = -k_2^2 e^{-ik_2x_1}$

Then the right-hand side of Eq. (\ref{Eq:phi4_counter}) is
\[
\frac{-i}{2} ( 2\pi)^4 \delta^4(k_1-k_2) Ak_2^2 + \frac{-i}{2} Bm^2(2\pi)^4 \delta^4(k_1-k_2) = \frac{-i}{2} (Ak_2^2 + Bm^2) (2\pi)^4 \delta^4(k_1-k_2)
\]

Because of $\delta^4(k_1-k_2)$, then the results are the same if labels 1,2, switched.  Thus, a $ \times 2$ counting factor.  

\[
\langle f | i \rangle = -i ( Ak_1^2 + Bm^2) (2\pi)^4 \delta^4(k_1 - k_2) + \dots
\]

\setlength{\unitlength}{0.5cm}

\begin{picture}(8,4)\label{Fi:phi4_counter}
\put(0,2.75){\vector(1,0){2}}
\put(2,2.75){\line(1,0){2}}
\put(4,2.75){\vector(1,0){2}}
\put(6,2.75){\line(1,0){2}}
\put(3.70,2.50){X}
\put(1.80,2.95){$k_1$}
\put(2.20,1.65){$-i (Ak_1^2 + Bm^2)$}
\end{picture}



%%---------------------------

The associated vertex factor is 
\[
\boxed{ -i Z_{\lambda} \lambda}
\]

We also need to sprinkle counterterm vertices for each propagator with a factor of $-i ( Ak_1^2 + Bm^2)$.  
\item[(b)]
\item[(c)] A clever way to see this (Arthur Lipstein, Nov. 17, 2008), that tadpoles are canceled, is to note that the symmetry $\varphi \to - \varphi$.  Then $\langle 0 | \varphi | 0 \rangle =0$.  

\hrulefill

Or, note that $E = 2 P - 4V$.  For $\langle 0 | \varphi(x) | 0 \rangle \neq 0$, $E=1$, but $1 = 2P - 4V$ $\Longrightarrow 4V  + 1 = 2P$ which is not possible, for $V,P \in \mathbb{Z}^+$.  So no counterterm linear in $\varphi$ is needed.  

\end{itemize}


\problemhead{9.3} Consider a complex scalar field (see problems 3.5, 5.1, and 8.7) with $\mathcal{L} = \mathcal{L}_0 + \mathcal{L}_1$, where
\[
\begin{aligned}
  \mathcal{L}_0 & = - \partial^{\mu} \varphi^{\dag} \partial_{\mu} \varphi - m^2 \varphi^{\dag} \varphi, \\
  \mathcal{L}_1 & = - \frac{1}{4} Z_{\lambda}\lambda (\varphi^{\dag} \varphi)^2 + \mathcal{L}_{ct} \\ 
  \mathcal{L}_{ct} & = -(Z_{\varphi} - 1)\partial^{\mu} \varphi^{\dag} \partial_{\mu} \varphi - (Z_m -1) m^2 \varphi^{\dag} \varphi 
\end{aligned}
\]
This theory has two kinds of sources, $J$ and $J^{\dag}$, and so we need a way to tell which is which when we draw the diagrams.  Rather than labeling the source blobs with a $J$ or $J^{\dag}$, we will indicate which is which by putting an arrow on the attached propagator that points \emph{towards} the source if it is a $J^{\dag}$, and \emph{away} from the source if it is a $J$.  
\begin{itemize}
\item[(a)] What kind of vertex appears in the diagrams for this theory, and what is the associated vertex factor?  Hint: your answer should involve those arrows!
\item[(b)] Ignoring the counterterms, draw all the connected diagrams with $1 \leq E \leq 4$ and $0\leq V \leq 2$, and find their symmetry factors.  Hint: the arrows are important!   
\end{itemize}

\solutionhead{9.3} 
\begin{itemize}
\item[(a)]Recall
\[
\begin{aligned}
  Z_0(J,J^{\dag}) & = \langle 0 | 0 \rangle_{J,J^{\dag}} = \int \mathcal{D}\phi \mathcal{D}\phi^{\dag} e^{i \int d^4x [ \mathcal{L}_0 + J \phi^{\dag} + J^{\dag} \phi ]} = \\
  & = \exp{ [ i \int d^4x d^4x' J(x) \Delta(x-x') J^{\dag}(x')]}
\end{aligned}
\]

Consider 
\[
\exp{ \left[ i \int d^4x \mathcal{L}_1 \right]} = \exp{ \left[ i \int d^4x \left( \frac{-1}{4} Z_{\lambda} \lambda \right) (\phi^{\dag} \phi)^2 \right]}
\]
Note that
\[
\phi^{\dag}(x) \phi(x) \phi^{\dag}(x) \phi(x) \rightarrow \frac{1}{i} \frac{\delta}{\delta J(x) } \frac{1}{i}\frac{ \delta}{ \delta J^{\dag}(x)} \frac{1}{i}\frac{ \delta}{ \delta J(x)} \frac{1}{i}\frac{ \delta}{ \delta J^{\dag}(x)} \xrightarrow{ \text{ notation}} \delta_{J_x} \delta_{J_x^{\dag}} \delta_{J_x} \delta_{J^{\dag}_x}
\]

\hrulefill

Consider, as an exercise, that's \emph{not directly} related to the problem at hand,
\[
\begin{gathered}
  \delta_{J_x^{\dag}} Z_0 = [ i \int d^4x' J_{x'} \Delta(x'-x)]Z_0 \xrightarrow{ \delta_{J_x}} i \Delta(x-x)Z_0 + [i \int d^4x' \Delta(x-x') J_{x'}^{\dag}] Z_0 \\ 
  \xrightarrow{ \delta_{J_x^{\dag}}} [ i \int d^4x' J_{x'} \Delta(x'-x) ]i \Delta(0) Z_0 + i \Delta(x-x) Z_0 + [i\int d^4x' \Delta(x-x') J_{x'}^{\dag}] [i \int d^4x' J_{x'} \Delta(x'-x)]Z_0 \\
  \xrightarrow{ \delta_{J_x}, J,J^{\dag} =0} -(\Delta (0))(1)
\end{gathered}
\]

In a sense, I say I can set to as an overall multiplicative factor $1$, diagrams that's closed as a loop.  

\hrulefill

Recall doing the double Taylor series expansion:
\[
\begin{gathered}
  \exp{ \left[ \frac{-i}{4} Z_{\lambda} \lambda \int d^4 x\left( \frac{1}{i} \frac{ \delta}{ \delta J(x)} \frac{1}{i} \frac{\delta}{\delta J^{\dag}(x)} \right)^2 \right] } Z_0(J) \\
  Z_1(J^{\dag},J) = \sum_{V=0}^{\infty} \frac{1}{V!} \left[ \frac{-i}{4} Z_{\lambda} \lambda \int d^4x \left( \frac{1}{i} \frac{\delta}{\delta J(x)} \frac{1}{i} \frac{ \delta }{ \delta J^{\dag}(x) } \right)^2 \right]^V  \sum_{P=0}^{\infty} \frac{1}{P!} \left[ i \int d^4y d^4z J(y) \Delta(y-z) J^{\dag}(z) \right]^P
\end{gathered}
\]

Consider 1 vertex.  $\frac{-i}{4} Z_{\lambda} \lambda \int d^4x \left( \frac{1}{i} \frac{ \delta}{\delta J(x)} \frac{1}{ i} \frac{ \delta}{ \delta J^{\dag}(x)} \right)^2$  

Consider $\frac{1}{P!} \left[ i \int d^4y d^4 z J(y) \Delta(y-z) J^{\dag}(z) \right]^P$.  For $P=1$, applying the vertex gives zero.  

For $P=2$, 
\[
\left( \frac{-i}{4} Z_{\lambda} \lambda \right) \left( \frac{1}{-1} \right) \left( \frac{1}{2} \right)(-1) (\Delta(0))^2 = \frac{-i}{8} Z_{\lambda} \lambda (\Delta(0))^2
\] 
is obtained.  Result: closed double loop, figure eight.  

For $P=3$, 
\[
\frac{1}{3!} [ i \int d^4 y_1 d^4z_1 J_{y_1} \Delta_{y_1 - z_1} J^{\dag}_{z_1} ] [ i \int d^4 y_2 d^4z_2 J_{y_2} \Delta_{y_2 - z_2} J^{\dag}_{z_2} ] [ i \int d^4 y_3 d^4z_3 J_{y_3} \Delta_{y_3 - z_3} J^{\dag}_{z_3} ]
\]
Consider only connected diagrams.
\[
\xrightarrow{\delta_{J_x^{\dag}}} \delta_{J_x} \delta_{J_x^{\dag}} \int d^4y_1 J_{y_1} \Delta_{y_1 -x} \int d^4z_2  \Delta_{x -z_2} J^{\dag}_{z_2} \int d^4y_3 J_{y_3} \Delta_{y_3 - x} \xrightarrow{ \delta_{J_x}} \Delta(0) \int d^4z_2 J^{\dag}_{z_2} \int d^4y_3 J_{y_3} \Delta_{x-z_2} \Delta_{y_3 - x}
\]

Result: closed loop and 2 propagator legs.  

For $P=4$, 
\[
\Longrightarrow  \int \left( \frac{-i}{4} Z_{\lambda} \lambda \right) d^4x \left(  \int d^4z_1 \Delta_{x-z_1} J^{\dag}_{x_1} \right) \left(  \int d^4 y_2 J_{y_2} \Delta_{y_2 - x} \right) \left(  \int d^4 z_3 \Delta_{x-z_3} J^{\dag}_{z_3} \right)\left(  \int d^4 y_4 J_{y_4} \Delta_{y_4 - x} \right)
\]
so we get the vertex, with 4 propagators attached in the following manner: \\

\setlength{\unitlength}{0.5cm}

\begin{picture}(5,5)
\put(2,2){\vector(1,1){1.3}}
\put(3,3){\line(1,1){1}}
\put(4,0){\vector(-1,1){1.1}}
\put(3,1){\line(-1,1){1}}
\put(0,4){\vector(1,-1){1}}
\put(1,3){\line(1,-1){1}}
\put(2,2){\vector(-1,-1){1.3}}
\put(1,1){\line(-1,-1){1}}
\put(2,2){\circle*{0.3}}
\put(0,0){\circle{0.7}}
\put(0,4){\circle*{0.7}}
\put(4,0){\circle*{0.7}}
\put(4,4){\circle{0.7}} 
\put(-1.0,-0.44){$J^{\dag}$}
\put(-1.0,4.44){$J$}
\put(4.6,-0.44){$J$}
\put(4.6,4.44){$J^{\dag}$}
\put(1,3){$k_1$} 
\put(0.25,1.1){$k_2$}
\put(3.1,1.1){$k_4$}
\put(3.4,3.0){$k_3$}
\end{picture}

\quad \\ 

Note that $E = 2P- 4V$ by induction.  

Note first the $4!$ factor, canceling $\frac{1}{4!}$ from the 4 propagators, due to matching functional derivatives to propagators.  Then another $2 \times 2$ factor when using the LSZ formula (see Problem 10.2) on the remaining propagators (``amputating the external legs'').  

The associated vertex factor is $\boxed{ -i Z_{\lambda} \lambda }$.  

\item[(b)] See diagram.  
\end{itemize}

\solutionhead{9.4} 
\begin{itemize}
\item[(a)] Given $\exp{ W(g,J)} = \frac{1}{\sqrt{2 \pi }} \int_{-\infty}^{+\infty} dx \exp{ \left[ \frac{-1}{2} x^2 + \frac{1}{6} gx^3 + Jx \right] }$  

The \emph{trick} is to do the following (Arthur Lipstein, Nov. 17, 2008).

\[
\boxed{
\exp{ W(g,J)} = \frac{1}{\sqrt{2\pi}} \exp{ \left( \frac{1}{6} g \left( \delta_J\right)^3 \right) } \int_{-\infty}^{\infty} dx e^{ -\frac{1}{2} x^2 + Jx } }
\] 

\small

Scratchwork:
\[
\int_{-\infty}^{\infty} dx e^{ - \frac{1}{2} x^2 + Jx } = \int_{-\infty}^{\infty} dx e^{ - \frac{1}{2} (x-J)^2 } e^{J^2/2} = \sqrt{ 2 \pi } e^{J^2/2}
\]
since 
\[
\left( \int_{-\infty}^{\infty} e^{-x^2} \right)^2 = \int_{-\infty}^{\infty} dx \int_{-\infty}^{\infty} dy e^{ - x^2 - y^2 } = \int_0^{\infty} r dr d\varphi e^{-r^2 } = 2\pi \left. \left( \frac{ -e^{-r^2} }{2} \right) \right|_0^{\infty} = \frac{ 2 \pi }{2} = \pi 
\]
\[
\int_{-\infty}^{\infty} e^{-x^2 } = \sqrt{ \pi }
\]
Note also, 
\[
\begin{aligned} y & = \frac{1}{\sqrt{2} } ( x- J)  \\ dy & = \frac{dx}{ \sqrt{2}} \end{aligned}
\]

\normalsize
\[
\exp{ W(g,J)} = \exp{ \left( \frac{1}{6} g(\delta_J)^3 \right) } e^{J^2/2 } = \sum_{V=0}^{\infty} \frac{1}{V!} \left( \frac{1}{6} g(\delta_J)^3 \right)^V \sum_{P=0}^{\infty} \frac{1}{P!} \left( \frac{J^2}{2} \right)^P
\]

\hrulefill

Recall Srednicki, Sec. 9.  Follow the development in Sec. 9.  

It begins with $\varphi^3$ theory.  
\[
Z_0(J) = \exp{ \left( \frac{i}{2} \int d^4y d^4z J(y) \Delta(y-z) J(z) \right) }
\]
\[
\begin{aligned}
  Z_1(J) & \propto \sum_{V=0}^{\infty} \frac{1}{V!} \left[ \frac{ i Z_g g}{6} \int d^4x \left( \frac{1}{i} \frac{ \delta}{ \delta J(x)} \right)^3 \right] \times \\ 
  & \phantom{ \propto } \times \sum_{P=0}^{\infty} \frac{1}{P!} \left[ \frac{ i}{2} \int d^4y d^4z J(y) \Delta(y-z) J(z) \right]^P  \quad \quad \quad \, (9.11) 
\end{aligned}
\]

Most general diagram contributing to $Z(J)$ consists of product of several connected diagrams.  

Let $C_I = $ particular connected diagram.  

Then general diagram is $D = \frac{1}{S_D} \prod_I (C_I)^{n_I}$.  

$n_I \in \mathbb{Z}$ counts number of $C_I$'s in $D$,  \\
$S_D$ is additional symmetry factor for $D$ (part of symmetry factor not already accounted for by symmetry factors included in each of the connected diagrams).  

Since propagator and vertex rearrangements \emph{within} each $C_I$ accounted for we only need to consider exchanged of propagators and vertices among different connected diagrams.  

These can leave the entire diagram $D$ unchanged only if (surely the following must be true):
\begin{enumerate}
\item exchange made among different but \emph{identical} connected diagrams.  
\item exchanges involve \emph{all} of the propagators and vertices in a given connected diagram (exchanged within a $C_I$ already accounted for)  
\end{enumerate}
\[
\Longrightarrow n_I!
\]
\[
S_D = \prod_I n_I!
\]

\[
\begin{gathered}
  Z_1(J) \propto \sum_{ \lbrace n_I \rbrace} D = \sum_{ \lbrace n_I \rbrace } \frac{1}{ S_D} \prod_I (C_I)^{n_I} = \sum_{ \lbrace n_I \rbrace } \prod_I \frac{1}{ n_I! } (C_I)^{n_I}   \propto \prod_I \sum_{n_I =0 }^{\infty} \frac{1}{ n_I!} (C_I)^{n_I} = \prod_I \exp{ ( C_I ) } = \exp{ \left( \sum_I C_i \right) }
\end{gathered}
\]

\hrulefill

So I claim that $Z_1(J) \propto \exp{ \left( \sum_{I} C_I \right) }$.  

Then if we identify \\ 
$\Delta(y-z) \to 1$,  \\
$i \int d^4y d^4z  J(y) J(z) \to J^2$ \\ 
$ iZ_g g \to g$ \\
$\int dx \left( \frac{1}{i} \delta_J \right)^3 \to \delta_J^3$

Then the result of $Z_1 \propto e^{ \sum_I C_I }$ applies here to, since the combinatorics are still the same, i.e., we still exchange derivatives $(\delta_J)^3$ at vertex $g$, we still can exchange $(2)$ sources, $J^2$.  

Thus in our case, 
\[
Z_1(J) = \exp{ W(g,J)} = \exp{ \left( \frac{1}{6} g (\delta_J)^3 \right)} \exp{ \left( \frac{J^2}{2} \right) } = \sum_{V=0}^{\infty} \frac{1}{V!} \left[ \frac{ g}{6} (\delta_J)^3 \right]^V \sum_{P=0}^{\infty} \frac{1}{P!} \left( \frac{J^2}{2} \right)^P \propto e^{\sum_I C_I }
\]
so $W(g,J) = \sum_I C_I$.  

A connected diagram is merely a term in the double Taylor expansion
\[
C_I = \frac{1}{S_I} g^{V_I} J^{E_I}
\]
So then
\[
W(g,J) = \sum_I \frac{1}{S_I} g^{V_I} J^{E_I}
\]

For each set of terms of some (specified) $V$ vertices, $E$ external sources, there's the same $g^V J^E$ factor for each of these terms.  Sum over all connected diagrams with such specified $V$, $E$.  

\[
\sum_K \frac{1}{S_K} g^{V_K} J^{E_K} = g^V J^E \sum_K \frac{1}{S_K}
\]

%But note that we sought connected diagrams with no tadpoles, or that the diagram has ample external sources, i.e. $E\geq 0$.  

%So impose $E = 2P - 3V$, $E\geq 0$.  

As in Problem 9.2(a), $V!$ ways to arrange $V$ $(\delta_J)^3$'s, $3!$ ways to arrange $3$ $\delta_J$'s, $2$ ways to arrange $J$'s, $P!$ ways to arrange $P$ $(J^2)$'s.  So the above result, saying that a connected diagram is simply a term of the form $\frac{1}{S_I} g^{V_I} J^{E_I}$ is justified.    

So we could impose the requirement of no tadpoles to rearrange our sum so that 
 
\[
W(g,J) = \sum_{V=0}^{\infty} \sum_{E=0}^{\infty} C_{V,E} g^V J^E
\]
where 
\[
\boxed{ C_{V,E} = \sum_K \frac{1}{S_K} }
\]

and $\sum_K$ is the sum over all connected diagrams with $V$ vertices and $E$ external sources.
\item[(b)]
\item[(c)] Now recall 
\[
exp{ W(g,J)} = \frac{1}{ \sqrt{ 2 \pi } } \int_{-\infty}^{\infty} dx \exp{ \left[ \frac{-1}{2} x^2 + \frac{1}{6} gx^3 +Jx \right] } \quad \quad \quad \, (9.27) 
\]
Adding a counterterm $Y$, 
\[
\exp{ W(g, J + Y) } = \frac{1}{ \sqrt{2 \pi }} \int_{-\infty}^{+\infty} dx \exp{ \left[ \frac{-1}{2} x^2 + \frac{1}{6} gx^3+ Jx + Y x \right] }
\]

Note
\[
\left. \frac{ \partial}{ \partial J} e^{W(g,J+Y) } \right|_{J=0} = e^{ W(g, J+Y) } \frac{ \partial }{ \partial J} \left. W(g,J+Y) \right|_{J=0} = 0 \Longleftrightarrow \frac{ \partial}{ \partial J} \left. W(g,J+Y) \right|_{J=0} = 0 
\]
are equivalent conditions.

Interpret $e^W$ as $Z(J)$
\[
Z(J) = e^{ W(g,J) } = e^{ \frac{1}{6} g ( \delta_J )^3 } e^{ J^2/ 2}  = \sum_{V=0}^{\infty} \frac{1}{V!} \left[ \frac{ g}{6} (\delta_J)^3 \right] \sum_{P=0}^{\infty} \frac{1}{P!} \left( \frac{J^2 }{2} \right)^P \propto e^{\sum_I C_I}
\]

Recall the development in Sec. 9.  Analogously, note
\[
Z(J) \simeq \left( 1 + \frac{g}{6} \delta_J^3 \right)\left( 1 + \frac{J^2}{2} + \frac{J^4}{4}  + \dots \right) + O(g^2) = 1 + \frac{J^2}{2} + \frac{J^4}{4} + gJ + O(J^3) 
\]
\[
\Longrightarrow \left. \delta_J Z_1 \right|_{J=0} = g \text{ or } \delta_J \left. W \right|_{J=0} \propto g 
\]
Thus a counterterm is necessary.  

Assuming $Y = O(g)$,
\[
Z_1(J) \simeq  \left( 1 + \frac{g}{6} \delta_J^3 \right) \left( 1 + \frac{ (J+Y)^2 }{2} + \frac{ (J+Y)^4}{4} \right) = 1 + \frac{ (J+Y)^2}{2} + \frac{ (J+Y)^4}{4} + g(J+Y)
\]
\[
 \left. \delta_J Z \right|_{J=0} = g + Y = 0 
\]
So $Y = - g$, assuming $Y = O(g)$.  

The ``no tadpole'' condition of $\frac{ \delta}{\delta J} \left. W(g,J+Y) \right|_{J=0} =0$ is equivalent to $\left. \frac{ \delta}{\delta J} e^{ W(g,J+Y) } \right|_{J= 0} = 0$ or 
\[
\delta_J \left. \exp{  W(g,J+Y) } \right|_{J=0} = \left. \frac{1}{ \sqrt{ 2 \pi }} \int_{-\infty}^{+\infty} dx \, x \exp{ \left[ \frac{-1}{2} x^2 + \frac{1}{6} gx^3 + Jx + Yx \right] } \right|_{J=0 } = 0 
\]
analogous to $\langle 0 | \varphi(x) | 0 \rangle = 0$.  

Interpreting connected Feynman diagrams as terms of the form $C_I = \frac{1}{S_I} g^{V_I} J^{E_I}$.  

For a $W(g,J + Y) = \sum_{V=0}^{\infty} \sum_{E=0}^{\infty} \widetilde{C}_{V,E} g^V J^E$  \quad \, (9.31) s.t. $ \frac{ \delta}{ \delta J}\left. W(g, J+Y) \right|_{J=0} = 0$, then no connected Feynman diagrams with tadpoles are included in the sum of connected Feynman diagrams $W(g,J+Y)$.  

\[
W(g,J+Y) = \sum_I \frac{1}{S_I} g^{V_I} J^{E_I} = \sum_{V=0}^{\infty} \sum_{E=0}^{\infty} \sum_K \frac{1}{ S_K} g^V J^E = \sum_{V=0}^{\infty} \sum_{E=0}^{\infty} \widetilde{C}_{V,E} g^V J^E
\]
where the sum $\sum_K$ is over all connected diagrams with no tadpoles of $V$ vertices, $E$ external sources.  

\item[(d)]
\end{itemize}


\section{Scattering amplitudes and the Feynman rules} 

\problemhead{10.2} Write down the Feynman rules for the complex scalar field of problem 9.3.  Remember that there are two kinds of particles now (which we can think of as positively and negatively charged), and that your rules must have a way of ditinguishing them.  Hint: the most direct approach requires two kinds of arrows: momentum arrows (as discussed in this section) and what we might call ``charge'' arrows (as discussed in problem 9.3).  Try to find a more elegant approach that requires only one kind of arrow.  

\solutionhead{10.2} With an expression for $Z(J) = \exp{ i W(J)}$, we can take functional derivatives to compute vacuum expectation values of time-ordered products of fields.  

Recall from Problem 9.3, 
\[
Z_0(J, J^{\dag} ) = \langle 0 | 0 \rangle_{J, J^{\dag}} = \int \mathcal{D}\varphi \mathcal{D} \varphi^{\dag} e^{ i \int d^4x [ \mathcal{L}_0 + J \varphi^{\dag} + J^{\dag} \varphi ] } = \exp{ \left[ i \int d^4x d^4x' J(x) \Delta(x-x') J^{\dag}(x') \right] }
\]
\[
\exp{ \left[ i \int d^4x \mathcal{L}_1  \right]} = \exp{ \left[ i \int d^4x \left( \frac{-1}{4} Z_{\lambda} \lambda \right) (\varphi^{\dag} \varphi)^2 \right] }
\]

Keep in mind the counterterms 
\[
\mathcal{L}_{ct} = - (Z_{\varphi} -  1 ) \partial^{\mu} \varphi^{\dag} \partial_{\mu} \varphi - (Z_m-1) m^2 \varphi^{\dag} \varphi = - A \partial^{\mu} \varphi^{\dag} \partial_{\mu} \varphi - B m^2 \varphi^{\dag} \varphi
\]

Recall doing the double Taylor series expansion:
\[
\exp{ \left[ \frac{-i }{4} Z_{\lambda} \lambda \int d^4x \left( \frac{1}{i} \delta_{J_x} \frac{1}{i} \delta_{J_x^{\dag}} \right)^2 \right] } Z_0(J) = Z(J,J^{\dag})
\]

So the path integral, with interaction, for our complex scalar field theory is 
\[
Z_1(J,J^{\dag}) = \sum_{V=0}^{\infty} \frac{1}{V!} \left[ \frac{i}{4} Z_{\lambda} \lambda \int d^4x \left( \frac{1}{i} \delta_{J_x} \frac{1}{i} \delta_{J_x^{\dag}} \right)^2 \right]^V \sum_{P=0}^{\infty} \frac{1}{P!} \left[ i \int d^4y d^4z J(y) \Delta(y-z) J^{\dag}(z) \right]^P 
\]

Recall, from Problem 8.7,

\[
\langle 0 |T \varphi(x_1) \varphi(x_2) | 0 \rangle = \langle 0 | T\varphi^{\dag}(x_1) \varphi^{\dag}(x_2) | 0 \rangle = 0 
\]
since, for example
\[
\begin{gathered}
  \langle 0 | T\varphi(x_1) \varphi(x_2) | 0 \rangle = \frac{1}{i} \delta_{J^{\dag}_{x_1} } \left[ \int d^4x J(x) \delta(x-x_1) \right] \left. Z_0(J,J^{\dag}) \right|_{J,J^{\dag}=0 } = \\
  = ( \int d^4x J(x) \Delta(x-x_2) ) (\int d^4x' J(x') \Delta(x'-x_1)) \left. Z_0(J,J^{\dag} )\right|_{J,J^{\dag} = 0} = 0 
\end{gathered}
\]
\[
\langle 0 | T\varphi^{\dag}(x_1) \varphi(x_2) | 0 \rangle = \frac{1}{i} \delta_{J_{x_1}} \left[ i \int d^4x J(x) \delta(x-x_2) \right] e^{ \frac{-i}{4} Z_{\lambda} \lambda \int d^4x \left( \frac{1}{i} \delta_{J_x} \frac{1}{i} \delta_{J_x^{\dag} } \right)^2 } Z_0(J,J^{\dag} ) = \Delta(x_1-x_2) + \mathcal{O}(\lambda)
\]

To study interaction in this theory, consider the following.  

\[
\begin{gathered}
  \langle 0 | T\varphi(x_3) \varphi^{\dag}(x_4) \varphi^{\dag}(x_1) \varphi(x_2) | 0 \rangle = \left. \left( \frac{1}{i} \delta_{J^{\dag}_{x_3}} \frac{1}{i} \delta_{J_{x_4}} \frac{1}{i} \delta_{J_{x_1}} \frac{1}{i} \delta_{J^{\dag}_{x_2} } \right) Z(J,J^{\dag}) \right|_{J,J^{\dag} = 0} \\
  = \left( \frac{1}{i} \delta_{J^{\dag}_{x_3}} \frac{1}{i} \delta_{J_{x_4}} \frac{1}{i} \delta_{J_{x_1}} \frac{1}{i} \delta_{J^{\dag}_{x_2} } \right) \sum_{V=0}^{\infty} \frac{1}{V!} \left( \frac{-i}{4} Z_{\lambda} \lambda \int d^4z \left( \frac{1}{i} \delta_{J_z} \frac{1}{i} \delta_{J_z^{\dag}} \right)^2 \right)^V \cdot \sum_{P=0}^{\infty} \frac{1}{P!} \left( i \int d^4x d^4x' J_x \Delta(x-x') J_{x'}^{\dag} \right)^P 
\end{gathered}
\]

% Note the first couple of terms of this double expansion.  

% \[
% \left( 1 + \frac{-1}{4} Z_{\lambda} \lambda \int d^4z \left( \frac{1}{i} \delta_{J_z} \frac{1}{i} \delta_{J^{\dag}_{z}} \right)^2 + \mathcal{O}(\lambda^2) \right) \cdot \left( 1 + i \int d^4x d^4x' J_x \Delta_{x-x'} J^{\dag}_{x'} + \frac{ ( i \int d^4x d^4x' J_x \Delta(x-x') J^{\dag}_{x'} )^2}{2} + \dots \right)
% \]

% Consider this particular term:
% \[
% \frac{1}{i} \delta_{J_{y_1}} \frac{1}{i} \delta_{J_{y_2}} \frac{1}{i} \delta_{J^{\dag}_{x_2}} \frac{1}{i} \delta_{J^{\dag}_{x_2}} \frac{1}{2} ( i \int d^4x d^4x' J_x \Delta_{x-x'} J_{x'}^{\dag} ) ( i \int d^4y d^4y' J_y \Delta_{y-y'} J_{y'}^{\dag} )
% \]
% Matching derivatives with propagators, there are $2*2 =4$ ways.  Those 4 ways divide into 2 different type of diagrams:
% \[
% \begin{gathered}
%  \frac{1}{i} \delta_{J_{y_1}} \frac{1}{i} \delta_{J_{y_2}} \frac{1}{i} \delta_{J^{\dag}_{x_2}} \frac{1}{i} \delta_{J^{\dag}_{x_2}} \frac{1}{2} ( i \int d^4x d^4x' J_x \Delta_{x-x'} J_{x'}^{\dag} ) ( i \int d^4y d^4y' J_y \Delta_{y-y'} J_{y'}^{\dag} ) = \\ 
%   = \Delta_{y_1 -x_2} \Delta_{y_2 - x_1} \text{ or } = \Delta_{y_2-x_2} \Delta_{y_1 - x_1}
% \end{gathered}
% \]


% \setlength{\unitlength}{0.5cm}

% \begin{picture}(4,4)\label{Fi:complexscalarfieldnoscattering}
% \put(0,1.25){\vector(1,0){2}}
% \put(2,1.25){\line(1,0){2}}
% \put(0,2.75){\vector(1,0){2}}
% \put(2,2.75){\line(1,0){2}}
% \end{picture}


% Observe the following diagrams, Figure (\ref{Fi:complexscalarfieldnoscattering}).  No scattering!  

Remember to restrict ourselves to fully connected diagrams.  

Consider, for $O(\lambda)$, $\lambda$ order,
\[
\left( \frac{1}{i} \delta_{J^{\dag}_{x_3}} \frac{1}{i} \delta_{J_{x_4}} \frac{1}{i} \delta_{J_{x_1}} \frac{1}{i} \delta_{J^{\dag}_{x_2} } \right) \frac{-i}{4} Z_{\lambda} \lambda \int d^4z \left( \frac{1}{i} \delta_{J_z} \frac{1}{i} \delta_{J_z^{\dag}} \right)^2 \cdot \frac{ ( i \int d^4x d^4x' J_x \Delta_{x-x'} J_{x'}^{\dag})^4}{ 4! }
\]

Matching up $\left( \delta_{J_z} \delta_{J_z^{\dag}} \right)^2$ with sources, there are $4!$ ways, if including only connected diagrams.  

Matching up $\delta_{J_{y_1}} \delta_{J_{y_2}} \delta_{J^{\dag}_{x_1}} \delta_{J_{x_2}^{\dag}}$, there are $2 \cdot 2 =4$ ways.  

Note that the factors $4 \cdot 4!$ match up to exactly cancel the factors in the denominators.  

\[
\begin{gathered}
  \left(  \frac{1}{i} \delta_{J^{\dag}_{x_3}} \frac{1}{i} \delta_{J_{x_4}} \frac{1}{i} \delta_{J_{x_1}} \frac{1}{i} \delta_{J^{\dag}_{x_2} } \right)  (-i) Z_{\lambda} \lambda \int d^4z ( \int d^4x J_x \Delta_{x-z} )( \int d^4y J_y \Delta_{y-z} ) ( \int d^4x' \Delta_{z-x'} J^{\dag}_{x'} ) ( \int d^4y' \Delta_{z-y'} J^{\dag}_{y'} )  \\
  = -i Z_{\lambda} \lambda \int d^4z \Delta_{x_4 - z} \Delta_{z -x_3 } \Delta_{x_1-z} \Delta_{z-x_2}
\end{gathered}
\]

\[
\Longrightarrow   \langle 0 | T\varphi(x_3) \varphi^{\dag}(x_4) \varphi^{\dag}(x_1) \varphi(x_2) | 0 \rangle = -i Z_{\lambda} \lambda \int d^4z \Delta_{x_4 - z} \Delta_{z -x_3 } \Delta_{x_1-z} \Delta_{z-x_2} + \mathcal{O}(\lambda^2)
\]

Recall Problem 5.1.  LSZ formula for complex scalar field.  

\[
\begin{gathered}
  \langle f | i \rangle = i^{n' +n } \int d^4x_{1'} e^{- i k_{1'} x_{1'} } ( - \partial_{1'}^2 + m_a^2 )  \dots  \int d^4x_{m'} e^{-ik_m x_{m'} } (- \partial^2_{m'} + m_a^2 ) \cdot \\ 
  \cdot \int d^4x_{m+1'} e^{ - i k_{m+1'} x_{m+1'} } ( - \partial^2_{m+1'} + m_b^2 )  \dots  \int d^4x_{n'} e^{ - i k_{n'} x_{n'} } (- \partial_{n'}^2 + m_b^2 ) \cdot \\ 
  \cdot \int d^4x_1 e^{ik_1 x_1} (- \partial_1^2 + m_a^2) \dots  \int d^4x_m e^{+ i k_m x_m } (- \partial^2_m + m_a^2 ) \cdot \\
  \cdot \int d^4x_{m+1} e^{i k_{m+1} x_{m+1} } ( - \partial_{m+1}^2 + m_b^2 ) \dots \int d^4x_n e^{ik_n x_n } (- \partial_n^2 + m_b^2 ) \cdot \\  
  \cdot \langle 0 | T\varphi(x_1') \dots \varphi(x_{m'} ) \varphi^{\dag}(x_{m+1}' ) \dots \varphi^{\dag}(x_{n'}) \varphi^{\dag}(x_1) \dots \varphi^{\dag}(x_m) \varphi(x_{m+1}) \dots \varphi(x_n) | 0 \rangle 
\end{gathered}
\]


In this case, for the vacuum expectation, $J, J^{\dag} =0$, so this condition forces ``a''-type particles, arising from $\varphi$ to be each paired up with ``b''-type particles, arising from $\varphi^{\dag}$.  Then, note that for the LSZ formula, there's the condition
\[
m' + (n-m) = n' - m' + m 
\]

Also, note that we are attempting a ``more elegant approach'' because the most direct approach requires two kinds of arrows: momentum arrows and what we might call ``charge'' arrows, that distinguish the sources $J$, $J^{\dag}$, or, in turn, distinguish between $\varphi$ and $\varphi^{\dag}$.  

%Note how in the above LSZ formula, the phase factors distinguish between the two type of particles, and between which particles are incident and which are final, i.e. a time ordering.  

Then I claim that because of the time ordering operator of the vacuum expectation, then $a$ particles could be distinguished from $b$ particles.  For incoming particles, $a$ particles correspond to $\varphi^{\dag}$ field and $b$ particles correspond to $\varphi$ field.  For outgoing particles, $a$ particles correspond to $\varphi$ field and $b$ particles correspond to $\varphi^{\dag}$ field.  

The term $\langle 0 | T \varphi(x_1') \dots \varphi^{\dag}(x'_{m+1}) \dots \varphi^{\dag}(x_1) \dots \varphi(x_{m+1} \dots | 0 \rangle$ make it very clear which particles are incoming and which are outgoing, compelled by the time ordering operator $T$.  


%then the integration part of the LSZ formula will also be time-ordered, appropriately.  

If in this case, $y_1, y_2$ come after $x_1, x_2$, then for this case, plug the above, very last result into the LSZ formula,

\[
\begin{gathered}
  \langle f | i \rangle = (i)^4 \int d^4x_3 e^{ - i k_{3} x_3 } (- \partial_{3}^2 + m_a^2) \int d^4x_4 e^{ -i k_{4} x_4 } ( - \partial_{4}^2 + m_b^2 ) \cdot \\ 
  \cdot \int d^4x_1 e^{i k_1 x_1} ( - \partial_1^1 + m_a^2) \int d^4x_2 e^{ik_2 x_2} (- \partial_2^2 + m_b^2)  ( -i) Z_{\lambda} \lambda \int d^4z \Delta_{x_4 - z} \Delta_{z -x_3 } \Delta_{x_1-z} \Delta_{z-x_2}
\end{gathered}
\]

Using $(-\partial_i^2 + m^2) \Delta(x_i - y) = \delta^4(x_i - y)$ (i.e. Klein-Gordon wave operator acts on propagator, i.e. propagator is a Green's function), \\
then for instance, 

\[
(-\partial_{4}^2 + m_b^2)\Delta_{x_4 - z} = \delta^4(x_4 - z) \xrightarrow{ \int d^4x_4 e^{-i k_{4} x_4} } e^{-ik_{4} z}
\]
\[
\Longrightarrow \langle f | i \rangle = -i Z_{\lambda} \lambda \int d^4z \exp{ \left[ ( - i k_3 - i k_4 + ik_1 + i k_2 ) z \right] } = -i Z_{\lambda} \lambda (2\pi)^4 \delta^{(4)}(k_1 + k_2 - k_3 - k_4 )
\]


\setlength{\unitlength}{0.5cm}

\begin{picture}(5,5)\label{Fi:complexscalar_abtoab_tree}
\put(2,2){\vector(1,1){1.2}}
\put(3,3){\line(1,1){1}}
\put(4,0){\vector(-1,1){1.2}}
\put(3,1){\line(-1,1){1}}
\put(0,4){\vector(1,-1){1}}
\put(1,3){\line(1,-1){1}}
\put(2,2){\vector(-1,-1){1.2}}
\put(1,1){\line(-1,-1){1}}
\put(2,2){\circle*{0.5}} 
\put(-0.7,0.5){$k_2$}
\put(-0.7,3.7){$k_1$}
\put(4,0.5){$k_4$}
\put(4,3.7){$k_3$}
\end{picture}

\quad \\ 

The above, Figure (\ref{Fi:complexscalar_abtoab_tree}), is for 1 $a$ and 1 $b$ particle coming in and a pair of $a$ and $b$ particles coming out, i.e., symbolically, $ab \to ab$.  

What about 2 other distinct cases, $aa \to bb$ and $bb \to aa$?  Are those processes allowed?  

At the tree level, no.  \\
For $aa \to bb$, we must consider $ \langle 0 | T \varphi^{\dag}(x_3) \varphi^{\dag}(x_4) \varphi^{\dag}(x_1) \varphi^{\dag}(x_2) | 0 \rangle$.   \\
For $bb \to aa$, we must consider $ \langle 0 | T \varphi(x_3) \varphi(x_4) \varphi(x_1) \varphi(x_2) | 0 \rangle$.  

The very first first order term, with no self-loops, with 4 propagators, that includes interaction, is \\
$ \frac{-i}{4} Z_{\lambda} \lambda \int d^4x \left( \int d^4y J_y \Delta_{y-x} \right)^2 \left( \int d^4z \Delta_{x-z} J_z^{\dag} \right)^2$ 

Then, when another four functional derivatves are to be applied, of the same kind, the result is $0$.  

Consider the counterterms $A,B$.  

\[
\begin{gathered}
  \mathcal{L}_{ct} = - (Z_{\varphi} - 1) \partial^{\mu} \varphi^{\dag} \partial_{\mu} \varphi - (Z_m - 1 ) m^2 \varphi^{\dag} \varphi = - A \partial^{\mu} \varphi \partial_{\mu} \varphi - Bm^2 \varphi^{\dag} \varphi \\ 
  e^{ i \int d^4x \mathcal{L}_{ct} } = \exp{ \left[ i \int d^4x \left( - A \partial^{\mu} \varphi^{\dag} \partial_{\mu} \varphi - B m^2 \varphi^{\dag} \varphi \right) \right] }  \\
  = \exp{ \left[ -i  \int d^4x \left( \delta_{J_x} \left(  - A \partial_x^2 + m^2 B \right) \delta_{J^{\dag}_x} \right) \right] }
\end{gathered}
\]
the last step by integration by parts.  

Consider the first order term in $e^{i \int \mathcal{L}_{ct}}$.  When this acts on 1 propagator, then the result is zero, after applying the LSZ formula. 

Go to 2 propagators.  
\[
-i \int d^4x \left( \left( - \delta_{J_x} A \partial_x^2 \right) \delta_{J_x^{\dag}} + m^2 B \delta_{J_x} \delta_{J_x^{\dag} } \right) \frac{1}{2} \left( i \int d^4y d^4y' J_y \Delta_{y-y'} J_{y'}^{\dag} \right)^2  
\]

For the term involving $B$,

\[
\Longrightarrow + i \int d^4x m^2 B \int d^4y \Delta_{x-y} J^{\dag}_y \int d^4z J_z \Delta_{z-x}
\]

For the term involving $A$,
\[
i d^4x ( - \delta_{J_x} A \partial_x^2 ) \int d^4y d^4y' J_y \Delta_{y-y'} J_{y'}^{\dag} \int d^4z J_z \Delta_{z-x} = (-i A) \int d^4x d^4y d^4z \Delta_{x-y} J^{\dag}_y J_z \partial_x^2 \Delta_{z-x}
\]

Consider 
\[
i^2 \int d^4x_2 e^{-ik_2 x_2} ( - \partial_2^2 + m_a^2) d^4x_1 e^{ik_1 x_1} ( - \partial_1^2 + m_a^2) \langle 0 | T \varphi(x_2) \varphi^{\dag}(x_1) | 0 \rangle
\]

So for the term involving $B$, applying the LSZ formula,
\[
\frac{i}{i^2} \int d^4x m^2 B \Delta_{x-x_2} \Delta_{x_1 -x} \Longrightarrow im^2 B (2\pi)^4 \delta^4(k_1 -k_2)
\]

Then for the term involving $A$, first applying the functional derivatives due to the LSZ formula,
\[
\frac{1}{i} \delta_{J_{x_2}^{\dag}} \frac{1}{i} \delta_{J_{x_1} } (-i A) \int d^4x d^4y d^4z \Delta_{x-y} J_y^{ \dag} J_z \partial_x^2 \Delta_{z-x} = i A \int d^4x \Delta_{x-x_2} \partial_x^2\Delta_{x_1- x}
\]
\[
\begin{gathered}
  \Longrightarrow i^2 \int d^4x_2 e^{-ik_2 x_2} ( - \partial_2^2 + m_a^2 ) d^4x_1 e^{ik_1 x_1} ( - \partial^2_1 + m_a^2) i A \int d^4x \Delta_{x-x_2} \partial_x^2 \Delta_{x_1 - x} = \\ 
   = - i A \int d^4x e^{-ik_2 x} \partial_x^2 e^{i k_1 x} = ( - i A) ( - k_1^2 ) ( 2\pi)^4 \delta^4(k_1 - k_2) = i A k^2 ( 2\pi)^4 \delta^4(k_1 - k_2) 
\end{gathered}
\]

So the counterterms must be accounted for in our diagrams by sprinkling vertices of factor $i (Ak^2 + Bm^2)$ on each propagator line.  

\quad \\ 

% \setlength{\unitlength}{0.5cm}

% \begin{picture}(4,4)
% \put(2,2){\vector(1,1){1.3}}
% \put(3,3){\line(1,1){1}}
% \put(2,2){\vector(1,-1){1.3}}
% \put(3,1){\line(1,-1){1}}
% \put(0,4){\vector(1,-1){1}}
% \put(1,3){\line(1,-1){1}}
% \put(0,0){\vector(1,1){1}}
% \put(1,1){\line(1,1){1}}
% \put(2,2){\circle*{0.5}} 
% \end{picture}


%Of course there are a total of 3 different kinds of diagrams, depending on the time-ordering, the 2 other kinds being diagrams for the following expressions:
% \[
% -i Z_{\lambda} \lambda (2\pi)^4 \delta^{(4)}(k_2' - k_1' + k_1 - k_2)
% \]




% and
% \[
% -i Z_{\lambda} \lambda ( 2 \pi)^4 \delta^{(4)}( -k_1' -k_2' -k_1 - k_2)
% \]


% \setlength{\unitlength}{0.5cm}

% \begin{picture}(4,4)
% \put(4,4){\vector(-1,-1){1.2}}
% \put(3,3){\line(-1,-1){1}}
% \put(4,0){\vector(-1,1){1.2}}
% \put(3,1){\line(-1,1){1}}
% \put(2,2){\vector(-1,1){1.2}}
% \put(1,3){\line(-1,1){1}}
% \put(2,2){\vector(-1,-1){1.2}}
% \put(1,1){\line(-1,-1){1.2}}
% \put(2,2){\circle*{0.5}} 
% \end{picture}



The Feynman rules would be the following:
\begin{enumerate}
\item Draw lines (called \emph{external lines}) for each incoming and each of the outgoing particles.
\item Leave one end of each external line free, and attach the other to a vertex at which exactly 4 lines meet.  Include extra \emph{internal lines} in order to do this.  In this ways, draw all possible diagrams that are \emph{topologically inequivalent}.  
\item Assign each line its own four-momentum.  The four-momentum of an external line should be the four-momentum of the corresponding particle.  
\item The time-ordering makes clear which particles are incoming (from the left) and which particles are outgoing (to the right).  

Now there are two type of particles, arising from the two different sources $J,J^{\dag}$, resulting in fields $\varphi^{\dag}, \varphi$.  Type ``a'' particle that is incoming is distinguished by the arrow flowing towards the vertex and type ``b'' particle that is incoming is distinguished by the arrow flowing away from the vertex.  

Again, because of the time-ordering operator, and if, on the diagram, we agree that time flows from left to right, then it's clear how to distinguish type ``a'' particle from type ``b'' particle just by the momentum flow.  The outgoing $a$ particles have an arrow flowing away from the vertex and the outgoing $b$ particles have an arrow flowing towards the vertex.  



Incoming type ``a'' particles have a 4-momentum of $k_a$; outgoing type ``a'' particles have a 4-momentum of $k'_a$.  Incoming type ``b'' particles have a 4-momentum of $k_b$; outgoing type ``b'' particles have a 4-momentum of $k'_b$.  
\item Thus so, conserve 4-momentum at each vertex.
\item The value of a diagram consists of the following factors: \\
for each external line, $1$: \\
for each vertex, $-iZ_{\lambda} \lambda$ \\
for each propagator, 
\end{enumerate}




\solutionhead{10.3} 

%Consider a complex scalar field $\varphi$ that interacts with a real scalar field $\chi$ via $\mathcal{L}_1 = g\chi \varphi^{\dag} \varphi$.  Use a solid line for the $\varphi$ propagator and a dashed line for the $\chi$ propagator.  Draw the vectex (remember the arrows!), and find the associated vertex factor.  

% \solutionhead{10.3} Consider $\mathcal{L}_1 = g \chi \varphi^{\dag} \varphi$ \\
% Consider utilizing $\frac{ \delta}{\delta J_{\chi}}$.  

% Analogous to previous developments, \\
% free scalar field: $Z_{0\chi}(J_{\chi}) = \exp{ \left[ \frac{i}{2} \int d^4x d^4x' J_{\chi}(x) \Delta_{\chi}(x-x') J_{\chi}(x') \right] }$  \\
% so it's clear that to get $\chi$ fields, apply the corresponding functional derivative: $\chi \to \frac{\delta}{\delta J_{\chi}}$.   \\

% Complex scalar field, free theory: $Z_{0 \varphi^{\dag}}(J^{\dag},J) = \exp{ \left[ \frac{i}{2} \int d^4x J^{\dag}(x) \Delta_{\phi \phi^{\dag}}(x) J(x) \right]}$, \\
% and recall, the source terms come from the following: $J^{\dag} \varphi + J\varphi^{\dag}$.   \\
% $\frac{\delta}{\delta J}$ to get $\varphi^{\dag}$. \\
% $\frac{\delta}{\delta J^{\dag}}$ to get $\varphi$.  

% \[
% \exp{ \left[ i \int d^4 x \mathcal{L}_1(\chi, \varphi^{\dag},\phi ) \right]} \Longrightarrow \exp{ \left[ i \int d^4 x \mathcal{L}_1\left( \frac{ \delta}{ \delta J_{\chi}}, \frac{ \delta}{\delta J}, \frac{\delta}{ \delta J^{\dag}} \right) \right] } = \exp{ [ i \int d^4x g \delta_{J_{\chi}} \delta_J \delta_{J^{\dag}} ]}
% \]

% Consider only connected diagrams.  \\
% Consider 1 vertex (first order in $g$).  

% \[
% ig \int d^4x \delta_{J_{\chi}} \delta_J \delta_{J^{\dag}} \Longrightarrow \left. \left( \frac{i}{2} \int d^4x' J_{\chi}(x') \Delta_{\chi}(x'-x) \right) W_{\chi} \frac{i}{2} \Delta_{\varphi \varphi^{\dag} }(x) W_{\varphi^{\dag}} \right|_{J^{\dag},J =0}
% \]

% Inevitably, a source for the real scalar field is necessary in this 1-vertex diagram.  

% Go to the second order:
% \[
% \begin{aligned}
%  &  \left( i g \int d^4x \delta_{J_{\chi}} \delta_J \delta_{J^{\dag}} \right)^2 \\ 
%  \Longrightarrow & (ig)^2 \int d^4 y d^z \left( \frac{i}{2} \Delta_{\chi}(y-z) \right) \frac{i}{2} \Delta_{\varphi^{\dag}(y)}(y) \frac{i}{2} \Delta_{\varphi^{\dag}}(z)
% \end{aligned}
% \]

Consider using this free field path integral for the real scalar field.  
\[
Z_{\chi_0}(K) = \exp{ \left[ \frac{i}{2} \int d^4x d^4x' K_x \rho_{x-x'} K_{x'} \right] }
\]
which is coming from a source term $K \chi_0$ and free-field Lagrangian.  

Consider using this free field path integral for the complex scalar field.
\[
Z_{\varphi^{\dag} 0}(J,J^{\dag}) = \exp{ \left[ \frac{i}{2} \int d^4x d^4x' J^{\dag}_x \Delta_{x-x'} J_{x'} \right] }
\]
which came originally from a source term $J^{\dag} \varphi + J \varphi^{\dag}$ and a Lagrangian. 

\[
\exp{ \left[ i \int d^4x \mathcal{L}_1( \chi, \varphi, \varphi^{\dag} ) \right] } \Longrightarrow \exp{ \left[ i \int d^4x \mathcal{L}_1\left( \frac{1}{i} \delta_K , \frac{1}{i} \delta_J , \frac{1}{i} \delta_{J^{\dag}} \right) \right] } = \exp{ \left[ i \int d^4x g \delta_K \delta_J \delta_{J^{\dag}} \right] }
\]
\[
\Longrightarrow \begin{gathered}
  Z = e^{i \int d^4x \mathcal{L}_1 }Z_{\chi 0} Z_{\varphi^{\dag} 0}   = \left( 1 + i \int d^4x g \frac{1}{i} \delta_{K_x} \frac{1}{i} \delta_{J_x} \frac{1}{i} \delta_{J_x^{\dag}} + \mathcal{O}(g^2) \right) ( Z_{\chi 0} Z_{\varphi^{\dag}0 } )
\end{gathered}
\]

Consider the term $i \int d^4x g \frac{1}{i} \delta_{K_z} \frac{1}{i} \delta_{J_z} \frac{1}{i} \delta_{J_z^{\dag}} Z_{\chi_0} Z_{\varphi^{\dag}_0 }$.  

\[
\begin{gathered}
  \int d^4z g \frac{1}{i} \delta_{K_z} \frac{1}{i} \delta_{J_z} \frac{1}{i} \delta_{J^{\dag}_z} Z_{\chi 0} Z_{\varphi^{\dag} 0 } = - \int d^4z g \delta_{K_z} \delta_{J_z} \delta_{J^{\dag}_z} \left( 1 + \frac{i}{2} \int d^4u d^4w K_u \rho_{u -w} K_w \right) \cdot \\ 
  \cdot \left( 1 + \frac{i}{2} \int d^4x d^4y J_x^{\dag} \Delta_{x-y} J_y + \frac{1}{2} \left( \frac{i}{2} \int d^4x d^4y J_x^{\dag} \Delta_{x-y} J_y \right)^2 + \dots \right) 
\end{gathered}
\]

If $\delta_{J_z} \delta_{J^{\dag}_z}$ acts on $\frac{i}{2} \int d^4x d^4y J_x^{\dag} \Delta_{x-y} J_y$, we'd get $\Delta(0)$ (i.e. closed loop in the diagram).  

\[
- \int d^4z g \delta_{K_z} \delta_{J_z} \delta_{J_z^{\dag}} \left( \frac{i}{2} \int d^4u d^4w K_u \rho_{u-w} K_w \right) \left( \frac{1}{2} \left( \frac{i}{2} \int d^4x_1 d^4y_1 J_{x_1}^{\dag} \Delta_{x_1- y_1} J_{y_1} \right) \left( \frac{i}{2} \int d^4x_1 d^4y_2 J_{x_2}^{\dag} \Delta_{x_2 - y_2} J_{y_2} \right) \right)
\]

\noindent 2 ways to match $\delta_{J_z} \delta_{J_z^{\dag}}$ with propagator.   \\
2 ways to match $\delta_{K_z}$ with propagator.  \\
Symmetry factor of $2^2 =4$ (as we'll see from the diagram).   \\

$2*2*4$ cancels exactly with the denominator.   \\

\[
i \int d^4z g \int d^4w \rho_{z-w} K_w \int d^4x_1 J_{x_1}^{\dag} \Delta_{x_1-z} \int d^4y_2 \Delta_{z-y_2} J_{y_2} 
\]

To compute out a probability amplitude, we'll ultimately need to strip off the sources.  
\[
\langle 0 | T\chi(x_1') \varphi^{\dag}(x_1) \varphi(x_2) | 0 \rangle = \frac{1}{i} \delta_{K_{x_1'}} \frac{1}{i} \delta_{J_{x_1}} \frac{1}{i} \delta_{J^{\dag}_{x_2}} Z
\]

Consider the following.
\[
\begin{aligned}
  \frac{1}{i} \delta_{K_{x_1'}} \frac{1}{i}  \delta_{J_{x_1}} \frac{1}{i} \delta_{J^{\dag}_{x_2}} i \int d^4z g \int d^4w \rho_{z-w}  K_w \int d^4y_1 J^{\dag}_{y_1} \Delta_{y_1 -z} \int d^4y_2 \Delta_{z-y_2} J_{y_2}  = - i \int d^4z g \rho_{z- x_1'} \Delta_{x_2 - z} \Delta_{z-x_1} 
\end{aligned}
\]

\[
\langle f | i \rangle = i^3 \int d^4x_1' e^{-ik_1' x_1'} (- \partial_{ \chi}^2 + m_{\chi}^2) \int d^4x_1 e^{ik_1 x_1} (-\partial_1^2 + m_{\varphi}^2 ) \int d^4x_2 e^{-ik_2 x_2} (- \partial_2^2 + m_{\varphi^{\dag}}^2 ) \langle 0 | T \chi(x_1') \varphi^{\dag}(x_1) \varphi(x_2) | 0 \rangle
\]

Noting that propagators are Green's functions satisfying Klein-Gordon,
\[
\Longrightarrow \langle f | i \rangle = i \int d^4z g e^{ -i k_1' z} e^{ i k_1 z} e^{-ik_2 z} = ig (2\pi)^4 \delta^{(4)}(k_1 -k_2 -k_1' )
\]



I had, as developed fully for my solution in Problem 10.2, denoted the antiparticle corresponding to $\varphi^{\dag}$ of the complex scalar field, with momentum $-k_2$.  However, this overall $(-1)$ factor, negative sign, doesn't change the fact that both incident particle and antiparticle corresponding to $\varphi$, $\varphi^{\dag}$ are moving forward in spacetime.  

\setlength{\unitlength}{0.5cm}

\begin{picture}(4,4)
\put(0,4){\vector(1,-1){1.2}}
\put(1,3){\line(1,-1){1}}
%\put(4,0){\vector(-1,1){1.2}}
%\put(3,1){\line(-1,1){1}}
%\put(2,2){\vector(-1,1){1.2}}
%\put(1,3){\line(-1,1){1}}
\put(2,2){\vector(-1,-1){1.2}}
\put(1,1){\line(-1,-1){1.2}}
%\put(2,2){\circle*{0.5}} 
\put(2,2){\line(1,0){3.2}}
\put(2.5,2.5){ig}
\end{picture}


\problemhead{10.4} Consider a real scalar field with $\mathcal{L}_1 = \frac{1}{2} g \varphi \partial^{\mu} \varphi \partial_{\mu} \varphi$.  Find the associated vertex factor.  

% 20100129

\solutionhead{10.4}

Given $\mathcal{L}_1 = \frac{1}{2} g \varphi \partial^{\mu} \varphi \partial_{\mu} \varphi$.   \\

Consider $\exp{ \left[ i \int d^4x \left( \mathcal{L}_0  + J \varphi + \mathcal{L}_1 \right) \right] }$
 
Recall the 4-dim. Fourier transforms: 
\[
\varphi(x) = \int \frac{d^4k}{(2\pi)^4} e^{ikx} \widetilde{\varphi}(k), \quad \quad \quad \, \widetilde{\varphi}(k) = \int d^4x e^{-ikx} \varphi(x) \quad \quad \quad \, (8.6) 
\]

Consider $\int d^4x \varphi \partial^{\mu} \varphi \partial_{\mu} \varphi$.  

Recall 
\[
\begin{aligned}
  \partial^{\mu} & = ( - \partial_t, \mathbf{\nabla} ) \\ 
  \partial_{\mu} & = ( \partial_t, \mathbf{\nabla} )
\end{aligned} \quad \quad \quad \, 
\begin{aligned}
  \partial^{\mu} e^{ikx} & = ( ik^0, i \mathbf{k} ) e^{ikx} \\ 
  \partial_{\mu} e^{ikx} & = ( -i k^0, i \mathbf{k} ) e^{ikx}
\end{aligned}
\]
\[
\begin{gathered}
  \Longrightarrow \int d^4x \int \frac{d^4k_a}{ (2\pi)^4} e^{ ik_a x} \widetilde{\varphi}_{k_a} \partial^{\mu} \int \frac{d^4 k_b}{ (2\pi)^4} e^{i k_b x} \widetilde{\varphi}_{k_b} \partial_{\mu} \int \frac{ d^4 k_c}{  (2\pi)^4} e^{ik_c x} \widetilde{\varphi}_{k_c} = \\
  \begin{aligned}
    & = \int \frac{d^4x }{(2\pi)^4} \int d^4k_a \int \frac{d^4k_b}{(2\pi)^4} \int \frac{d^4k_c}{(2\pi)^4} \widetilde{\varphi}_{k_a} \widetilde{\varphi}_{k_b} \widetilde{\varphi}_{k_c} ( - k_b k_c) e^{ i ( k_a + k_b + k_c) x} = \\
    & = \frac{1}{ (2\pi)^8 } \int d^4k_a d^4k_b d^4k_c \delta^{(4)}(k_a + k_b + k_c) (-k_b k_c) \widetilde{\varphi}_{k_a} \widetilde{\varphi}_{k_b} \widetilde{\varphi}_{k_c}
  \end{aligned}
\end{gathered}
\]

Note that we can do $\widetilde{\chi}(k) = \widetilde{\varphi}(k) - \frac{ \widetilde{J}(k)}{k^2 + m^2}$ change of integration variable for $\mathcal{L}_0 + J \varphi$ only, separately; doesn't concern us for $\mathcal{L}_1$.  

\[
e^{ i \int d^4x \mathcal{L}_1} = e^{ i \int d^4x \frac{1}{2} g \varphi \partial^{\mu} \varphi \partial_{\mu} \varphi} = \exp{ \left[ \frac{ig}{2} \frac{1}{(2\pi)^ 8 } \int d^4k_a d^4k_b d^4k_c \delta^{(4)}(k_a + k_b +k_c)(-k_b k_c) \widetilde{\varphi}_{k_a} \widetilde{\varphi}_{k_b} \widetilde{\varphi}_{k_c} \right] }
\]

Looking at the form of $S_0$, and that $Z_0(J) = \int \mathcal{D} \varphi e^{iS_0}$.  

\[
S_0 = \frac{1}{2} \int \frac{d^4k}{(2\pi)^4} \left[ - \widetilde{\varphi}(k) ( k^2 + m^2) \widetilde{\varphi}(-k) + \widetilde{J}(k) \widetilde{\varphi}(-k) + \widetilde{J}(-k) \widetilde{\varphi}(k) \right]
\]

Consider this functional derivative, but for momentum space: $\frac{1}{i} \frac{ \delta}{ \delta J(-k) } = \frac{1}{i} \delta_{J_{-k} }$

So to $O(g)$, 

\[
\frac{ig}{2} \frac{1}{(2\pi)^8} \int d^4k_a d^4k_b d^4k_c \delta^{(4)}(k_a + k_b + k_c) (-k_b k_c) \frac{1}{i} \delta_{J_{-k_a}} \frac{1}{i} \delta_{J_{-k_b} } \frac{1}{i} \delta_{J_{-k_c}}
\]
operate on $Z_0$.  

Need 3 propagators:
\[
\frac{1}{3!} \left( \frac{i}{2} \right)^3 \int \frac{d^4k_1}{ (2\pi)^4} \frac{ \widetilde{J}(k_1) \widetilde{J}(-k_1) }{ k_1^2 + m^2 - i \epsilon } \int \frac{d^4k_2}{ (2\pi)^4} \frac{ \widetilde{J}(k_2) \widetilde{J}(-k_2) }{ k_2^2 + m^2 - i \epsilon } \int \frac{d^4k_3}{ (2\pi)^4} \frac{ \widetilde{J}(k_3) \widetilde{J}(-k_3) }{ k_3^2 + m^2 - i \epsilon } 
\]

$3!$ ways to match $\delta_J$ to propagator, $2^3$ ways for the 3 times to choose which $\widetilde{J}$ of each propagator to act on 2 ways to choose which is $k_b$ or $k_c$ (indistinguishable).  

\[
\Longrightarrow ig \frac{1}{(2\pi)^8} \int d^4k_a d^4k_b d^4k_c \delta^{(4)}(k_a + k_b + k_c) (-k_b k_c) \frac{1}{(2\pi)^{12}} \frac{ \widetilde{J}_{k_a} }{ k_a^2 + m^2 - i \epsilon} \frac{ \widetilde{J}_{k_b} }{ k_b^2 + m^2 - i \epsilon} \frac{ \widetilde{J}_{k_c} }{ k_c^2 + m^2 - i \epsilon}
\]

\hrulefill 

Let's examine the LSZ formula for momentum space.  

\[
\langle f | i \rangle = i^{n + n' } \int d^4x_1 e^{ik_1 x_1} ( - \partial_1^2 + m^2 ) \dots d^4x_1' e^{-ik_1' x_1'} ( - \partial_{1'}^2 + m^2 ) \dots \times \langle 0 | T \varphi(x_1) \dots \varphi(x_1') \dots | 0 \rangle 
\]

Use $\varphi(x) = \int \frac{ d^4k}{(2\pi)^4} e^{ikx} \widetilde{\varphi}(k)$ when considering
\[
\begin{gathered}
  \int d^4x_1 e^{ik_1 x_1} ( - \partial_1^2 + m^2 ) \int \frac{d^4k_a}{(2\pi)^4} e^{i k_a x_1} \dots \langle 0 | T \widetilde{\varphi}(k_a) \dots | 0 \rangle = \\
   = \int d^4x_1 \int \frac{d^4k_a}{(2\pi)^4} e^{ik_1 x_1} ( k_a^2 + m^2 ) e^{ik_a x_1} \dots \langle 0 | T\widetilde{\varphi}(k_a) \dots | 0 \rangle
\end{gathered}
\] 

\noindent For incoming particles, $ \frac{1}{ (2\pi)^4} \int d^4x_1 e^{ i (k_1 + k_a) x_1} = \delta^{(4)}(k_1 + k_a)$ \\ 
\phantom{For } outgoing particles, $ \frac{1}{ (2\pi)^4} \int d^4x_1'^2 e^{ -i (k_1' - k_a')x_1' } = \delta^{(4)}(k_1' - k_a')$

\[
\Longrightarrow \boxed{ \langle f | i \rangle = i^{n+n'} (k_1^2 + m_1^2) \dots (k_{1'}^2 + m_{1'}^2 ) \dots \times \langle 0 | T \widetilde{\varphi}(-k_1) \dots \widetilde{\varphi}(k_1') \dots | 0 \rangle }
\]

\hrulefill

So for our case,
\[
\langle f | i \rangle = i^3 ( k_1^2 + m^2)(k_2^2 + m^2)(k_3^2 + m^2) \langle 0 | T \widetilde{\varphi}(-k_1) \widetilde{\varphi}(-k_2) \widetilde{\varphi}(k_3) | 0 \rangle
\]

So for our $Z_1$, we need to further apply $\frac{1}{i} \delta_{J_{-k_3} } \frac{1}{i} \delta{J_{k_2} } \frac{1}{i} \delta_{J_{k_1} }$ to bring down $T \widetilde{\varphi}_{-k_1} \widetilde{\varphi}_{-k_2} \widetilde{\varphi}_{k_3}$.  \\
There are a number of ways to match these functional derivatives to the external sources, external legs, i.e. ``amputate the legs,'' distinguishable by $(-k_bk_c)$ factor.   \\

\begin{tabular}{ l r c c r c c r c c }
  a & \, & 1 & 1 & \, & 2 & 2 & \, & 3 & 3  \\ 
  b & \, & 2 & 3 & \, & 1 & 3 & \, & 1 & 2  \\ 
  c & \, & 3 & 2 & \, & 3 & 1 & \, & 2 & 1    \\
  & & & $-(-k_2)k_3$  & & & $-(-k_1)k_3$  & & & $-(-k_1)(-k_2)$ 
\end{tabular} \\

There are actually 3 distinguishable arrangements, not $3!$, as the particles associated with $k_b,k_c$ should be indistinguishable.  

\[
\Longrightarrow ig \delta^{(4)}(k_1 + k_2 - k_3)(k_2 k_3 + k_1 k_3 - k_1 k_2) = \boxed{ ig ( k_1^2 + k_2^2 + k_1k_2) }
\]


%%%%%%%%%%%% 20100128 -------------




\problemhead{10.5} The scattering amplitudes should be unchanged if we make a \emph{field redefinition}.  Suppose, for example, we have
\begin{equation*}
  \mathcal{L} = -\frac{1}{2} \partial^{\mu} \varphi \partial_{\mu} \varphi - \frac{1}{2} m^2 \varphi^2  \quad \quad \quad \, (10.15)
\end{equation*}
and we make the field redefinition
\begin{equation*}
  \varphi \to \varphi + \lambda \varphi^2 \quad \quad \quad \, (10.16)
\end{equation*}
Work out the lagrangian in terms of the redefined field, and the corresponding Feynman rules.  Compute (at tree level) the $\varphi \varphi \to \varphi \varphi$ scattering amplitude.  You should get zero, because this is a free-field theory in disguise.  (At the loop level, we also have to take into account the transformation of the functional measure $\mathcal{D} \varphi$; see section 85.)

\solutionhead{10.5}

\[
\mathcal{L} = - \frac{1}{2} \partial^{\mu} \varphi \partial_{\mu} \varphi - \frac{1}{2} m^2 \varphi^2
\]
\[
\begin{aligned}
  & \varphi \to  \varphi + \lambda \varphi^2 \\ 
  & \varphi^2 \to  (\varphi + \lambda \varphi^2)(\varphi + \lambda \varphi^2) = \varphi^2 + 2 \lambda \varphi^3 + \lambda^2 \varphi^4
\end{aligned}
\]
\[
\begin{aligned}
  & \partial^{\mu} \varphi  \to & \partial^{\mu} \varphi + 2 \lambda \varphi \partial^{\mu} \varphi \\ 
  & \partial_{\mu} \varphi \to & \partial_{\mu} \varphi + 2 \lambda \varphi \partial_{\mu} \varphi 
\end{aligned}
\]

\[
\begin{gathered}
  (\partial^{\mu} \varphi + 2 \lambda \varphi \partial^{\mu} \varphi) ( \partial_{\mu} \varphi + 2 \lambda \varphi \partial_{\mu} \varphi )  = \partial^{\mu} \varphi \partial_{\mu} \varphi + 4 \lambda \varphi \partial^{\mu} \varphi \partial_{\mu} \varphi + 4 \lambda^2 \varphi^2 \partial^{\mu} \varphi \partial_{\mu} \varphi = \\ 
  = ( 1 + 4 \lambda \varphi + 4 \lambda^2 \varphi^2 ) \partial^{\mu} \varphi \partial_{\mu} \varphi
\end{gathered}
\]
\[
\mathcal{L} \to \mathcal{L}_0 + \mathcal{L}_3 + \mathcal{L}_4
\]
where
\[
\begin{aligned}
  & \mathcal{L}_0 = - \frac{1}{2} \partial^{\mu} \varphi \partial_{\mu} \varphi - \frac{1}{2} m^2  \varphi^2 \\
  & \mathcal{L}_3 = \lambda \left( -2 \varphi \partial^{\mu} \varphi \partial_{\mu} \varphi - m^2 \varphi^3 \right) \\ 
  & \mathcal{L}_4 = \lambda^2 \left( -2 \varphi^2 \partial^{\mu} \varphi \partial_{\mu} \varphi - \frac{1}{2} m^2 \varphi^4 \right)
\end{aligned}
\]
So then
\[
Z(J) \equiv \langle 0 | 0 \rangle_J = \int \mathcal{D} \varphi e^{ i \int d^4x \left[ \mathcal{L}_0 + \mathcal{L}_3 + \mathcal{L}_4 + J \varphi \right] }
\]

Now

\[
e^{ i \int d^4x \mathcal{L}_3 } = \exp{ \left[ - i 2 \lambda \frac{ 1 }{ (2\pi)^8 } \int d^4k_a d^4k_b d^4k_c \delta^{(4)}(k_a + k_b +k_c)( -k_b k_c) \widetilde{\varphi}_{k_a} \widetilde{\varphi}_{k_b} \widetilde{\varphi}_{k_c} - i \lambda m^2 \int d^4x \varphi^3 \right] }
\]
to first-order in $\lambda$, i.e. $O(\lambda)$, 
\[
\lambda \left( - 2 i \frac{1}{ ( 2\pi)^8 } \int d^4k_a d^4k_b d^4k_c \delta^{(4)}(k_a + k_b + k_c) (- k_b k_c) \widetilde{\varphi}_{k_a} \widetilde{\varphi}_{k_b} \widetilde{\varphi}_{k_c} - i m^2 \int d^4x \varphi^3 \right)
\]

Acting on 3 propagators,
\[
\begin{gathered}
  \frac{1}{3!} \left( \frac{ i }{2} \int \frac{ d^4k_{ \alpha} }{ (2\pi)^4 } \frac{ \widetilde{ J}_{k_{\alpha} } \widetilde{ J}_{-k_{\alpha}} }{ k_{\alpha}^2 + m^2 - i \epsilon } \right) \left( \frac{ i }{2} \int \frac{ d^4k_{ \beta} }{ (2\pi)^4 } \frac{ \widetilde{ J}_{k_{\beta} } \widetilde{ J}_{-k_{\beta}} }{ k_{\beta}^2 + m^2 - i \epsilon } \right) \left( \frac{ i }{2} \int \frac{ d^4k_{ \gamma} }{ (2\pi)^4 } \frac{ \widetilde{ J}_{k_{\gamma} } \widetilde{ J}_{-k_{\gamma}} }{ k_{\gamma}^2 + m^2 - i \epsilon } \right) = \\ 
  = \frac{1}{3!} \left( \frac{i }{2} \int d^4x d^4x' J_x \Delta_{x-x'} J_{x'} \right) \left( \frac{i }{2} \int d^4y d^4y' J_y \Delta_{y-y'} J_{y'} \right) \left( \frac{i }{2} \int d^4z d^4z' J_z \Delta_{z-z'} J_{z'} \right)
\end{gathered}
\]

$3!$ ways to match derivatives $\delta_{J_{-k_i}}$ to propagators.  $2^3$ total ways to pick which of 2 sources.  

\[
\begin{gathered}
  \lambda \left( - 2i \frac{1}{(2\pi)^{20}} \int d^4k_a d^4k_b d^4k_c \delta^{(4)}(k_a + k_b + k_c)(-k_b k_c) \frac{ \widetilde{J}_{k_a} }{ k_a^2 + m^2 - i\epsilon } \frac{ \widetilde{J}_{k_b}}{ k_b^2 + m^2 - i \epsilon}  \frac{ \widetilde{J}_{k_c}}{ k_c^2  +m^2 - i \epsilon }  + \right. \\ 
\left.   + -i m^2 \int d^4x \int d^4y J_y \Delta_{y-x} \int d^4z J_z \Delta_{z-x} \int d^4w J_w \Delta_{w-x} \right)
\end{gathered}
\]

Considering 2 incoming particles and 1 intermediary particle, \\
\quad Recall the LSZ formula in momentum space (See Problem 10.4).  For this case,

\[
\langle f | i \rangle = i^3 (k_1^2 + m^2 )(k_2^2 + m^2 ) ( k_{in}^2 + m^2 ) \langle 0 | T\widetilde{\varphi}(-k_1) \widetilde{\varphi}(-k_2) \widetilde{ \varphi}(k_{in}) | 0 \rangle
\]

To obtain $ \langle 0 | T\widetilde{\varphi}(-k_1) \widetilde{\varphi}(-k_2) \widetilde{ \varphi}(k_{in}) | 0 \rangle$, we must apply functional derivatives $\frac{1}{i} \delta_{\widetilde{J}_{k_1}} \frac{1}{i} \delta_{\widetilde{J}_{k_2}} \frac{1}{i} \delta_{\widetilde{J}_{-k_{in}} }$ to our $O(\lambda)$ term above.  

Consider all the different ways to uniquely match functional derivatives to propagators.  

\begin{tabular}{ l r c c r c c r c c }
  a & \, & $k_{in}$ & $k_{in}$ & \, & 1 & 1 & \, & 2 & 2  \\ 
  b & \, & 1 & 2 & \, & $k_{in}$ & 2 & \, & 1 & $k_{in}$  \\ 
  c & \, & 2 & 1 & \, & 2 & $k_{in}$ & \, & $k_{in}$ & 1    \\
  & &  $(k_1)k_2$ & $k_2 k_1$  & & $-k_{in} k_2$  & $-(k_2)k_{in}$  & & $-k_1 k_{in} $ & $-(k_{in})(k_1)$ 
\end{tabular} \\

So we obtain 
\[
\lambda ( + 2 i ) ( 2) ( k_1 k_2 - k_{in} ( k_1 + k_2) ) = - \lambda 4 i (k_1^2 + k_2^2 + k_1 k_2) 
\]
where there's a factor of 2 for each unique $k_b k_c$.  

Now for the second term in the sum, \\
\phantom{Now} Recall the LSZ formula,

\[
\langle f | i \rangle = i^{n + n' } \int d^4x_1 e^{ik_1 x_1} (- \partial_1^2 + m_1^2 ) \dots \int d^4x_1' e^{-ik_1' x_1' } (-\partial_{1'}^2 + m_{1'}^2 ) \dots \times \langle 0 | T \varphi(x_1) \dots \varphi (x_1') \dots | 0 \rangle 
\]

In this case,
\[
\langle f | i \rangle = i^3 \int d^4x_1 e^{ik_1 x_1} ( - \partial_1^2 + m^2 ) \int d^4x_2 e^{ik_2 x_2 } ( - \partial_2^2 + m^2 ) \int d^4x' e^{-ik'x'} ( - \partial_{1'}^2 + m^2 ) \langle 0 | T\varphi(x_1) \varphi(x_2) \varphi(x') | 0 \rangle 
\]

To get $\langle 0 | T\varphi(x_1) \varphi(x_2) \varphi(x') | 0 \rangle $, we must apply $\frac{1}{i} \delta_{J_{x_1} } \frac{1}{i} \delta_{J_{x_2}} \frac{1}{i} \delta_{J_{x'}}$ to the 3 propagators.  

There are $3!$ ways to match functional derivative $\delta_{J_{x_i}}$ to propagator (and all 6 yield the same result)
\[
\Longrightarrow - i 3! \int d^4x_1 e^{ik_1 x_1} ( - \partial_1^2 + m^2 ) \int d^4x_2 e^{ik_2 x} ( - \partial_2^2 + m^2 ) \int d^4x' e^{-ik'x' } ( - \partial_{1'}^2 + m^2 ) \cdot m^2 \int d^4x \Delta_{x_1-x} \Delta_{x_2 - x} \Delta_{x' - x}
\]

Using $( - \partial_i^2  +m^2 ) \Delta(x_i - y) = \delta^{(4)}(x_i - y )$, 
\[
-i 6\lambda  m^2 \int d^4x e^{ i (k_1 + k_2 - k') x}  = -i \lambda 6 m^2 (2\pi)^4 \delta^{(4)}(k_1 + k_2 - k')
\]

So for the cubic vertex, for a single vertex, $\langle f | i \rangle$ has a factor of  
\[
\boxed{ -4 \lambda i ( k_1^2 + k_2^2  + k_1 k_2 + \frac{3}{2} m^2 ) }
\]

Before going on with the contribution due to cubic vertices to  $\varphi \varphi \to \varphi \varphi$ scattering, consider the quartic vertex.  

Similar to what was done for $\mathcal{L}_3$, 
\[
e^{ i \int d^4x \mathcal{L}_4 } = \exp{ \left[ -i 2 \lambda^2 \frac{1}{ (2\pi)^{12} } \int d^4k_a d^4k_b d^4k_c d^4k_d \delta^{(4)}(k_a + k_b +k_c + k_d ) (- k_c k_d) \widetilde{\varphi}_{k_a} \widetilde{\varphi}_{k_b} \widetilde{\varphi}_{k_c} \widetilde{\varphi}_{k_d} - i \frac{ \lambda^2}{2} m^2 \int d^4x \varphi^4 \right] }
\]


To first order in $\lambda^2$, 
\[
\lambda^2 \left[ \frac{ -2 i }{ (2\pi)^{12} } \int d^4k_a d^4k_b d^4k_c d^4k_d \delta^{(4)}(k_a + k_b +k_c + k_d) (- k_c k_d) \widetilde{\varphi}_{k_a} \widetilde{\varphi}_{k_b} \widetilde{\varphi}_{k_c}\widetilde{\varphi}_{k_d} -  \frac{i}{2} m^2 \int d^4x \varphi^4 \right] 
\]

\noindent Let this term act on 4 propagators.  \\
\phantom{Let } Then consider the LSZ formula, with 2 incoming particles and 2 outgoing particles.   \\

\noindent For the first term, a factor of 4 for each unique way to match $\frac{1}{i} \delta_{J_{k_i}}$, $i = 1,2,3,4$ with 4 propagators.  \\
For the second term, a factor of $4!$ for $4!$ ways to match $\frac{1}{i} \delta_{J_{x_i}}$, $i = 1,2,3,4$ with propagators.  

\[
\Longrightarrow -i \lambda^2 \left[ 8 ( - k_1 k_2 + k_1 k_3 + k_1 k_4 + k_2 k_3 + k_2 k_4 - k_3 k_4) + \frac{ 4! m^2 }{2} \right]
\]
Using momentum conservation, 
\[
\begin{aligned}
  & k_1( k_3 +k_4) = k_1(k_1 + k_2) = k_1^2 + k_1 k_2 \\ 
  & k_2(k_3 +k_4) = k_2(k_1 + k_2) = k_1 k_2 + k_2^2 
\end{aligned}
\]
So the quartic vertex contributes
\[
 \boxed{ -8i\lambda^2 ( k_1^2 + k_2^2 + k_1 k_2 - k_3 k_4 + \frac{3}{2} m^2 ) }
\]

\section{Cross sections and decay rates}


% ------------------------ 20100204 --------------------

\solutionhead{11.1} \begin{itemize}
\item[(a)] Given $\mathcal{L}_1 = g AB^2$.  If source terms are $J(x) A(x)$ and $K(x) B(x)$, \\

Consider first order term.  
\[
ig \int d^4x \frac{1}{i} \delta_{J(x)} \left( \frac{1}{i} \delta_{K(x)} \right)^2 
\]

Consider 2 $B$ propagators, 1 $A$ propagator.  

\[
\left( \frac{i}{2} \int d^4x d^4x' J(x) \Delta(x-x') J(x') \right) \frac{1}{2} \left( \frac{i}{2} \int d^4y d^4y' K(y) \Delta(y-y') K(y') \right)^2 
\]

Include $2! \times 2^3$ combinatoric factor.  

\[
ig \int d^4x \int d^4x' J_{x'} \Delta_{x'-x} \int d^4y K_y \Delta_{y-x} \int d^4z K_z \Delta_{z-x}
\]
\[
\Longrightarrow \langle f|i \rangle = i^3 \int d^4x_1 e^{ik_1 x_1} ( - \partial_1^2 + m_A^2 ) \int d^4x_2 e^{-ik_2 x_2} ( - \partial_2^2 + m_B^2 ) \int d^4x_3 e^{-ik_3 x_3} (- \partial_3^2 + m_B^2 ) \langle 0 | T \varphi(x_1) \varphi(x_2) \varphi(x_3) | 0 \rangle
\]

So apply $\frac{1}{i} \delta_{K_{x_2}} \frac{1}{i} \delta_{K_{x_3} } \frac{1}{i} \delta_{J_{x_1}}$.  
\[
\begin{gathered}
  \langle f | i \rangle = \int d^4x_1 e^{ik_1 x_1} (-\partial_1^2 + m_A^2) \int d^4x_2 e^{-ik_2x_2} (- \partial_2^2 +m_B^2 ) \int d^4x_3 e^{-ik_3 x_3} ( - \partial_3^2 + m_b^2 )   \cdot \\ 
  \cdot (2) ig \int d^4x \Delta_{x_1-x} \Delta_{x_2 -x} \Delta_{x_3 - x} = 2 ig (2\pi)^4 \delta^4(k_1 -k_2 -k_3)
\end{gathered}
\]
Now
\[
\langle f | i \rangle = (2\pi)^4 \delta^4(k_{in} - k_{out} ) i \mathcal{T}
\]
So 
\[
\mathcal{T} = 2g
\]
$dLIPS(k_1)$ is Lorentz invariant.  Use CM frame.  

Recall 
\[
dLIPS_{n'}(k) \equiv (2\pi)^4 \delta^4(k- \sum_{i=1}^{n'} k_i') \prod_{j=1}^{n'} \widetilde{dk}'_j
\]
where 
\[
\widetilde{dk} \equiv \frac{d^3k}{ (2\pi)^3 2 \omega }
\]

Recalling the steps in evaluating $dLIPS_2(k_1) = (2\pi)^4 \delta^4(k_1 - k_2 - k_3 ) \widetilde{dk}_2 \widetilde{dk}_3$ for 2 outgoing particles, then in the CM frame,
\[
dLIPS_2(k) = (2\pi)^4 \delta(E-E_1-E_2) \delta^3(\mathbf{k} - \mathbf{k}_1 \mathbf{k}_2 ) \frac{ d^3k_1}{ (2\pi)^3 2 E_1 } \frac{ d^3k_2 }{ (2\pi)^3 2E_2 } = \frac{1}{16 \pi^2 } \delta(E-E_1 - E_2) \delta^3(\mathbf{k} - \mathbf{k}_1  - \mathbf{k}_2 ) \frac{ d^3k_1 d^3k_2 }{ E_1 E_2} 
\]
For $\mathbf{k} = 0$ (CM frame),
\[
\xrightarrow{ \int d^3(\mathbf{k} - \mathbf{k}_1 - \mathbf{k}_2 ) d^3k_2} \frac{1}{16\pi^2} \delta(E- 2E_1) \frac{ d^3k_1}{ E_1 ( E-E_1) } = \frac{1}{16 \pi^2 } \delta(E-2E_1) \frac{ p_1^2 dp_1 d\Omega_{cm} }{ E_1 ( E-E_1) }
\]

Now $\frac{ \partial ( m_A - 2 \sqrt{ p_1^2 + m_B^2  } ) }{ \partial p_1 } = \frac{-2 p_1 }{ E_1 } $

Recalling $\int dx \delta(f(x) ) = \sum_i | f'(x_i) |^{-1}$ where $x_i$ s.t. $f(x_i) = 0$.  

\[
\Longrightarrow \frac{1}{16 \pi^2 } \frac{ p_1 d\Omega_{cm} }{ 2 (E - E_1) } = \frac{1}{16 \pi^2 } \frac{ \sqrt{ \frac{ m_A^2 }{4 } - m_B^2 } d\Omega_{cm } }{ 2 \left( \frac{m_A}{2} \right) } = \frac{1}{16 \pi^2 } \frac{ \sqrt{ \frac{ m_A^2 }{ 4} - m_B^2 } }{ m_A } d\Omega_{cm}
\]

Since by kinematics, $E = m_A$.  $E_1 + E_2 = 2E_1 = E$ so, $E_1 = \frac{m_A}{2}$.  

So
\[
\Gamma = \frac{1}{S} \int d\Gamma = \frac{1}{S} \int \frac{1}{2 E_1 } | \mathcal{T} |^2 dLIPS_{n'}(k_1) = \frac{1}{2} \int \frac{1}{ 2m_A} 4 g^2 \frac{1}{16 \pi^2 } \frac{ \sqrt{ \frac{ m_A^2 }{ 4 } - m_B^2 } }{ m_A} d\Omega_{cm}
\]
Now $\int d\Omega_{cm} = 4\pi$.  
\[
\boxed{ \Gamma = \frac{ g^2 }{ 8 \pi^2 } \frac{ \sqrt{ m_A^2 - 4 m_B^2 }}{ m_A^2 } }
\]

\item[(b)] Now consider
\[
Z_{\varphi}(K) = \exp{ \left[ \frac{i}{2} \int d^4x d^4x' K_x \rho_{x-x'} K_{x'} \right] }
\]
\[
Z_{\chi^{\dag}, 0}(J,J^{\dag}) = \exp{ \left[ \frac{i}{2} \int d^4x d^4x' J_x^{\dag} \Delta_{x-x'} J_{x'} \right] }
\]
Source term $J^{\dag} \chi + J \chi^{ \dag}$.  

Consider first order
\[
ig \int d^4x \frac{1}{i} \delta_{K_x} \frac{1}{i} \delta_{J_x} \frac{1}{i} \delta_{J_x^{\dag}}
\]

Consider 2 $\chi$ propagators, 1 $\varphi$ propagator.  

\[
\left[ \frac{i}{2} \int d^4y d^4y' K_y \rho_{y-y'} K_{y'} \right] \frac{1}{2} \left( \frac{i}{2} \int d^4z d^4z' J_z^{\dag} \Delta_{z-z'} J_{z'} \right) \left( \frac{i}{2} \int d^4x d^4w' J_x^{\dag} \Delta_{w-w'} J_{w'} \right)
\]
include $2! \cdot 2 = 4$ counting factor.  Another counting factor of $2^2$, because eventually the outgoing particles are distinguishable.  

\[
ig \int d^4x \int d^4y K_y \rho_{y - x} \int d^4z J_z^{\dag} \Delta_{z-x} \int d^4w \Delta_{x-w} J_w
\]

Recall LSZ formula, for complex field $\chi$ and scalar field $\varphi$.  
\[
\langle f | i \rangle = i^3 \int d^4x_1 e^{ik_1 x_1} (- \partial_1^2 + m_1^2 ) \int d^4x_1' e^{-ik_1' x_1' } ( - \partial_{1'}^2 + m_2^2 ) \int d^4x_2' e^{ i k_2' x_2' } ( - \partial_{2'}^2 + m_2^2  ) \cdot \langle 0 | T \varphi(x_1) \chi(x_1') \chi^{\dag}(x_2') | 0 \rangle 
\]
so do $\frac{1}{i} \delta_{K_{x_1}} \frac{1}{i} \delta_{J^{\dag}_{x_1'}} \frac{1}{i} \delta_{J_{x_2'} }$.  
\[
\begin{gathered}
  \langle f | i \rangle = \int d^4x_1 e^{ik_1 x_1 }( - \partial_1^2 + m_1^2 ) \int d^4x_1' e^{-ik_1' x_1' } ( - \partial_{1'}^2 + m_2^2 ) \int d^4x_2' e^{ik_2' x_2'} ( - \partial_{2'}^2 + m_2^2 ) \cdot ig \int d^4x \rho_{x_1-  x} \Delta_{x_1' - x} \Delta_{x-x_2'} = \\ 
  = ig (2\pi)^4 \delta^4(k_1 - k_1' + k_2')
\end{gathered}
\]

\[
\mathcal{T} = ig 
\]

Since the particles in the final state are distinguishable, there's no symmetry factor of 2 to divide by.

\[
\Longrightarrow \boxed{ \Gamma = \frac{g^2 }{ 16\pi^2} \frac{ \sqrt{ m_{\varphi}^2 - 4 m_{\chi}^2 } }{ m_{\varphi}^2 } }
\]

\end{itemize}

\solutionhead{11.2}
\begin{itemize}
\item[(a)] If, for 4-vectors,  \\
$(\omega, \mathbf{p})$ is the incident photon, $\omega^2 = \mathbf{p}^2$,  \\
$(m,0)$ is the initial electron at rest \\ 
$(\omega_2, \mathbf{q})$ is the final photon, $\omega_2^2 = \mathbf{q}^2$, \\
$(E, \mathbf{q}_e)$ is the final electron, $E^2 - \mathbf{q}_e^2 = m^2$.  

Then momentum-energy conservation says
\[
\begin{gathered}
  \mathbf{p} = \mathbf{q} + \mathbf{q}_e \\ 
  \omega + m = \omega_2 + E
\end{gathered}
\] 
or 
\[
( \omega, \mathbf{p} ) + ( m, 0 ) = ( \omega_2, \mathbf{q}) + (E, \mathbf{q}_e) = (\omega + m, \mathbf{p} )
\]

Recall that 
\[
\begin{aligned}
  & s = - (k_1 + k_2)^2 \\ 
  & t = - (k_1 - k_1')^2 \\ 
  & u = - (k_1 - k_2')^2 
\end{aligned}
\]
then for our case
\[
\begin{aligned}
  & s = ( \omega + m)^2 - \mathbf{p}^2 = 2 \omega m + m^2 = m ( m + 2 \omega )  \\
  & t = - ((\omega - \omega_2, \mathbf{p} - \mathbf{q}))^2 = \omega^2 - 2 \omega \omega_2 + \omega_2^2 - ( \mathbf{p}^2 - 2 \mathbf{p} \cdot \mathbf{q} + \mathbf{q}^2 ) = -2 \omega \omega_2 ( 1 - \cos{\theta}) \\ 
  & u = - ( ( \omega - E, \mathbf{p} - \mathbf{q}_e ) )^2 = ( \omega_2 - m)^2 - \omega_2^2 = m (m - 2 \omega_2)
\end{aligned}
\]
\item[(b)] Starting from $\mathbf{p} - \mathbf{q} = \mathbf{q}_e$, 
\[
\begin{gathered}
  (\mathbf{p} - \mathbf{q})^2 = \omega^2 + \omega_2^2 - 2 \mathbf{p} \cdot \mathbf{q} = E^2 - m^2 = \mathbf{q}_e^2 = ( \omega + m - \omega_2)^2 - m^2 = \\ 
  = \omega^2 + m^2 + \omega_2^2 + 2\omega m - 2 \omega \omega_2 - 2 m \omega_2 - m^2 = \omega^2 + \omega_2^2 + 2 \omega m - 2 \omega \omega_2 - 2 m \omega_2
\end{gathered}
\]
\[
\begin{gathered}
  \Longrightarrow -2 \mathbf{p} \cdot \mathbf{q} = 2 m ( \omega - \omega_2) - 2 \omega \omega_2 \\ 
   -\omega \omega_2 \cos{\theta} = m ( \omega - \omega_2) - \omega \omega_2 \\ 
   \Longrightarrow \cos{\theta} = 1 - \frac{ m ( \omega - \omega_2) }{ \omega \omega_2 }
\end{gathered}
\]
\item[(c)] Consider first the $|\mathcal{T}|^2$ term.  
\[
\begin{aligned}
  \frac{ m^4 + m^2 ( 3s + u ) - su }{ (m^2 - s)^2 } & = \frac{ m^4 + m^2 ( 3m^2 + 6 m \omega + m^2 - 2 m \omega_2 ) - m^2 ( m+ 2\omega) ( m-2\omega_2)  }{ (m^2 - m^2 - 2 \omega m )^2 } = \\
  & = \frac{ m^4 + 3m^4 + 6 m^3 \omega + m^4 - 2 m^3 \omega_2 - m^2 ( m^2 - 2 \omega_2 m + 2 \omega m - 4 \omega \omega_2 ) }{ 4 \omega^2 m^2 }  = \\ 
  & = \frac{ 4 m^4 + 4 m^3 \omega + 4 m^2 \omega \omega_2 }{ 4 \omega^2 m^2 } = \frac{ m^2 + m\omega + \omega \omega_2 }{ \omega^2 } 
\end{aligned}
\]

\[
\begin{aligned}
  \frac{ m^4 + m^2 ( 3 u + s ) - su }{ (m^2 - u)^2 } & = \frac{ m^4 + m^2 ( 3 ( m^2 - 2 \omega_2 m ) + m^2 + 2 m \omega )  - m^2 ( m^2 - 2 \omega_2 m + 2 \omega m - 4 \omega \omega_2 ) }{ ( m^2 - m^2 + 2 m \omega_2)^2 } = \\ 
  & = \frac{ m^4 + m^2 ( 4 m^2 - 6 \omega_2 m + 2 m \omega) - m^4 + 2 \omega_2 m^3 - 2 \omega m^3 + 4 m \omega \omega_2  }{ 4 m^2 \omega_2^2 } = \frac{ 4 m^4 + -4 \omega_2 m^3 + 4 m^2 \omega \omega_2 }{ 4 m^2 \omega_2^2 } = \\
  & = \frac{ m^2 - \omega_2 m + \omega \omega_2 }{ \omega_2^2 } 
\end{aligned}
\]
\end{itemize}

\[
\begin{aligned}
  \frac{ 2 m^2 ( s + u + 2 m^2 ) }{ (m^2 - s) (m^2 - u ) } = \frac{ 2 m^2 ( m^2 + 2 m \omega + m^2-  2 m \omega_2 + 2m^2 ) }{ ( -2 \omega m ) (2 m \omega_2 ) } = \frac{ 2 m^2 ( 4 m^2 + 2 m ( \omega- \omega_2 )  ) }{ -2^2 m^2 \omega \omega_2  }  = \frac{ m (2m + (\omega - \omega_2)) }{ - \omega \omega_2  } 
\end{aligned}
\]
Summing those 3 terms up above,
\[
\begin{gathered}
  ( m^2 \omega_262 + m\omega \omega_2^2 + \omega \omega_2^3 + m^2 \omega^2 - \omega_2 m \omega^2 + \omega^3 \omega_2 + - 2 m^2 \omega \omega_2 - m \omega^2 \omega_2 + m \omega_2^2 \omega )/ ( \omega \omega_2)^2  = \\
  = ( m^2 ( \omega_2 - \omega)^2 + \omega \omega_2 (\omega_2^2 + \omega^2 ) + 2 m \omega \omega_2 ( \omega_2 - \omega) )/ ( \omega \omega_2)^2 
\end{gathered}
\]

Now 
\[
-(1 - \cos^2{\theta} ) = \frac{ -2 m ( \omega - \omega_2) }{ \omega \omega_2} + \frac{ m^2 ( \omega - \omega_2)^2 }{ ( \omega \omega_2)^2 } = -\sin^2{\theta}
\]

Then we have
\[
|\mathcal{T}|^2 = 32 \pi^2 \alpha^2 \left( \frac{ \omega }{ \omega_2 } + \frac{ \omega_2}{ \omega } - \sin^2{\theta_{FT} } \right)
\]

Now in the lab frame, 
\[
d\sigma = \frac{ 1 }{ 4 | \mathbf{k}_1 | m_2 } | \mathcal{T} |^2 dLIPS_2(k_1 + k_2) = \frac{1}{ 4 \omega m } | \mathcal{T} |^2 dLIPS_2(k_1 + k_2) 
\]

Now, evaluated in the lab frame,
\[
\begin{gathered}
  dLIPS_2(k_1 + k_2) = (2\pi)^4 \delta^4(k_1 + k_2 -k_3 - k_4 ) \frac{ d^3 q }{ (2\pi)^3 2 \omega_2 } \frac{ d^3 q_e }{ (2\pi)^3 2 E } = \\ 
  = (2\pi)^4 \delta(\omega_2  + E - \omega - m ) \delta^3(\mathbf{p} - \mathbf{q} - \mathbf{q}_e ) \frac{ d^3q }{ (2\pi)^3 2 \omega_2 } \frac{ d^3q_e}{ (2\pi)^3 2 E } = \frac{1}{16 \pi^2 } \delta(E + \omega_2 - \omega - m ) \frac{ d^3q }{ \omega_2 E }
\end{gathered}
\]
With $d^3q = \omega_2^2 d\omega_2 d\Omega$, 
\[
dLIPS_2 = \frac{1}{16 \pi^2 } \delta( E + \omega_2 - m - \omega )  \frac{ \omega_2 d\omega_2 d\Omega }{ E } 
\]
Now
\[
\begin{gathered}
  \frac{ \partial ( -\omega - m + \omega_2 + E ) }{ \partial \omega_2  } = 1 + \left( \frac{1}{ 2 E } \left( 2 \omega_2 - 2 \omega \cos{\theta} \right) \right) = 1 + \frac{ ( \omega_2 - \omega \cos{\theta}  ) }{ E } = \\ 
   = \frac{  ( E + (\omega_2 - \omega \cos{\theta} )  ) }{ E } 
\end{gathered}
\]
\[
\begin{gathered}
  \Longrightarrow dLIPS_2(k_1 + k_2) = \frac{ 1 }{16 \pi^2 } \frac{ E}{ ( E + ( \omega_2 - \omega \cos{\theta} ) ) } \frac{ \omega_2 d\Omega }{ E } = \frac{ 1 }{ 16 \pi^2 } \frac{ \omega_2 d\Omega }{ ( \omega ( 1 - \cos{\theta} ) + m ) } = \\ 
  = \frac{1 }{ 16 \pi^2 } \frac{ \omega_2 d\Omega }{ \left( \omega \left( \frac{ m ( \omega - \omega_2  ) }{ \omega \omega_2 } \right) + m \right) } = \frac{ 1 }{ 16 \pi^2 } \frac{ \omega_2 d\Omega}{  \frac{ m \omega - m \omega_2 + m \omega_2 }{ \omega_2 } } = \frac{ 1 }{ 16 \pi^2 } \frac{ \omega_2^2 d\Omega }{ m \omega}
\end{gathered}
\]
\[
\Longrightarrow \boxed{ \frac{ d \sigma }{ d\Omega} = \frac{ \alpha^2}{2m^2 } \frac{ \omega_2^2 }{ \omega^2 }  \left( \frac{ \omega }{ \omega_2 } + \frac{ \omega_2}{ \omega } - \sin^2{\theta_{FT} } \right)}
\]


% ------------------------ 20100204 -------------end----




\problemhead{11.3} Consider the process of \emph{muon decay}, $\mu^- \to e^- \overline{\nu}_e \nu_{\mu}$.  In section 88, we will compute $|\mathcal{T}|^2$ for this process (summed over the possible spin states of the decay products, and averaged over the possible spin states of the initial muon), with the result
\[
|\mathcal{T}|^2 = 64 G_F^2 (k_1 \cdot k_2')(k_1' \cdot k_3'), \quad \quad \quad (11.52)
\]
where $G_F$ is the \emph{Fermi constant}, $k_1$ is the four-momentum of the muon, and $k_{1,2,3}'$ are the four-momenta of the $\nu_e$, $\nu_{\mu}$, and $e^-$, respectively.  In the rest frame of the muon, its decay rate is therefore
\[
\Gamma = \frac{32 G_F^2}{m} \int (k_1 \cdot k_2')(k_1'\cdot k_3') dLIPS_3(k_1), \quad \quad \quad (11.53)
\]
where $k_1 = (m,\mathbf{0})$, and $m$ is the muon mass.  The neutrinos are massless, and the electron mass is 200 times less than the muon mass, so we can take the electron to be massless as well.  To evaluate $\Gamma$, we perform the following analysis.  
\begin{itemize}
\item[(a)] Show that 
\[
\Gamma = \frac{32 G_F^2 }{m} \int \widetilde{dk}'_3 k_{1\mu} k'_{3\nu} \int k^{'\mu}_2 k_1^{' \nu} dLIPS_2(k_1-k_3') \quad \quad \quad (11.54)
\]
\item[(b)] Use Lorentz invariance to argue that, for $m_{1'} = m_{2'} = 0$, 
\[
\int k_1^{'\mu} k_2^{'\nu} dLIPS_2(k) = Ak^2 g^{\mu \nu} + Bk^{\mu} k^{\nu}, \quad \quad \quad (11.55)
\]
where $A$ and $B$ are numerical constants.
\item[(c)] Show that, for $m_{1'} = m_{2'} =0$,
\[
\int dLIPS_2(k) = \frac{1}{8\pi} . \quad \quad \quad  (11.56)
\]
\item[(d)] By contracting both sides of eq. (11.55) with $g_{\mu \nu}$ and with $k_{\mu} k_{\nu}$, and using eq. (11.56), evaluate $A$ and $B$.  
\item[(e)] Use the results of parts (b) and (d) in eq. (11.54).  Set $k_1 = (m,\mathbf{0})$, and compute $d\Gamma/dE_e$; here $E_e \equiv E_3'$ is the energy of the electron.  Note that the maximum value of $E_e$ is reached when the electron is emitted in one direction, and the two neutrinos in the opposite direction; what is the maximum value?  
\item[(f)] Perform the integral over $E_e$ to obtain the muon decay rate $\Gamma$.  
\item[(g)] The measured liftime of the muon is $2.197 \times 10^{-6}$ s.  The muon mass is $105.66 \, MeV$.  Determine the value of $G_F$ in $GeV^{-2}$.  (Your answer is too low by about $0.2\%$, due to loop corrections to the decay rate.)
\item[(h)] Define the \emph{energy spectrum of the electron} $P(E_e) \equiv \Gamma^{-1} d\Gamma/dE_e$.  Note that $P(E_e) dE_e$ is the probability for the electron to be emitted with energy between $E_e$ and $E_e + dE_e$.  Draw a graph of $P(E_e)$ versus $E_e/m_{\mu}$.  
\end{itemize}

\solutionhead{11.3} 
\begin{itemize}
\item[(a)] Recall $\widetilde{dk} \equiv \frac{d^3 k}{ (2\pi)^3 2k^0}$.  
\[
\begin{aligned}
  dLIPS_3(k_1) & = (2\pi)^4 \delta^4(k_1 - (k_1' + k_2' + k_3' ) ) \prod_{j=1}^3 \widetilde{dk}_j' = \\
   & = (2\pi)^4 \delta^4(k_1 - k_3' - (k_1' + k_2'))\prod_{k=1}^2 \widetilde{dk}_k' \widetilde{dk}_3'
\end{aligned}
\]
Now
\[
dLIPS_2(k_1-k_3') = (2\pi)^4 \delta^4(k_1 -k_3' - (k_1' + k_2') ) \prod_{j=1}^2 \widetilde{dk}_j'
\]
\[
\begin{aligned}
  \Gamma & = \frac{32 G_F^2}{m} \int (k_1 \cdot k_2')(k_1' \cdot k_3') dLIPS_3(k_1) = \frac{32 G_F^2}{m} \int k_{1\mu} k_{3\nu}' k_{2'}^{\mu} k_{1'}^{\nu} dLIPS_2(k_1 -k_3')\widetilde{dk}_{3'} = \\
  & = \frac{32 G_F^2}{m} \left( \int \widetilde{dk}_{3'} k_{1\mu} k_{3\nu}' \right) \int k_{2'}^{\mu} k_{1'}^{\nu} dLIPS_2(k_1 - k_3')
\end{aligned}
\]
\item[(b)] $dLIPS_2(k)$ is Lorentz invariant.  \\
There are 2 upper indices $\mu,\nu$.  \\
All Lorentz-invariant objects with 2 upper indices $\mu,\nu$ must be a linear combination of $k_{1'}^2 g^{\mu \nu}$, $k_{2'}^2 g^{\mu \nu}$, $k^2 g^{\mu \nu}$ and $k^{\mu}k^{\nu}$.  Note $k^2$, $k_{1'}^2$, $k_{2'}^2$, $g^{\mu \nu}$ are Lorentz invariant.  Note $k_{1'}^2 = k_{2'}^2 =0$.  

Thus
\[
\int k_{1'}^{\mu} k_{2'}^{\nu} dLIPS_2(k) = Ak^2 g^{\mu \nu} + Bk^{\mu} k^{\nu} \quad \quad \quad (11.55)
\]
\item[(c)] Note that $k_1' +k_2' = k_1 - k_3'$ by momentum conservation and that  \\
$(k_1' + k_2')^2 = 2k_1'k_2' = -2(E_{1'} E_{2'})(1- \cos{\theta})$ \\ 
$(k_1 - k_3')^2 = -m^2$

Then of course we can always boost to a frame where $\mathbf{k} = \mathbf{k}_1 - \mathbf{k}'_3 = 0$.  Then $E_{1'} = E_{2'}$.  


It's \textbf{important} to see that we can use (11.30) from Srednicki (as Arthur Lipstein points out in his homework solutions):
\[
dLIPS_2(k) = \frac{ |\mathbf{k}_1|}{ 16 \pi^2 \sqrt{s}} d\Omega_{CM} \quad \quad \quad (11.30)
\]
For this frame, where $\mathbf{k}_{1'} + \mathbf{k}_{2'} =0$,
\[
\begin{aligned}
  s & = -(k_1 + k_2)^2 = - (k_{1'} + k_{2'})^2 = -(2k_{1'} k_{2'}) = (-2)(-E_{1'} E_{2'} - E_{1'} E_{2'} ) = \\
  & = 4E_{1'} E_{2'}
\end{aligned}
\]
and again, note that $E_{1'} = E_{2'}$ in this frame (since, again, $m_{1'} = m_{2'} =0$.  
\[
\int dLIPS_2(k) = \int \frac{ E_{1'}}{ 16 \pi^2 2 \sqrt{ E_{1'} E_{2'}} } d\Omega_{CM} = \boxed{ \frac{1}{8\pi}}
\]
\item[(d)] 
Note that
\[
\begin{aligned}
  k & = k_1' + k_2' \\ 
  k^2 & = k_{1'}^2 + k_{2'}^2 + 2k_{1'}k_{2'} = 2k_{1'} k_{2'}
\end{aligned}
\]

Applying $g_{\mu \nu}$ onto (11.55), 
\[
\begin{gathered}
  g_{\mu \nu} \int k_{1'}^{\mu} k_{2'}^{\nu} dLIPS_2(k) = Ak^2 ( 4) + B g_{\mu \nu} k^{\mu} k^{\nu} = \\
= \int k_{1'\nu} k_{2'}^{\nu} dLIPS_2(k) = 4Ak^2 + Bk^2 = k^2 ( 4A + B) = \frac{k^2}{2} \int dLIPS_2(k) = k^2 ( 4A + B) \\ \text{ or } \frac{1}{16\pi} = 4A + B
\end{gathered}
\]

Note that 
\[
\begin{aligned}
  k_1' k & = k_{1'}^2 + k_1' k_2' \\ 
  kk_2' & = k_1'k_2' + k_{2'}^2 
\end{aligned}
\]

Applying $k_{\mu} k_{\nu}$ on both sides to (11.55), 
\[
\begin{gathered}
  \int k_{\mu} k_{1'}^{\mu} k_{\nu} k_{2'}^{\nu} dLIPS_2(k) = Ak^2 k^2 + B(k^2)^2 \\ 
  \Longrightarrow \int (k_1' k_2')(k_1' k_2') dLIPS_2(k) = \int \frac{(k^2)^2}{4} dLIPS_2(k) = (A+B)(k^2)^2 \\ 
  \text{ or } \frac{1}{4} \frac{1}{8\pi} = \frac{1}{32 \pi} = A+B
\end{gathered}
\]
\[
\Longrightarrow \boxed{ \begin{aligned} A & = \frac{1}{96 \pi}  \\ B & = \frac{1}{48 \pi} \end{aligned} }
\]
\item[(e)] \[
\begin{aligned}
  \int k_{1'}^{\mu} k_{2'}^{\nu} dLIPS_2(k_1-k_2') & = \frac{1}{96 \pi} ((k_1-k_3')^2 g^{\mu \nu} + 2(k_1-k_3')^{\mu}(k_1-k_3')^{\nu} ) = \\
& =   \frac{1}{96 \pi} ((-m^2 - 2k_1k_3')g^{\mu \nu} + 2(k_1 - k_3')^{\mu}(k_1-k_3')^{\nu})
\end{aligned}
\]
Recall $\widetilde{dk}_3' = \frac{ d^3 k_{3'}}{ (2\pi)^3 2 E_{3'}}$

\[
\begin{aligned}
  \xrightarrow{ k_{1\mu}k_{3\nu}' } & \frac{1}{96 \pi} ((-m^2 - 2k_1 k_{3'}) k_{1\mu} k_{3'}^{\mu} + 2(-m^2 -k_1 k_{3'} )(k_1 k_{3'} - 0 )) = \\
  & = \frac{1}{96 \pi} ( - m^2 k_1 k_{3'} - 2(k_1 k_{3'} )^2 + -2m^2 k_1 k_{3'} - 2(k_1 k_{3'})^2 ) =  \\
  & = \frac{-k_1 k_{3'}}{ 96 \pi} (3m^2 + 4(k_1 k_{3'}) )
\end{aligned}
\]
Now $k_1 k_{3'} = -m E_e$.  

\[
\Longrightarrow \frac{ mE_e}{96 \pi} (3m^2 + - 4m E_e) = \frac{ m^3 E_e}{32 \pi} \left( 1 - \frac{4}{3} \frac{E_e}{m} \right)
\]
\[
\begin{aligned}
  \Gamma & = \frac{32 G_F^2}{m} \int \frac{ E_{3'}^2 dE_{3'} d\Omega_{3'} }{ (2\pi)^3 2 E_{3'}} \frac{ m^3 E_e}{32 \pi} \left( 1 - \frac{4}{3} \frac{E_e}{m} \right) = \\ 
  & = \frac{ G_F^2 m^2 }{(2\pi)^2 \pi } \int E_e^2 \left( 1 - \frac{4}{3} \frac{E_e}{m} \right) dE_e = \frac{ G_F^2 m^2 }{ \pi^3} \int E_e^2 \left( \frac{1}{4} - \frac{E_e}{3m } \right) dE_e
\end{aligned}
\]
\[
\boxed{ \frac{d\Gamma}{dE_e}   = \frac{ G_F^2 m^2 }{ \pi^3} E_e^2 \left( \frac{1}{4} - \frac{E_e}{3m } \right) }
\]
\item[(f)] \[
\begin{aligned}
  \Gamma & = \frac{G_F^2 m^2 }{ \pi^3} \left( \frac{1}{4} \frac{1}{3} \left( \frac{m}{2} \right)^3 - \frac{1}{3m} \frac{1}{4} \left( \frac{m}{2} \right)^4 \right) = \frac{ G_F^2 m^2 }{ \pi^3 12} \left( \frac{1}{8} - \frac{1}{16} \right)  = \\
  & = \boxed{ \frac{ G_F^2 m^5}{ 192 \pi^3 } }
\end{aligned}
\]
\item[(g)] $\Gamma = \frac{1}{\tau}$.  \\
$G_F^2 = \frac{192 \pi^3}{m^5} \frac{1}{\tau} = \frac{ 192 \pi^3}{m^5} \frac{\hbar}{\tau}$
\[
G_F = \sqrt{ \frac{ 192 \pi^3 (6.582 \times 10^{-22} \, MeV\cdot s ) }{ (105.66 \, MeV)^5 (2.197 \times 10^{-6} \, s ) } } = \frac{ 1.1637 \times 10^{-11}}{ (MeV)^2} \left( \frac{ 10^3 \, MeV}{ 1 \, GeV} \right)^2 = \boxed{ 1.1637 \times 10^{-5} \, (GeV)^{-2} }
\]
vs. $1.16637 \times 10^{-5} \, (GeV)^{-2}$ from physics.nist.gov.  $0.229 \%$ too low.  
\item[(h)] \quad \\
$P(E_e) \equiv \Gamma^{-1} \frac{d\Gamma}{dE_e}$ \\
$\begin{aligned} P(E_e) & = \frac{ 192 \pi^3}{ G_F^2 m^5} \frac{G_F^2 m^2 }{\pi^3} E_e^2 \left( \frac{1}{4} - \frac{E_e}{3m} \right) = \frac{192}{m^3} E_e^2 \left( \frac{1}{4} - \frac{E_e}{3m} \right) = \\
  & = \frac{192}{m} \left( \frac{E_e}{m} \right)^2 \left( \frac{1}{4} - \frac{E_e}{3m} \right) = 1.8171 \, (MeV)^{-1} \left( \frac{E_e}{m} \right)^2 \left( \frac{1}{4} - \frac{E_e}{3m} \right)
\end{aligned}$


with $\frac{E_e}{m}$ ranging from $0$ to $\frac{1}{2}$, since maximum $E_e$ is $\frac{m}{2}$.  See Figure (\ref{Fi:1103h}).  

\begin{figure}
  \scalebox{.70}{\includegraphics{Fig_Srednicki_1103h}}
  \caption{$P(E_e) \, (MeV)^{-1}$ vs. $\frac{E_e}{m}$ from $0$ to $\frac{1}{2}$ (maximum $E_e$ physically allowed)}\label{Fi:1103h}
\end{figure}
\end{itemize}






\section{Dimensional analysis with $\hbar = c = 1$ }


\problemhead{12.1} Express $\hbar c$ in GeV fm, where $1 \, fm = 1 \, Fermi = 10^{-13} \, cm$.  

\solutionhead{12.1} 
\[
\begin{aligned}
  \hbar c & = (6.582 \times 10^{-22} \, MeV \cdot s)(3\times 10^8 \, m/s)\left( \frac{ 10^{15} \, fm }{ 1 \, m } \right)\left( \frac{ 1 \, GeV }{ 10^3 \, MeV} \right) = \\
  & = 19.746 \times 10^{-2} \, GeV \cdot fm = 1.9746 \times 10^{-2} \, GeV \cdot fm
\end{aligned}
\]

\problemhead{12.2} Express the masses of the proton, neutron, pion, electron, muon, and tau in GeV.  

\solutionhead{12.2} proton $0.9383 \, GeV$ \\ 
neutron $0.9396 \, GeV$ \\ 
pion $\pi^{\pm}$: $0.13957 \, GeV$ \\ 
\phantom{pion }$\pi^0$: $0.13498$  \\
electron: $0.000511 \, GeV$ \\ 
muon: $0.1057 \, GeV$ \\
tau: $0.001777 \, GeV$

\section{ The Lehmann-K\"{a}ll\'{e}n form of the exact propagator}

\solutionhead{13.1} Given $\mathcal{L} = \frac{-1}{2} Z_{\varphi} \partial^{\mu} \varphi \partial_{\mu} \varphi - \frac{1}{2} Z_m m^2 \varphi^2 - \mathcal{L}_1(\varphi)$, 

Recall the canonical commutation relation, 
\[
[ \varphi(\mathbf{x}, t), \Pi(\mathbf{x}', t) ] = i \delta^3(\mathbf{x} - \mathbf{x}')
\]

\textbf{Note carefully that the time is the same for each operator}, above.  

Then note that 
\[
\left. [ \varphi(x), \Pi(y)] \right|_{x^0 = y^0} =  [ \varphi(\mathbf{x},t), \Pi,(\mathbf{x}',t) ]
\]
which is a \textbf{big math trick}, here.  

Recall this definition,
\[
\Pi(x) = \frac{  \partial \mathcal{L}}{ \partial \dot{\varphi(x) }}
\]
Since $\mathcal{L}_1(\varphi)$ is not a function of its derivatives
\[
\frac{ \partial \mathcal{L}}{ \partial \dot{\varphi}(x)} = Z_{\varphi} \dot{\varphi}(x) = \Pi(x)
\]

\[
\Longrightarrow  \left. [ \varphi(x) , \Pi(y) ] \right|_{x^0 = y^0 } = \left. [ \varphi(x), Z_{\varphi} \dot{\varphi}(y) ] \right|_{x^0 = y^0 } = i \delta^3(\mathbf{x} -  \mathbf{y} )
\]
or  
\[
\left. [ \varphi(x), \dot{\varphi}(y) ] \right|_{x^0 = y^0 } = i Z_{\varphi}^{-1} \delta^3(\mathbf{x}- \mathbf{y} )
\]

Now
\[
\langle 0 | \varphi(x) \varphi(y) | 0 \rangle = \int \widetilde{dk} e^{ik(x-y) } + \int_{4m^2}^{\infty} ds \rho(s) \int \widetilde{dk} e^{ik(x-y)}  \quad \quad \quad \, (13.12)
\]
where $\widetilde{dk} \equiv \frac{ d^{d-1} k }{ (2\pi)^{d-1} 2 \omega } \quad \quad \, (13.5)$, and where $k^0 = \omega = \sqrt{ \mathbf{k}^2 + m^2 }$ for the first term, above, and $k^0 = \sqrt{ \mathbf{k}^2 + M^2}$ for the second term, where $M$ is a \emph{parameter}.  

\textbf{ Note} how $k^0$ is \textbf{different} for each of these terms.  

\[
\begin{aligned}
  \left. \langle 0 | \varphi(x) \dot{\varphi}(y) |0  \rangle \right|_{x^0 = y^0} & = \left. \langle 0 | \varphi( \mathbf{x}, x^0 ) \partial_{y^0} \varphi(\mathbf{y}, y^0 )|0 \rangle \right|_{x^0 = y^0 } =  \\ 
  & = \left. \partial_{y^0} \langle  0 | \varphi(\mathbf{x}, x^0 ) \varphi(\mathbf{y}, y^0) | 0 \rangle \right|_{x^0 = y^0 }
\end{aligned}
\]

Then
\[
\begin{gathered}
  \left.  \partial_{y^0} \left( \int \widetilde{dk} e^{ik(x-y)} + \int_{4m^2 }^{\infty} ds \rho(s) \int \widetilde{dk} e^{ik(x-y) } \right) \right|_{x^0 = y^0 } = \left. \left( \int \widetilde{dk} ik^0 e^{i k (x-y) } + \int_{4m^2 }^{\infty} ds \rho(s) \int \widetilde{dk} ik^0 e^{ik(x- y) } \right) \right|_{x^0 = y^0 } = \\ 
  = \left. \left( \int \frac{d^dk }{ (2\pi)^{d-1}2\omega } ik^0 e^{ik(x-y) } + \int_{4m^2}^{\infty} ds \rho(s) \int \frac{d^dk }{ (2\pi)^{d-1}2\omega} ik^0 e^{ik(x-y) }  \right) \right|_{x^0 = y^0 } = \\
   = \frac{i}{2} \left( \int \frac{d^{d-1}k }{ (2\pi)^{d-1} } e^{i \mathbf{k}\cdot (\mathbf{x} - \mathbf{y} ) } + \int_{4m^2 }^{\infty} ds \rho(s) \int \frac{d^{d-1}k}{ (2\pi)^{d-1}} e^{i \mathbf{k} \cdot ( \mathbf{x} - \mathbf{y}) } \right) = \\
  = \frac{i}{2} \left( 1 + \int_{4m^2}^{\infty} ds \rho(s) \right) \delta^{d-1}(\mathbf{x} - \mathbf{y} )
\end{gathered}
\]

Now
\[
\left. \langle 0 | \dot{\varphi}(x) \varphi(y) | 0 \rangle  \right|_{x^0 = y^0 } = \frac{-i}{2} \left( 1 + \int_{4m^2}^{\infty} ds \rho(s) \right) \delta^{d-1}(\mathbf{x}- \mathbf{y})
\]
So 
\[
\left. \langle 0 | \dot{\varphi}_y \varphi_x | 0\rangle \right|_{x^0 = y^0} = \frac{-i}{2} \left( 1+ \int_{4m^2}^{\infty} ds \rho(s) \right) \delta^{d-1}(\mathbf{x} -  \mathbf{y} ) 
\]

\[
\left. \langle 0 | [ \varphi(x), \dot{\varphi}(y) ]  | 0 \rangle \right|_{x^0 = y^0 } = i Z_{\varphi}^{-1} \delta^{d-1}(\mathbf{x} - \mathbf{y} )  = i \left( 1 + \int_{4m^2 }^{\infty} ds \rho(s) \right) \delta^{d-1}(\mathbf{x} - \mathbf{y} )
\]

\[
\Longrightarrow \boxed{ Z_{\varphi}  = \frac{1}{ 1 + \int_{4m^2 }^{\infty} ds \rho(s) } }
\]

\section{Loop corrections to the propagator}

\solutionhead{14.1} 

Want: 
\[
\frac{1}{ A_1^{\alpha_1} \dots A_n^{\alpha_n} } = \frac{ \Gamma( \sum_i \alpha_i ) }{ \prod_i \Gamma( \alpha_i ) } \frac{1}{ ( n-1)! } \int dF_n \frac{ \prod_i x_i^{\alpha_i - 1 } }{ \left( \sum_i x_i A_i \right)^{ \sum_i \alpha_i } }  \quad \quad \quad \, (14.49) 
\]

Let's convert this to make it easier.  

\[
\frac{ \prod_i \Gamma(\alpha_i) }{ \prod_i A_i^{\alpha_i } }  = \prod_i \frac{ \Gamma( \alpha_i ) }{ A_i^{\alpha_i } } = \Gamma( \sum_i \alpha_i ) \prod_i^n \int_0^1 dx_i \delta(\sum_j x_j - 1 ) \frac{ x_i^{ \alpha_i -1 } }{ \left(  \sum_j x_j A_j \right)^{x_i } } 
\]
where $\int dF_n = (n-1)! \int_0^1 dx_1 \dots dx_n \delta(x_1 + \dots + x_n - 1 )$ and $\int dF_n 1 =$.  

Start with, as the hint suggests, 
\[
\frac{ \Gamma(\alpha) }{ A^{\alpha} } = \int_0^{\infty} dt t^{\alpha- 1 } e^{-At}
\]
\[
 \prod_i \frac{ \Gamma(\alpha_i ) }{ A^{\alpha_i} } =\prod_i \int_0^{\infty} dt_i t_i^{\alpha_i -1 } e^{ - A_i t_i }
\]

Multiply the right side by $1 = \int_0^{\infty} ds \delta( s - \sum_j t_j )$.  

\[
\prod_i \frac{ \Gamma(\alpha_i ) }{ A^{\alpha_i }} = \prod_i \int_0^{\infty} dt_i t_i^{\alpha_i -1 } e^{ - A_i t_i} \int_0^{\infty} ds \delta( s-  \sum_i t_i  )
\]
Note
\[
\begin{aligned} t_i & = sx_i \\ dt_i & = sdx_i \end{aligned} \quad \quad \, \frac{t_i }{s} = x_i 
\]
\begin{equation}\label{Eq:condition_ts}
  1 - \sum_i \frac{ t_i }{s} = 0 
\end{equation}
$t_j$'s are all positive. \\
$s$ is positive.  \\
Then by Eq. (\ref{Eq:condition_ts}), $\frac{t_i}{s}$'s must at least be less than 1.  Thus $x_i$'s must be less than 1.  

Also 
\[
\delta(s - \sum_i t_i ) = \frac{1}{s} \delta( 1 - \sum_j \frac{ t_j}{s} )  =\frac{1}{s} \delta(\sum_j x_j - 1 ) 
\]
\[
\begin{gathered}
  \prod_i \frac{ \Gamma(\alpha_i ) }{ A^{\alpha_i }} = \prod_i \int_0^{\infty} ds \int_0^1 s dx_i s^{\alpha_i - 1 } x_i^{\alpha_i -1 } e^{- A_i x_i s } \delta( s - \sum_j t_j ) = \\ 
  = \prod_i \int_0^1 dx_i x_i^{\alpha_i  -1 } \left( \int_0^{\infty} ds s^{ \sum_i \alpha_j - 1 } e^{ - \sum_j A_j x_j } \right) \delta(\sum_j x_j - 1 ) = \\ 
  = \left( \prod_i \int_0^1 dx_i x_i^{\alpha_i - 1 } \right) \left( \frac{ \Gamma( \sum_i \alpha_i ) }{ \left( \sum_j A_j x_j \right)^{\sum_i \alpha_i } } \right) \delta(\sum_j x_j -1 ) = \frac{ \Gamma(\sum_i \alpha_i ) }{ (n-1)! } \int dF_n \frac{ \prod_i x_i^{\alpha_i -1 } }{ \left( \sum_i x_i A_i \right)^{\sum_i \alpha_i } }
\end{gathered}
\]

\solutionhead{14.2} Recall $\int dx e^{-x^2 } = \sqrt{ \pi}$.  

\[
\int d^n x e^{- \mathbf{x}^2 } = \int dx_1 dx_2 \dots dx_n e^{-x_1^2 } e^{-x_2^2 } \dots e^{-x_n^2 } = \pi^{n/2} = \int r^{d-1} dr d\Omega_d e^{-r^2 } = \int r^{d-1} e^{-r^2 } \int d\Omega_d
\]
\[
\int r^{d-1} e^{-r^2} = \frac{ r^{d-2} e^{ - r^2} }{ -2} -\int (d-2) r^{d-3} \frac{ e^{-r^2} }{ -2} = \frac{d-2}{2} \int r^{d-3} e^{-r^2} 
\]

If $d=2n$, 
\[
\frac{2n-2}{2} \int r^{2n -3} e^{-r^2} = \left( \frac{ 2n - 2 }{2} \right) \cdot \lbrace r^{2n-4} \frac{e^{-r^2}}{ -2} - \int (2n -4) r^{2n-5} \frac{ e^{-r^2} }{ -2} \rbrace = (n-1)! \int re^{-r^2} = \frac{ (n-1)! }{ 2} = \frac{ \Gamma(d/2) }{2} 
\]
So that $\left( \int d\Omega_d \right) \cdot \frac{ \Gamma(d/2) }{2} = \pi^{d/2}$ or $\boxed{ \Omega_d = \frac{ 2\pi^{d/2}}{ \Gamma(d/2) } }$.  

Note that $\Gamma(n+1) = n!$ so $\Gamma\left( \frac{ d}{2} \right) = \Gamma\left( \frac{2n}{2} \right) = \Gamma(n) = (n-1)!$ and $2n-1-2(n-1) = 1$.  

If $d= 2n + 1$
\[
\begin{aligned}
  \int r^{d-1} e^{-r^2} & = \frac{ 2n - 1}{2} \int r^{2n-2} e^{-r^2} = \left( \frac{2n-1}{2} \right) \lbrace r^{2n-3} \frac{ e^{-r^2}}{-2} - \int (2n-3) r^{2n-4} \frac{ e^{-r^2}}{ -2 } \rbrace  = \\
  & = \left( \frac{2n-1}{2} \right) \left( \frac{2n-3}{2} \right) \dots \frac{1}{2} \int e^{-r^2}  = \frac{ (2n -1)!! }{ 2^{n+1} } \sqrt{ \pi }
\end{aligned}
\]

Note $2n-2(k)$ so for $k=n$, $2n - 2n=0$ and $2n+1 - 2k$, so for $k=n$, $2n+1 - 2n = +1$.  

Recall $\Gamma\left( n+\frac{1}{2} \right) = \frac{ (2n)! }{ 2^{2n}  n!} \sqrt{\pi} = \frac{ (2n-1)!!}{ 2^n} \sqrt{ \pi}$.  

So $\pi^{d/2} = \Gamma\left(n + \frac{1}{2} \right)\frac{1}{2} \int d\Omega_d$ or $\boxed{ \Omega_d = \frac{ 2 \pi^{d/2}}{ \Gamma(d/2) } }$.  












\solutionhead{14.4} Recall the development in Sec. 14, Srednicki, so to understand where $\kappa_A$, $\kappa_B$ comes from, that due to this renormalization scheme, to make $g$ dimensionless, through parameter $\mu$, we calculate the necessary counterterms, $A$, and $B$, but they're dependent on renormalization scheme, as we'll see.  

Starting from $\Pi(k^2) = \frac{1}{2} g^2 I(k^2) - Ak^2 - Bm^2 + \mathcal{O}(g^2)$ where $I(k^2) \equiv \int_0^1 dx \int \frac{d^d \overline{q}}{ (2\pi)^d} \frac{1}{ ( \overline{q} + D)^2 }$ using (14.27) and following Srednicki Sec. 14.  

\[
\Pi(k^2) = \frac{ \alpha}{2} \int_0^1 dx D \ln{ \left( \frac{D}{m^2} \right) } - \lbrace \frac{ \alpha}{6} \left[ \frac{1}{ \epsilon } + \ln{ \left( \frac{ \mu}{m} \right)} + \frac{1}{2} \right] + A \rbrace k^2 - \lbrace \alpha \left[ \frac{1}{ \epsilon }+ \ln{ \left( \frac{\mu}{m} \right) } + \frac{1}{2} \right] + B \rbrace m^2 + \mathcal{O}(\alpha^2) 
\]
where $ \begin{aligned} D & = x(1-x)k^2 + m^2 \\ \mu & = \sqrt{ 4\pi} e^{-\gamma/2} \widetilde{\mu} \end{aligned}$ \quad \quad \, $\begin{aligned} \gamma & = 0.5772 \\ \alpha & = \frac{ g^2}{ (4\pi)^3 } \end{aligned}$.

Take
\[
\begin{aligned}
  A & = \frac{-1}{6} \alpha \left[ \frac{1}{3} + \ln{ \left( \frac{ \mu}{ m} \right) } + \frac{1}{2} + \kappa_A \right] + \mathcal{O}(\alpha^2)  \\ 
  B & = - \alpha \left[ \frac{1}{\epsilon} + \ln{ \left( \frac{ \mu}{ m} \right) } + \frac{1}{2} + \kappa_B \right] + \mathcal{O}(\alpha^2)
\end{aligned}
\]
\[
\Pi(k^2) = \frac{ \alpha }{2} \int_0^1 dx D \ln{ \left( \frac{ D}{m^2} \right) } + \alpha \left( \frac{1}{6} \kappa_A k^2 + \kappa_B m^2\right) + \mathcal{O}(\alpha^2)
\]

 \quad \\ 

Impose  
$\begin{aligned} \Pi(-m^2) & = 0 \\ \Pi'(-m^2) & = 0 \end{aligned}$.  

Now $D_0 \equiv \left. D \right|_{k^2 = - m^2} = (1 - x(1-x))m^2$.  

\[
\Pi(-m^2) = \frac{ \alpha }{2} \int_0^1 dx m^2 ( 1 - x(1-x) ) \ln{ (1- x(1-x) ) } + \alpha m^2 \left( \frac{ - \kappa_A}{6} + \kappa_B \right) + \mathcal{O}(\alpha^2)  = 0 
\]
\[
\Longrightarrow \kappa_B = \frac{ \kappa_A}{6} - \frac{1}{2} \int_0^1 dx ( 1 - x(1-x) ) \ln{ (1-x(1-x) ) }
\]

Now
\[
\ln{ \left( \frac{D}{D_0} \right)} = \ln{ \left( \frac{D}{m^2} \right) } - \ln{ (1 - x(1-x) )}
\]
\[
\frac{ \partial }{ \partial k^2} ( D \ln{ \left( \frac{ D}{D_0 } \right) } ) = x(1-x) \ln{ \frac{ D}{D_0} } + \frac{ D}{ D/D_0 } \frac{ x(1-x)}{ D_0} = x(1-x) \left( \ln{ \frac{D}{D_0 } } + 1 \right)
\]
\[
\Pi(k^2) = \frac{ \alpha}{2} \int_0^1 dx \left( D \ln{ \left( \frac{D}{D_0} \right)} + D \ln{ (1 - x(1-x) )} \right) + \alpha \left( \frac{1}{6} \kappa_A k^2 + \kappa_B m^2 \right) + \mathcal{O}(\alpha^2)
\]
\[
\frac{  \partial}{ \partial k^2} \Pi(k^2) = \frac{ \alpha}{2} \int_0^1 dx \left( x(1-x) \left( \ln{ \frac{D}{D_0} } + 1 \right) + x(1-x) \ln{ (1 - x(1-x)) } \right) + \frac{ \alpha \kappa_A}{6} 
\]
\[
\Pi'(-m^2) = \frac{ \alpha}{2} \int_0^1 dx ( x(1-x) + x(1-x) \ln{ ( 1- x(1-x) ) } ) + \frac{ \alpha \kappa_A}{6} = \alpha \left( \frac{1}{12} + \frac{ \kappa_A}{6} + \frac{1}{2} \int_0^1 x (1-x) \ln{ (1-x(1-x)  ) } \right) = 0 
\]

Notice that 
\[
\kappa_B = -\frac{1}{12} - \frac{1}{2} \int_0^1 x(1-x) \ln{ (1-x(1-x) ) } - \frac{1}{2} \int_0^1 (1-x+x^2) \ln{ (x^2 - x+1 ) } = \frac{-1}{12} - \frac{1}{2} \int_0^1 dx \ln{ ( x^2 -x +1 ) } 
\]

\[
\int_0^1 dx \ln{ (1- x(1-x) )} = \left( x \ln{ ( 1 -  x + x^2)  } - (2x + \frac{1}{2} \ln{ ( x^2 - x + 1 ) } - \sqrt{3} \left. \arctan{ \left( \frac{3x - 1 }{ \sqrt{3}} \right) } ) \right) \right|_0^1 = \frac{ \pi}{ \sqrt{3}} - 2 
\]
\[
\boxed{ \kappa_B = \frac{11}{12} - \frac{ \pi}{ 2 \sqrt{3}} = 0.009767 }
\]
Now $\kappa_A = \frac{-1}{2} - 3 \int_0^1 x(1-x) \ln{ (1-x(1-x) ) }$.  
\[
\begin{gathered}
  \int_0^1 dx ( x^2 - x) \ln{ ( x^2 - x+1) } = \\
  = \left. \left( \frac{x^3}{3} \ln{ ( 1 - x + x^2) } - \frac{x^2}{2} \ln{ ( x^2 - x + 1 ) } - \frac{ 2 x^3}{9} + \frac{x^2}{3} + \frac{5x}{6} + \frac{1}{12} \ln{ ( x^2 - x  + 1 ) } - \frac{ \sqrt{3}}{2} \arctan{ \left( \frac{ 2 x - 1 }{ \sqrt{3}} \right) } \right) \right|_0^1 = \\ 
  = \frac{ 51}{54} - \frac{ \pi \sqrt{3}}{6} 
\end{gathered}
\]
\[
\kappa_A = \frac{-1}{2} + 3 \left( \frac{51}{54} - \frac{ \pi \sqrt{3}}{6} \right) = \boxed{ \frac{7}{3} - \frac{ \pi \sqrt{3}}{2} \simeq -0.3874 }
\]

\solutionhead{14.5} Work directly in momentum space.

I will develop $\varphi^4$ theory here.  
\[
\begin{gathered}
  i \int d^4x \frac{-1}{4!} Z_{\lambda} \lambda ( \varphi(x))^4 = \frac{-i}{ 4!} Z_{\lambda} \lambda \int d^4x \left( \int \frac{ d^4k }{ (2\pi)^4} e^{ikx} \widetilde{\varphi}(k) \right)^4 = \\
  \frac{-i}{4!} Z_{\lambda} \lambda \left( \prod_{i=a}^d \int \frac{ d^4k_i}{ (2\pi)^4} \widetilde{\varphi}(k) \right) (2\pi)^4 \delta^4(k_a + k_b + k_c + k_d)
\end{gathered}
\]
So

\begin{align}
  Z_1  & \simeq \frac{- i }{4!} Z_{\lambda} \lambda \left( \prod_{i=a}^d \int \frac{ d^4k_i}{ (2\pi)^4} \frac{ (2\pi)^4}{ i } \delta_{J_{-k_i} } \right) (2\pi)^4 \delta^4\left( \sum_{i=a}^d k_i \right) \cdot \frac{1}{3!} \left( \frac{i }{2} \int \frac{ d^4k}{ (2\pi)^4} \frac{ \widetilde{J}_k \widetilde{J}_{-k} }{ k^2 + m^2 - i \epsilon} \right)^3 = \label{Eq:14_5-phi4_counting_01} \\  
    & = \frac{- i}{4} Z_{\lambda} \lambda ( 2\pi)^4 \delta^4(k_c + k_d) \left( \prod_{i=c}^d \int \frac{d^dk_i}{ (2\pi)^d} \frac{ \widetilde{J}_{k_i} }{ k_i^2 + m^2 - i \epsilon } \right) \int \frac{ d^dk_a }{ ( 2\pi)^4} \frac{ 1/i }{ k_a^2 + m^2 - i \epsilon } \label{Eq:14_5-phi4_counting_02} 
\end{align}


Now $\frac{1}{i} \mathbf{\Delta}(x_1-x_2) \equiv \langle 0 | T\varphi(x_1) \varphi(x_2) | 0 \rangle$, i.e. the definition of the exact propagator.  

Note about the counting factor from going from Eq. (\ref{Eq:14_5-phi4_counting_01}) to Eq. (\ref{Eq:14_5-phi4_counting_02}): there is a $\binom{4}{2} = 3!$ factor for choosing which of the functional derivatives, which 2 legs of the $\phi^4$ theory, will loop into itself.  The functional derivatives are indistinguishable.  Then there is the usual $3! \cdot 2^{3}$ factor, canceling the denominators in the propagators, from matching 3 functional derivatives to propagators.  

Indeed, the counting factor could be checked by this \textbf{very educational diagram} (cf. Arthur Lipstein and Peskin and Schroeder, \textbf{An Introduction to Quantum Field Theory}, 1995)


\[
\langle \underbracket{ k_1 | \phi_x }_2 \underbracket{ \phi_x \phi_x }_{ \binom{4}{2} } \underbracket{ \phi_x | k_2 }_1 \rangle
\]



So apply $ \frac{ (2\pi)^4}{i} \delta_{\widetilde{J}_k } \frac{ (2\pi)^4}{ i } \delta_{\widetilde{J}_{-k} }$ to obtain the first order term.  

\[
\left( \frac{ 1 / i }{ k^2 + m^2 - i \epsilon } \right) \frac{ Z_{\lambda} \lambda}{2 i } \int \frac{ d^d l }{ (2\pi)^d} \frac{ 1 / i }{ l^2 + m^2 - i \epsilon } \frac{ 1 /i }{ k^2 + m^2 - i \epsilon } = \widetilde{\Delta}(k) i \Pi(k^2 ) \widetilde{\Delta}(k)
\]



I'll include the $A,B$ counterterms soon.

Before evaluating $\int \frac{ d^d l}{ (2\pi)^d} \frac{ 1 }{ l^2 + m^2 - i \epsilon }$, renormalize $\lambda$ to make it dimensionless.  

So before, we had, 
\[
\begin{gathered}
  \left[ d^d x \right] = - d \\ 
  [ \varphi ] = \frac{1}{2} ( d-2) \\ 
  \Longrightarrow - d + 4 \cdot \frac{1}{2} ( d-2) + [ \lambda ] = [ \lambda ] + - d + 2d - 4 = [ \lambda ] + d -4 = 0 \text{ or } [ \lambda ] = 4 - d
\end{gathered}
\]

So let $\epsilon = 4 -d$.  

$\lambda \to \lambda \widetilde{\mu}^{\epsilon}$, which will force $\lambda$ to be dimensionless.  

\textbf{Remember} to do the Wick rotation before applying (14.27).  

\[
\int \frac{ d^d \overline{q}}{ (2\pi)^d} \frac{ ( \overline{q}^2)^2 }{ ( \overline{q}^2 + D )^b} = \frac{ \Gamma(b-a- \frac{d}{2} ) \Gamma(a+ d/2) }{ (4\pi)^{d/2} \Gamma(b) \Gamma(d/2) } D^{ - (b-a-\frac{d}{2} ) }
\]
\[
\begin{gathered}
  \int \frac{ d^d l}{ (2\pi)^d} \frac{1}{ l^2 + m^2 - i \epsilon } = i \int \frac{ d^d q}{ (2\pi)^d} \frac{1}{ q^2 + m^2 - i \epsilon } = i \frac{ \Gamma( 1 - 0 - \frac{d}{2} ) \Gamma( 0 + \frac{d}{2} ) }{ (4\pi)^{d/2} \Gamma(1) \Gamma(d/2) } ( m^2)^{- (1 - 0 - \frac{d}{2} ) } = i \frac{ \Gamma\left( 1 - \frac{d}{2} \right) }{ (4\pi)^{d/2} } (m^2)^{ \left( \frac{d}{2} -1 \right) } = \\ 
   =i \frac{ \Gamma\left( \frac{ \epsilon-2}{2} \right) }{ (4\pi)^{(4-\epsilon)/2} } m^{ 2 \left( \frac{ 2 - \epsilon }{2} \right) }  = \frac{ i \Gamma\left( \frac{ \epsilon}{2} - 1 \right) }{ (4\pi)^2} ( 4\pi)^{\epsilon/2} \frac{ m^2}{ m^{\epsilon }}
\end{gathered}
\]

now include the counterterms $-i (Ak^2 + Bm^2)$.  

\[
i \Pi(k^2) = i \frac{ \lambda \widetilde{\mu}^{\epsilon} }{2} \Gamma\left( \frac{ \epsilon}{2} - 1 \right) \left( \frac{ m}{4\pi} \right)^2 \left( \frac{ 4\pi }{ m^2 } \right)^{\epsilon/2} -i ( Ak^2 + Bm^2)
\]

$\Pi'(k^2 = - m^2) = 0$ yields $A=0$.  

\noindent Now recall $\Gamma(-n +x ) = \frac{ (-1)^n}{n!} \left[ \frac{1}{x} - \gamma + \sum_{k=1}^n k^{-1} + O(x) \right]$, \quad \, (14.26), so \\ 
\phantom{Now recall} $\Gamma(-1 + \frac{ \epsilon}{2} ) = (-1) \left[ \frac{1}{\epsilon/2} - \Gamma + 1 + \mathcal{O}(\epsilon) \right]$,
and recall $A^{\epsilon/2} = 1 + \frac{\epsilon}{2} \ln{A} + \mathcal{O}(\epsilon^2)$ \quad \, (14.33) and $\mu \equiv \sqrt{ 4\pi} e^{-\gamma/2} \widetilde{\mu}$  \quad \, (14.35).  

\[
\begin{gathered}
  \frac{ \lambda}{2} (-1) \left( \frac{1}{ \epsilon/2} -\Gamma +1 + \mathcal{O}(\epsilon) \right) \left( 1 + \frac{ \epsilon}{2} \ln{ \frac{ 4\pi \widetilde{\mu}^2 }{ m^2} }+ \mathcal{O}(\epsilon^2) \right) \left( \frac{m}{4\pi} \right)^2  = \\ 
  = \frac{ -\lambda}{2} \left( \frac{ m}{4\pi} \right)^2 \left( \frac{1}{ \epsilon/2} - \Gamma + 1 + \ln{ \frac{ 4\pi \widetilde{\mu}^2}{ m^2 }} + \mathcal{O}(\epsilon) \right) = \frac{ - \lambda}{2} \left( \frac{m}{4\pi} \right)^2 \left( \frac{1}{ \epsilon/2} + 1 + 2 \ln{ \frac{ \mu}{m}} \right) + \mathcal{O}(\epsilon) 
\end{gathered}
\]
\[
\Pi(k^2 =  - m^2 ) = 0 = - Bm^2 + \frac{ -\lambda}{2} \left( \frac{m}{4\pi} \right)^2 \left( \frac{1}{ \epsilon/2} + 1 + 2 \ln{ \frac{ \mu}{m }  } \right) + \mathcal{O}(\epsilon) 
\]
\[
\boxed{ B = \frac{ \lambda}{  ( 4\pi)^2 } \left( \frac{1}{ \epsilon} + \ln{ \frac{ \mu}{m} } + \frac{1}{2} \right) + \mathcal{O}(\epsilon) }
\]











\solutionhead{14.6} 

The answer ends up being the same as for $\phi^4$ theory, Problem 14.5.  I will develop complex scalar theory here in case it is not clear where the Feynman rules come from.  

Consider 1 vertex and 3 propagators.  
\[
\frac{-i}{4} \lambda \int d^4x \left( \frac{1}{i} \delta_{J_x} \frac{1}{i} \delta_{J_x^{\dag}} \right)^2 \frac{1}{3!} \left( i \int d^4y d^4z J_y \Delta_{y-z} J_z^{\dag} \right)^3 =  -\lambda \int d^4x \int d^4a J_a \Delta_{a-x} \Delta(0) \int d^4b \Delta_{x-b} J_b^{\dag}
\]
\[
\Longrightarrow \langle 0 | T \varphi(x_2) \varphi^{\dag}(x_1) | 0 \rangle \simeq  - \lambda \int d^4x \Delta_{x_1-x} \Delta(0) \Delta_{x-x_2} \frac{1}{i^2}
\]

It's easier to evaluate $\Delta(0)$ in momentum space, so let's redo all this in momentum space.  

 \begin{gather}
   i \int d^4x \frac{-1}{4}  \lambda ( \varphi^{\dag} \varphi)^2  = \frac{ - i \lambda}{4} \int d^4x \left( \int \frac{ d^4k }{(2\pi)^4} e^{-ikx} \widetilde{\varphi}^{\dag}_k \frac{ \int d^4l }{ (2\pi)^4} e^{i lx} \widetilde{\varphi}_{l } \right)^2 = \notag \\ 
    = \frac{- i \lambda}{4} \left( \prod_{\alpha = a }^d \int \frac{ d^4k_{\alpha}}{ (2\pi)^4} \right) \delta^4( - k_a + k_b - k_c + k_d) \frac{ (2\pi)^4}{i} \delta_{J_a} \frac{ (2\pi)^4}{i } \delta_{J^{\dag}_b} \frac{ (2\pi)^4}{ i } \delta_{J_c} \frac{ (2\pi)^4}{i } \delta_{J_d^{\dag}} \label{Eq:complexscalar_3prop}
 \end{gather}


Referencing Problem 8.7, $Z_0(J^{\dag},J) = \exp{ \left[ i \int \frac{ d^4k}{ (2\pi)^4} \frac{ \widetilde{J}_k \widetilde{J}_k^{\dag}}{ k^2 + m^2 - i \epsilon } \right] }$ and source terms are $ i \int \frac{d^4k}{(2\pi)^4} \left( \widetilde{J}_k^{ \dag} \varphi_k + \widetilde{J}_k \varphi^{\dag}_k \right)$

Before applying the first order term, Eq. (\ref{Eq:complexscalar_3prop}), to 3 propagators, multiply by $2 \cdot 2$ counting factor to determine which of the functional derivatives will be involved with the loop.  This could also be clearly seen, as in Problem 14.5, by this diagram:

\[
\langle \underbracket{ k_1 | \phi_x }_2 \underbracket{ \phi_x^{\dag} \phi_x }_{ 2 } \underbracket{ \phi_x^{\dag} | k_2 }_1 \rangle
\]



So applying the first order term, Eq. (\ref{Eq:complexscalar_3prop}), to 3 propagators, $\frac{1}{3!} \left( i \int \frac{ d^4k}{(2\pi)^4} \frac{ \widetilde{J}_k \widetilde{J}_k^{\dag}}{  k^2 + m^2 - i \epsilon } \right)^3$, 

\[
\begin{gathered}
   - \lambda  \left( \prod_{\alpha = a}^c \int d^4k_{\alpha} \right) \int d^4k_d \delta^4(-k_a + k_b - k_c + k_d) \delta_{J_d^{\dag}} \frac{ \widetilde{J}_a^{\dag}}{ k_a^2 + m^2 } \frac{ \widetilde{J}_c^{\dag}}{ k_c^2 + m^2 } \frac{ \widetilde{J}_b}{ k_b^2 + m^2 } \frac{1}{ (2\pi)^{12}} = \\ 
  =  - \lambda \left( \prod_{\alpha = a}^c \int d^4k_a \right) \delta^4(k_b - k_c) \frac{1}{k_a^2 + m^2 } \frac{ \widetilde{J}_c^{\dag} }{ k_c^2 + m^2 } \frac{ \widetilde{J}_b }{ k_b^2 + m^2 } \frac{1}{(2\pi)^{12}}  
  \end{gathered}
\]
\[
\langle 0 | T \widetilde{\varphi}(k_2) \widetilde{\varphi}^{\dag}(k_1) | 0 \rangle \simeq  \lambda \int d^4k_a \frac{1}{k_a^2 + m^2 } \frac{1}{(2\pi)^4} \frac{1}{ k_2^2 + m^2 } \frac{1}{ k_1^2 + m^2} \delta^4(k_1-k_2)
\]

\[
\Longrightarrow \left( \frac{1/i}{ k^2 + m^2 } \right)  (\lambda/i)  \int \frac{d^4l }{ (2\pi)^4} \frac{ 1 /i }{  l^2 + m^2 } \frac{ 1 /i }{ k^2 + m^2 } = \widetilde{\Delta}(k) i \Pi(k^2) \widetilde{\Delta}(k)
\]

So since $i \Pi(k^2) =  \lambda/i \int  \frac{ d^4l}{(2\pi)^4} \frac{ 1 /i }{ l^2 + m^2 }$, and the counterterms $A$, $B$ are exactly the same as in the case of $\varphi^4$ theory (cf. Problem 14.5), then the steps are exactly the same as with Problem 14.5.  

\[
\boxed{ B =  \frac{\lambda}{  ( 4\pi)^2 } \left( \frac{1}{ \epsilon} + \ln{ \frac{ \mu}{m} } + \frac{1}{2} \right) + \mathcal{O}(\epsilon) }
\]
\[
\boxed{ A = 0}
\]

Note that there are 2 kinds of propagators, for type $a$ particles and type $b$ particles, but the factors and results are the same for each.  

\section{The one-loop correction in Lehmann-K\"all\'en form}

%%----------------------20100211-begin-%%%

\solutionhead{15.1}
\begin{itemize}
\item[(a)] Now given $\Pi(k^2) = \frac{1}{2} \alpha \Gamma(-1 + \frac{ \epsilon}{2} ) \int_0^1 dx D \left( \frac{ 4\pi \widetilde{\mu}^2 }{ D } \right)^{\epsilon/2} - Ak^2 - Bm^2 + \mathcal{O}(\alpha^2)$ \quad \quad \,  (14.32), where $D = x(1-x) k^2 + m^2$.  

\[
\begin{aligned}
  & \Pi_{loop}(k^2) = \frac{1}{2} \alpha \Gamma( -1 + \frac{ \epsilon}{2} ) \int_0^1 dx ( 4\pi \widetilde{\mu}^2 )^{\epsilon/2} D^{1- \epsilon/2}  \\
  & \Pi'_{loop}(k^2) = \frac{ \alpha}{2} \Gamma\left( \frac{ \epsilon }{2} - 1 \right) \int_0^1 dx (4\pi \widetilde{\mu}^2 )^{\epsilon/2} \left( 1-  \frac{ \epsilon}{2} \right) D^{-\epsilon/2} x(1-x) \\ 
  & \Pi'_{loop}(-m^2) = \frac{ \alpha}{2} \Gamma\left( \frac{ \epsilon}{2} -1 \right) (4\pi \widetilde{\mu}^2)^{\epsilon/2} \left( 1 - \frac{ \epsilon}{2} \right) \int_0^1 dx (m^2)^{-\epsilon/2} (1 - x(1-x))^{- \epsilon/2} x(1-x)
\end{aligned}
\]

For $x \in [0,1]$, $x^2 - x+1 \geq \frac{3}{4}$ (note $(x^2 - x+1)'=0$ when $x= \frac{1}{2}$, and thus, $\frac{3}{4}$).  

So using $A^{\epsilon/2} = 1 + \frac{ \epsilon}{2} \ln{ A} + \mathcal{O}(\epsilon^2)$, $(x^2 - x+1)^{-\epsilon/2} = 1 + \frac{ - \epsilon }{2} \ln{ (x^2 - x +1) } + \mathcal{O}(\epsilon^2)$.  

Also, using $\Gamma(-n + x ) = \frac{ ( -1)^n}{n!} \left[ \frac{1}{x} - \gamma + \sum_{k=1}^n k^{-1} + \mathcal{O}(x) \right]$  \quad \, (14.26), we get

\[
\begin{gathered}
  \Pi'_{loop}(-m^2) = \frac{ \alpha}{2} ( -1) \left[ \frac{ 1}{ \epsilon/2} - \Gamma + 1 + \mathcal{O}(\epsilon) \right] \left( \frac{ 4\pi \widetilde{\mu}^2 }{ m^2} \right)^{\epsilon/2} \left( 1 - \frac{ \epsilon}{2} \right) \int_0^1 dx \left( 1 - \frac{ \epsilon}{2} \ln{ ( 1 - x(1-x) ) } + \mathcal{O}(\epsilon^2) \right) x( 1 - x) = \\ 
  = \frac{ - \alpha}{2} \left( 1 + \frac{ \epsilon}{2} \ln{ \frac{ 4\pi \widetilde{\mu}^2 }{ m^2 } } \right) \left( \frac{1}{ \epsilon/2} - \Gamma + 1 + \mathcal{O}(\epsilon) - 1 \right) \left( \frac{1}{6} - \frac{ \epsilon}{2} \int_0^1 dx x ( 1-x) \ln{ ( 1- x(1-x)) } + \mathcal{O}(\epsilon^2) \right) = \\
  = \frac{ - \alpha}{2} \left( 1 + \frac{ \epsilon}{2} \ln{ \frac{ 4\pi \widetilde{\mu}^2 }{ m^2 } } \right) \left( \frac{1}{6} \left( \frac{ 1 }{ \epsilon/2} - \gamma \right) - \int_0^1 dx x(1-x) \ln{ (1 - x(1-x) )} + \mathcal(\epsilon) \right) = \\ 
  = \frac{- \alpha}{2} \left( \frac{1}{6} \left( \frac{2}{ \epsilon } - \gamma \right) - \int_0^1 dx x(1-x) \ln{ ( 1 - x(1-x)) } + \frac{1}{6} \ln{ \frac{ 4 \pi \widetilde{\mu}^2 }{m^2} } + \mathcal{O}(\epsilon) \right) = \\
  = \frac{ - \alpha}{6} \left( \frac{1}{ \epsilon } - 3 \int_0^1 dx x(1-x) \ln{ ( 1 - x(1-x)) } + \ln{ \frac{ \mu}{m} } \right) = \frac{ - \alpha}{6} \left[ \frac{1}{\epsilon} + \ln{ \left( \frac{ \mu}{m} \right) } + \frac{1}{2} + \kappa_A \right] + \mathcal{O}(\alpha^2) = A
\end{gathered}
\]
where the substitution $\mu = \sqrt{ 4 \pi } \widetilde{\mu} e^{-\gamma/2}$ was used and the result from Problem 14.4, that $\kappa_A = -\frac{1}{2} - 3 \int_0^1 x ( 1 - x) \ln{ (1-x(1-x) ) } = -0.3874$ was used.  

So $\boxed{ \Pi'_{loop}(-m^2) = A}$.  

Recall the Cauchy integral formula $f^{(n)}(z) = \frac{n!}{ 2\pi i} \int \frac{ f(s) ds}{ (s-z)^{n+1} }$.  

Then 
\[
A = \Pi'_{loop}(-m^2) = \frac{1}{2 \pi i} \int \frac{ \Pi_{loop}(w) }{ (w + m^2)^2 }
\]
\item[(b)]
\item[(c)]
\end{itemize}










%%----------------------20100211-end-%%%







\section{Loop corrections to the vertex}

\solutionhead{16.1} 

I will develop $\varphi^4$ theory here in case it is not familiar.  

Recall in $\varphi^4$ theory,
\[
\begin{aligned}
  & \mathcal{L}_0 = - \frac{1}{2} \partial^{\mu} \varphi \partial_{\mu} \varphi - \frac{1}{2} m^2 \varphi^2 \\ 
  & \mathcal{L}_1 = - \frac{1}{4!} Z_{\lambda} \lambda \varphi^4 \\ 
  & \mathcal{L}_{ct} =  \frac{-1}{2} (Z_{\varphi}- 1) \partial^{\mu} \varphi \partial_{\mu} \varphi - \frac{1}{2} ( Z_m - 1) m^2 \varphi^2 = -\frac{1}{2} A \partial^{\mu} \varphi \partial_{\mu} \varphi - \frac{1}{2} Bm^2 \varphi^2
\end{aligned}
\]

For $\mathcal{O}(\lambda)$ order,

\[
Z_1 \simeq \frac{-i}{4!} Z_{\lambda} \lambda \int d^4x \left( \frac{1}{i} \delta_{J_x} \right)^4 \frac{1}{4!} \left[ \frac{i}{2} \int d^4y d^4y' J_y \Delta_{y-y'} J_y' \right]^4 = \frac{-Z_{\lambda} \lambda i }{ 4!} \int d^4x \left( \int dy J_y \Delta_{y-x} \right)^4
\]
$4! \cdot 2^4$ counting factor for ways to match $(\delta_{J_x})^4$ to propagators and sources.  

Apply LSZ formula.  
\[
\begin{gathered}
  \langle f | i \rangle   = \int d^4x_1 e^{ik_1 x_1} d^4x_2 e^{ i k_2 x_2} d^4x_3 e^{ -i k_3 x_3} d^4x_4 e^{ -i k_4 x_4} (-Z_{\lambda} \lambda i ) \int d^4x \Delta_{x_1- x} \Delta_{x_2 - x} \Delta_{x_3 - x} \Delta_{x_4 - x} = \\ 
  = -Z_{\lambda} \lambda i (2\pi)^4 \delta^4(k_1 + k_2 - k_3 - k_4)
\end{gathered}
\]

For $\mathcal{O}(\lambda^2)$ order, assuming $Z_{\lambda} = 1 + \mathcal{O}(\lambda^2)$,

\begin{gather}
  \left( \frac{ -i}{4!} \lambda \right)^2 \left( \int d^4x \left( \frac{1}{i} \delta_{J_x} \right)^4 \right)^2 \frac{1}{6!} \left[ \frac{i}{2} \int d^4y d^4x' J_y \Delta_{y-y'} J_{y'} \right]^6 = \notag \\
  = \left( \frac{ -i}{4!} \lambda \right)^2 \int d^4x d^4y \frac{1}{i} \delta_{J_x} \frac{1}{i} \delta_{J_y} \left( \int d^4x' J_{x'} \Delta_{x'-x} \right)^3 \left( \int d^4y' J_{y'}  \Delta_{y'-y} \right)^3     \label{Eq:phi4theory-counting01} \\
  = \frac{ \lambda^2}{(4!)^2} \int d^4x d^4y (3)^2 \left( \int d^4x' J_{x'} \Delta_{x'-x} \right)^2 \left( \int d^4y' J_{y'} \Delta_{y' - y} \right)^2 \left( \Delta_{y-x} \right)^2  \label{Eq:phi4theory-counting02}
\end{gather}
  
where I used a $6! \cdot 2^6$ counting factor matching $\delta_{J_x}$'s to propagators in Eq. (\ref{Eq:phi4theory-counting01}) and $3^2$ counting factor used in Eq. (\ref{Eq:phi4theory-counting02}) to match last 2 functional derivatives to sources.  There's also a $4 \cdot 4$ counting factor from distinguishing 1 of the functional derivatives, among 4, for each vertex, to be the one involved with the internal propagator.  

Apply LSZ formula,
\[
\begin{gathered}
  \int d^4x_1 e^{ik_1x_1}  d^4x_2 e^{ik_2x_2}  d^4x_3 e^{-ik_3x_3}  d^4x_4 e^{-ik_4x_4} \frac{ \lambda^2}{ 4} \int d^4x d^4y \Delta_{x_1 -x}  \Delta_{x_2 -x}  \Delta_{x_3 -y}  \Delta_{x_4 -y} ( \Delta_{y-x} )^2 = \\
  = \frac{ \lambda^2}{ 4 } \int d^4x d^4y e^{ i (k_1 + k_2)x} e^{-i (k_3 + k_4)y } \int \frac{ d^4k}{(2\pi)^4} \frac{ e^{ik (y-x) }}{ k^2 + m^2 - i\epsilon } \int \frac{ d^4l}{(2\pi)^4} \frac{ e^{il (y-x) }}{ l^2 + m^2 - i\epsilon } = \\
  = \frac{ \lambda^2 }{4} \int d^4k \frac{ 1 }{ k^2 + m^2} \int d^4l \frac{1}{ l^2 + m^2 } \delta^4(k_1 + k_2 - k - l) \delta^4(-k_3 -k_4 + k+ l ) = \\ 
  = \frac{ \lambda^2 }{ 4 } \int d^4l \frac{ 1}{  ( k_1 + k_2 - l)^2 + m^2 } \frac{1}{ l^2 + m^2 } \delta^4(k_1 + k_2 - k_3 -k_4)
\end{gathered}
\]

Note that there are 3 cases:  \\
$\begin{aligned}
  & e^{i (k_1 + k_2)x}    e^{-i (k_3 + k_4)y}  \\
  & e^{i (k_1 - k_3)x}    e^{-i (k_4 - k_2)y}  \\
  & e^{i (k_1 - k_4)x}    e^{-i (k_3 - k_2)y}  
\end{aligned}$

And for each case, there's $2$ counting factor for matching $\delta_{J_{x_i}}$'s to sources for each case.   

\[
\Longrightarrow \frac{ (i \lambda)^2 }{ 2 } \int d^4l \frac{ 1/i}{ l^2 +  m^2 } \left( \frac{ 1 /i }{ (k_1 + k_2 - l)^2 + m^2 } + \frac{ 1 /i }{ (k_1 - k_3 - l)^2 + m^2 }  + \frac{ 1 /i }{ (k_1 - k_4 - l)^2 + m^2 }  \right) 
\]

\hrulefill

Calculating the vertex $i \mathbb{V}_4 = i \mathbb{V}_4(k_1,k_2,k_3,k_4)$, with $k_1 + k_2 = k_3 + k_4$ (statement of momentum conservation), with \\
$-i \lambda$ being the original, tree-level vertex,

So to second order in $\lambda$,

\[
i \mathbb{V}_4 = -i \lambda + \frac{1}{2} \lambda^2 \int \frac{d^4l}{ (2\pi)^d} \mathbf{\widetilde{\Delta}}(l^2) \lbrace \mathbf{\widetilde{\Delta}}(k_1 + k_2 - l )^2 + \mathbf{\widetilde{\Delta}}(l + k_1 - k_3)^2 + \mathbf{\widetilde{\Delta}}(l + k_1 - k_4)^2 \rbrace
\]
with
\[
\mathbf{ \widetilde{\Delta}}(l^2) \mathbf{\widetilde{\Delta}}(k_1 + k_2 - l )^2 = \frac{1}{ l^2 + m^2 } \frac{ 1 }{ (l - (k_1 + k_2)  )^2 + m^2 } = \int dF_2 ( x_1 l^2 + x_2(l - (k_1 + k_2))^2 + m^2 )^{-2}
\]
where
\[
dF_2 = 1 \int_0^1 dx_1 dx_2 \delta(x_1+ x_2 - 1 )
\]

\[
\begin{gathered}
  \Longrightarrow \int_0^1 dx ((1-x)l^2 + x(l^2 - s -  l(k_1 + k_2)) + m^2 )^{-2} = \\ 
  = \int_0^1 dx (l^2 - 2lx(k_1 + k_2) - xs +m^2)^{-2} = \int_0^1 dx ((l - x(k_1 + k_2))^2 - x(1-x)s + m^2)^{-2}
\end{gathered}
\]
Likewise,
\[
\begin{aligned}
  & \mathbf{ \widetilde{\Delta}}(l^2) \mathbf{\widetilde{\Delta}}( l +k_1 - k_3)^2 = \int_0^1 dx ((l+ x(k_1- k_3))^2 - x(1-x)t + m^2)^{-2}  \\
  &   \mathbf{ \widetilde{\Delta}}(l^2) \mathbf{\widetilde{\Delta}}( l +k_1 - k_4)^2 = \int_0^1 dx ((l+ x(k_1- k_4))^2 - x(1-x)u + m^2)^{-2} 
\end{aligned}
\]

In each $\mathcal{O}(\lambda^2)$ diagram, $q = (l \pm xv)$, $v = s,t,u$, $D = -x(1-x)v + m^2$.  

Do a Wick rotation for each diagram: $\int \frac{d^dl}{(2\pi)^d} = i \int \frac{d^d \overline{q}}{ (2\pi)^d}$.  

\[
\frac{ \mathbb{V}_4}{ \lambda} = - Z_{\lambda} + \frac{ \lambda}{2} \int \frac{ d^d \overline{q}}{ (2\pi)^d} \int_0^1 dx \lbrace (\overline{q}^2 - x(1-x)s + m^2)^{-2} + (\overline{q}^2 - x(1-x)t + m^2)^{-2} + (\overline{q}^2 - x(1-x)u + m^2)^{-2} \rbrace
\]

Do the renormalization (make interaction constant dimensionless).  \\
$\lambda \to \lambda \widetilde{\mu}^{\epsilon}$; \quad \, $\epsilon = 4 -d$; \quad \, $d= 4-\epsilon$.  

\[
\begin{gathered}
  \Longrightarrow \frac{ \Gamma(2 - \left( \frac{ 4 - \epsilon}{2} \right) ) }{ (4\pi)^{2- \epsilon/2} } D^{-2 + \frac{ 4- \epsilon}{2} } = \frac{ \Gamma\left( \frac{ \epsilon}{2} \right)}{ (4\pi)^{2 - \epsilon/2} } D^{-\epsilon/2} \xrightarrow{ \times \widetilde{\mu}^{\epsilon}} \frac{ \Gamma\left( \frac{ \epsilon}{2} \right)}{ (4\pi)^2} \left( \frac{ 4 \pi \widetilde{\mu}^2}{ D} \right)^{\epsilon/2} \simeq \\ 
  \simeq \frac{ \left( \frac{2}{\epsilon} - \gamma + \mathcal{O}(\epsilon) \right) }{ (4 \pi)^2} \left( 1 + \frac{\epsilon}{2} \ln{ \left( \frac{ 4 \pi \widetilde{\mu}^2 }{D} \right) } + \mathcal{O}(\epsilon^2) \right) =  \frac{ \frac{2}{\epsilon} + \ln{ \left( \frac{ 4 \pi \widetilde{\mu}^2 }{ D} \right) } - \gamma }{ (4\pi)^2 }
\end{gathered}
\]

With $\ln{ \frac{ 4 \pi \widetilde{\mu}^2 }{ e^{\gamma}} } = \ln{ \mu^2}$, 
\[
\begin{gathered}
  \frac{  \mathbb{V}_4}{\lambda} = - Z_{\lambda} + \frac{  \lambda}{2} \int_0^1 \frac{dx}{ (4\pi)^2} \left( 3 \left( \frac{2}{\epsilon } - \gamma + \ln{4\pi \widetilde{\mu}^2 } \right) - \ln{ \left( m^2 - x(1-x)s \right) } - \ln{ (m^2 - x(1-x)t) } - \ln{ (m^2 - x(1-x) u ) } \right) = \\ 
  = - Z_{\lambda} + \frac{ \lambda }{2} \int_0^1 \frac{dx}{ (4\pi)^2} \left( 3 \left( \frac{2}{\epsilon} + 2 \ln{ \frac{ \mu}{ m} } \right) - \ln{ \left( \frac{m^2 - x(1-x) s}{ m^2 } \right) } - \ln{ \left( \frac{ m^2 - x(1-x)t }{ m^2} \right) } - \ln{ \left( \frac{ m^2 - x(1-x) u }{ m^2 } \right) } \right)
\end{gathered}
\]

For $\mathbb{V}_4 = \lambda$, \\
and setting $Z_{\lambda} = 1 - C$,

\[
\begin{gathered}
  C =   \frac{\lambda}{2} \int_0^1 \frac{ dx}{ (4\pi)^2 } \lbrace 3 \left( \frac{2}{ \epsilon} + 2 \ln{ \frac{\mu}{m} } \right) - \ln{ \left( \frac{ m^2 - x(1-x)s}{ m^2 } \right) } - \ln{ \left( \frac{ m^2 - x(1-x)t }{ m^2 } \right) } - \ln{ \left( \frac{m^2 - x(1-x)u }{m^2} \right) } \rbrace = \\ 
  =  \frac{\lambda}{2}  \frac{ 1}{ (4\pi)^2 } \lbrace  \frac{6}{ \epsilon} + 6 \ln{ \frac{\mu}{m} }  -\int_0^1 dx \left( \ln{ \left(  1 - \frac{ x(1-x)s}{ m^2 } \right)} - \ln{ \left( 1 - \frac{ x(1-x)t }{ m^2 } \right) } - \ln{ \left( 1 - \frac{ x(1-x)u }{m^2} \right) } \right)   
\end{gathered}
\]
and remember that for $\mathbb{V}_4 = \lambda$, $(k_1,k_2,k_3,k_4) = ((m_1,0), (m_2,0), (m_3,0), (m_4,0) ) = (0,0,0,0)$.  

So
\[
\begin{aligned}
  & s = - (k_1 +k_2)^2 = ( m_1 +m_2)^2 = 4 m^2 \\ 
  & t = -(k_1 -k_3)^2 = 0 \\ 
  & u = - (k_1-k_4)^2 =0  
\end{aligned}
\]

Now $1-  x(1-x)4 = 4 (x^2 -x + \frac{1}{4} ) = 4(x-\frac{1}{2} )^2$.  
\[
\begin{gathered}
  \int_0^1 dx \ln{ (1 - x(1-x)4) } = \int_0^1 dx \ln{ ( 2 ( x- \frac{1}{2} ) )^2 } = \int_{-1}^1 dy \ln{ y^2} \left( \frac{1}{2} \right) = \frac{1}{2} \lbrace \int_0^1 dy \ln{y^2} + \int_{-1}^0 dy \ln{ y^2} \rbrace = \\ 
  = \frac{1}{2} \left( \int_0^1 dy \ln{y^2} + \int_0^1 dy \ln{y^2} \right) = 2 \int_0^1 \ln{y} = 2 \left. ( y\ln{y} - y ) \right|_0^1 = 2(-1) \\ 
  \text{ having used the substitution } \\
  \begin{aligned}
    y & = 2(x - \frac{1}{2} )  \\ 
    dy & = 2 dx
\end{aligned}
\end{gathered} 
\]
\[
C = \frac{ \lambda}{2} \left( \frac{1}{4\pi } \right)^2 \lbrace \frac{ 6}{\epsilon} + 6\ln{ \frac{ \mu}{m }} + 2 \rbrace  = \boxed{ \frac{ \lambda}{ ( 4\pi)^2} \left( \frac{3}{\epsilon} + 3 \ln{ \frac{ \mu}{ m}} + 1 \right)} 
\]

\solutionhead{16.2} 
In case it is unfamiliar, I will develop the Feynman rules for the complex scalar field theory here.  

Recall that 
\[
\begin{aligned}
  & \mathcal{L}_0 = - \partial^{\mu} \varphi^{\dag} \partial_{\mu} \varphi - m^2 \varphi^{\dag} \varphi \\  
  & \mathcal{L}_1 = -\frac{1}{4} Z_{\lambda} \lambda ( \varphi^{\dag} \varphi )^2 + \mathcal{L}_{ct} \\ 
  & \mathcal{L}_2 = - (Z_{ \varphi} - 1) \partial^{\mu} \varphi^{\dag} \partial_{\mu} \varphi - (Z_m - 1) m^2 \varphi^{\dag} \varphi
\end{aligned}
\]
Recall also 

\[
Z_0(J,J^{\dag}) = \left. \langle 0 | 0 \rangle \right|_{J,J^{\dag}} = \int \mathcal{D}\varphi \mathcal{D} \varphi^{\dag} e^{ i \int d^4x [ \mathcal{L}_0 + J \varphi^{\dag} + J^{\dag} \varphi ] } = \exp{ \left[ i \int d^4x d^4x' J(x) \Delta(x-x') J^{\dag}(x') \right] }
\]
and that 
\[
\exp{ \left[ i \int d^4x \mathcal{L}_1 \right]} = \exp{ \left[ i \int d^4x \left( \frac{-1}{4} Z_{\lambda} \lambda \right)( \varphi^{\dag} \varphi)^2 \right] }
\]
so that 
\[
Z_1 = \exp{ \left[ \frac{ -i}{4} Z_{\lambda} \lambda \int d^4x \left( \frac{1}{i} \delta_{J(x)} \frac{1}{i} \delta_{J^{\dag}(x)} \right)^2 \right]}Z_0(J)
\]

The first order term is $-i Z_{\lambda} \lambda$.  

Consider the second order term in momentum space.  

\begin{gather}
  \frac{1}{2} \left( \frac{ -i  \lambda}{4} \left( \prod_{\alpha= a }^d \int \frac{ d^4k_{\alpha}}{ (2\pi)^4} \frac{ (2\pi)^4}{i } \right) \delta_{J_a} \delta_{J_b^{\dag}} \delta_{J_c} \delta_{J_d^{\dag} } \delta^4(-k_a + k_b - k_c + k_d) \right)^2 \cdot \frac{1}{6!} \left( i \int \frac{d^4k}{ (2\pi)^4} \frac{ \widetilde{J}_k \widetilde{J}_k^{\dag} }{ k^2 + m^2 - i \epsilon } \right)^6 = \label{Eq:16_2-complexscalar_loopvertex01} \\ 
   \begin{gathered} = \frac{1}{2} ( - \lambda^2 ) \prod_{\alpha = a}^d \int  d^4k_{\alpha} \delta^4(-k_a+k_b  - k_c + k_d)  \frac{1}{ (2\pi)^8 } \frac{ \widetilde{J}_{ d} }{ k_d^2 + m^2 } \frac{ \widetilde{J}_{ c}^{\dag} }{ k_c^2 + m^2 } \cdot \\ 
 \cdot     \prod_{\beta = e}^h \int d^4k_{\beta} \delta^4( - k_e + k_f - k_g + k_h ) \frac{1}{(2\pi)^ 8 } \frac{ \widetilde{J}_h}{ k_h^2 + m^2 } \frac{ \widetilde{J}_g^{\dag}}{ k_g^2 + m^2 }  \delta_{J_a} \delta_{J_b^{\dag}} \delta_{J_e } \delta_{J_f^{\dag}} \frac{1}{2} (-1) \int \frac{ d^4k}{ (2\pi)^4} \frac{ d^4l}{ (2\pi)^4} \frac{ \widetilde{J}_k \widetilde{J}_k^{\dag}}{ k^2 + m^2 }  \frac{ \widetilde{J}_l \widetilde{J}_l^{\dag}}{ l^2 + m^2 }   = \end{gathered} \label{Eq:16_2-complexscalar_loopvertex02} \\
   \begin{gathered} = \lambda^2  \int d^4k d^4l d^4k_c d^4k_d \delta^4(-l + k - k_c + k_d) \frac{1}{(2\pi)^8} \frac{ \widetilde{J}_d}{ k_d^2 + m^2 } \frac{ \widetilde{J}_c^{\dag}}{ k_c^2 + m^2 } \cdot \\ 
     \quad \quad \quad \,      \cdot d^4k_g d^4k_h \delta^4(-k + l - k_g + k_h ) \frac{1}{ (2\pi)^8 } \frac{ \widetilde{J}_h \widetilde{J}_g^{\dag}}{ (k_h^2 + m^2 )(k_g^2 + m^2 ) } \frac{1}{ l^2 + m^2 } \frac{1}{ k^2 + m^2 } \frac{ 1 }{ (2\pi)^8 } \end{gathered} \notag
\end{gather}

Note the counting factor of $(4)^2$ from going from Eq. (\ref{Eq:16_2-complexscalar_loopvertex01}) to Eq. (\ref{Eq:16_2-complexscalar_loopvertex01}), due to choosing which functional derivative, for each vertex, will participate in the internal propagator.  

$\langle 0 | T \varphi_{k_3}  \varphi_{k_4}^{\dag} \varphi_{k_1} \varphi_{k_2} | 0 \rangle$ to second order in $\lambda$ is   

\[
\begin{gathered}
   \lambda^2 \int \frac{ d^4 k}{ (2\pi)^4} \frac{1 }{ k^2 + m^2 } \frac{ d^4l}{ (2\pi)^4} \frac{1}{ l^2 + m^2 } \delta^4( - l + k - k_2 + k_1 ) \frac{1}{ k_2^2 + m^2 } \frac{1}{ k_1^2  + m^2 } \delta^4(-k + l - k_3 + k_4) \frac{ 1 }{ k_4^2 + m^2 } \frac{1}{ k_3^2 + m^2 } = \\ 
  = \lambda^2  \int \frac{ d^4 l }{ (2\pi)^4} \frac{ 1 }{ l^2 + m^2 } \frac{ 1 }{ (l+k_1 - k_2)^2  + m^2 } \frac{1}{ (2\pi)^4 } \delta^4(k_1 - k_2 - (k_3 - k_4)  ) \frac{1}{ k_1^2 + m^2 }  \frac{1}{ k_2^2 + m^2 } \frac{1}{ k_3^2 + m^2 } \frac{1}{ k_4^2 + m^2 }   
\end{gathered}
\]

Applying the LSZ formula to amputate the external legs, we have,
\[
\left(  \lambda  \right)^2 \int \frac{ d^4l }{(2\pi)^4} \mathbf{\widetilde{\Delta}}(l^2 ) \mathbf{\widetilde{\Delta}}(l + k_1 - k_2)^2 
\]
There is only 1 other distinct diagram that is allowed: for $\lambda^2$ order terms
\[
\left(  \lambda \right)^2 \int \frac{ d^4l }{(2\pi)^4} \mathbf{ \widetilde{\Delta}}(l^2 ) \mathbf{\widetilde{\Delta} }( l + k_1-  k_3)^2
\]

% This is different from $\phi^4$ theory.  

\hrulefill

So
\[
i \mathbf{V}_4 = -i Z_{\lambda} \lambda + \left(  - i \lambda \right)^2 \left( \frac{ 1 }{i} \right)^2 \int \frac{ d^4l}{ (2\pi)^d} \widetilde{\mathbf{\Delta}}(l^2 ) \left( \widetilde{\mathbf{\Delta}}(l + k_1  - k_2)^2 + \widetilde{\mathbf{\Delta}}( l +k_1 - k_3)^2 \right)
\]
\[
\begin{gathered}
  \frac{ \mathbf{V}_4}{\lambda} = - Z_{\lambda} + \lambda \int \frac{ d^d\overline{l}}{ (2\pi)^d} \int_0^1 dx \lbrace ( \overline{l}^2 - x(1-x) s + m^2 )^{-2}  + (\overline{l}^2 - x(1-x) t + m^2 )^{-2} \rbrace = \\
  = - Z_{\lambda} + \lambda \int_0^1 \frac{dx}{ (4\pi)^2 } \lbrace 2 \left( \frac{2}{ \epsilon }+ 2\ln{ \frac{ \mu}{ m } } \right) - \ln{ \left( \frac{ m^2 - x(1-x)s}{ m^2 } \right) } - \ln{ \left( \frac{ m^2 - x(1-x)t }{ m^2 } \right) } \rbrace
\end{gathered}
\]

For $\mathbf{V}_4 = \lambda$, $Z_{\lambda} = 1 + C$,

\[
\begin{aligned}
  C & = \lambda  \frac{ 1}{ (4\pi)^2 } \lbrace \frac{4}{ \epsilon} + 4 \ln{ \frac{ \mu}{ m} } - \int_0^1 dx \left( \ln{ \left( 1 - \frac{ x ( 1-x) s }{ m^2 } \right) } + \ln{ \left( 1 - \frac{ x(1-x) t}{ m^2 } \right) } \right) \rbrace = \\ 
  & =  \lambda  \frac{1}{ (4\pi)^2 } \lbrace \frac{4}{ \epsilon } + 4 \ln{ \frac{ \mu}{ m }} +2 \rbrace = \boxed{  \lambda  \frac{1}{ (4\pi)^2 } \lbrace \frac{1}{ \epsilon} + \ln{ \frac{ \mu}{ m} } + \frac{1}{2} \rbrace } 
\end{aligned}
\]

% $\varphi \varphi \to \varphi \varphi$ scattering.  
% \[
% \begin{gathered}
%  \left( \frac{-i}{4} Z_{\lambda} \lambda \right)^2 \left( \int d^4x \left( \frac{1}{i} \delta_{J_x} \frac{1}{i} \delta_{J^{\dag}_x} \right)^2 \right)^2 \frac{1}{6!} \left( i \int d^4x d^4x' J_x \Delta_{x-x'} J^{\dag}_{x'} \right)^6 = \\ 
%   = \left( \frac{-i}{4} Z_{\lambda} \lambda \right)^2 \int d^4x d^4y \frac{1}{i} \delta_{J_x} \frac{1}{i} \delta_{J_y} \left( \int d^4x' J_{x'} \Delta_{x' - x} \right)^2 \left( \int d^4y' J_{y'} \Delta_{y'- y} \right)^2 \left( \int d^4x' \Delta_{x-x'} J_{x'}^{\dag} \right) \left( \int d^4y' \Delta_{y-y'} J_{y'}^{ \dag } \right) = \\ 
%  = \frac{ ( - i Z_{\lambda} \lambda)^2}{ 4 i^2 } \int d^4x d^4y \int d^4x' J_{x'} \Delta_{x'-x} \int d^4y' J_{y'} \Delta_{y' - y} \int d^4x'' \Delta_{x-x''} J_{x''}^{\dag} \int d^4y'' \Delta_{y-y''} J_{y''}^{\dag} \Delta_{y-x} \Delta_{y-x}
% \end{gathered}
% \]
% \[
% \Longrightarrow \langle 0 | T\varphi_{x_1} \varphi^{\dag}_{x_2} \varphi_{x_3} \varphi^{\dag}_{x_4} | 0 \rangle = \frac{ (Z_{\lambda} \lambda )^2}{ 4 i^4} \int d^4x d^4y \Delta_{x_1-x}  \Delta_{x_3-y} \Delta_{x-x_2} \Delta_{y-x-4} \Delta_{y-x} \Delta_{y-x}     
% \]

% Apply the LSZ formula developed in Problem 5.1, which was 
% \[
% \begin{gathered}
%  \langle f | i \rangle  = i^{n' + n} \int d^4x_1' e^{-ik_1'x_1'} ( - \partial_{1'}^2 + m_a^2 ) \dots \int d^4x_{m+1'} e^{ + i k_{m+1'} x_{m+1'} } ( - \partial^2_{m+1'} + m_b^2) \dots \\ 
%  \dots \int d^4x_1 e^{ik_1 x_1}( -\partial_1^2 + m_a^2) \dots \int d^4x_{m+1} e^{-ik_{m+1} x_{m+1} } ( -\partial_{m+1}^2 + m_b^2 ) \dots \cdot  \langle 0 | T \varphi(x_1') \dots \varphi^{\dag}(x'_{m+1} ) \dots \varphi^{\dag}(x_1) \dots \varphi(x_{m+1} \dots |  0 \rangle
% \end{gathered}
% \]  

% So then
% \[
% \begin{gathered}
%   \langle f | i \rangle = i^4 \int d^4x_3 e^{-ik_3x_3} ( - \partial_3^2 + m^2) \int d^4x_4 e^{ik_4x_4} ( - \partial_4^2 + m^2) \cdot \int d^4x_2 e^{ik_2x_2} (- \partial_2^2 + m^2 ) \int d^4x_1 e^{-ik_1 x_1 } ( - \partial_1^2 + m^2)  \cdot \\ 
%   \cdot \frac{ (Z_{\lambda} \lambda)^2}{ 4i^4} \int d^4x d^4y \Delta_{x_1 - x} \Delta_{x_3 - y} \Delta_{x - x_2} \Delta_{y - x_4} \Delta_{y - x} \Delta_{y - x}   
% \end{gathered}
% \]

% $x_1,x_3$ can be exchanged and $x_2,x_4$ can be exchanged among this pair.  So there are 3 cases to consider:
% \begin{align}
%  &  e^{-i (k_3 - k_4) y} e^{i (k_2 - k_1) x} \label{Eq:complexscalar_combi_counting} \\ 
%  &  e^{-i (k_1 - k_4) y} e^{i (k_2 - k_3) x} \notag \\ 
%  &  e^{-i (k_3 - k_2) y} e^{i (k_4 - k_1) x} \notag
% \end{align}

% There's a $\times 2$ counting factor if the ``other pair(s)'' are exchanged (e.g. for the 1st. term, Eq. (\ref{Eq:complexscalar_combi_counting}), that's above, $k_1 \leftrightarrow k_3$ and $k_2 \leftrightarrow k_4$ are exchanged simultaneously, and the result is the same).  

% \[
% \begin{gathered}
%  \Longrightarrow \frac{ ( Z_{\lambda} \lambda)^2}{2} \int d^4x d^4y e^{-i(k_3 - k_4)y } e^{i(k_2 -k_1) x} \int \frac{ d^4k}{(2\pi)^4} \frac{ e^{ik(y-x)}}{ k^2 + m^2 - i \epsilon } \int \frac{ d^4l}{ (2\pi)^4} \frac{ e^{i l (y-x)} }{ l^2 + m^2 - i \epsilon }  = \\ 
%  = \frac{ ( Z_{\lambda} \lambda)^2}{ 2} \int d^4k d^4l \frac{1}{k^2 + m^2 - i \epsilon } \frac{ 1 }{ l^2 + m^2 - i \epsilon } \delta^4(k_3- k_4 + k+ l ) \delta^4(k_2 - k_1 -k - l ) = \\ 
%  = \frac{ (Z_{\lambda} \lambda)^2 }{ 2 } \int d^4l \frac{1}{ (k_2 - k_1- l)^2 + m^2 } \frac{ 1 }{ l^2 + m^2} \delta^4(k_2- k_1 + k_3 - k_4) 
% \end{gathered}
% \]

%So indeed, the usual Feynman rules work: there are two internal propagators, momentum is conserved, there is a factor of $(-i \lambda)^2$ for the 2 vertices.  



\section{Other 1PI vertices}

\solutionhead{17.1} Want: 
\[
D_{1234} = x_1 x_4 k_1^2 + x_2 x_4 k_2^2 + x_2 x_3 k_3^2 + x_1 x_3 k_4^2 + x_1 x_2 ( k_1 + k_2)^2 + x_3 x_4 (k_2 + k_3)^2 + m^2 \quad \quad \, (17.3) 
\]
Rewrite using $x_1 + x_2 + x_3 + x_4 = 1$, $k_1 + k_2 + k_3 + k_4 = 0$.  

\[
\begin{gathered}
  \Longrightarrow D_{1234} = x_1 ( 1 - x_1 - x_2 - x_3) k_1^2  + x_2 ( 1 - x_1 - x_2 - x_3) k_2^2 + x_2 x_3 k_3^2 + \\
  + x_1 x_3(-k_1 - k_2 - k_3)^2 + x_1 x_2 (k_1 + k_2)^2 + x_3 ( 1 - x_1 - x_2 - x_3) (k_2 + k_3)^2 + m^2 = \\
  = x_1( 1 - x_1) k_1^2 - x_1(x_2 + x_3) k_1^2 + x_2 (1 - x_2) k_2^2 + x_2( x_1 + x_3) ( - k_2^2) + x_2 x_3 k_3^2 + \\
  + x_1 x_3(k_1^2 + (k_2+k_3)^2 + 2 k_1(k_2 + k_3)) + x_1 x_2( k_1^2 + k_2^2 + 2k_1 k_2) + x_3 (1-x_3)(k_2+k_3)^2 - x_3(x_1 + x_2) (k_2^2 + k_3^2 + 2k_2 k_3) = \\
  = x_1 ( 1 - x_1 )k_1^2 + x_2(1-x_2)k_2^2 - x_2 x_3 (k_2^2 - k_3^2)  + 2 x_1 x_3 k_1 (k_2 + k_3) + 2 x_1 x_2 k_1 k_2 + \\ 
  + x_3(1-x_3)(k_2 + k_3)^2 - x_2 x_3 ( k_2^2 + k_3^2 + 2k_2 k_3) = \\
  = x_1 ( 1 - x_1 )k_1^2 + x_2 (1-x_2)k_2^2 + x_3(1-x_3)(k_2 + k_3)^2  + 2 x_1 x_3 k_1(k_2 + k_3) + 2x_1 x_2 k_1 k_2 - 2x_2 x_3 (k_2 (k_2 + k_3))
\end{gathered}
\]

Calculating $i \mathbf{V}_4$, which includes 3 distinct terms at the $g^4$ order,
\[
i \mathbf{V}_4 = g^4 \int \frac{ d^6l}{ (2\pi)^6} \mathbf{\widetilde{\Delta}}(l - k_1)^2 \mathbf{\widetilde{\Delta}}(l+k_2)^2 \mathbf{\widetilde{\Delta}}(l+k_2 + k_3)^2 \mathbf{\widetilde{\Delta}}(l^2) + (k_3 \leftrightarrow k_2 ) + ( k_3 \leftrightarrow k_4) + \mathcal{O}(g^6) \quad \quad \, (17.1)
\] 

\[
\mathbf{\widetilde{\Delta}}(l - k_1)^2 \mathbf{\widetilde{\Delta}}(l+k_2)^2 \mathbf{\widetilde{\Delta}}(l+k_2 + k_3)^2 \mathbf{\widetilde{\Delta}}(l^2) = \frac{1}{ ( l-k_1)^2 + m^2 }  \frac{1}{ ( l+k_2)^2 + m^2 } \frac{1}{ ( l+k_2 +k_3)^2 + m^2 } \frac{1}{ l^2 + m^2 } 
\]

Recall Feynman's formula
\[
 = \int dF_4 [ x_1 (l-k_1)^2 + x_2 (l+k_2)^2 + x_3(l+k_2 + k_3)^2 + x_4 l^2 + m^2 ]^{-4} 
\]
where
\[
\int dF_4 = 3! \int_0^1 dx_1 dx_2 dx_3 dx_4 \delta(x_1 + x_2 + x_3 + x_4 -1 )
\]

\begin{align}
  \Longrightarrow [ x_1 ( l-k_1)^2 + x_2 ( l+k_2)^2 + x_3(l+k_2 + k_3)^2 + (1-x_1 - x_2 - x_3)l^2 + m^2 ] = \notag \\
  = l^2 + x_1 ( k_1^2 - 2lk_1) + x_2 (k_2^2 + 2lk_2) + x_3( (k_2+ k_3)^2 + 2l (k_2 +k_3) ) + m^2 \label{Eq:17.1-1pIvertex}
\end{align}

Now
\[
\begin{aligned}
  q & = l - x_1 k_1 + x_2 k_2 + x_3 (k_2 + k_3) \\
  q^2 & = l^2 +  x_1^2 k_1^2 + x_2^2 k_2^2 + x_3^2 (k_2+ k_3)^2 + \\
  & \phantom{ = } + 2l (-x_1 k_1 + x_2 k_2 + x_3 (k_2 + k_3) ) + -2 x_1 x_2 k_1 k_2 - 2x_1 x_3 (k_2 + k_3) + 2x_2 k_2 x_3 (k_2 +k_3)
\end{aligned}
\]

Then the left-hand side, or the entire expression, of Eq. (\ref{Eq:17.1-1pIvertex}), is equal to 

\[
q^2 + x_1(1-x_1) k_1^2 + x_2(1-x_2) k_2^2 + x_3(1-x_3) (k_2 + k_3)^2 + 2x_1 x_2 k_1 k_2 + 2 x_1 x_3 k_1 (k_2 + k_3) - 2x_2 x_3 k_2 (k_2 + k_3)
\]
\[
\Longrightarrow \boxed{ D_{1234} = x_1 ( 1 - x_1 ) k_1^2 + x_2 ( 1 - x_2 )k_2^2 + x_3(1-x_3) (k_2 + k_3)^2 + 2 x_1 x_2 k_1 k_2 + 2x_1 x_3 k_1(k_2+ k_3) - 2 x_2 x_3 k_2 ( k_2 + k_3) }
\]

\section{Higher-order corrections and renormalizability}

\solutionhead{18.1} 
\begin{itemize}
\item[(a)] Recall that the action must be dimensionless since it's in the exponential.  Recall the action: $S \sim i \int d^dx i \overline{\Psi} \gamma^{\mu} \partial_{\mu} \Psi$.  where $\overline{\Psi} = \Psi^{\dag} \gamma^0$. 

Recall also the following  mass dimensions: \\
$ [ d^d x ] = - d $ \\
$ [ \overline{\Psi} \gamma^{\mu} \partial_{\mu} \Psi ] = 2 [ \Psi ] +1$

So
\[
-d + 2 [ \Psi ] + 1 = 0 \text{ or }  \boxed{ [ \Psi ] = \frac{d - 1}{2} }
\]
\item[(b)]  Then to make the action dimensionless,
  \[
    [ g_n ( \overline{\Psi} \Psi )^n ] = [ g_n ] + ( 2 [ \Psi ]) n = [g_n ] + 2n \frac{d- 1}{2} = [g_n ] + n (d-1) = d
    \]
So 
\[
\boxed{ [ g_n ] = n + d(1- n) }
\]
  \item[(c)]  Likewise,
\[
[ g_{m,n} \phi^m ( \overline{\Psi} \Psi )^n  ] = [ g_{m,n} ] + m \frac{1}{2} ( d-2)+ n (2) \left( \frac{d-1}{2} \right)  = [ g_{m,n} ] + \frac{ m (d-2) }{2} + n (d-1) =d
\]
\[
\Longrightarrow [ g_{m,n} ] = d ( 1 - n - \frac{m}{2} ) + (n+m)
\]
  \item[(d)] For $d=4$, 
\[
[ g_n ] = n + 4(1- n) \geq 0 \text{ or } -3n + 4 \geq 0 \text{ or } \frac{4}{3} \geq n \text{ or } \boxed{ n =1 }
\]
\[
[ g_{m,n} ] = 4 ( 1 - n  - \frac{m}{2} ) + n + m = 4 -3n - m \geq 0  \text{ so } \boxed{ n = m = 1 }
\]
\end{itemize}



\section{Perturbation theory to all orders}

\section{Two-particle elastic scattering at one loop}

\subsection*{Problems}

\problemhead{20.1} Verify eq. (20.17).

\problemhead{20.2} Compute the $O(\alpha)$ correction to the two-particle scattering amplitude \emph{at threshold}, that is, for $s= 4m^2$ and $t=u=0$, corresponding to zero three-momentum for both the incoming and outcoming particles.  

\solutionhead{20.2} 
What we want to calculate is 
\[
i \mathcal{T}_{1-loop} = \frac{1}{i} \left( [ i \mathbf{V}_3(s) ]^2  \mathbf{\widetilde{\Delta}}(-s) + [ i \mathbf{V}_3(t)]^2 \mathbf{\widetilde{\Delta}}(-t) + [ i \mathbf{V}_3(u)]^2 \widetilde{\mathbf{\Delta}}(-u) \right) + i \mathbf{V}_4(s,t,u)
\]
for $s = 4m^2$, $t= u =0$.  

Recall 
\[
 \widetilde{\mathbf{\Delta}}(-s) = \frac{1}{ - s + m^2 - \Pi(-s) }
\]
\[
\Longrightarrow \widetilde{\mathbf{\Delta}}(-4m^2 ) = \frac{1}{ -3 m^2 - \Pi( -4m^2) }, \quad \quad \, \widetilde{\mathbf{\Delta}}(0) = \frac{1}{ m^2 - \Pi(0) }
\]



%\emph{at threshold}: $s = -(k_1+k_2)^2 = - ( -(E_1 + E_2)^2) = (m + m)^2 = 4m^2$  \\
%\phantom{ at threshold: } $t= u =0$

From (20.4), Srednicki,
\[
\Pi(-s) = \Pi ( -4m^2) = \frac{1}{2} \alpha \int_0^1 dx D_2(4m^2) \ln{ \left( \frac{ D_2(4m^2)}{D_0} \right)} - \frac{1}{12} \alpha (-4m^2 + m^2)
\]

From (20.8), (20.9), Srednicki, the definitions
\[
\begin{aligned}
  & D_2(s) = -x(1-x)s+m^2 \\ 
  & D_2(4m^2) = m^2 (-4x + 4x^2 + 1) = m^2 ( 2x-1)^2 \\ 
  & D_0 = m^2 [1 - x(1-x)] = m^2 [1-x +x^2] = m^2 \left( (x-\frac{1}{2} )^2 +  \frac{3}{4} \right) = \frac{1}{4} m^2 ((2x-1)^2 + 3)
\end{aligned}
\]

\[
\Pi(0)  = \frac{ m^2 \alpha}{2} \int_0^1 dx \ln{ \left( \frac{1}{ 1 - x + x^2} \right) } - \frac{ \alpha m^2}{ 12} = -\frac{ \alpha m^2}{2} \left( \frac{1}{6} + \int_0^1 dx \ln{ (x^2 - x+1) } \right)= \boxed{ \frac{ \alpha m^2}{2} \left( \frac{ 11 }{6} - \frac{ \pi }{\sqrt{3}} \right) }
\]

\[
\begin{gathered}
  \Pi(-4m^2) = \frac{ \alpha m^2}{2} \left( \frac{1}{2} + \int_0^1 dx (2x-1)^2 \ln{ \left( \frac{ 4(2x-1)^2 }{ (2x-1)^2 + 3 } \right) } \right) = \frac{ \alpha m^2}{2} \left( \frac{1}{ 2} + \int_{-1}^1 \frac{du}{2} u^2 \ln{ \left( \frac{ 4 u^2 }{ u^2 + 3 } \right) } \right) = \\ 
  = \frac{ \alpha m^2}{ 4 } \left( 1 + \frac{2}{3} \ln{4} - \int_{-1}^1 du u^2 \ln{ \left( 1 + \frac{3}{ u^2 } \right) } \right) = \frac{ \alpha m^2 }{4} \left( 1 + \frac{2}{3} \ln{4} - \left. \left( \frac{ u^2 }{3 }\ln{ \left( 1 + \frac{3}{u^2 } \right) } + 2 u -2 \sqrt{3} \arctan{ \left( \frac{u}{\sqrt{3}} \right) } \right) \right|_{-1}^1 \right) = \\
  = \frac{  \alpha m^2 }{4} \left( 1 - 4 + 4\sqrt{3} \frac{ \pi}{6} \right) = \boxed{ \alpha m^2 \left( \frac{-3}{4} + \frac{ \pi \sqrt{3}}{6} \right) }
\end{gathered}
\]
since 
\[
\left( \frac{u^3}{3} \ln{ \left( 1 + \frac{ 3}{u^2 } \right) } \right)' = u^2 \ln{ \left( 1 + \frac{3}{u^2 } \right) } + \frac{u^3}{3} \frac{1}{ 1 + 3/u^2} \left( \frac{ -6}{u^3} \right) = u^2 \ln{ \left( 1 + \frac{3}{u^2 } \right) } + -2 \left( 1 - \frac{ 3}{ u^2 + 3} \right) 
\]




Consider the 3 point vertices.  

Recall that \[
\frac{ \mathbf{V}_3(s)}{ g} = 1 - \frac{ \alpha}{2} \int dF_3 \ln{ \left( \frac{D_3(s)}{ m^2} \right) }  \quad \quad \quad \, (20.5)
\]
\[
D_3(s) = -x_1 x_2 s + [ 1 - (x_1 + x_2)x_3] m^2 \quad \quad \quad \, (20.10)
\]

For $x_3 = 1 - (x_1 + x_2)$,
\[
\begin{aligned}
  & D_3(0) = (1 - (x_1 + x_2) + ( x_1 + x_2)^2)m^2 \\ 
  & D_3(4m^2) = -4 m^2 x_1 x_2 + (1 - (x_1 + x_2) + (x_1 + x_2)^2)m^2 = m^2 ( 1 - (x_1 + x_2) + (x_1-x_2)^2 )
\end{aligned}
\]
Then consider
\[
\begin{gathered}
  \int_0^1 dx \int_0^{1-x} dy \ln{ ( 1 - x-  y + (x+y)^2) } = \int_0^1 dx \int_0^{1-x} dy \ln{ ( y^2 + y(2x-1) + x^2 - x + 1 ) } = \\
  = \int_0^1 dx \int_0^{1-x} dy \ln{ \left( \left( y - \left( \frac{2x - 1 }{2} \right) \right)^2 + \frac{3}{4} \right) }
\end{gathered}
\]
\[
\begin{gathered}
  \int_0^1 dx \int_0^{1-x} dy \ln{ ( 1 - x - y + x^2 - 2xy + y^2) } = \int_0^1 dx \int_0^{1-x} dy \ln{ (y^2 - y (2x+  1 ) + x^2 - x+1 ) } = \\ 
  = \int_0^1 dx \int_0^{1-x} dy \ln{ \left( \left( y - \left( \frac{2x + 1 }{2} \right) \right)^2 - 2x + \frac{3}{4} \right) }
\end{gathered}
\]








\section{The quantum action}

\subsection*{Problems}

\problemhead{21.1} Show that 
\[
\Gamma(\varphi) = W(J_{\varphi}) - \int d^dx J_{\varphi} \varphi, \quad \quad \quad \, (21.20)
\]
where $J_{\varphi}(x)$ is the solution of 
\[
\frac{ \delta }{ \delta J(x) } W(J) = \varphi(x) \quad \quad \quad \, (21.21)
\]
for a specified $\varphi(x)$.  

\solutionhead{21.1}

From the main result of the section, by stationary phase method and $W(J) = W_{\Gamma,L=0}(J)$,
\[
W(J) = \Gamma(\varphi_J) + \int d^dxJ \varphi_J \quad \quad \, (21.12) 
\]
where 
\[
\phantom{W(J) = } \frac{ \delta \Gamma(\varphi) }{ \delta \varphi(x)} = - J(x) \quad \quad \, (21.10)
\]
and $J(x)$ is a specified source.  

Then applying $\frac{ \delta}{ \delta J(x)}$, since we could vary $J(x)$ through varying $\varphi$,
\[
\begin{gathered}
  \frac{ \delta W(J)}{ \delta J(x) } = \frac{ \delta \Gamma(\varphi_J)}{ \delta J(x)} + \int d^dy \frac{ \delta }{\delta J(x)} ( J(y) \varphi_J(y) ) = \\ 
  = \int d^d y \frac{ \delta \Gamma(\varphi_J)}{ \delta \varphi_J(y)} \frac{ \delta \varphi_J(y)}{ \delta J(x)} + \int d^dy \lbrace \delta(x-y)\varphi_J(y) + J(y) \frac{ \delta \varphi_J(y)}{ \delta J(x)} \rbrace = \\
  =\int d^4y \left( - J(y) \frac{ \delta \varphi_J(y)}{\delta J(x) } + J(y) \frac{ \delta \varphi_J(y)}{ \delta J(x)} \right) + \varphi(x) = \varphi(x)
\end{gathered}
\]
So 
\[
\frac{ \delta W(J)}{ \delta J(x)} = \varphi(x)
\]
If so, we can say
\[
\boxed{ \Gamma(\varphi) = W(J_{\varphi}) - \int d^dx J_{\varphi}\varphi }
\]
for 
\[
\phantom{\Gamma(\varphi)} \boxed{ \frac{ \delta W(J)}{ \delta J(x)} = \varphi(x) }
\]
so that $\varphi$ is specified and $J = J_{\varphi}(x)$ is a solution to $\frac{ \delta W(J)}{ \delta J(x)} = \varphi(x)$.  

\solutionhead{21.2} 
\begin{itemize}
\item[(a)]
From the main result, that the sum of connected diagrams with sources $iW(J)$, 
\[
W(J) = \Gamma(\varphi_J) + \int d^dx J \varphi_J
\]
The action was given to be invariant.  

Given how the field transforms, \\
$ \varphi_a(x) \to \int d^dy R_{ab}(xy) \varphi_b(y)$

and how the current transforms, then
\[
J_a \varphi_J \to \int d^d y_1 J_b(y_1) R_{ba}(y_1,x)  \int d^dy_2 R_{Ac}(x,y_2) \varphi_c(y_2) = \int d^dy_1 \int d^dy_2 J_b(y_1) R_{ba}(y_1,x) R_{ac}(x,y_2) \varphi_c(y_2) = J_b(x) \varphi_b(x)
\]
since $\delta^d(x-y_1) \delta^d(x-y_2)$.  
\item[(b)] $\Gamma(\varphi) = W(J_{\varphi}) - \int d^4x J_{\varphi}\varphi$
\[
\int d^dx J_{\varphi} \varphi \to \int d^dx \int d^dy_1 J_{b\varphi}(y_1) R_{ba}(y_1,x) \int d^4y_2 R_{ac}(x,y_2) \varphi_c(y_2) = \int d^dx J_{b} \varphi_b(x)
\]
$W(J_{\varphi})$ was shown to be invariant under $J_a(x) \to \int d^dy J_b(y) R_{ba}(y,x)$  (21.23).  Thus $\Gamma(\varphi)$ is invariant.  
\end{itemize}








\section{Continuous symmetries and conserved currents}


\[
\begin{aligned}
  & \mathcal{L} = \mathcal{L}( \varphi, \partial_{\mu} \varphi) \\ 
  & \mathcal{L} = \frac{ \partial \mathcal{L}}{ \partial \varphi_a} \delta \varphi_a + \frac{ \partial \mathcal{L} }{ \partial ( \partial_{\mu} \varphi_a ) } \delta  \partial_{\mu} \varphi_a = \frac{ \partial \mathcal{L}}{ \partial \varphi_a} \delta \varphi_a + \frac{ \partial \mathcal{L}}{ \partial ( \partial_{\mu} \varphi_a ) } \partial_{\mu} \delta \varphi_a
\end{aligned}
\]
\[
\begin{aligned}
   S & = \int d^4x \mathcal{L} \\ 
   \delta S & = \int d^4x \lbrace \frac{ \partial \mathcal{L} }{ \partial \varphi_a } \delta \varphi_a + \frac{ \partial \mathcal{L} }{ \partial ( \partial_{\mu} \varphi_a ) } \delta \partial_{\mu} \varphi_a \rbrace = \int d^dx \lbrace \frac{ \partial \mathcal{L}}{ \partial \varphi_a } \delta \varphi_a - \left( \partial_{\mu} \frac{ \partial \mathcal{L} }{ \partial ( \partial_{\mu} \varphi_a ) } \right) \delta \varphi_a \rbrace + \int_{\partial V} d\sigma \frac{ \partial \mathcal{L}}{ \partial ( \partial_{\mu} \varphi_a  ) } \delta \varphi_a 
\end{aligned}
\]

So

\[
\mathcal{L} = \frac{ \partial \mathcal{L}}{ \partial \varphi_a} \delta \varphi_a + \partial_{\mu} \left( \frac{ \partial \mathcal{L} }{ \partial ( \partial_{\mu} \varphi_a ) } \delta \varphi_a \right) - \left( \partial_{\mu} \frac{ \partial \mathcal{L} }{ \partial ( \partial_{\mu} \varphi_a  )} \right) \delta \varphi_a = \partial_{\mu} \left( \frac{ \partial \mathcal{L}}{ \partial ( \partial_{\mu} \varphi_a ) } \delta \varphi_a \right) + \frac{ \delta S }{ \delta \varphi_a } \delta \varphi_a
\]

\textbf{ \emph{Complex} Scalar field }

cf. Srednicki.  
\[
\begin{aligned}
  & \mathcal{L} = - \partial^{\mu} \varphi^{\dag} \partial_{\mu} \varphi - m^2 \varphi^{\dag} \varphi - \frac{1}{4} \lambda (\varphi^{\dag} \varphi)^2  \quad \quad \quad \, & (22.9) \\
  & \mathcal{L} = - \frac{1}{2} \partial^{\mu} \varphi_1 \partial_{\mu} \varphi_1 - \frac{1}{2} \partial^{\mu} \varphi_2 \partial_{\mu} \varphi_2 - \frac{1}{2} m^2 ( \varphi_1^2 + \varphi_2^2) - \frac{1}{16} \lambda ( \varphi_1^2 + \varphi_2^2)^2 
\end{aligned}
\]
\[
\begin{gathered}
  \frac{ \partial \mathcal{L} }{ \partial A_{\mu} } - \partial^{\nu} \frac{ \partial \mathcal{L} }{ \partial ( \partial^{\nu} A_{\mu} ) } = 0 \\
  \Longrightarrow - \partial^2 \varphi = - m^2 \varphi - \frac{1}{2} \lambda |\varphi|^2 \varphi \text{ or } \partial^2 \varphi  -m^2 \varphi = \frac{\lambda}{2} |\varphi |^2 \varphi
\end{gathered}
\]
For $\varphi  = \frac{1}{\sqrt{2}} ( \varphi_1 + i \varphi_2 )$ 
\[ \Longrightarrow
\begin{aligned}
  - \frac{1}{2} \partial^2 \varphi_1 & = - m^2 \varphi_1 - \frac{\lambda}{4} ( \varphi_1^2 + \varphi_2^2) \varphi_1 \\ 
  \frac{1}{2} \partial^2 \varphi_2 & = m^2 \varphi_2 + \frac{\lambda}{4} ( \varphi_1^2 + \varphi_2^2) \varphi_2
\end{aligned}
\]
Note $|\varphi|^2 = \frac{1}{2} ( \varphi_1^2 + \varphi_2^2 )$.   \\

Note, possibly this is correct and not what's in the text:
\[
\mathcal{L} = -\frac{1}{2} \partial^{\mu} \varphi_i \partial_{\mu} \varphi_i - \frac{1}{4}m^2 (\varphi_i \varphi_i ) - \frac{1}{16} \lambda ( \varphi_i \varphi_i)^2 = -\frac{1}{2} \partial^{\mu} \varphi_1 \partial_{\mu} \varphi_1  -\frac{1}{2} \partial^{\mu} \varphi_2 \partial_{\mu} \varphi_2 - \frac{1}{4} m^2 (\varphi_1^2 + \varphi_2^2) - \frac{1}{16} \lambda ( \varphi_1^2 + \varphi_2^2)^2
\]



\problemhead{22.1} For the Noether current of eq. (22.6), and assuming that $\delta \varphi_a$ does not involve time derivatives, use the canonical commutation relations to show that 
\[
[\varphi_a, Q ] = i \delta \varphi_a  \quad \quad \quad \, (22.41)
\]
where $Q$ is the Noether charge.  

\solutionhead{22.1} Now 
\[
Q \equiv \int d^3 x j^0(x) \quad \quad \quad \, \text{ Noether charge }
\]
\[
\begin{aligned}
  \left[\varphi_a, Q \right]  & = \int d^3 x [ \varphi_a, j^0(x) ] = \int d^3 x [ \varphi_a, \frac{ \partial \mathcal{L}(x) }{ \partial (\partial_0 \varphi_a) } \delta \varphi_a ] = \\
  & = \int d^3 x [ \varphi_a(x'), \Pi(x) ] \delta \varphi_a = \int d^3 x i \delta^3(\varphi{x}' - \varphi{x}) \delta \varphi_a =  i \delta \varphi_a
\end{aligned}
\]

\problemhead{22.2} Use the canonical commutation relations to verify eq. (22.38).  

\solutionhead{22.2}  Recall from (22.35) Srednicki,
\[
P^{\mu} = \int d^3x T^{0\mu}(x)
\]
and recall
\[
T^{0 \mu} = - \frac{ \partial \mathcal{L}}{ \partial (\dot{\varphi}_a)} \partial^{\mu} \varphi_a = -\Pi_a \partial^{\mu} \varphi_a
\]

\[
\begin{aligned}
  \left[ \varphi_a(x), P^{\mu} \right] & = \int d^3x' [ \varphi_a(x), - \Pi_a \partial^{\mu} \varphi_a(x') ] = + \int d^3x' [ \Pi_a \partial^{\mu} \varphi_a(x'), \varphi_a(x) ] = \\
  & = \int d^3 x' \lbrace \Pi_a(x') [ \partial_{x'}^{\mu} \varphi_a(x') , \varphi_a(x) ] + [ \Pi_a, \varphi_a] \partial^{\mu} \varphi_a \rbrace 
\end{aligned}
\]
Now 
\[
\left[ \partial_{x'}^{\mu} \varphi_a(x'), \varphi_a(x) \right] = \partial_{x'}^{\mu} [ \varphi_a(x'), \varphi_a(x)] = 0 
\]
so then
\[
\boxed{ \left[ \varphi_a(x), P^{\mu} \right]  = -i \partial^{\mu} \varphi_a(x)  }
\] 

\problemhead{22.3} \begin{itemize} \item[(a)] With $T^{\mu \nu}$ given by eq. (22.31), compute the equal-time $(x^0 = y^0)$ commutators $[T^{00}(x), T^{00}(y)]$, $[T^{0i}(x), T^{00}(y)]$, and $[T^{0i}(x), T^{0j}(y)]$.
\item[(b)] Use your results to verify eqs. (2.17), (2.19), and (2.20). 
\end{itemize}

\solutionhead{22.3} \begin{itemize}
\item[(a)]  (20090511) INCOMPLETE.  

Recall that 
\[
\begin{aligned}
  & T^{\mu \nu} = \partial^{\mu} \varphi_a \partial^{\nu} \varphi_a + g^{\mu \nu} \mathcal{L} \\ 
  & \mathcal{L}  = - \frac{1}{2} \partial^{\mu} \varphi_a \partial_{\mu} \varphi_a - V(\varphi) \\ 
  & \Pi_a = \partial_0 \varphi_a \\ 
  & T^{00} = \mathcal{H} \\ 
  & T^{0j} = \partial^0 \varphi_a \partial^j \varphi_a = - \Pi_a \nabla^j \varphi_a
\end{aligned}
\]

So then, for $g^{00} =-1$,
\[
\begin{aligned}
  \mathcal{H} & = T^{00} = + \Pi_a^2 + \mathcal{L} = + \Pi_a^2 + \frac{1}{2} ( - \Pi_a^2 + (\nabla \varphi_a)^2 ) + V(\varphi) = \\
 & = \frac{1}{2} \Pi_a^2 + \frac{1}{2} ( \nabla \varphi_a)^2 + V(\varphi)
\end{aligned}
\]

Assume $i,j \neq 0$.  
\[
\begin{aligned}
  \left[ T^{i0}(x), T^{0j}(y) \right] & = [ - \Pi_a(x) \nabla^i \varphi_a(x), -\Pi_a(y) \nabla^j \varphi_a(y) ] = \Pi_x [ \nabla^i \varphi_x , \Pi_y \nabla^j \varphi_y ] + [ \Pi_x, \Pi_y \nabla^j \varphi_y] \nabla^i \varphi_x = \\ 
  & = - \Pi_x \left( \Pi_y [ \nabla^j \varphi_y, \nabla^i \varphi_x] + [ \Pi_y, \nabla^i \varphi_x ] \nabla^j \varphi_y  \right) - \left( \Pi_y [ \nabla^j \varphi_y, \Pi_x ] + [ \Pi_y, \Pi_x ] \nabla^j \varphi_y \right) \nabla^i \varphi_x = \\
  &= \Pi_x ( \Pi_y [ \nabla^i \varphi_x, \nabla^j \varphi_y] + [ \nabla^i \varphi_x, \Pi_y] \nabla^j \varphi_y ) + \Pi_y [ \Pi_x, \nabla^j \varphi_y ] \nabla^i \varphi_x
\end{aligned}
\]

What to do with the spatial derivatives on the fields $\varphi$?

Consider 

\[
\begin{aligned}
  \left[ \partial^i \varphi_x, \Pi_y \right] f & = ( \partial^i \varphi_x) \Pi_y f - (\Pi_y \partial^i \varphi_x) f = \partial^{x_i} ( \varphi_x \Pi_y - \Pi_y \varphi_x ) f = \\ 
  & = \partial^{x_i} [ \varphi_x, \Pi_y] f
\end{aligned}
\]
($f$ is just a single or multiple particle state, labeled by quantum numbers number of particles and total momentum $k$, and not a function of position $x$).  

\[
\Longrightarrow \left[ \nabla^i \varphi_x, \Pi_y \right] = \nabla^i i \delta^3(x-y)
\]

Likewise,
\[
\begin{aligned}
  [\nabla^i \varphi_x, \nabla^j \varphi_y ]f & = \nabla^i \varphi_x \nabla^j \varphi_y f - \nabla^j \varphi_y \nabla^i \varphi_x f = \left( \partial^{x_i} \varphi_x \partial^{y_j} \varphi_y - \partial^{y_j} \varphi_y  \partial^{x_i} \varphi_x \right) f = \\
  & = \partial^{x_i} \partial^{y_j} [\varphi_x, \varphi_y] f = 0 
\end{aligned}
\]
\[
\begin{gathered}
  \Longrightarrow \Pi_x ( (\nabla^i i \delta^3(x-y) ) \nabla^j \varphi_y ) + \Pi_y ( \nabla^{y_j} (-i \delta^3(x-y)) ) \nabla^i \varphi_x = \\
  = i \lbrace \Pi_x ( \nabla^{x_i} \delta^3(x-y) ) \nabla^j \varphi_y + - \Pi_y (\nabla^{y_j} \delta^3(x-y) ) \nabla^i \varphi_x \rbrace = \\
  = i \lbrace \Pi_x \nabla^j \varphi_y \nabla^{x_i} \delta^3(x-y) - \Pi_y \nabla^i \varphi_x \nabla^{y_j} \delta^3(x-y) \rbrace
\end{gathered}
\]
\item[(b)]
\end{itemize} 




\section{Discrete symmetries: $P$, $T$, $C$ and $Z$ }

\section{Nonabelian symmetries}

\problemhead{24.1} Show that $\theta_{ij}$ in eq. (24.4) must be antisymmetric if $R$ is orthogonal.  

\solutionhead{24.1} $R$ orthogonal.  
\[
\Longrightarrow 
\begin{aligned}
  R_{ik} R^T_{kj} & = \delta_{ij} = \\
  & = R_{ik} R_{jk} = (\delta_{ik} + \theta_{ik} + O(\theta^2) )( \delta_{jk} + \theta_{jk} + O(\theta^2) ) = \\
  & = \delta_{ij} + \delta_{ik} \theta_{jk} + \theta_{ik} \delta_{jk} + O(\theta^2)
\end{aligned}
\]
\[
\Longrightarrow \boxed{ - \theta_{ji} = \theta_{ij} }
\]

\problemhead{24.2} By considering the $SO(N)$ transformation $R^{'-1} R^{-1} R'R$, where $R$ and $R'$ are independent infinitesimal $SO(N)$ transformations, prove eq. (24.7).  

\solutionhead{24.2} By the property of $SO(N)$ being a group, $R^{'-1} R^{-1} R'R$ is also an orthogonal $N \times N$ matrix of determinant 1: 
\[
R^{'-1} R^{-1} R'R = R'' = \delta + \theta''
\]
Recall that 
\[
\begin{aligned}
  & R_{ij} = \delta_{ij} + \theta_{ij} \\ 
  & (R^{-1})_{ij} = (R^T)_{ij} = R_{ji} = \delta_{ji} + \theta_{ji} = \delta_{ij} - \theta_{ij} \\ 
  & \begin{aligned} (R^{'-1} R^{-1} R' R)_{ij} & = R^{'-1}_{ik} (R^{-1} R' R)_{kj} = R^{'-1}_{ik} R^{-1}_{kl} (R'R)_{lj}  = R^{'-1}_{ik} R^{-1}_{kl} R'_{lm} R_{mj} = R'_{ki} R_{lk} R'_{lm} R_{mj}
\end{aligned}
\end{aligned}
\]

So now
\[
(R^{'-1} R^{-1} R'R)_{ij} = R''_{ij} = \delta_{ij} + \theta''_{ij} + O(\theta^{''2})
\]
with $\theta''_{ij} = -i \theta^c(T^c)_{ij}$.  

Keeping terms of order $\theta^2$ (which I found in retrospect in first attempting the problem)
\[
\begin{gathered}
(R^{'-1} R^{-1} R' R)_{ij} = R^{'-1}_{ik} R^{-1}_{kl} R'_{lm} R_{mj} = R'_{ki} R_{lk} R'_{lm} R_{mj} = \\
    = (\delta_{ik} + \theta'_{ki} + O(\theta^{'2}) )(\delta_{kl} + \theta_{lk} + O(\theta^2))(\delta_{lm} + \theta'_{lm} + O(\theta^2))(\delta_{mj} + \theta_{mj} + O(\theta^2)) =  \\
    = (\delta_{il} + \theta_{li} + \theta'_{li} + \theta_{ki}' \theta_{lk} )( \delta_{lj} + \theta_{lj} + \theta'_{lj} + \theta'_{lm} \theta_{mj} )  = \\
    = \delta_{ij} + \theta_{ij} + \theta'_{ij} + \theta'_{im} \theta_{mj} + \theta_{ji} + \theta_{li} \theta_{lj} + \theta_{li} \theta'_{lj} + \theta_{ji}' + \theta_{li}' \theta_{lj} + \theta'_{li} \theta'_{lj} + \theta'_{ki} \theta_{jk} = \\
    = \delta_{ij} + \theta_{li} \theta_{lj} + \theta_{li} \theta_{lj}' + \theta_{ki}' \theta_{jk} + \theta'_{li} \theta'_{lj}
\end{gathered}
\]

So $O(\theta^2)$ terms were kept to equate to $\theta''_{ij}$.  

Since it was defined that 
\[
\theta_{jk} = -i \theta^a (T^a)_{jk}
\]
so then
\[
\begin{aligned}
  & \theta_{li} \theta_{lj}' = - \theta_{il} \theta'_{lj} = - - i \theta^a(T^a)_{il} -i \theta^b (T^b)_{lj} = (\theta^a \theta^b) T^a_{il} T^b_{lj} = (\theta^a \theta^b) (T^a T^b)_{ij}  \\ 
  & \theta_{ki}' \theta_{jk} = + \theta'_{ik} \theta_{kj} = - \theta^b(T^b)_{ik} \theta^a (T^a)_{kj} = - \theta^a \theta^b (T^b T^a)_{ij}
\end{aligned}
\]
\[
\theta_{li} \theta_{lj} = - \theta_{il} \theta_{lj} = - - i \theta^a (T^a)_{il} - i \theta^a (T^a)_{lj} = (\theta^a)^2 (T^a)^2_{ij}
\]
\[
\Longrightarrow \theta_{li} \theta_{lj}' + \theta_{ki}' \theta_{jk} = \theta^a \theta^b [T^a, T^b]_{ij}
\]

Better yet, if we included all $\theta^2$ order terms from the beginning, 

\[
\begin{gathered}
  R'_{ki} R_{lk} R'_{lm} R_{mj} \\
\begin{aligned}
  & = \left( \delta_{ki} + \beta_{ki} + \frac{ \beta_{kn} \beta_{ni} }{ 2} \right) \left( \delta_{lk} + \alpha_{lk} + \frac{ \alpha_{lp} \alpha_{pk} }{2} \right)\left( \delta_{lm} + \beta_{lm} + \frac{ \beta_{lq} \beta_{qm} }{ 2} \right) \left( \delta_{mj} + \alpha_{mj} + \frac{ \alpha_{mr} \alpha_{rj} }{2} \right) =  \\
  & =  \left( \delta_{li} + \alpha_{li} + \frac{ \alpha_{lp} \alpha_{pi} }{2} + \beta_{li} + \beta_{ki}\alpha_{lk} + \frac{  \beta_{ln} \beta_{ni} }{2} \right) \left( \delta_{lj} + \alpha_{lj} + \frac{ \alpha_{lr} \alpha_{rj} }{2} + \beta_{lj} + \beta_{lm}\alpha_{mj} + \frac{  \beta_{lq} \beta_{qj} }{2} \right) = \\
  & = \delta_{ij} + \alpha_{ij} + \frac{ \alpha_{ir} \alpha_{rj} }{2}  + \beta_{ij} + \beta_{im}\alpha_{mj} + \frac{ \beta_{iq} \beta_{qj}}{2} + \alpha_{ji} + \alpha_{li} \alpha_{lj} + \alpha_{li} \beta_{lj} + \frac{ \alpha_{jp} \alpha_{pi} }{2} + \beta_j + \beta_{li} \alpha_{lj} + \beta_{li} \beta_{lj} + \beta_{ki} \alpha_{jk} + \frac{ \beta_{jn} \beta_{ni}}{2} = \\
  & = \delta_{ij} + \beta_{im} \alpha_{mj} -\alpha_{il} \beta_{lj} - \beta_{il}\alpha_{lj} + \beta_{ik} \alpha_{kj}
\end{aligned}
\end{gathered}
\]

\[
\begin{gathered} \Longrightarrow \beta_{im} \alpha_{mj} - \alpha_{il} \beta_{lj} = - i \beta(T^b)_{im}(-i \alpha(T^a)_{mj} ) - (-i \alpha(T^a)_{il} )(-i \beta (T^b)_{lj}) = \alpha \beta (T^a)_{il} (T^b)_{lj} - \alpha\beta (T^b)_{im} (T^a)_{mj} = \\
  = (\alpha \beta) (T^a T^b - T^b T^a)_{ij} = (\alpha \beta) [ T^a, T^b]_{ij} 
\end{gathered}
\]

Now since a group is closed under group multiplication, 
\[
  R'_{ki} R_{lk} R'_{lm} R_{mj} = \delta_{ij} + \gamma_{ij} + \delta_{ij} + - i \theta^c (T^c)_{ij}
\]

\[
\Longrightarrow \boxed{ [ T^a, T^b]_{ij} = i \frac{ (- \theta^c)}{\alpha \beta} T^c = i f^{abc} T^c }
\]

\solutionhead{24.3}
\begin{enumerate}
\item[(a)] Recall \\
$\phi_i(x) \to R_{ij} \phi_j(x)$ \quad (24.3) $N$ real scalar fields $\varphi_i$ \\
$R_{ij} = \delta_{ij} + \theta_{ij} + O(\theta^2)$ \quad (24.4) \\
$\theta_{jk} = -i \theta^a (T^a)_{jk}$  \quad (24.6) $\theta^a$ \, $\frac{1}{2} N ( N-1)$ real, infinitesimal parameters.   \\

So then 
\[
R_{ij} \varphi_j \simeq (\delta_{ij} + \theta_{ij} ) \varphi_j = \phi_i + -i \theta^a (T^a)_{ij} \varphi_j = \varphi_i + -i \theta^a ( T^a \varphi)_i
\]
Recall the Noether current
\[
j^{\mu}(x) = \frac{ \partial \mathcal{L}(x) }{ \partial ( \partial_{\mu} \varphi_i(x)) } \delta \varphi_i(x)
\]
and the Lagrangian
\[
\mathcal{L} = \frac{-1}{2} \partial^{\mu} \varphi_i \partial_{\mu} \varphi_i - \frac{1}{2} m^2 \varphi_i \varphi_i - \frac{1}{16} \lambda ( \varphi_i \varphi_i)^2
\]
so that 
\[
\frac{ \partial \mathcal{L} }{ \partial ( \partial_{\mu} \varphi_i ) } = - \partial^{\mu} \phi_i 
\]
and that 
\[
j^{\mu} = - \partial^{\mu} \varphi_i ( -i  \theta^a (T^a)_{ij} \varphi_j ) = i \theta^a \partial^{\mu} \varphi_i T^a_{ij} \varphi_i 
\]
so that 
\[
\theta^a j^{\mu a} =i \theta^a \partial^{\mu} \varphi_i T^a_{ij} \varphi_j \text{ or } \boxed{ j^{\mu a} = i \partial^{\mu} \varphi_i T^a_{ij} \varphi_j  }
\]

\item[(b)] \[
Q^a = \int d^3x j^{0 a}(x) = - \int d^3x i \dot{\varphi}_i T_{ij}^a \varphi_j = - i \int d^3x \pi_i T_{ij}^a \varphi_j
\]
so
\[
[\varphi_i, Q^a ] = [ \varphi_i, -i \int d^3x \pi_j T^a_{jk} \varphi_k ] = -i \int d^3x \left( [ \varphi_i, \pi_j ] T_{jk}^a \varphi_k + \pi_j T^a_{jk} [ \varphi_i, \varphi_k ] \right) = -i \int d^3x \left( i \delta^3(x'-x) \delta_{ij} T_{jk}^a \varphi_k \right) = T_{ij}^a \varphi_j(x)
\]
\[
\Longrightarrow \boxed{ [\varphi_i, Q^a] = T^a_{ij}\varphi_j }
\]
\item[(c)] Recall $[T^a, T^b] = i f^{abc} T^c$  \quad (24.7)  and the Jacobi identity,
\[
[[A,B],C ] + [[B,C] ,A] + [[C,A] ,B] = 0 
\]
\[
\begin{gathered}
  [[Q^a, Q^b], \varphi_i ] = [ Q^a, [Q^b, \varphi_i ]] + [ Q^b , [\varphi_i, Q^a]] = - [ Q^a , T^b_{ij} \varphi_j ] + [Q^b , T^a_{ij} \varphi_j ] = T^b_{ij} T^a_{jk} \varphi_k + - T_{ij}^a T^b_{jk} \varphi_k = -i f^{abc} T^c_{ik} \varphi_k = \\
  = -i f^{abc}  [ \varphi_i, Q^c] = [if^{abc} Q^c, \varphi_i ]
\end{gathered}
\]
$\varphi_i$ arbitrary so $[Q^a, Q^b] = i f^{abc} Q^c$
\end{enumerate}

\solutionhead{24.4} 

Now 
\[
S_{ij} S_{kl} \eta_{jl} = \eta_{ik}
\]
If 
\[
S_{ij} = \delta_{ij} + \theta_{ij} + O(\theta^2)
\]
then
\[
S_{ij} S_{kl} \eta_{jl} \simeq ( \delta_{ij} \delta_{kl} + \delta_{ij} \theta_{kl} + \theta_{ij} \delta_{kl} ) \eta_{jl} = \eta_{ik} + \theta_{kl} \eta_{il} + \theta_{ij} \eta_{jk} = \eta_{ik}
\]
\[
\Longrightarrow \theta_{kl} \eta_{il} + \theta_{ij} \eta_{jk} = 0 \text{ or } \theta_{ik} \eta_{kj} = - \eta_{ik} \theta_{jk}
\]
which is the same as, which is nice to note and rewrite as, $(\theta \eta )_{ij} = - (\eta \theta^T)_{ij} \text{ or } \theta \eta = - \eta \theta^T$   \\

$\eta$ is antisymmetric (and now include this fact):
\[
\theta_{ik} \eta_{kj} = - \theta_{ik} \eta_{jk} = - \eta_{jk} \theta_{ik} = \theta_{jk} \eta_{ki}
\]
Immediately, the number of independent equations from above is 
\[
\frac{ 2N (2N+1) }{ 2 } = \frac{ (2N)^2}{ 2} + N
\]
so that the number of \emph{(independent)} generators is 
\[
(2N)^2  - \left( \frac{ (2N)^2}{2} + N \right) = \frac{(2N)^2}{2} - N = 2N^2 - N
\]



\section{Unstable particles and resonnaces }

\section{Infrared divergences }

\section{Other renormalization schemes }

\solutionhead{27.1} 

Given \\
$\begin{aligned} 
  \beta(\alpha) & = b_1 \alpha^2 + \mathcal{O}(\alpha^3) \quad \quad \quad \, & (27.27) \\ 
  \gamma_m(\alpha) & = C_1 \alpha + \mathcal{O}(\alpha^2) \quad \quad \quad \, & (27.28) 
\end{aligned}$ \\

and recalling the development in Sec. 27, 

\[
\begin{gathered}
  \frac{ d \alpha }{ d\ln{ \mu} } = b_1 \alpha^2 + \mathcal{O}(\alpha^3) = \beta(\alpha) \quad \quad \quad \, \frac{ dm}{ d\ln{ \mu} } = \gamma_m(\alpha) m = (C_1 \alpha + \mathcal{O}(\alpha^2) ) m
\end{gathered}
\]

Using chain rule, for $\mu = \mu(\alpha)$, 
\[
\frac{dm}{d\alpha}  = \frac{ dm}{ d\ln{\mu}} \frac{ d\ln{ \mu }}{ d\alpha} = (C_1 \alpha + \mathcal{O}(\alpha^2) ) m \left( \frac{1}{ b_1 \alpha^2 + \mathcal{O}(\alpha^3)}  \right) = \frac{ (C_1 \alpha + \mathcal{O}(\alpha^2)) m}{ b_1 \alpha^2 + \mathcal{O}(\alpha^3) } \simeq \left( \frac{ C_1}{b_1} \frac{1}{ \alpha} + \mathcal{O}(\alpha) \right) m 
\]

\[
\frac{dm}{m} = \frac{C_1}{b_1} \frac{1}{\alpha} d\alpha \Longrightarrow \ln{ \left( \frac{ m_2}{m_1} \right) } = \frac{ C_1}{b_1} \ln{ \frac{ \alpha_2}{\alpha_1} } \text{ or } \frac{m_2}{m_1} = \left( \frac{ \alpha_2}{\alpha_1} \right)^{C_1/b_1}
\]
\[
\boxed{ m(\mu_2) = \left( \frac{ \alpha(\mu_2) }{ \alpha(\mu_1) } \right)^{C_1/b_1} m(\mu_1) }
\]

\section{The renormalization group}

\solutionhead{28.1} 

Consider $\varphi^4$ theory.  

$\mathcal{L} = -\frac{1}{2} Z_{\varphi} \partial^{\mu} \varphi \partial_{\mu} \varphi - \frac{1}{2} Z_m m^2 \varphi^2 - \frac{1}{24} Z_{\lambda} \lambda \widetilde{\mu}^{\epsilon} \varphi^4$  (28.41), in $d= 4 - \epsilon$.  

For bare fields and parameters,
\[
\mathcal{L} = -\frac{1}{2} \partial^{\mu} \varphi_0 \partial_{\mu} \varphi_0  - \frac{1}{2} m_0^2 \varphi_0^2 - \frac{1}{4!} \lambda_0 \varphi_0^4
\]

Comparing the renormalized and bare Lagrangian,

\[
\begin{aligned}
  & \varphi_0(x) = Z_{\varphi}^{1/2} \varphi(x) \\ 
  & m_0 = Z_\varphi^{-1/2} Z_m^{1/2} m \\ 
  & \lambda_0 = Z_{\lambda} \lambda \widetilde{\mu}^{\epsilon} Z_{\varphi}^{-2}
\end{aligned}
\]

In $\overline{MS}$ scheme, \\
$\begin{aligned}
  Z_{\varphi} & = 1 + \sum_{n=1}^{\infty} \frac{ a_n(\lambda)}{ \epsilon^n } \\ 
  Z_m & = 1 + \sum_{n=1}^{\infty} \frac{b_n(\lambda) }{  \epsilon^n } \\ 
  Z_{\lambda} & = 1 + \sum_{n=1}^{\infty} \frac{ c_n(\lambda) }{ \epsilon^n }
\end{aligned}$

Recall, from Prob. 14.5,

\[
\Pi(k^2) = \alpha \left( \frac{1}{ \epsilon}  + \ln{\mu} + \frac{1 - \gamma}{2}  + \mathcal{O}(\epsilon)  \right) m^2 - (Ak^2 + Bm^2) + \mathcal{O}(\alpha)
\]
where $\alpha = \frac{ \lambda }{ (4\pi)^2}$;  \quad \, $\mu = \frac{ \sqrt{4 \pi } \widetilde{\mu}}{ e^{\gamma}}$,  and  \\
$\begin{aligned} 
 A & = Z_{\varphi} - 1 = 0 \\ 
 B & = \alpha \left( \frac{1}{ \epsilon } + \ln{ \frac{ \mu}{ m }} + \frac{1}{2} + \mathcal{O}(\epsilon) \right) = Z_m - 1
\end{aligned}$.  

From Prob. 16.1, to get $\mathbf{V}_4(0,0,0,0) = \lambda$, $Z_{\lambda} = 1 + C$ and so
\[
C = \alpha \left( \frac{3}{\epsilon} + 3\ln{ \frac{\mu}{m} } + 1 \right)
\]

\[
\lambda_0 = Z_{\lambda} \lambda \widetilde{\mu}^{\epsilon} Z_{\varphi^{-2}} \Longrightarrow \alpha_0 = \frac{ \lambda_0}{ (4\pi)^2} = Z_{\lambda } \alpha \widetilde{\mu}^{\epsilon} Z_{\varphi}^{-2}
\]
\[
\ln{ \alpha_0} = \ln{Z_{\lambda}} + \ln{\alpha}+ \epsilon \ln{ \widetilde{\mu}} + - 2 \ln{ Z_{\varphi}}
\]
\[
\begin{aligned}
  \frac{d \ln{ \alpha_0} }{ d\ln{ \mu} }  = 0  = \frac{d}{d\ln{\mu}} ( + \alpha \left( \frac{3}{ \epsilon } + 3\ln{ \frac{\mu}{m} } + 1 \right) ) + \frac{ 1 }{ \alpha} \frac{d\alpha}{ d\ln{ \mu} } + \epsilon + 0  = + \frac{d\alpha}{ d\ln{ \mu }} \left( \frac{3}{\epsilon } + 3 \ln{ \frac{ \mu}{ m}} + 1 \right) - \alpha(3) + \frac{1}{ \alpha} \frac{d\alpha}{ d\ln{ \mu} } + \epsilon 
\end{aligned}
\]
\[
\frac{d\alpha}{ d\ln{ \mu}} \left(  \frac{3}{\epsilon} + \frac{1}{ \alpha} + 3 \ln{ \frac{\mu}{m}} + 1 \right) =3 \alpha - \epsilon 
\]
\[
\begin{aligned}
  \frac{d\alpha}{ d\ln{\mu}} & = \frac{3\alpha - 3}{ \frac{1}{ \alpha} + \frac{3}{\epsilon}  + 3\ln{ \frac{\mu}{m}} + 1 } = \frac{ 3\alpha^2 - \epsilon \alpha}{ 1 + \alpha \left( \frac{3}{\epsilon } + 3\ln{ \frac{\mu}{m} } + 1 \right) }  \simeq (3\alpha^2 - \epsilon \alpha) ( 1 - \alpha \left( \frac{3}{\epsilon} + 3 \ln{ \frac{\mu}{m}} + 1 \right) + \mathcal{O}(\alpha^2) ) = \\
   & = - \epsilon \alpha + 3\alpha^2 + \mathcal{O}(\alpha^3)
\end{aligned}
\]
\[
\Longrightarrow \boxed{ \beta(\alpha) = 3\alpha^2 + \mathcal{O}(\alpha^3) }
\]

\noindent Now $m_0 = Z_{\varphi}^{-1/2} Z_m^{1/2} m$  \\
\phantom{Now} $Z_{\varphi} =1$ \\ 
\phantom{Now} $Z_m = 1 + \alpha \left( \frac{1}{ \epsilon } \right) + \mathcal{O}(\epsilon^0)$ \\
$\ln{ m_0} = \frac{1}{2} \alpha \left( \frac{1}{ \epsilon } \right) + \mathcal{O}(\epsilon^0) + \ln{m}$ 

\[
\frac{d\ln{m_0} }{ d \ln{\mu}} = 0 = \frac{1}{2} \frac{ 1 }{ \epsilon } \frac{ d\alpha}{ d\ln{ \mu}} + \frac{1}{m } \frac{dm}{ d\ln{ \mu }}   
\]
\[
  \Longrightarrow \frac{1}{m} \frac{dm}{ d\ln{ \mu}} = \frac{-1}{2\epsilon } \frac{d\alpha}{ d\ln{ \mu}} = \frac{-1}{2 \epsilon } ( - \epsilon \alpha + \mathcal{O}(\alpha^2) ) = \frac{\alpha}{2} + \mathcal{O}(\alpha^2)
    \]



\[
\boxed{ \gamma_{ \varphi}(\alpha) = \frac{1}{2} \frac{ d\ln{ Z_{\varphi}} }{ d\ln{ \mu}} = 0 }
\]






\solutionhead{28.3}

Since for $\chi \to - \chi$, $\mathcal{L}$ is invariant, then physically, $\langle 0 | \chi | 0 \rangle =\langle 0 | - \chi | 0 \rangle = - \langle 0 | \chi | 0 \rangle$.  

So $\langle 0 | \chi | 0 \rangle = 0$.  

Consider the interaction term $\frac{1}{2} Z_h h \widetilde{\mu}^{\epsilon/2} \varphi \chi^2$.  I will develop the theory and Feynman diagrams here.  Work directly in momentum space.  

\[
\begin{gathered}
  i \int d^4x \frac{1}{2} Z_h h \widetilde{\mu}^{3/2} \varphi \chi^2 = \\
  = \frac{i}{2} Z_h h \widetilde{ \mu}^{\epsilon/2} \int d^4x \int \frac{d^4k_a}{(2\pi)^4} e^{ik_a x} \widetilde{\varphi}(k_a) \int \frac{d^4k_b}{(2\pi)^4} \int \frac{ d^4k_c}{(2\pi)^4} e^{ik_b x} \widetilde{\varphi}(k_b) e^{ik_c x} \widetilde{\varphi}(k_c)   = \\
  = \frac{i}{2} (2\pi)^4 Z_h h \widetilde{\mu}^{\epsilon/2} \left( \prod_{\alpha = a}^c \int \frac{d^4k_{\alpha}}{ (2\pi)^4} \right) \delta^4(k_a+k_b+k_c) \varphi_{k_a} \chi_{k_b} \chi_{k_c} = \\
  = \frac{i}{2} (2\pi)^4 Z_h h\widetilde{\mu}^{\epsilon/2} \left( \prod_{\alpha =a}^c \int \frac{d^4k_{\alpha}}{ (2\pi)^4} \right) \delta^4(k_a+k_b+k_c) \left( \frac{ (2\pi)^4}{ i } \right)^3 \delta_{J_{-k_a}} \delta_{K_{-k_b}} \delta_{K_{-k_c}}
\end{gathered}
\]

\noindent Consider $P$ $\varphi$ propagators.  \\
\phantom{Consider} $Q$ $\chi$ propagators.  

\noindent Let $2P - V = E_{\varphi}$ \\
\phantom{Let} $2Q - 2V = 2(Q-V) = E_{\chi}$.

Now $V \in \mathbb{Z}^+$.  Consider $V=1$.  Consider $Q=1$.  For $P=1$, there's the term $\langle 0 | \varphi | 0 \rangle =0$, fixed by $Y$.  

For $P=2$,
\[
\frac{1}{2} \left( \int \frac{ d^4l }{ (2\pi)^4} \frac{ \widetilde{J}_l \widetilde{J}_{-l} }{ l^2 + m^2 } \right)^2  \left( \frac{i}{2} \int \frac{ d^4q }{ (2\pi)^4} \frac{ \widetilde{K}_q \widetilde{K}_{-q} }{ q^2 + M^2 } \right) 
\]


Applying $\delta_{K_{-k_b}} \delta_{K_{-k_a}}$ then 1 vertex on the propagators leads to \\
$\delta_{K_{-k_b}} \frac{ \widetilde{K}_{k_c}}{ k_c^2  + M^2 } = \frac{1}{ k_b^2 + M^2 } \delta^4(k_c - (-k_b))$ \\
so that $k_b+k_c =0$.  So there results a factor of $\delta^4(k_a)$, forcing $k_a = 0$.  

But that means that for one of the particles of $\varphi$ field, $k=k_a =0$.  No particle.  This is not physical.  

Then $Q \geq 2$.   \\
Consider $P =1, Q=2$ (tree-level vertex).  

\[
\begin{gathered}
  \frac{i}{2} (2\pi)^4 Z_h h \widetilde{\mu}^{\epsilon/2} \left( \prod_{\alpha = a}^c \int \frac{d^4k_{\alpha}}{ (2\pi)^4} \right) \delta^4(k_a+k_b +k_c)  \left( \frac{ (2\pi)^4}{ i} \right)^3 \delta_{J_{-k_a}} \delta_{K_{-k_b} } \delta_{K_{-k_c}} \cdot \\ 
  \cdot \left( \frac{i}{2} \int \frac{d^4l}{(2\pi)^4} \frac{ \widetilde{J}_l \widetilde{J}_{-l}}{ l^2 + m^2 } \right) \frac{1}{2} \left( \frac{i}{2} \int \frac{d^4q}{ (2\pi)^4} \frac{ \widetilde{K}_q \widetilde{K}_{-q} }{ q^2 + M^2 } \right)^2 = \\
  = \frac{i}{2} (2\pi)^4 Z_h h \widetilde{\mu}^{\epsilon/2} \left( \prod_{\alpha = a}^c \int d^4k_{\alpha} \right) \delta^4(k_a+k_b+k_c) \frac{1}{ i^{3V} } i^{P +Q} \frac{ \widetilde{J}_{k_a}}{ k_a^2  +m^2} \frac{ \widetilde{K}_{k_b}}{ k_b^2  +M^2} \frac{ \widetilde{K}_{k_c}}{ k_c^2  +M^2} \left( \frac{1}{ (2\pi)^4} \right)^{P+Q}
\end{gathered}
\]

Applying LSZ formula,
\[
i^3(k_2^2 + M^2) (k_3^2+ M^2) ( k^2+m^2) \langle 0 | T\varphi_{-k_1} \chi_{k_2} \chi_{k_3} | 0 \rangle = i (2\pi)^4 Z_h h \widetilde{\mu}^{\epsilon/2} \delta^4(k_1 - k_2 - k_3)
\]


For each vertex is assigned a factor of $iZ_h h \widetilde{\mu}^{\epsilon/2}$.  

%\setlength{\unitlength}{10mm}
% \begin{picture}(6,6)
%  \multiput(3,2)(0.5,0.1){5}{ \line(0,1){0.3} }
%  \multiput(3,3)(0.5,0.5){5}{\line(1,1){0.3} }
%  \multiput(3,3)(0.5,-0.5){5}{\line(1,-1){0.3} }
%  \linethickness{ 7.5mm}
%  \put(3,3){ \vector(1,1){1.4}  }
%  \put(3,3){ \vector(1,-1){1.4} }
% \linethickness{ 15mm }
%  \put(3,3){ \line(-1,0){3} }
%  \put(3,3,5){ $i Z_h h \widetilde{\mu}^{\epsilon/2}$ } 
% \end{picture}

The counting factor of 2 is correct; indeed
\[
\langle \overbracket{ k_2 \underbracket{ k_3 | \chi_x}_2 \chi_x } \underbracket{ \varphi_x | k_1 \rangle }
\]

Consider $P = 2$, $Q=2$.  
\[
\begin{gathered}
  \frac{1}{2} \left( \frac{i}{2} (2\pi)^4 Z_h h \widetilde{\mu}^{\epsilon/2} \left( \prod_{\alpha = a}^c \int d^4k_{\alpha} \right) \delta^4(k_a+k_b + k_c) \delta_{J_{-k_a}} \delta_{K_{-k_b}} \delta_{K_{-k_c} } \right)^2 \frac{1}{i^{3V}} \cdot \\ 
  \cdot \frac{1}{2} \left( \frac{i}{2} \int  \frac{ d^4l}{ (2\pi)^4} \frac{ \widetilde{J}_l \widetilde{J}_{-l} }{ l^2 + m^2 } \right)^2 \frac{1}{2}  \left( \frac{i}{2} \int  \frac{ d^4q}{ (2\pi)^4} \frac{ \widetilde{K}_q \widetilde{K}_{-q} }{ q^2 + M^2 } \right)^2 = \\ 
  = \frac{1}{2} ( i (2\pi)^4 Z_h h \widetilde{\mu}^{\epsilon/2} )^2 \int d^4k_a d^4k_d d^4k_e d^4k_f \delta^4(k_a - k_e - k_f) \delta^4(k_d + k_e +k_f) \left( \frac{ -1}{ (2\pi)^4} \right)^{ P+Q} \cdot \\
  \cdot \frac{ \widetilde{J}_{k_a} }{ k_a^2 + m^2} \frac{ \widetilde{J}_{k_d} }{ k_d^2 + m^2} \left( \frac{1}{ k_f^2 + M^2 } \right)  \left( \frac{1}{ k_e^2 + M^2 } \right)  
\end{gathered}
\]

Applying LSZ formula, and while there is a counting factor of $2$ from matching functional derivatives from $\langle 0 | \widetilde{\varphi}_{k_1} \widetilde{\varphi}_{k_2} | 0 \rangle$, there is a $\frac{1}{2}$ factor from symmetry (the 2 internal propagators of $\chi$ can be exchanged).  

\[
(i Z_h h \widetilde{\mu}^{\epsilon/2} )^2 \frac{1}{2} \int \frac{ d^4l}{(2\pi)^d} \frac{ 1 /i }{ (k_1 + l)^2 + M^2 } \frac{ 1/ i }{ l^2 + M^2 } (2\pi)^4 \delta^4(k_1 - k_2)
\]

The counting is correct, a factor of $2^3 = 8$, as can be checked by the following:
\[
\langle \underbracket{ k_2 | \varphi_y }_2 \underbracket{ \chi_y \overbracket{ \chi_y \chi_x}^2 \chi_x }_2 \underbracket{ \varphi_x | k_1 } \rangle 
\]

The consider another term (Feynman diagram).

\begin{gather}
  \begin{gathered}
    \frac{1}{2} \left( \frac{i}{2} (2\pi)^4 Z_h h \widetilde{\mu}^{\epsilon/2} \left( \prod_{\alpha = a}^c  \int d^4k_{\alpha } \right) \delta^4( k_a + k_b + k_c) \delta_{J_{-k_a}} \delta_{K_{-k_b}} \delta_{K_{-k_c}} \right)^2 \frac{1}{i^{3V}} \cdot \\ 
    \cdot \left( \frac{i}{2} \int \frac{d^4l}{(2\pi)^4} \frac{ \widetilde{J}_l \widetilde{J}_{-l} }{ l^2 + m^2 } \right) \cdot \frac{1}{3!} \left( \frac{i}{2} \int \frac{ d^4 q}{ (2\pi)^4} \frac{ \widetilde{K}_q \widetilde{K}_{-q} }{ q^2 + M^2 } \right)^3 = 
  \end{gathered} \label{Eq:28_3-chi_propagator_1} \\
  \begin{gathered}
    = \frac{1}{2} ( (2\pi)^4 Z_h h \widetilde{\mu}^{\epsilon/2} )^2 \left( \prod_{\alpha = c}^f \int d^4k_{\alpha} \right) \delta^4(-k_d +  - k_e + k_c) \delta^4(k_d + k_e + k_f) \left( \frac{1}{ (2\pi)^4} \right)^{P+Q} \cdot \\
    \cdot \frac{ \widetilde{K}_c }{ k_c^2 + M^2 }   \frac{ \widetilde{K}_f }{ k_f^2 + M^2 }  \frac{1}{ k_d^2 + m^2 } \frac{1}{ k_e^2 + M^2 }  = 
  \end{gathered} \label{Eq:28_3-chi_propagator_2} \\
  \begin{gathered}
    = (i Z_h h \widetilde{\mu}^{\epsilon/2} )^2 \int d^4k \frac{ 1 / i }{ k^2 + m^2 } \frac{ 1 /i}{ (l+k)^2 + M^2 } (2\pi)^4 \delta^4(k_1 +k_2  - k_3 - k_4) 
    \end{gathered} \label{Eq:28_3-chi_propagator_3} 
\end{gather}

where a counting factor of $2 \cdot 2 \cdot 2 \cdot 3! \cdot 2^3$ was needed going from Eq. (\ref{Eq:28_3-chi_propagator_1}) to Eq. (\ref{Eq:28_3-chi_propagator_2}) and factor of 2 was included, going from Eq. (\ref{Eq:28_3-chi_propagator_2}) to Eq. (\ref{Eq:28_3-chi_propagator_3}).  




\hrulefill





\begin{itemize}
\item[(a)] So to calculate the 1-loop contributions to each of the $Z$'s in the $\overline{MS}$ renormalization scheme, be careful to include all contributing Feynman diagrams.  

For the $\varphi$ propagator, there are 3 distinct diagrams, including 1 for the counterterms $A$, $B$.  

\[
\begin{gathered}
  \Pi_{\varphi}(k^2) = \frac{1}{2} \frac{ (Z_g g)^2 + (Z_h h)^2 }{ (4\pi)^3} \int_0^1 dx D\ln{ \frac{D}{\mu^2 }} - \left( \frac{1}{12} \frac{ (Z_g g)^2 + (Z_h h)^2 }{ (4\pi)^3} \left( \frac{2}{ \epsilon } + 1 \right) + A \right) k^2 - \\
  - \left( \frac{1}{2} \frac{ (Z_g g)^2 m^2 + (Z_h h)^2 M^2 }{ (4\pi)^3} \left( \frac{2}{ \epsilon } + 1 \right) + B m^2 \right)
\end{gathered}
\]

Assuming $Z_g = 1 + \mathcal{O}(g^2)$, $Z_h = 1 + \mathcal{O}(h^2)$, \\
$A = Z_{\varphi}  -1$, $B = Z_m - 1$
\[
\boxed{ A = \frac{-1}{6} \frac{ g^2 + h^2 }{ (4\pi)^3} \frac{1}{ \epsilon }, \quad \, B = -1 \frac{ g^2 + h^2 \frac{M^2}{m^2}}{ (4\pi)^3} \frac{1}{\epsilon }  }
\]

For $\chi$ propagator, there's only 2 distinct diagrams, including 1 for counterterms $A'$, $B'$.  

Note that because the 2 internal propagators, 1 $\chi$, 1 $\varphi$, are distinct, there's no symmetry factor of $2$ in the denominator.  

\[
i \Pi_{\chi}(k^2 ) = ( i Z_h h \widetilde{\mu}^{\epsilon/2})^2 \int \frac{d^4l}{(2\pi)^d} \frac{ 1 /i }{ (l + k_1)^2 + m^2 } \frac{ 1/i }{ l^2 + M^2 } 
\]
\[
\begin{gathered}
  ( (l+k)^2 + m^2 )^{-1} (l^2 + M^2)^{-1} = \int_0^1 dx \left( x ( (l^2 + 2lk + k^2 ) + m^2 )  + ( 1 - x) (l^2 + M^2 ) \right)^{-2} = \\
   = \int_0^1 dx ( l^2 + 2lkx + x(k^2 + m^2) + (1-x)M^2 )^{-2} = \int_0^1 dx (( l+kx)^2 + x(1-x) k^2 + x(m^2 - M^2) + M^2 )^{-2}
\end{gathered}
\]
here $D = x(1-x)k^2 + x(m^2 - M^2) + M^2$.  

Note
\[
\int_0^1 dx D  = \frac{1}{6} k^2 + \frac{1}{2} ( m^2 - M^2) + M^2 = \frac{1}{6} k^2 + \frac{1}{2} (m^2 + M^2 )
\]
\[
\Pi_{\chi}(k^2 ) = \frac{ (Z_h h)^2}{(4\pi)^3} \lbrace \int_0^1 dx D \ln{ \frac{D}{ \mu^2 } } + - \left( \frac{2}{ \epsilon } +1 \right) \left( \frac{k^2 }{6} + \frac{1}{2} (m^2 + M^2 ) \right) \rbrace - A' k^2 - B' M^2 
\]
matching terms of $k^2$ and $m^2, M^2$,
\[
\boxed{
  \begin{aligned}
    A' & = - \frac{ (Z_h h)^2 }{ (4\pi)^3} \frac{ 1/3 }{ \epsilon } \\ 
    B' & = - \frac{ (Z_h h)^2 }{ (4\pi)^3 } \frac{1}{ \epsilon } \left( 1 + \frac{m^2 }{M^2 } \right) 
  \end{aligned} }
\]

For $\varphi^3$ vertex, there are 3 diagrams.  

\[
\mathbb{V}_{3\varphi^3}  = Z_g g \widetilde{\mu}^{\epsilon/2} + \frac{ g^3}{(4\pi)^3} \left( \frac{1}{ \epsilon } + \frac{1}{2} 2 \int_0^1 dx \int_0^{1-x} dy \ln{ \left( \frac{ \mu^2}{ D} \right) } \right) + \frac{h^3}{ (4\pi)^3} \left( \frac{1}{ \epsilon } + \frac{1}{2} 2 \int_0^1 dx \int_0^{1-x} dy \ln{ \frac{\mu^2}{ D} } \right) 
\]
For $Z_g = 1 + C$, for the tree-level term, first term, above,
\[
C = \frac{-1}{ \epsilon} \frac{ (g^3 + h^3) }{ (4\pi)^3} / g\widetilde{\mu}^{\epsilon/2} \to \boxed{ C = \frac{-1}{\epsilon} \frac{ (g^3 + h^3) }{ g(4\pi)^3 } }
\]  

For $\varphi \chi^2$ vertex, there are 3 diagrams.  

\[
\mathbb{V}_{3\varphi \chi^2} = Z_h h \widetilde{\mu}^{\epsilon/2} + \frac{ h^3}{ (4\pi)^3} \left( \frac{1}{ \epsilon} + \frac{1}{2} 2 \int_0^1 dx \int_0^{1-x} dy \ln{ \frac{ \mu^2}{ D} } \right) + \frac{h^2 g}{ (4\pi)^3} \left( \frac{1}{ \epsilon} + \frac{1}{2} 2 \int_0^1 dx \int_0^{1-x} dy \ln{ \frac{\mu^2 }{ D} } \right)
\]
For $Z_h = 1 + C'$,
\[
C' = \frac{-1}{\epsilon}\frac{ h (h^2 + hg) }{ (4\pi)^3 h} = \boxed{ \frac{-1}{\epsilon} \frac{ (h^2 + hg) }{ (4\pi)^3} }
\]
\item[(b)] \[
\begin{gathered}
  \ln{ g_0 } = G(g,h,\epsilon) + \ln{g} + \frac{\epsilon}{2} \ln{\widetilde{\mu}} \\ 
  \frac{ d\ln{ g_0}}{ d\ln{ \mu}} = 0 = \frac{ \partial G}{ \partial g} \frac{dg}{ d\ln{\mu} } + \frac{ \partial G}{ \partial h} \frac{dh}{ d\ln{ \mu }} + \frac{1}{g} \frac{ dg}{ d\ln{ \mu}} + \frac{\epsilon}{2} \\ 
  \Longrightarrow \left( g \frac{ \partial G}{ \partial g} + 1 \right) \frac{dg}{ d\ln{ \mu}} + g \frac{\partial G}{ \partial h } \frac{ dh}{ d\ln{ \mu}} + \frac{ \epsilon g}{ 2}  =0 
\end{gathered}
\]

\[
\begin{gathered}
  \ln{ h_0 } = H(g,h,\epsilon) + \ln{h} + \frac{\epsilon}{2} \ln{\widetilde{\mu}} \\ 
  \frac{ d\ln{ h_0}}{ d\ln{ \mu}} = 0  \Longrightarrow \left( h \frac{ \partial H}{ \partial h} + 1 \right) \frac{dh}{ d\ln{ \mu}} + h \frac{\partial H}{ \partial g } \frac{ dg}{ d\ln{ \mu}} + \frac{ \epsilon h}{ 2}  =0 
\end{gathered}
\]
Now
\[
\begin{aligned}
  & G(g,h,\epsilon) = \sum_{n=1}^{\infty} G_n(g,h) \epsilon^{-n} = \frac{ G_1(g,h)}{\epsilon} + \mathcal{O}\left( \frac{1}{\epsilon^2}\right) \\ 
  & H(g,h,\epsilon) = \sum_{n=1}^{\infty} H_n(g,h) \epsilon^{-n} = \frac{ H_1(g,h)}{\epsilon} + \mathcal{O}\left( \frac{1}{\epsilon^2}\right) 
\end{aligned}
\]
Then
\[
\begin{aligned}
  & \left( 1 + g \left( \frac{ \partial_g G_1 }{ \epsilon} + \dots \right) \right) \frac{dg}{ d\ln{ \mu}} + g \left( \frac{ \partial_h G_1}{ \epsilon} + \dots \right) \frac{dh}{ d\ln{ \mu}} = \frac{- \epsilon g}{2} \\ 
  & \left( 1 + h \left( \frac{ \partial_h H_1 }{ \epsilon} + \dots \right) \right) \frac{dh}{ d\ln{ \mu}} + h \left( \frac{ \partial_g H_1}{ \epsilon} + \dots \right) \frac{dg}{ d\ln{ \mu}} = \frac{- \epsilon h}{2}  
\end{aligned}
\]
$ \frac{dg}{ d\ln{ \mu}}$, $\frac{dh}{d \ln{\mu}}$ finite, so $\frac{dg}{d\ln{ \mu}}$, $\frac{dh}{ d \ln{\mu}}$ at ``most'' $\mathcal{O}(\epsilon^0)$, no powers of $\frac{1}{ \epsilon}$.  

\[
\left( \begin{matrix} 1 + g \frac{ \partial_g G_1}{\epsilon} & g \frac{ \partial_h G_1}{\epsilon} \\ h \frac{ \partial_g H_1}{ \epsilon } & 1 + h \frac{ \partial_h H_1}{ \epsilon } \end{matrix} \right) \left( \begin{matrix} \frac{dg}{d\ln{\mu}} \\ \frac{dh}{d\ln{\mu}} \end{matrix} \right) = \frac{-\epsilon}{2} \left( \begin{matrix} g \\ h \end{matrix} \right)
\]
\[
\left( \frac{1}{ 1 + h \frac{ \partial_h H_1}{\epsilon} + g \frac{ \partial_g G_1}{\epsilon } } \right) \left( \begin{matrix} 1 + h \frac{ \partial_h H_1}{\epsilon } & - g \frac{ \partial_h G_1}{\epsilon } \\ - h \frac{\partial_g H_1}{\epsilon } & 1 + g \frac{ \partial_g G_1}{ \epsilon } \end{matrix} \right) \left( \frac{ -\epsilon}{2} \right) \left( \begin{matrix} g \\ h \end{matrix} \right) \simeq \left( \begin{matrix} \frac{dg}{d\ln{ \mu} } \\ \frac{dh}{ d\ln{ \mu}} \end{matrix} \right)
\]
Doing long division with the determinant, for instance,
\[
\left( 1 + h \frac{ \partial_h H_1}{\epsilon} \right)/ \left( 1 + h \frac{ \partial_h H_1}{ \epsilon } + g \frac{ \partial_g G_1}{ \epsilon} \right) = 1 - g \frac{ \partial_g G_1}{\epsilon} + \dots
\]

\[
\begin{gathered}
  \Longrightarrow \frac{ - \epsilon}{2} \left( \begin{matrix} 1 - g \frac{ \partial_g G_1}{ \epsilon} & - g \frac{ \partial_1 G_1}{\epsilon} \\ -h \frac{ \partial_g H_1}{ \epsilon } & 1 - h \frac{ \partial_h H_1 }{\epsilon } \end{matrix} \right) \left( \begin{matrix} g \\ h \end{matrix} \right)  \simeq \\
  \simeq \boxed{ \begin{aligned}
      & \frac{ - \epsilon g}{2} + \frac{g}{2} \left( g \frac{ \partial G_1}{ \partial g} + h \frac{ \partial G_1}{ \partial h } \right) = \frac{dg}{ d\ln{ \mu }} \\ 
      & \frac{ - \epsilon h}{2} + \frac{h}{2} \left( g \frac{ \partial H_1}{ \partial g} + h \frac{  \partial H_1}{ \partial h} \right) = \frac{ dh}{ d\ln{ \mu}}
\end{aligned} }
\end{gathered}
\]
\item[(c)] Using $\ln{(1-x)} = -x - x^2 + \dots$ and going to $\mathcal{O}\left( \frac{1}{\epsilon} \right)$ order,
\[
\begin{aligned}
  G & = -\frac{3}{2} \ln{Z_{\varphi}} + \ln{ Z_g} = \\ 
  & = - \frac{3}{2} - \frac{1}{6} \frac{ g^2 + h^2 }{ (4\pi)^3 } \frac{1}{ \epsilon } + \frac{-1}{\epsilon} \frac{ g^3 + h^3 }{ g(4\pi)^3} = \frac{1}{ (4\pi)^3 \epsilon } \left( \frac{-3}{4} g^2 - \frac{h^3}{g} + \frac{h^2 }{4} \right)
\end{aligned}
\]
\[
H = \frac{-1}{2} \ln{ Z_{\varphi}} - \ln{ Z_{\chi}} + \ln{ Z_h } = \frac{1}{ (4\pi)^3 \epsilon } \left( \frac{1}{12} (g^2 + h^2 ) + \frac{h^2 }{3} + - h^2  -hg \right) = \frac{1}{ (4\pi)^3 \epsilon } \left( \frac{ - 7h^2 }{12} -hg + \frac{g^2}{12} \right)
\]
\[
\begin{aligned}
  \frac{ \partial G_1}{\partial g} & = \frac{1}{(4\pi)^3} \left( \frac{-3g}{2} + \frac{h^3}{g^2} \right) \\
  \frac{ \partial G_1}{ \partial h} & = \frac{1}{(4\pi)^3} \left( \frac{ - 3h^2 }{g}  + \frac{h}{2} \right)
\end{aligned}
\quad \quad \quad \, 
\begin{aligned}
  \frac{ \partial H_1}{ \partial h } & = \frac{1}{(4\pi)^3} \left( \frac{ - 7h }{6} -g \right) \\ 
  \frac{ \partial H_1}{ \partial g} & = \frac{1}{(4\pi)^3 } \left( - h + \frac{g}{6} \right)
\end{aligned}
\]
\[
\Longrightarrow \boxed{ \begin{aligned} & \beta_g = \frac{g}{2} \left( \frac{1}{(4\pi)^3} \right) \left( \frac{-3g^2}{2} + \frac{h^3}{g} - \frac{3h^3}{g} + \frac{h^2}{2} \right) = \frac{g}{2 (4\pi)^3} \left( \frac{-3g^2}{2} + \frac{ -2h^3}{g} + \frac{h^2}{2} \right)  \\ 
    & \beta_h = \frac{h}{2} \frac{1}{(4\pi)^3} \left( -gh + \frac{g^2}{6}  - \frac{7h^2}{6} -gh \right) = \frac{h}{2} \frac{1}{(4\pi)^3} \left( \frac{-7h^2}{6} - 2gh + \frac{g^2}{6} \right) 
\end{aligned} }
\]
\item[(d)] Now
\[
\begin{aligned}
  & \beta_g = \frac{g^3}{ 2(4\pi)^3} \left( \frac{-3}{2} + \frac{1}{2} \left( \frac{h}{g} \right)^2 - 2\left( \frac{h}{g} \right)^3 \right) \\ 
  & \frac{\beta_h}{h}  = \frac{g^2}{2(4 \pi)^3}  \left( \frac{-7}{6} \left( \frac{h}{g} \right)^2 - 2\left( \frac{h}{g} \right) + \frac{1}{6} \right) 
\end{aligned}
\]

If $\beta_g$ is negative, $\mu \frac{dg}{d\mu}$ is negative.  If $\mu$ increases (and assuming $\mu$ positive), then coupling constant decreases.  Recall in $\overline{MS}$ scheme $\mu \sim s^2$ (or just by dimensional analysis) so at higher energy, coupling constant decreases: \emph{asymptotic freedom}.   \\

\noindent For $g >0$, $\beta_g < 0$ when $\frac{h}{g} > -0.8324$.  \\
\phantom{For $g$ } For $-0.8324 < \frac{h}{g} < 0$, $h<0$.  \\
\phantom{For $g$ } \quad when $h < 0$, $\frac{\beta_h}{h} >0$, then $\beta_h < 0$, then $\beta_h <0$, for \\
\phantom{For $g$ } \quad when $-1.7936 < \frac{h}{g} < 0.07979$ 

So $\beta_g$ and $\beta_h$ is negative when $-0.8324 < \frac{h}{g} < 0$.  

For $\frac{h}{g} > 0$ and $h>0$,  \\
 $\frac{\beta_h}{h} < 0$ so that $\beta_h < 0$ when $\frac{h}{g} > 0.07979$.  

\[
\Longrightarrow \boxed{ -0.8324 < \frac{h}{g} < 0 \text{ or } \frac{h}{g} > 0.07979 } 
\]
for both $\beta_g$ and $\beta_h$ to be negative.  

\end{itemize}

\section{ Effective field theory}

\section{ Spontaneous symmetry breaking }

\section{ Broken symmetry and loop corrections}

\section{ Spontaneous breaking of continuous symmetries }

\section*{Part II  Spin One Half }


\section{ Representations of the Lorentz group }

\subsection*{ Problems} 

\problemhead{33.1} Express $A^{\mu \nu}(x)$ \, $S^{\mu \nu}(x)$, and $T(x)$ in terms of $B^{\mu \nu }(x)$.  

\solutionhead{33.1} Given $B^{\mu \nu}(x) = A^{\mu \nu}(x) + S^{\mu \nu }(x) + \frac{1}{4} g^{\mu \nu} T(x)$, (33.6), 

\[
B^{\nu \mu}(x) = A^{\nu \mu}(x) + S^{\nu \mu} (x) + \frac{1}{4} g^{\nu \mu} T(x) = - A^{\mu \nu}(x) + S^{\mu \nu} + \frac{1}{4} g^{\mu \nu} T(x)
\]

Now
\[
\begin{gathered}
\boxed{  \frac{B^{\mu \nu} + B^{\nu \mu} }{2} = S^{\mu \nu } } \\
\boxed{  g_{\mu \nu} B^{\mu \nu}(x) = \frac{1}{4}(4) T(x) = T(x) }
\end{gathered}
\]
Note $g_{\mu \nu}g^{\mu \nu} = 4$ because (proto\v{z}e, perch\'{e}), $g_{\mu \nu} g^{\mu \nu} =  \left( \begin{matrix} 1 & & & \\ & -1 & & \\ & & -1 & \\ & & & -1 \end{matrix} \right)\left( \begin{matrix} 1 & & & \\ & -1 & & \\ & & -1 & \\ & & & -1 \end{matrix} \right) = 4$ \\
\phantom{ Note } and $g_{\mu \nu} S^{\mu \nu} =0$  (traceless).  

\[
\boxed{ \frac{B^{\mu \nu } - B^{\nu \mu }}{2} =A^{\mu \nu }(x) }
\]


\problemhead{33.2} Verify that eqs. (33.18)-(33.20) follow from eqs. (33.11)-(33.13).  

\solutionhead{33.2} Recall, 
\[  \begin{aligned}
& \left[ J_i, J_j \right] = + i \epsilon_{ijk} J_k \quad \quad \, & (33.11)\\  
&  [ J_i , K_j ] = i \epsilon_{ijk} K_k \quad \quad \, & (33.12 ) \\
& [K_i, K_j ] = - i \epsilon_{ijk} J_k \quad \quad \, & (33.13) 
\end{aligned} \quad \quad \, \begin{aligned} N_i &= \frac{1}{2} (J_i - i K_i ) \quad \quad \, & (33.16) \\ N_i^{\dag} &= \frac{1}{2} (J_i + i K_i ) \quad \quad \, & (33.16) 
\end{aligned}
\]
Then we want
\[
\begin{aligned}
& \left[ N_i, N_j \right] = + i \epsilon_{ijk} N_k \quad \quad \, & (33.18)\\  
&  [ N_i^{\dag} , N_j^{\dag} ] = i \epsilon_{ijk} N_k^{\dag} \quad \quad \, & (33.19) \\
& [N_i, N_j^{\dag} ] =0 \quad \quad \, & (33.13) 
\end{aligned}
\]

\[
\begin{aligned}
  \left[ N_i, N_j \right] & = \frac{1}{4} \lbrace [J_i - i K_i, J_j - iK_j ] \rbrace = \frac{1}{4} \lbrace [ J_i, J_j] - i[ J_i, K_j] - i [ K_i, J_j] - [ K_i, K_j]  \rbrace = \\
  & = \frac{1}{4} \lbrace i \epsilon_{ijk} J_k + \epsilon_{ijk} K_k - \epsilon_{jik} K_k + i \epsilon_{ijk} J_k \rbrace = \frac{i \epsilon_{ijk}}{2} N_k 
\end{aligned}
\]

\[
\begin{aligned}
  \left[ N_i^{\dag}, N_j^{\dag} \right] & = \frac{1}{4} \lbrace [J_i + i K_i, J_j + iK_j ] \rbrace = \frac{1}{4} \lbrace [ J_i, J_j] + i[ J_i, K_j] + i [ K_i, J_j] - [ K_i, K_j]  \rbrace = \\
  & = \frac{1}{4} \lbrace i \epsilon_{ijk} J_k - \epsilon_{ijk} K_k + \epsilon_{jik} K_k + i \epsilon_{ijk} J_k \rbrace = \frac{i \epsilon_{ijk}}{2} N_k^{\dag} 
\end{aligned}
\]
\[
\left[ N_i, N_j^{\dag} \right] = \frac{1}{4} \lbrace [J_i - i K_i , J_j + iK_k ] \rbrace = \frac{1}{4} \lbrace i \epsilon_{ijk} J_k - \epsilon_{ijk} K_k - \epsilon_{jik} K_k - i \epsilon_{ijk} J_k \rbrace = 0 
\]




\section{Left- and right-handed spinor fields}

\subsection*{Problems}

\problemhead{34.1} Verify that eq. (34.6) follows from eq. (34.1).

\solutionhead{34.1} Given $U(\Lambda)^{-1} \psi_a(x) U(\Lambda) = L_a^b(\Lambda) \psi_b(\Lambda^{-1} x)$  \quad \, (34.1), 
\[
\begin{aligned}
&  \Lambda^{\mu}_{\nu} = \delta^{\mu}_{\nu} + \delta \omega^{\mu}_{\nu} \\ 
&  (\Lambda^{-1})^{\rho}_{\nu} = \Lambda_{\nu}^{\rho} = \delta_{\nu}^{\rho} + \delta \omega_{\nu}^{\rho} \\ 
&  (\Lambda^{-1})^{\rho}_{\nu} = x^{\rho} + \delta \omega_{\nu}^{\rho} x^{\nu} 
\end{aligned}
\]
\[
\begin{aligned}
  \psi_b(\Lambda^{-1} x) & \simeq \psi_b(x^{\rho}) + \delta \omega_{\nu}^{\rho} x^{\nu} \partial_{\rho} \psi_b = \psi_b(x^{\rho} ) + g^{\rho k} \delta \omega_{\nu k} x^{\nu} g_{\rho l} \partial^l \psi_b = \\ 
  & = \psi_b(x) + \delta \omega_{\nu k} x^{\nu} \partial^k \psi_b
\end{aligned}
\]
\[
(I + - \frac{i}{2} \delta \omega_{\mu \nu} M^{\mu \nu} ) ( \psi_a + \frac{i}{2} \delta \omega_{\mu \nu} \psi_a M^{\mu \nu} ) = \psi_a + \frac{i}{2} \delta \omega_{\mu \nu} \psi_a M^{\mu \nu} - \frac{i}{2} \delta \omega_{\mu \nu} M^{\mu \nu} \psi_a + \frac{1}{4} \delta \omega_{ab} \delta \omega_{\mu \nu} M^{ab} M^{\mu \nu} 
\]
Consider first-order terms.  
\[
\Longrightarrow i \delta \omega_{\mu \nu} \left[ \psi_a, M^{\mu \nu} \right]
\]
(when summing over $\mu, \nu$, we obtain 2 terms, each), i.e. 
\[
\delta \omega_{\mu \nu} \psi_a M^{\mu \nu} + \delta \omega_{\nu \mu} \psi_a M^{\nu \mu } = \delta \omega_{\mu \nu} \psi_a M^{\mu \nu} + (-\delta \omega_{\mu \nu} \psi_a (-M^{\mu \nu })
\]

\[
\begin{aligned}
  & L_a^b(\Lambda) = L_a^b(1 + \delta \omega) = \delta_a^b + \frac{i}{2} \delta \omega_{\mu \nu} (S_L^{\mu \nu})_a^b \\ 
  & L_a^b(\Lambda)\psi_b(\Lambda^{-1} x) \xrightarrow{ \delta \omega_{\mu \nu} } \frac{i}{2} \delta \omega_{\mu \nu} (S_L^{\mu \nu})_a^b \psi_b(x) + \delta \omega_{\mu \nu} (x^{\mu} \partial^{\nu} - x^{\nu} \partial_{\mu} )\psi_a + \frac{i}{2} \delta \omega_{\nu \mu} (S_L^{\nu \mu} )_a^b \psi_b(x)
\end{aligned}
\]

\[
\Longrightarrow \boxed{ \left[ \psi_a, M^{\mu \nu} \right] = \mathcal{L}^{\mu \nu} \psi_a + (S_L^{\mu \nu})_a^b \psi_b(x) }
\]
with $(S_L^{\mu \nu})_a^b = - (S_L^{\nu \mu})_a^b$  

\solutionhead{34.2} 

Recall that 
\[
\begin{aligned}
  & [ S^{\mu \nu}_L, S^{\rho \sigma}_L ] = i (g^{\mu \rho} S_L^{\nu \sigma} - ( \mu \leftrightarrow \nu ) ) - (\rho \leftrightarrow \sigma ) \quad \quad \quad \quad \, & (34.4) \\
  & (S_L^{ij})_a^b = \frac{1}{2} \epsilon^{ijk} \sigma_k  \quad \quad \quad \quad \, & (34.9) \\
  & (S_L^{k0})_a^b = \frac{1}{2} i \sigma_k \quad \quad \quad \quad \, & (34.10)
\end{aligned}
\]

For (34.9) obeying (34.4), without loss of generality, assume only one pair of indices of (34.4) are the same, since $\mu, \nu, \rho, \sigma = 1,2,3$, and if there are two matching pairs, then because $S_L$ is antisymmetric, indices could be exchanged, $S^{\mu \nu}_L$ and $S^{\rho \sigma}_L$ are the same, and the commutator is 0.  
\[
[S^{\mu \nu}_L, S^{\mu \sigma}_L] = i S_L^{\nu \sigma}
\]
\[
[ \frac{1}{2} \epsilon^{\mu \nu j } \sigma_j , \frac{1}{2} \epsilon^{\mu \sigma k } \sigma_k ] = \frac{1}{2} \epsilon^{\mu \nu j } \epsilon^{ \mu \sigma k } [ \sigma_j, \sigma_k ] = \frac{1}{2} \left( \delta^{\nu \sigma} \sigma^{jk } - \delta^{\nu k} \delta^{j \sigma} \right) i \sigma_l \epsilon^{jkl} = \frac{-i}{2} \sigma_l \epsilon^{ \sigma \nu l } = + i S_L^{\nu \sigma}
\]

\[
[S^{j0}_L, S_L^{k0} ] = [ \frac{1}{2} i \sigma_j, \frac{1}{2} i \sigma_k ] = \frac{-1}{4} 2 i e^{ jkl } \sigma^l = \frac{-i }{2} \epsilon^{jkl} \sigma^l = -i S_L^{jk}
\]
Indeed,
\[
[ S^{j 0}_L, S_L^{k0} ] = i ( 0 - 0 ) - ( 0 - S_L^{jk} ) = i S_k^{jk}
\]





\problemhead{34.3}




\solutionhead{34.3} $\epsilon^{\mu \nu \rho \sigma} \epsilon_{\alpha \beta \gamma \sigma} = \epsilon^{\sigma \mu \nu \rho} \epsilon_{\sigma \alpha \beta \gamma}$.  

Consider all permutations of $\mu \nu \rho$ to $\alpha \beta \gamma$.  



\section{Manipulating spinor indices}

\subsection*{Problems}





\solutionhead{35.1}

$\overline{\sigma}^{\mu \dot{a} a} \equiv \epsilon^{ab} \epsilon^{\dot{a} \dot{b} } \sigma^{\mu}_{b\dot{b}} = \epsilon^{\dot{a} \dot{b}} \epsilon^{ab} \sigma^{\mu}_{b\dot{b}}$ since we're given (3.19).  

What does it mean, "written in `matrix multiplication'"?  

Recall $\epsilon^{12} = + 1$, \, $\epsilon^{21} = -1$.  

To begin, consider writing everything explicitly:

\[
\epsilon^{ab} \sigma^0_{b\dot{b}} = -\delta_{a2} \delta_{\dot{b} 1} + \delta_{a1} \delta_{\dot{b} 2} = i ( \sigma^2)_{a\dot{b}}.  \text{Equivalently}, \left( \begin{matrix} & 1 \\ -1 & \end{matrix} \right) \left( \begin{matrix} 1 & \\ & 1 \end{matrix} \right) = \left( \begin{matrix} & 1 \\ -1 & \end{matrix} \right) 
\]

Likewise, for example, $\mu =3$, 
\[
\epsilon^{ab} \sigma^3_{b\dot{b}} = -\delta_{a2} \delta_{\dot{b} 1} - \delta_{a1} \delta_{\dot{b} 2} = i^2 ( \sigma^1)_{a\dot{b}}.  \text{Equivalently}, \left( \begin{matrix} & 1 \\ -1 & \end{matrix} \right) \left( \begin{matrix} 1 & \\ & -1 \end{matrix} \right) = \left( \begin{matrix} 0 & -1 \\ -1 & 0 \end{matrix} \right) 
\]

So, being careful about the spinor indices, that we can't cavalierly say "$\epsilon^{ab} \sigma^{\mu}_{b\dot{b}} = \sigma^{\mu a}_{\dot{b}}$", and keeping in mind the definition $\sigma^{\mu}_{a\dot{b}} = (I, \vec{\sigma})$, I conclude that (can be shown explicitly for $\mu = 1,2$ as well), 
\[
\epsilon^{ab} \sigma^{\mu}_{b\dot{b}} = \begin{cases} i (\sigma^2)_{a\dot{b}} & \text{ for } \mu =0 \\ 
i (\sigma^2 \sigma^{\mu})_{a\dot{b}} = i ( i \epsilon_{2\mu \rho} \sigma^{\rho} + \delta^{2\mu } 1 ) _{a\dot{b}} = ( - \epsilon_{2\mu \rho} \sigma^{\rho} + i \delta^{2\mu} 1 )_{a\dot{b}} & \text{ if } \mu \neq 0 
\end{cases}
\]
where, recall $\sigma^i \sigma^k =  i \epsilon^{ikl} \sigma^l + \delta^{ik}$.  

\[
\xrightarrow{ \epsilon^{\dot{a} \dot{b}} = - \epsilon^{\dot{b}\dot{a}} } \text{ if $\mu =0$}, \begin{gathered}
i ( \sigma^2)_{a\dot{b}}(-i) (\sigma^2_{\dot{b} \dot{a}} ) = \quad \, \text{(regular matrix multiplication)} \\ 
= (1)_{a\dot{a}} = (1)_{\dot{a} a} \end{gathered}
\]

If $\mu \neq 0$, 
\[
\begin{gathered}
-\epsilon^{\dot{b} \dot{a}} \epsilon^{ab} \sigma^{\mu}_{b\dot{b}} = - (- \epsilon_{2\mu \rho} \sigma^{\rho} + i \delta^{2\mu} 1)_{a\dot{b}} i (\sigma^2)_{\dot{b}\dot{a}} = i \epsilon_{2\mu \rho} (\sigma^{\rho} \sigma^2)_{a\dot{a}} + \delta^{2\mu} \sigma^2_{a\dot{a}} =  \\ 
	= - \epsilon_{2\mu \rho} \epsilon_{\rho 2 \nu} \sigma^{\nu}_{a\dot{a}}+ \delta^{2\mu} \sigma^2_{a\dot{a}} = - \sigma^{\mu \neq 2}_{a\dot{a}} + \delta^{2\mu} \sigma^2_{a\dot{a}} = - \sigma^{\mu \neq 2}_{\dot{a} a} + \delta^{2\mu} (- \sigma^2_{\dot{a} a} )
\end{gathered}
\]
Recall,
$\sigma^1 = \left( \begin{matrix} & 1 \\ 1 & \end{matrix} \right), \quad \, \sigma^2 = \left( \begin{matrix} & -i \\ i & \end{matrix} \right), \quad \, \sigma^3 = \left( \begin{matrix} 1 &  \\  & -1 \end{matrix} \right)$.  Note that $\begin{aligned} (\sigma^1)^T &  = \sigma^1 \\ (\sigma^3)^T & = \sigma^3 \end{aligned}$ but $(\sigma^2)^T = - \sigma^2$,  !

\[
\Longrightarrow \boxed{ \overline{\sigma}^{\mu \dot{a} a} = -\epsilon_{\dot{b}\dot{a}} \epsilon^{ab} \sigma^{\mu}_{b\dot{b}} = (I, - \vec{\sigma}) }
\]








\solutionhead{35.2} Given
\[
\begin{aligned}
  (S_L^{ij})_a^b & = \frac{1}{2} \epsilon_{ijk} \sigma_k \quad \quad & (34.9) \\ 
  (S_L^{k0})_a^b & = \frac{1}{2} i \sigma_k \quad \quad & (34.10) 
\end{aligned}
\]
we want $(S_L^{\mu \nu})_a^b = \frac{i}{4} ( \sigma^{\mu} \overline{\sigma}^{\nu} - \sigma^{\nu } \overline{\sigma}^{\mu} )_a^b$ \quad \, (35.21).  

Now $(S_L^{ij})_a^b = \frac{i}{4} ( \sigma^i \overline{\sigma}^j - \sigma^j \overline{\sigma}^i )_a^b$ \\
If $i =j$, \\
$\begin{aligned} \sigma^{i} \overline{\sigma}^i &= 1 + -4 = -3  \\  \sigma^{j} \overline{\sigma}^j &= 1 + -4 = -3  \end{aligned} \quad \, \Longrightarrow (S_L^{ij})_a^b = 0$.

\[
\begin{aligned}
  (S_L^{\ij})_a^b & = \frac{i}{4} ( -i \epsilon_{ijk} \sigma^k + i \epsilon_{jik} \sigma^k )_a^b = \frac{1}{2} \epsilon_{ijk} \sigma^k \\ 
  (S_L^{k0})_a^b & = \frac{i}{4} ( \sigma^{\mu} \overline{\sigma}^0 - \sigma^0 \overline{\sigma}^{\mu} )_a^b = \frac{i}{2} (\sigma^{\mu})_a^b
\end{aligned}
\]



\solutionhead{35.3} 
\[
\begin{aligned}
(S_L^{\mu \nu} )_a^b & = \frac{i}{4} ( \sigma^{\mu} \overline{\sigma}^{\nu} - \sigma^{\nu} \overline{\sigma}^{\mu})_a^b = \frac{i}{4} ( (\sigma^{\mu} \overline{\sigma}^{\nu})_a^b - (\sigma^{\nu} \overline{\sigma}^{\mu} )_a^b )  \\
((S_L^{\mu \nu})_a^b)^* & = \frac{-i}{4} ( ((\sigma^{\mu} \overline{\sigma}^{\nu})_a^b)^* - ((\sigma^{\nu} \overline{\sigma}^{\mu} )_a^b)^*) 
\end{aligned}
\]
Now recall that Hermetian conjugation swaps the 2 SU(2) Lie algebras.  
\[
\begin{aligned}
  ((\sigma^{\mu})_{ak}(\overline{\sigma}^{\nu})^{kb} )^* = (\sigma^{\mu})^*_{ak}((\overline{\sigma}^{\nu})^{kb})^* = (\sigma^{\mu})_{\dot{k}\dot{a}}(\overline{\sigma}^{\nu})^{\dot{b}\dot{k}} = (\overline{\sigma}^{\nu} \sigma^{\mu})_{\dot{a}}^{\dot{b}} \\ 
  ((\sigma^{\nu})_{ak}(\overline{\sigma}^{\mu})^{kb} )^* = (\sigma^{\nu})^*_{ak}((\overline{\sigma}^{\mu})^{kb})^* = (\overline{\sigma}^{\mu})_{\dot{b}\dot{k}}(\sigma^{\nu})^{\dot{k}\dot{a}} = (\overline{\sigma}^{\mu} \sigma^{\nu})_{\dot{a}}^{\dot{b}} 
\end{aligned}
\]
\[
\Longrightarrow ((S_L^{\mu \nu})_a^b)^* = -\frac{i}{4} ( (\overline{\sigma}^{\nu} \sigma^{\mu})_{\dot{a}}^{\dot{b}} - (\overline{\sigma}^{\mu} \sigma^{\nu} )_{\dot{a}}^{\dot{b}} ) = - ( -\frac{i}{4} ( \overline{\sigma}^{\mu} \sigma^{\nu} - \overline{\sigma}^{\nu} \sigma^{\mu})_{\dot{a}}^{\dot{b}} ) = - (S_R^{\mu \nu})_{\dot{a}}^{\dot{b}}
\]





\solutionhead{35.4}  Given $\epsilon^{ab} \epsilon^{\dot{a} \dot{b}} \sigma_{a\dot{a}}^{\mu} \sigma^{\nu}_{b\dot{b}} = (\sigma^{\mu}_{a\dot{a}} \epsilon^{\dot{a} \dot{b}} )(\epsilon^{ab} \sigma^{\nu}_{b\dot{b}} )$,  \\
from previous problem 35.1, 
\[
\epsilon^{ab} \sigma^{\nu}_{b\dot{b}} = \begin{cases} i (\sigma^2)_{a\dot{b}} & \text{ for $\nu =0$ } \\ (-\epsilon_{2 \nu \rho} \sigma^{\rho} + i \delta^{2\nu} 1 )_{a\dot{b}} & \text{ if } \nu \neq 0 \end{cases}
\]
Likewise,
\[
\sigma^{\mu}_{a\dot{a}} \epsilon^{\dot{a} \dot{b}} = i (\sigma^{\mu} \sigma^2 )_{a\dot{b}} = \begin{cases} i (\sigma^2)_{a\dot{b}} & \text{ if } \mu = 0 \\ i ( i \epsilon_{\mu 2 \rho } \sigma^{\rho} + \delta^{\mu 2} 1 )_{a\dot{b}} = -\epsilon_{\mu 2 \rho} \sigma^{\rho}_{a \dot{b}} + i \delta^{\mu 2 } 1_{a\dot{b}} & \text{ if } \mu \neq 0 \end{cases}
\]

If $\mu = \nu =0$, $i (\sigma^2)_{a\dot{b}} i (\sigma^2)_{a\dot{b}} = - (\sigma^2)_{a\dot{b}} (\sigma^2)_{\dot{b} a}(-1) = (1) 1_{aa} = 2$ (it's the trace).  

If $\mu \neq 0$, $\nu =0$,
\[
\begin{gathered}
(-\epsilon_{\mu 2 \rho} \sigma^{\rho}_{a\dot{b}} + i \delta^{\mu 2} 1_{a\dot{b}} )i ( \sigma^2)_{a\dot{b}} = + i \epsilon_{\mu 2 \rho} \sigma^{\rho}_{a\dot{b}} \sigma^2_{\dot{b} a} - \delta^{\mu 2} 1_{a\dot{b}} (- \sigma^2_{\dot{b} a} ) = \\
	= i \epsilon_{\mu 2 \rho} \epsilon_{\rho 2 \alpha} \sigma^{\alpha}_{aa} + \delta^{\mu 2} \sigma^2_{aa} = 0 \quad \, \text{(trace of Pauli matrices is 0)}
\end{gathered}
\]

Likewise, if $\nu \neq 0$, $\mu = 0$, 
\[
i(\sigma^2)_{a\dot{b}} (-\epsilon_{2\nu \rho} \sigma^{\rho}_{a\dot{b}} + i \delta^{2\nu} 1_{a\dot{b}} ) = i \epsilon_{2\nu \rho} \sigma^2_{\dot{b} a} \sigma^{\rho}_{a\dot{b}} + \delta^{2\nu} \sigma^2_{\dot{b} \dot{b}} = i \epsilon_{2\nu \rho} \epsilon_{2\rho \alpha} \sigma^{\alpha}_{\dot{b} \dot{b}} + \delta^{2\nu} \sigma^2_{\dot{b} \dot{b}} = 0 
\]

If $\mu, \nu \neq 0 ,2$, 
\[
\begin{gathered}
	(- \epsilon_{\mu 2 \rho} \sigma_{a\dot{b}}^{\rho} )( - \epsilon_{2\nu \alpha} \sigma^{\alpha}_{a\dot{b}} ) = - \epsilon_{2\mu \rho} \epsilon_{2 \nu \alpha} \sigma^{\rho}_{a\dot{b}} \sigma^{\alpha}_{\dot{b} a} = - \lbrace \delta_{\mu \nu} 1_{aa} + - \sigma^{\nu}_{a\dot{b}} \sigma^{\mu}_{\dot{b} a} \rbrace = \\
	= - \lbrace \delta_{\mu \nu} (2) - i \epsilon_{\mu \nu \alpha} \sigma^{\alpha}_{aa} \rbrace = - 2\delta_{\mu \nu} 
\end{gathered}
\]

If $\mu \neq 0,2$, $\nu =2$, 
\[
(- \epsilon_{\mu 2 \rho} \sigma^{\rho}_{a\dot{b}} )( i 1_{a\dot{b}}) = - i \epsilon_{\mu 2 \rho} \sigma^{\rho}_{aa} = 0 
\]

Likewise, for $\nu \neq 0,2$, $\mu =2$, 
\[
(-\epsilon_{2\nu \rho} \sigma^{\rho} )_{a\dot{b} }( i 1_{a\dot{b}} ) = - i\epsilon_{2\nu \rho} \sigma^{\rho}_{a\dot{b}} 1_{\dot{b} a} = 0 
\]

If $\mu = \nu = 2$, 
\[
i 1_{a\dot{b} } i 1_{a\dot{b}} = -1_{aa} = -2 
\]

In my convention, $g_{\mu \nu} = \left( \begin{matrix} 1 & & & \\ & -1 & & \\ & & -1 & \\ & & & -1 \end{matrix} \right)$.  So I obtain 
\[
\boxed{ \epsilon^{ab} \epsilon^{\dot{a} \dot{b}} \sigma^{\mu}_{a\dot{a}} \sigma^{\nu}_{b \dot{b}} = +2 g^{\mu \nu} }
\]



\section{Lagrangians for spinor fields}

\subsection*{Problems}

\problemhead{36.1}



\problemhead{36.2} Verify that eq. (36.46) is consistent with eq. (36.43).

\solutionhead{36.2}

Recall that $\gamma^{\mu} = \left( \begin{matrix} 0 & \sigma^{\mu}_{e\dot{a}} \\ \overline{\sigma}^{\mu \dot{e} a} & 0 \end{matrix} \right) \quad \, (36.39)$, $\gamma_5 = \left( \begin{matrix} - \delta_a^c & \\ & +\delta^{\dot{a}}_{\dot{c}}  \end{matrix} \right) \quad \, (36.43) $.  

Rewrite the gamma matrices as tensor products:
\[
\begin{aligned}
  & \gamma^0 = \left( \begin{matrix} & 1 \\ 1 & \end{matrix} \right) = 1 \otimes \sigma^1 \\ 
  & \gamma^{\mu} = \left( \begin{matrix} & \sigma^{\mu} \\ -\sigma^{\mu} & \end{matrix} \right) = \sigma^{\mu} \otimes i \sigma^2 \quad \, \mu = 1,2,3
\end{aligned}
\]
Then multiply tensor products:
\[
\begin{aligned}
  \gamma^1 \gamma^2 & = (\sigma^1 \otimes i \sigma^2 )( \sigma^2 \otimes i \sigma^2) = \sigma^1 \sigma^2 \otimes - \sigma^2 \sigma^2 = i \sigma^3 \otimes -1 \\ 
  \gamma^1 \gamma^2 \gamma^3 & = - i (\sigma^3 \otimes 1)(\sigma^3 \otimes i \sigma^2) = (+1)(1 \otimes \sigma^2) \\ 
  \gamma^0 \gamma^1 \gamma^2 \gamma^3 & = i (1 \otimes i \sigma^3)= - (1 \otimes \sigma^3) = \left( \begin{matrix} -1 & \\ & 1 \end{matrix} \right)
\end{aligned}
\]




\problemhead{36.3} \begin{itemize} \item[(a)] Prove the \emph{Fierz identities}  
\[
\begin{aligned}
(\chi^{\dag}_1 \overline{\sigma}^{\mu} \chi_2)(\chi_3^{\dag} \overline{\sigma}_{\mu} \chi_4) & = -2 (\chi_1^{\dag} \chi_3^{\dag})(\chi_2 \chi_4), \quad \quad \quad \quad & (36.58) \\ 
(\chi^{\dag}_1 \overline{\sigma}^{\mu} \chi_2)(\chi_3^{\dag} \overline{\sigma}_{\mu} \chi_4) & (\chi_1^{\dag} \overline{\sigma}^{\mu} \chi_4)(\chi_3^{\dag} \overline{\sigma}_{\mu} \chi_2) , \quad \quad \quad \quad & (36.59)
\end{aligned}
\]
\item[(b)] Define fields 
\[
\Psi_i \equiv \left( \begin{matrix} \chi_i \\ \xi_i^{\dag} \end{matrix} \right), \quad \quad \quad \Psi_i^C \equiv \left( \begin{matrix} \xi_i \\ \chi_i^{\dag} \end{matrix} \right)
\]
\end{itemize}


\solutionhead{36.3} \begin{itemize}
\item[(a)] Consider $(\chi^{\dag}_1 \overline{\sigma}^{\mu} \chi_2)(\chi_3^{\dag} \overline{\sigma}_{\mu} \chi_4)$.  Note that $\chi_1^{\dag} \overline{\sigma}^{\mu} \chi_2$ is a scalar, i.e. $\text{column matrix} \cdot \text{row} = \text{ number }$.  
\[
\begin{aligned}
	(\chi^{\dag}_1 \overline{\sigma}^{\mu} \chi_2)(\chi_3^{\dag} \overline{\sigma}_{\mu} \chi_4) & 	= \chi_{1\dot{a}}^{\dag} \overline{\sigma}^{\mu \dot{a} a} \chi_{2a} \chi^{\dag}_{3\dot{c}} \overline{\sigma}_{\mu}^{\dot{c} c} \chi_{4c} = \chi_{1\dot{a}}^{\dag} \chi_{3\dot{c}}^{\dag} \chi_{2a} \chi_{4c} \epsilon^{ab} \epsilon^{\dot{a}\dot{b}} \epsilon^{cd} \epsilon^{\dot{c} \dot{d}} \sigma^{\mu}_{b\dot{b}} \sigma_{\mu d \dot{d}} = \\
	& = \chi_{1\dot{a}}^{\dag} \chi_{3\dot{c}}^{\dag} \chi_{2a} \chi_{4c} \epsilon^{ab} \epsilon^{\dot{a} \dot{b}} \epsilon^{cd} \epsilon^{\dot{c} \dot{d}} (-2\epsilon_{bd} \epsilon_{\dot{b} \dot{d}} )  \\
		& \text{ since, recall that $\sigma_{a\dot{a}}^{\mu} \sigma_{\mu b\dot{b}} = -2 \epsilon_{ab} \epsilon_{\dot{a} \dot{b}}$  (35.4) } \\ 
			& = -2 \chi_{1\dot{a}}^{\dag} \chi_{3\dot{c}}^{\dag} \chi_{2a} \chi_{4c} \epsilon^{ab} \epsilon^{\dot{a} \dot{b}} (-\delta_b^c )(-\delta_{\dot{b}}^{\dot{c}}) = \\
				& = -2(\chi_{1\dot{a}}^{\dag} \chi_{3\dot{c}}^{\dag} (- \epsilon^{\dot{c} \dot{a}} ) )(\chi_{2a} \chi_{4c} (-\epsilon^{ca}) ) = -2(\chi_{1\dot{a}}^{\dag} \chi_3^{\dag \dot{a}} )( \chi_{2a} \chi_4^a ) = \\
					& = -2 (\chi_1^{\dag} \chi_3^{\dag} )(\chi_2 \chi_4)  
\end{aligned}
\]


$\begin{aligned} (\chi_1^{\dag} \overline{\sigma}^{\mu} \chi_2)(\chi_3^{\dag} \overline{\sigma}_{\mu} \chi_4) & = -2 (\chi_1^{\dag} \chi_3^{\dag} )(\chi_2 \chi_4) \\ (\chi_1^{\dag} \overline{\sigma}^{\mu} \chi_4)(\chi_3^{\dag} \overline{\sigma}_{\mu} \chi_2) & = -2 (\chi_1^{\dag} \chi_3^{\dag} )(\chi_4 \chi_2) \end{aligned}$
\[
\begin{aligned}
	\chi_4 \chi_2 & = \chi_{4a} \chi_2^a = \epsilon_{ab} \chi_4^b \epsilon^{ac} \chi_{2c} = -\epsilon_{ba} \epsilon^{ac} ( - \chi_{2c} \chi_4^b ) = \\
		& = \delta_b^c \chi_{2c} \chi_4^b = \chi_{2b} \chi_4^b \text{ or } \chi_{2a} \chi_4^a = \chi_2 \chi_4
\end{aligned}
\]
Components of Weyl fields themselves anticommute, so $(\chi_4^b \chi_{2c} = - \chi_{2c} \chi_4^b)$.  

\emph{Pedagogically}, I figured the last identity by doing a \emph{simple example}:

\[
\begin{aligned} \chi_1 & = \left( \begin{matrix} a \\ b \end{matrix} \right) \\ \chi_1^{\dag} & = (\begin{matrix} \overline{a} & \overline{b} \end{matrix} ) \end{aligned} \quad \quad \quad \begin{aligned} \chi_2 & = \left( \begin{matrix} c \\ d \end{matrix} \right) \\ \chi_3^{\dag} & = ( \begin{matrix} \overline{e} & \overline{f} \end{matrix} ) \\ \chi_4 & = \left( \begin{matrix} g \\ h \end{matrix} \right) \end{aligned}
\] 

For the LHS of the identity:

For $\sigma^0 =1$, 
\[
\begin{aligned}
  (\begin{matrix} \overline{a} & \overline{b} \end{matrix}) \left( \begin{matrix} c \\ d \end{matrix}  \right) = (\overline{a}c + \overline{b} d ) \\ 
  (\begin{matrix} \overline{e} & \overline{f} \end{matrix}) \left( \begin{matrix} g \\ h \end{matrix}  \right) = (\overline{e}g + \overline{f} h ) \\ 
\end{aligned}
\]
\[
(\overline{a} c + \overline{b} d )(\overline{e} g + \overline{f} h) = \overline{a} c \overline{e} g + \overline{b} d \overline{e} g + \overline{a} c \overline{f} h + \overline{b} d\overline{f} h 
\]

For $\sigma^1 =\left( \begin{matrix} & 1 \\ 1 & \end{matrix} \right)$, 
\[
\begin{aligned}
  (\begin{matrix} \overline{a} & \overline{b} \end{matrix}) \left( \begin{matrix} d \\ c \end{matrix}  \right) = (\overline{a}d + \overline{b} c ) \\ 
  (\begin{matrix} \overline{e} & \overline{f} \end{matrix}) \left( \begin{matrix} h \\ g \end{matrix}  \right) = (\overline{e}h + \overline{f} g ) \\ 
\end{aligned}
\]
\[
(\overline{a} d + \overline{b} c )(\overline{e} h + \overline{f} g) = \overline{a} d \overline{e} h + \overline{b} c \overline{e} h + \overline{a} d \overline{f} g + \overline{b} c\overline{f} g 
\]

For the RHS of the identity, 

\[
\begin{gathered}
  ( \begin{matrix} \overline{a} & \overline{b} \end{matrix} ) \left( \begin{matrix} g \\ h \end{matrix} \right) ( \begin{matrix} \overline{e} & \overline{f} \end{matrix} ) \left( \begin{matrix} c \\ d \end{matrix} \right) = \\ 
  = (\overline{a} g + \overline{b} h )(\overline{e} c + \overline{f} d) = \overline{a} g \overline{e} c +\overline{b} h \overline{e} c + \overline{a} g \overline{f} d + \overline{b} h \overline{f} d
\end{gathered}
\]

\[
\begin{gathered}
  ( \begin{matrix} \overline{a} & \overline{b} \end{matrix} ) \left( \begin{matrix} h \\ g \end{matrix} \right) ( \begin{matrix} \overline{e} & \overline{f} \end{matrix} ) \left( \begin{matrix} d \\ c \end{matrix} \right) = \\ 
  = (\overline{a} h + \overline{b} g )(\overline{e} d + \overline{f} c) = \overline{a} h \overline{e} d +\overline{b} g \overline{e} d + \overline{a} h \overline{e} d + \overline{b} g \overline{f} c
\end{gathered}
\]


\item[(b)] Now $\Psi_i = \left( \begin{matrix} \chi_i \\ \xi_i^{\dag} \end{matrix} \right) = \left( \begin{matrix}  \chi_{ic} \\ \xi_i^{\dag \dot{c}} \end{matrix} \right)$.  (Dirac field)  \\

Consider $(\overline{\Psi}_1 \gamma^{\mu} P_L \Psi_2)(\overline{\Psi}_3 \gamma_{\mu} P_L \Psi_4)$.  

Recall the left and right projection matrices.  
\[
\begin{aligned}
  & P_L \equiv \frac{1}{2} ( 1 - \gamma_5)   = \left( \begin{matrix} \delta_a^c & \\ & \end{matrix} \right) \xrightarrow{ \text{ for a Dirac field } } P_L \Psi = \left( \begin{matrix} \chi_c \\ 0 \end{matrix} \right) \\ 
  & P_R \equiv \frac{1}{2} ( 1 + \gamma_5 )  = \left( \begin{matrix} & \\ & \delta^{\dot{a}}_{\dot{c}} \end{matrix} \right)
\end{aligned}
\]

Recall $\gamma$-matrix definition, $\gamma^{\mu}  = \left( \begin{matrix} & \sigma^{\mu}_{a\dot{c}} \\ \overline{\sigma}^{\mu \dot{a} c} & \end{matrix} \right)$.  
\[
\Longrightarrow \gamma^{\mu} P_L \Psi_2 = \left( \begin{matrix} 0 \\ \overline{\sigma}^{\mu \dot{a} c} \chi_{2c} \end{matrix} \right)
\]

Now recall $\overline{\Psi} = \Psi^{\dag} \beta = ( \xi^a, \chi_{\dot{a}}^{\dag} )$, where $\beta \equiv \left( \begin{matrix} 0 & \delta^{\dot{a}}_{\dot{c}}  \\ \delta_a^c & 0 \end{matrix} \right)$.  \\

$\begin{aligned} 
  \overline{\Psi}_1 \gamma^{\mu} P_L \Psi_2 & = \chi_{1\dot{a}}^{\dag} \overline{\sigma}^{\mu \dot{a} c } \chi_{2c}  \\
  \Longrightarrow \overline{\Psi}_3 \gamma_{\mu} P_L \Psi_4 & = \chi_{3\dot{a}}^{\dag} \overline{\sigma}_{\mu}^{\dot{a} c} \chi_{4c} \text{ where } \gamma_{\mu} = \left( \begin{matrix} & \sigma_{\mu a\dot{c}} \\ \overline{\sigma}_{\mu}^{\dot{a} c} & \end{matrix} \right)
\end{aligned}$ 

Recall charge conjugate of $\Psi$ $= \Psi^c \equiv \mathcal{C} \overline{\Psi}^T = \left( \begin{matrix} \xi_a \\ \chi^{\dag \dot{a}} \end{matrix} \right)$.  Notice how $\chi, \xi$ are switched.  
\[
\Longrightarrow P_R \Psi^c = \left( \begin{matrix} 0 \\ \chi^{\dag \dot{a}} \end{matrix} \right)
\]
$\overline{\Psi}_1 P_R \Psi_3^c = \chi_{1\dot{a}}^{\dag} \chi_3^{\dag \dot{a}} = \chi_1^{\dag} \chi_3^{\dag}$

Consider $\overline{\Psi}^c_4 P_L \Psi_2$
\[
P_L \Psi_2 = \left( \begin{matrix} \chi_{2c} \\ 0 \end{matrix} \right)
\]

What is $\overline{\Psi}^c_4$?  We cannot simply apply the definition as outlined in Srednicki (2007), Sec. 36, as it was for $\Psi$: the order of operations change.  But we know that we want to switch $\chi, \xi$ in the end, i.e., \\
\phantom{what} for $\overline{\Psi} = (\xi^a, \chi_{\dot{a}}^{\dag})$, $\overline{\Psi}^c = (\chi^a, \xi_{\dot{a}}^{\dag} )$ switch $\chi, \xi$ in the end, i.e.

\[
\overline{ (\overline{\Psi})} = \beta^{\dag} (\overline{\Psi})^{\dag} = \left( \begin{matrix} & \delta^a_c \\ \delta_{\dot{a}}^{\dot{c}} & \end{matrix} \right) \left( \begin{matrix} \xi^{\dag \dot{a}} \\ \chi_a \end{matrix} \right) = \left( \begin{matrix} \chi_c \\ \xi^{\dag \dot{c}} \end{matrix} \right) = \Psi; \quad \quad \quad \, \overline{ (\overline{\Psi})^T } = (\chi_c, \xi^{\dag \dot{c}}) 
\]

vs. 

\[
(\overline{\Psi}) = \Psi^{\dag} \beta  = (\chi^{\dag}_{\dot{a}}, \xi^a ) \left( \begin{matrix} & \delta^{\dot{a}}_{\dot{c}} \\ \delta_a^c & \end{matrix} \right) = ( \begin{matrix} \xi^c, & \chi^{\dag}_{\dot{c}} \end{matrix} ); \quad \quad \quad \, (\overline{\Psi})^T = \left( \begin{matrix} \xi^c \\ \chi^{\dag}_{\dot{c}} \end{matrix} \right)
\]

\[
\overline{ (\overline{\Psi})^T} \mathcal{C'} = ( \begin{matrix} \chi_c, & \xi^{\dag \dot{c}} \end{matrix} ) \left( \begin{matrix} \epsilon^{ac} & \\  & \epsilon_{\dot{a} \dot{c}} \end{matrix} \right) = ( \begin{matrix} \chi^a, & \xi^{\dag}_{\dot{a}} \end{matrix} ) = \overline{\Psi}^c \quad \quad  \text{ vs. } \quad \quad \mathcal{C}(\overline{\Psi})^T = \left( \begin{matrix} \epsilon_{ac} & \\ & \epsilon^{\dot{a} \dot{c}} \end{matrix} \right) \left( \begin{matrix} \xi^c \\ \chi^{\dag}_{\dot{c}} \end{matrix} \right) = \left( \begin{matrix} \xi_a \\ \chi^{\dag \dot{a}}  \end{matrix} \right) = \Psi^c
\]

Notice how for $\overline{\Psi}^c$, I had to multiply $\Psi^T = \overline{(\overline{\Psi})^T}$ with $\mathcal{C}^1$ to the right.  

Nevertheless, $\overline{\Psi}^c_4 P_L \Psi_2 = \chi_{4a}^c \chi_{2c} = \chi_4 \chi_2 =\chi_2 \chi_4$.  (by (35.25), for left-handed Weyl fields).  

Using the Fierz identity $(\chi_1^{\dag} \overline{\sigma}^{\mu} \chi_2)(\chi_3^{\dag} \overline{\sigma}_{\mu} \chi_4) = -2 (\chi_1^{\dag} \chi_3^{\dag})(\chi_2 \chi_4)$.  
\[
\Longrightarrow (\overline{\Psi}_1 \gamma^{\mu} P_L \Psi_2)(\overline{\Psi}_3 \gamma_{\mu} P_L \Psi_4) = -2( \overline{\Psi_1} P_R \Psi_3^c)(\overline{\Psi}_4^c P_L \Psi_2 )
\]

As before,
\[
\begin{aligned}
  \overline{\Psi}_1 \gamma^{\mu} P_L \Psi_4 & = \chi_{1\dot{a}}^{\dag} \overline{\sigma}^{\mu \dot{a}c} \chi_{4c} & = \chi_1^{\dag} \overline{\sigma}^{\mu} \chi_4 \\ 
  \overline{\Psi}_3 \gamma^{\mu} P_L \Psi_2 & = \chi_{3\dot{a}}^{\dag} \overline{\sigma}^{\mu \dot{a}c} \chi_{2c} & = \chi_3^{\dag} \overline{\sigma}^{\mu} \chi_2 \\ 
\end{aligned}
\]
\[
\Longrightarrow (\overline{\Psi}_1 \gamma^{\mu} P_L \Psi_2 )( \overline{\Psi}_3 \gamma_{\mu} P_L \Psi_4) = (\overline{\Psi}_1 \gamma^{\mu} P_L \Psi_4 )(\overline{\Psi}_3 \gamma_{\mu} P_L \Psi_2 )  \quad \quad \, (36.62)
\]

\item[(c)] \[
\begin{aligned}
  \overline{\Psi}_1 \gamma^{\mu} P_R \Psi_2 & = ( \begin{matrix} \xi_1^a, & \chi_{1\dot{a}}^{\dag} \end{matrix}) \gamma^{\mu} \left( \begin{matrix} 0 \\ \xi_{2}^{\dag \dot{c}} \end{matrix} \right) = (\begin{matrix} \xi_1^a, & \chi_{1\dot{a}}^{\dag} \end{matrix} ) \left( \begin{matrix}  \overline{\sigma}^{\mu}_{ a \dot{c}} \xi_2^{\dag \dot{c}} \\ 0 \end{matrix} \right) = \xi_1^a \sigma^{\mu}_{a\dot{c}} \xi_2^{\dag \dot{c}} = \xi_1 \sigma^{\mu} \xi_2^{\dag} \\ 
  \overline{\Psi}_2^c \gamma^{\mu} P_L \Psi_1^c & = ( \begin{matrix} \chi_2^a, & \xi_{2\dot{a}}^{\dag} \end{matrix}) \gamma^{\mu} \left( \begin{matrix}  \xi_{1c} \\ 0  \end{matrix} \right) = (\begin{matrix} \chi_2^a, & \xi_{2\dot{a}}^{\dag} \end{matrix} ) \left( \begin{matrix} 0 \\ \overline{\sigma}^{\mu \dot{a} c} \xi_{1c} \end{matrix} \right) = \xi_{2\dot{a}}^{\dag} \overline{\sigma}^{\mu \dot{a}c} \xi_{1c} = \\
    & = \xi_{2\dot{a}}^{\dag} \epsilon^{cb} \epsilon_{\dot{a}\dot{b}} \sigma^{\mu}_{b\dot{b}} \xi_{1c}   = \\
    & = (-\epsilon^{\dot{b} \dot{a}} \xi_{2\dot{a}}^{\dag})\sigma^{\mu}_{b\dot{b}} (- \epsilon^{bc} \xi_{1c} ) = \xi_{2}^{\dag \dot{b}} \sigma^{\mu b\dot{b}} \xi_1^b = \xi_1^b \sigma^{\mu}_{b\dot{b}} \xi_2^{\dag \dot{b}} = -\xi_1 \sigma^{\mu} \xi_2^{\dag}
\end{aligned}
\]
($\xi_2^{\dag \dot{b}}$, $\xi_1^b$ will anticommute with each other)
\[
\Longrightarrow \boxed{ \overline{\Psi}_1 \gamma^{\mu} P_R \Psi_2 = - \overline{\Psi}_2^c \gamma^{\mu} P_L \Psi_1^c }
\]


$
\begin{aligned}
  \overline{\Psi}_1 P_L \Psi_2 & = ( \begin{matrix} \xi_1^a , & \chi_{1\dot{a}}^{\dag} \end{matrix} ) \left( \begin{matrix} \chi_{2a} \\ 0 \end{matrix} \right) = \xi_1 \chi_2 \\ 
  \overline{\Psi}_2^c P_L \Psi_1^c & = ( \begin{matrix} \chi_1^a , & \xi_{2\dot{a}}^{\dag} \end{matrix} ) \left( \begin{matrix} \xi_{1a} \\ 0 \end{matrix} \right) = \chi_2^a \xi_{1a} = \chi_2 \xi_1 = \xi_1 \chi_2 \quad \quad \, \text{($\xi_1, \chi_2$ are left-handed Weyl fields and by (35.25))}
\end{aligned}
$
\[
\Longrightarrow \boxed{ \overline{\Psi}_1 P_L \Psi_2 = \overline{\Psi}_2^c P_L \Psi_1^c }
\]
\[
\begin{aligned}
  \overline{\Psi}_1 P_R \Psi_2 & = (\begin{matrix} \xi_1^a, & \chi_{1\dot{a}}^{\dag} \end{matrix} ) \left( \begin{matrix} 0 \\ \xi_2^{\dag \dot{a}} \end{matrix} \right) = \chi^{\dag}_{1\dot{a}} \xi_2^{\dag \dot{a}} = \chi_1^{\dag} \xi_2^{\dag} = \xi^{\dag}_2 \chi_1^{\dag}  \\ 
  \overline{\Psi}_2^c P_R \Psi_1^c & = (\begin{matrix} \xi_2^a, & \chi_{2\dot{a}}^{\dag} \end{matrix} ) \left( \begin{matrix} 0 \\ \chi_1^{\dag \dot{a}} \end{matrix} \right) = \xi^{\dag}_{2\dot{a}} \chi_1^{\dag \dot{a}} = \xi_2^{\dag} \chi_1^{\dag} 
\end{aligned}
\]
\[
\Longrightarrow \boxed{ \overline{\Psi}_1 P_R \Psi_2 = \overline{\Psi}_2^c P_R \Psi_1^c }
\]
\end{itemize}





\problemhead{36.5} \emph{Symmetries of fermion fields}.  (Prerequisite: 24.)  Consider a theory with $N$ massless Weyl fields $\psi_j$, 
\[
\mathcal{L} = i \psi_j^{\dag} \sigma^{\mu} \partial_{\mu} \psi_j \quad \quad \quad \, (36.74)
\]
where the repeated index $j$ is summed.  This lagrangian is clearly invariant under the $U(N)$ transformation, 
\[
\psi_j \to U_{jk} \psi_k, \quad \quad \quad \,  (36.75)
\]
where $U$ is a unitary matrix.  State the invariance group for the following cases:
\begin{itemize}
\item[(a)] $N$ Weyl fields with a common mass $m$, 
\[
\mathcal{L} = i \psi_j^{\dag} \sigma^{\mu} \partial_{\mu} \psi_j - \frac{1}{2} m (\psi_j \psi_j + \psi_j^{\dag} \psi_j^{\dag} ).  \quad \quad \quad \, (36.76)
\]
\item[(b)] $N$ massless Majorana fields,
\[
\mathcal{L} = \frac{1}{2} \Psi_j^T \mathcal{C} \gamma^{\mu}  \partial_{\mu} \Psi_j \quad \quad \quad \, (36.77)
\]
\item[(c)] $N$ Majorana fields with a common mass $m$, 
\[
\mathcal{L} = \frac{i}{2} \Psi_j^T \mathcal{C} \gamma^{\mu} \partial_{\mu} \Psi_j - \frac{1}{2} m \Psi_j^T \mathcal{C} \Psi_j.  \quad \quad \quad \, (36.78) 
\]
\item[(d)] $N$ massless Dirac fields, 
\[
\mathcal{L} = i \overline{\Psi}_j \gamma^{\mu} \partial_{\mu} \Psi_j \quad \quad \quad \, (36.79)  
\]
\item[(e)] $N$ Dirac fields with a common mass $m$,
\[
\mathcal{L} = i \overline{\Psi}_j \gamma^{\mu} \partial_{\mu} \Psi_j - m \overline{\Psi}_j \Psi_j \quad \quad \quad \, (36.80)   
\]
\end{itemize}

\solutionhead{36.5} 
\begin{itemize}
\item[(a)] Consider $SO(N)$ transformation.  Note that $\psi_j \psi_j = \psi^T \psi = 1$.  

\[
\begin{gathered}
  \psi_j' \psi_j' = \psi^{'T} \psi' = (R\psi)^T R\psi = \psi^T R^T R \psi = \psi^T \psi
\end{gathered}
\]

\[
\begin{gathered}
  \psi^{'\dag}_j \psi^{'\dag}_j = (\psi^{' \dag})^T \psi^{'\dag} = ((R\psi)^{\dag})^T (R\psi)^{\dag} = (\psi^{\dag} R^T)^T \psi^{\dag} R^{\dag} = R\psi^{\dag T} \psi^{\dag} R^T = (\psi^{\dag T} \psi^{\dag}) R R^T = \psi^{\dag T} \psi^{\dag} = \psi^{\dag}_j \psi_j^{\dag}
\end{gathered}
\]

Note that for column vectors, 

\[
\begin{gathered}
\psi'_i = (R\psi)_i  = R_{ij} \psi_j \text{ or }  
 \psi'_{i1} = R_{ij} \psi_{j1}
\end{gathered}
\]

but for row vectors, one must be careful about the indices:

\[
\begin{gathered}
  \psi^{ ' \dag}_i = (R\psi)_i^{\dag} = (\psi^{\dag} R^{\dag})_i = \psi_k^{\dag} R^T_{ki} = \psi^{\dag}_{k} R_{ik} \text{ or }  
  \psi^{' \dag}_{1i} = (R\psi)_{1i}^{\dag} = (\psi^{\dag} R^{\dag})_{1i} = \psi_{1k}^{\dag} R^T_{ki} = \psi^{\dag}_{1k}R_{ik}
\end{gathered}
\]

\[
i \psi^{' \dag}_i \sigma^{\mu} \partial_{\mu} \psi_i' = i \psi^{\dag}_k R_{ik} \sigma^{\mu} \partial_{\mu} R_{ij} \psi_j = i R_{ik} R_{ij} \psi^{\dag}_k \sigma^{\mu} \partial_{\mu} \psi_j = i \psi^{\dag}_j \sigma^{\mu} \partial_{\mu} \psi_j
\]
$\Longrightarrow SO(N)$ symmetry.  

Also, $N$ Weyl fields with a common mass $m$ has $Z_2$ symmetry because the $Z_2$ transformation $\psi \to -\psi$, and $\psi^{\dag} \to - \psi^{\dag}$ still leaves $\mathcal{L}$ invariant.  

With $Z_2$ symmetry, $SO(N)$ is ``extended'' to $O(N)$ symmetry, ``overall,'' for the Lagrangian of $N$ Weyl fields.  
\item[(b)] 
\item[(c)]
\item[(d)] Given $\mathcal{L} = i \overline{\Psi}_j \gamma^{\mu} \partial_{\mu} \Psi_j$ for $N$ massless Dirac fields, recall that for the Dirac field, $\Psi = \left( \begin{matrix} \chi_c \\ \xi^{\dag \dot{c}} \end{matrix} \right)$.  

Now 
\[
\begin{gathered}
  \overline{\Psi} \gamma^{\mu} \partial_{\mu} \Psi = ( \xi^a, \chi_{\dot{a}}^{\dag} ) \left( \begin{matrix} & \sigma^{\mu}_{a\dot{c}} \\ \overline{\sigma}^{\mu \dot{a} c} & \end{matrix} \right) \left( \begin{matrix} \partial_{\mu} \chi_{c} \\ \partial_{\mu} \xi^{\dag \dot{c}} \end{matrix} \right) =  \\
    = (\chi_{\dot{a}}^{\dag} \overline{\sigma}^{\mu \dot{a} c} , \chi^a \sigma^{\mu}_{a\dot{c}} ) \left( \begin{matrix} \partial_{\mu} \chi_c \\ \partial_{\mu} \xi^{\dag \dot{c}} \end{matrix} \right) = \chi^{\dag}_{\dot{c}} \overline{\sigma}^{\mu \dot{a} c} \partial_{\mu} \chi_c + \xi^a \sigma^{\mu}_{a\dot{c} } \partial_{\mu} \xi^{\dag \dot{c}}
\end{gathered}
\]

Now 
\[
\xi^a \sigma^{\mu}_{a\dot{c}} \partial_{\mu} \xi^{\dag \dot{c}} = \xi^c \sigma^{\mu}_{c\dot{c}} \partial_{\mu} \xi^{\dag \dot{c}} = \xi^c \epsilon_{ca} \epsilon_{\dot{c} \dot{a}} \overline{\sigma}^{\mu \dot{a} a} \partial_{\mu} \xi^{\dag \dot{c}} = \xi_a \overline{\sigma}^{\mu \dot{a} a} \partial_{\mu} \xi^{\dag \dot{a}} = \xi \overline{\sigma}^{\mu} \partial_{\mu} \xi^{\dag}
\]

and $(\partial_{\mu} \xi) \overline{\sigma}^{\mu} \xi^{\dag} = \xi^{\dag} \overline{\sigma}^{\mu} (\partial_{\mu} \xi)$ since $\xi$, $\xi^{\dag}$ anticommute (since $\xi$ is spin-$\frac{1}{2}$ field)

\[
\Longrightarrow \xi \overline{\sigma}^{\mu} \partial_{\mu} \xi^{\dag} = \xi^{\dag} \overline{\sigma}^{\mu} \partial_{\mu} \xi + \partial_{\mu} (\xi \overline{\sigma}^{\mu} \xi^{\dag} )
\]
So
\[
\overline{\Psi} \gamma^{\mu} \partial_{\mu} \Psi = \chi^{\dag} \overline{\sigma}^{\mu} \partial_{\mu} \chi + \xi^{\dag} \overline{\sigma}^{\mu} \partial_{\mu} \xi 
\]
(up to a total divergence, that will be integrated out).  

So then

\[
\mathcal{L} = i \overline{\Psi}_j \gamma^{\mu} \partial_{\mu} \Psi_j = i \chi^{\dag}_j \overline{\sigma}^{\mu} \partial_{\mu} \chi_j + i \xi^{\dag}_j \sigma^{\mu} \partial_{\mu} \xi_j
\]

Now let 
\[
\begin{aligned}
  \psi_{\alpha} & = \chi_j \text{ for } \alpha = j = 1, \dots N \text{ and } 
  \psi_{\alpha} & = \xi_j \text{ for } \alpha - N = j, \text{ for } \alpha = N+1, \dots 2N
\end{aligned}
\]

I was encouraged to do this by the form of $\mathcal{L}$ where we seemed to have 2 sets of $N$ Weyl fields.  

Then 
\[
\begin{aligned}
  \mathcal{L} & = i \chi_j^{\dag} \overline{\sigma}^{\mu} \partial_{\mu} \chi_j + i \xi^{\dag}_j \sigma^{\mu} \partial_{\mu} \xi_j = \\
  & = i \psi_{\alpha}^{\dag} \overline{\sigma}^{\mu} \partial_{\mu} \psi_{\alpha}
\end{aligned}
\]  

For $\psi \to U \psi$,
\[
\begin{aligned}
  & \psi^{'\dag}_{\alpha} = (U\psi)^{\dag}_{1\alpha} = (\psi^{\dag}U^{\dag})_{1\alpha} = \psi^{\dag}_{1k} U^{\dag}_{k\alpha} = \psi^{\dag}_{\alpha} U^{\dag} \\
  &  \psi'_{\alpha} = (U\psi)_{\alpha} = U_{\alpha j}\psi_j = U \psi_j
\end{aligned}
\]

Then 
\[
\psi_{\alpha}^{'\dag} \overline{\sigma}^{\mu} \partial_{\mu} \psi'_{\alpha} = \psi_{1k}^{\dag} U^{\dag}_{k\alpha} \overline{\sigma}^{\mu} \partial_{\mu} ( U_{\alpha j} \psi_j ) = \psi_{1k}^{\dag} U^{\dag}_{k\alpha} U_{\alpha j} \overline{\sigma}^{\mu} \partial_{\mu} \psi_j = \psi_{1k}^{\dag} \delta_{kj} \overline{ \sigma}^{\mu} \partial_{\mu} \psi_j = \psi^{\dag}_j \overline{\sigma}^{\mu} \partial_{\mu} \psi_j
\]

Now $det{\widetilde{U}} = 1$ was not specified.  Then $\mathcal{L}$ has $U(2N)$ symmetry.  

\item[(e)]
Now $\overline{\Psi} \Psi = \xi \chi + \xi^{\dag} \chi^{\dag}$.  So an $U(2N)$ or $O(2N)$ transformation cannot be applied directly since the $\overline{\Psi} \Psi$ terms ``mixes together'' $\xi$, $\chi$.  $\xi$ and $\chi$ are considered to be 2 different Weyl fields and transform each in their own way.  

Let
\[
\begin{aligned}
  \xi_j & = \frac{  \varphi_{j1} - i \varphi_{j2} }{\sqrt{2}} \\
  \chi_j & = \frac{ \varphi_{j1} + i \varphi_{j2} }{\sqrt{2}}
\end{aligned}
\]
(inspired by Lipstein's solution to Problem Set 3, Winter term 2009, Ph205b) where $\varphi_{j1}$, $\varphi_{j2}$ are real, scalar fields.  

Then
\[
\xi_j \chi_j = \left( \frac{ \varphi_{j1} - i \varphi_{j2} }{ \sqrt{2}} \right) \left( \frac{ \varphi_{j1} + i \varphi_{j2} }{ \sqrt{2}} \right)  = \frac{ \varphi_{j1} \varphi_{j1} + \varphi_{j2} \varphi_{j2} }{ 2} 
\]
likewise,
\[
\xi_j^{\dag} \chi_j^{\dag} = \left( \frac{ \varphi_{j1} + i \varphi_{j2} }{ \sqrt{2}} \right) \left( \frac{ \varphi_{j1} - i \varphi_{j2} }{ \sqrt{2}} \right)  = \frac{ \varphi_{j1} \varphi_{j1} + \varphi_{j2} \varphi_{j2} }{ 2} 
\]

\[
\begin{aligned}
  \chi_j^{\dag} \overline{\sigma}^{\mu} \partial_{\mu} \chi_j & =  \left( \frac{ \varphi_{j1} - i \varphi_{j2} }{ \sqrt{2}} \right) \overline{\sigma}^{\mu} \partial_{\mu} \left( \frac{ \varphi_{j1} + i \varphi_{j2} }{ \sqrt{2}} \right)  = \\
  & = \frac{1}{2} \left( \varphi_{j1} \overline{\sigma}^{\mu} \partial_{\mu} \varphi_{j1} + i \varphi_{j1} \overline{\sigma}^{\mu} \partial_{\mu}\varphi_{j2} - i \varphi_{j2} \overline{\sigma}^{\mu} \partial_{\mu} \varphi_{j1} + \varphi_{j2} \overline{\sigma}^{\mu} \partial_{\mu} \varphi_{j2} \right) \\
  \xi_j^{\dag} \overline{\sigma}^{\mu} \partial_{\mu} \xi_j & =   \frac{1}{2} \left( \varphi_{j1} \overline{\sigma}^{\mu} \partial_{\mu} \varphi_{j1} - i \varphi_{j1} \overline{\sigma}^{\mu} \partial_{\mu}\varphi_{j2} + i \varphi_{j2} \overline{\sigma}^{\mu} \partial_{\mu} \varphi_{j1} + \varphi_{j2} \overline{\sigma}^{\mu} \partial_{\mu} \varphi_{j2} \right)
\end{aligned}
\]

\[
\Longrightarrow   \chi_j^{\dag} \overline{\sigma}^{\mu} \partial_{\mu} \chi_j     +   \xi_j^{\dag} \overline{\sigma}^{\mu} \partial_{\mu} \xi_j  = \\
\varphi_{j1} \overline{\sigma}^{\mu} \partial_{\mu} \varphi_{j1} + \varphi_{j2} \overline{\sigma}^{\mu} \partial_{\mu} \varphi_{j2}
\]

Then
\[
\mathcal{L} = i (\varphi_{j1} \overline{\sigma}^{\mu} \partial_{\mu} \varphi_{j1}  + \varphi_{j2} \overline{\sigma}^{\mu} \partial_{\mu} \varphi_{j2} ) - m (\varphi_{j1} \varphi_{j1} + \varphi_{j2} \varphi_{j2} ) 
\]

In this form, as previously shown, for $\varphi_j \varphi_j$ terms, $\mathcal{L}$ is invariant under $O(2N)$ transformations (because we now have $2N$ real, scalar fields).  

Explicitly:
\[
\begin{aligned}
  \varphi_j' \varphi_j' & = \varphi^{' T}_j \varphi_j' = (R\varphi)^T_j (R\varphi)_j = (\varphi^T R^T)_j R_{jk} \varphi_k = \varphi_{1l}^T R_{lj}^T R_{jk} \varphi_k = \varphi_{l1} R_{jl} R_{jk} \varphi_k = \\
  & = \varphi_k \varphi_k
\end{aligned}
\]


\end{itemize}



\hrulefill

Complete solution to Problem 36.5; 20100302 

Recall definitions of unitary and orthogonal matrices.  
\[
\begin{aligned}
  U^{\dag} U & = 1 \\ 
  T^T T & =1 
\end{aligned} \quad \quad \quad \, \begin{aligned}
  U^{\dag}_{jk} U_{kl} & = \delta_{jl} & = U^*_{kj} U_{kl} \\ 
  T^{T}_{jk} T_{kl} & = \delta_{jl} & = T_{kj} T_{kl} \\ 
\end{aligned}
\]
I have to be careful about indices and the difference between column vectors and row vectors.  

For a column vector, $\psi_j = \psi_{j1}$.  Index 1 denoting only 1 column.  

Then for a row vector, such as $\psi^{\dag}$, 
\phantom{Then} $\psi_j^{\dag} = (\psi_{j1})^{\dag} = \psi_{1j}^*$.  If there's no confusion, denote $\psi_{1j}^*$ row vector as $\psi_j^{\dag}$.  

Now $\psi^{\dag} U$ should also be a row vector, of components
\[
(\psi^{\dag} U)_{1j} = \psi^{\dag}_{1k} U_{kj}
\]

So from the introduction to the problem, 
\[
\begin{aligned}
  & \psi_j \to \psi_j' = U_{jk} \psi_k \\ 
  & \psi_j^{\dag} \to \psi^{' \dag }_j = \psi^{\dag}_k U^{\dag}_{jk} \\
  & \psi_j^{\dag} \sigma^{\mu} \partial_{\mu} \psi_j \to \psi_k^{\dag} U_{kj}^{\dag} \sigma^{\mu} \partial_{\mu} U_{jl} \psi_l = \psi_k^{\dag} \sigma^{\mu}  U_{jk}^* U_{jl} \partial_{\mu} \psi_l = \psi_k^{\dag} \sigma^{\mu} \partial_{\mu} \psi_k 
\end{aligned}
\]

So for a theory of $N$ massless Weyl fields $\psi_j$, $\mathcal{L} = i \psi_j^{\dag} \sigma^{\mu} \partial_{\mu} \psi_j$, invariance group is $U(N)$.  

\begin{itemize}
\item[(a)] $N$ Weyl fields with a common mass $m$, $\mathcal{L} = i \psi_j^{\dag} \sigma^{\mu } \partial_{\mu} \psi_j - \frac{1}{2} m (\psi_j \psi_j + \psi_j^{\dag} \psi_j^{\dag})$.  \\
  \phantom{ Weyl fields with } (be careful that $\psi_j \psi_j$ implies $\psi_j$(row) $\psi_j$(column))
\[
\psi_j \psi_j \to \psi_k T^T_{jk} T_{jl} \psi_l = \psi_k T_{jk} T_{jl} \psi_l = \psi_k \psi_k
\]

Be careful in considering $\psi_j^{\dag} \psi_j^{\dag}$ term.  

For $\psi_j^{\dag}$, a column of $N$ Weyl fields, complex conjugated, $\psi_j^{\dag} = \left( \begin{matrix} \psi_1^* \\ \vdots \\ \psi_N^* \end{matrix} \right)$, 
\[
\psi_j^{\dag} \to \psi_j^{' \dag } = T_{jk} \psi_k^{\dag}
\]

For row vector $\psi_j^{\dag}$, $\psi_j^{\dag} \to \psi_j^{' \dag} = \psi_k^{\dag} T_{kj}^T$ (where, for instance if $T_{ij} x_j = x_{i1}$, $(x_{i1})^T = x_{1i} = (Tx)^T = x^T T^T = x_{1j} T_{jk}^T$).  

\[
\psi_j^{' \dag} \psi_j^{' \dag} = \psi_k^{\dag} T^T_{kj} T_{jl} \psi_l^{\dag} = \psi_k^{\dag} T_{jk} T_{jl} \psi_l^{\dag} = \psi_k^{\dag} \psi_k^{\dag} 
\]
So due to mass term, only $O(N)$ symmetry of Weyl fields.  
\item[(b)] $N$ massless Majorana fields, $\mathcal{L} = \frac{i}{2} \Psi_j^T \mathcal{C} \gamma^{\mu} \partial_{\mu} \Psi_j$.  

Now
\[
\begin{aligned}
  & \Psi_j^T = \left( \begin{aligned} \psi_{jc} \\ \psi_j^{\dag \dot{c}} \end{aligned} \right)^T  = \left( \begin{aligned} \epsilon_{ca} \psi_j^a \\ \epsilon^{\dot{c} \dot{a}} \psi_{j \dot{a}}^{\dag} \end{aligned} \right)^T = ( \epsilon_{ac} \psi_j^a, \epsilon^{\dot{a} \dot{c}} \psi_{j \dot{a}}^{\dag} ) \\ 
  & \Psi_j^T \mathcal{C} = ( \epsilon_{bc} \Psi_j^b, \epsilon^{\dot{b} \dot{c}} \psi_{j \dot{b}}^{\dag} ) \left( \begin{matrix}  - \epsilon^{ac} & \\ & - \epsilon_{\dot{a} \dot{c}} \end{matrix} \right)  = ( \epsilon_{bc} \epsilon^{ca} \psi_j^b, \epsilon^{\dot{b} \dot{c}} \epsilon_{\dot{c} \dot{a}} \psi_{j\dot{b}}^{\dag} ) = (\psi_j^a, \psi_{j\dot{a}}^{\dag} )
\end{aligned}
\]

\[
\Psi_j^T \mathcal{C} \gamma^{\mu} \partial_{\mu} \Psi_j = ( \psi_j^a, \psi_{j \dot{a}}^{\dag} ) \left( \begin{matrix} & \sigma^{\mu}_{ a \dot{c}} \\ \overline{\sigma}^{\mu \dot{a} c} & \end{matrix} \right) \partial_{\mu} \left( \begin{matrix} \partial_{\mu} \psi_{jc} \\ \partial_{\mu} \psi_j^{\dag \dot{c}} \end{matrix} \right) =   ( \psi_j^a, \psi_{j \dot{a}}^{\dag} ) \left( \begin{matrix} \sigma^{\mu}_{a \dot{c}} \partial_{\mu} \psi_j^{\dag \dot{c}} \\ \overline{\sigma}^{\mu \dot{a} c} \partial_{\mu} \psi_{jc} \end{matrix} \right) = + \psi_j \sigma^{\mu} \partial_{\mu} \psi_j^{\dag} + \psi_j^{\dag} \overline{\sigma}^{\mu} \partial_{\mu} \psi_j
\]

where $\begin{aligned} \epsilon^{12} & = \epsilon_{21} = + 1 \\  \epsilon^{21} & = \epsilon_{12} = - 1 \end{aligned}$ and likewise $\begin{aligned} \epsilon^{\dot{1}\dot{2}} & = \epsilon_{\dot{2}\dot{1}} = + 1 \\  \epsilon^{\dot{2}\dot{1}} & = \epsilon_{\dot{1}\dot{2}} = - 1 \end{aligned}$.  Also contract the second index to keep (+) sign.  

\[
\mathcal{L} = \frac{+i}{2} ( \psi_j \sigma^{\mu}  \partial_{\mu} \psi_j^{\dag} + \psi_j^{\dag} \overline{\sigma}^{\mu} \partial_{\mu} \psi_j )
\]

As given in the problem, for $\psi_j' \to U_{jk} \psi_k$, 

\[
\begin{aligned}
  & \psi_j(\text{row}) \to \psi_k U_{kj}^T \\ 
  & \psi_j^{\dag} \to U_{jk}^* \psi_k \\ 
  & \psi_j^{\dag}(\text{row}) \to \psi_k U_{kj}^{* T} 
\end{aligned}
\]

\[
\begin{aligned}
  & \psi_j \sigma^{\mu} \partial_{\mu} \psi_j^{\dag} \to \psi_k U_{kj}^T  \sigma^{\mu} \partial_{\mu} U_{jl}^* \psi_l^{\dag} = \psi_k U_{jk} U_{jl}^* \sigma^{\mu} \partial_{\mu} \psi_l^{\dag} = \psi_k \sigma^{\mu} \partial_{\mu} \psi_k^{\dag} \\ 
  & \psi^{\dag}_k U_{kj}^{* T} \overline{\sigma}^{\mu} \partial_{\mu} U_{jk} \psi_l = \psi^{\dag}_k U_{jk}^* U_{jl} \overline{\sigma}^{\mu} \partial_{\mu} \psi_l = \psi^{\dag}_k \overline{\sigma}^{\mu} \partial_{\mu} \psi_k 
\end{aligned}
\]

\[
\Longrightarrow \boxed{ U(N) \text{ symmetry}  }
\]


\item[(c)] $N$ Majorana fields with a common mass $m$, $\mathcal{L} = \frac{i}{2} \Psi_j^T \mathcal{C} \gamma6{\mu} \partial_{\mu} \Psi_j - \frac{1}{2} m \Psi_j^T \mathcal{C} \Psi_j$  (36.78).  

\[
\Psi_j^T \mathcal{C} \Psi_j = ( \psi_j^a , \psi_j^{\dag \dot{a}} ) \left( \begin{matrix} \psi_{jc} \\ \psi_j^{\dag \dot{c} } \end{matrix} \right) = ( \psi_j \psi_j + \psi_j^{\dag} \psi_j^{\dag} )
\]
as shown in (a), only $O(N)$ symmetry for $N$ Weyl fields with this form of the mass term.  
\item[(d)] $N$ massless Dirac fields, $\mathcal{L} = i \overline{\Psi}_j \gamma^{\mu} \partial_{\mu} \Psi_j$ (36.79).  

\[
\begin{aligned}
  \overline{\Psi}_j \gamma^{\mu} \partial_{\mu} \Psi_j & = (\xi_j^a, \chi_{j \dot{a}}^{\dag} ) \left( \begin{matrix} & \sigma_{a\dot{c}}^{\mu} \\ \overline{\sigma}^{\mu \dot{a} c} & \end{matrix} \right) \left( \begin{matrix} \partial_{\mu} \chi_{ja} \\ \partial_{\mu} \xi_j^{\dag \dot{a}} \end{matrix} \right) = ( \chi_{j\dot{a}}^{\dag} \overline{\sigma}^{\mu \dot{a} c} , \xi^a_j \sigma^{\mu}_{a\dot{c}} ) \left( \begin{matrix} \partial_{\mu} \chi_{ja} \\ \partial_{\mu} \xi_j^{\dag \dot{a}} \end{matrix} \right) = \chi_{j \dot{c}}^{\dag} \overline{\sigma}^{\mu \dot{c} a} \partial_{\mu} \chi_{ja} + \xi_j^c \sigma_{c\dot{a}}^{\mu} + \partial_{\mu} \xi^{\dag \dot{a}} = \\
  & = \chi_j^{\dag} \overline{\sigma}^{\mu} \partial_{\mu} \chi_j + \xi_j \sigma^{\mu} \partial_{\mu} \xi^{\dag}
\end{aligned}
\]

There are $N$ $\chi$ fields and $N$ $\xi$ fields.  

As in the development in Srednicki (36.26),
\[
( \xi_j \sigma^{\mu} \partial_{\mu} \xi_j^{\dag} ) = \partial_{\mu} ( \xi_j \sigma^{\mu} \xi_j^{\dag}) - (\partial_{\mu} \xi_j ) \sigma^{\mu} \xi_j^{\dag}
\]
So for
\[
(\partial_{\mu} \xi_j^c) \sigma_{c\dot{a}}^{\mu} \xi_j^{\dag \dot{a}} = - \xi_j^{\dag \dot{a} } \sigma^{\mu}_{c\dot{a}} \partial_{\mu} \xi_j^c = - \xi_j^{\dag \dot{a}} \epsilon_{ca} \epsilon_{\dot{a} \dot{b}} \overline{\sigma}^{\mu \dot{b} \dot{a} } \partial_{\mu} \xi_j^c = - ( - \epsilon_{\dot{b} \dot{a}} \xi_j^{\dag \dot{a}} ) \overline{\sigma}^{\mu \dot{b} a} ( - \epsilon_{ac} \partial_{\mu} \xi_j^c ) = - \xi_{j \dot{b}}^{\dag} \overline{\sigma}^{\mu \dot{b} a}  \partial_{\mu} \xi_{ja} = - \xi_j^{\dag} \overline{\sigma}^{\mu} \partial_{\mu} \xi_j 
\]
where in the first equality, the anticommutation of Grassmann \emph{numbers} was used.  
\[
\mathcal{L} = i \overline{\Psi}_j \gamma^{\mu} \partial_{\mu} \Psi_j = i ( \chi_j^{\dag} \overline{\sigma}^{\mu} \partial_{\mu} \chi_j + \xi_j^{\dag} \overline{\sigma}^{\mu} \partial_{\mu} \xi_j - \partial_{\mu} ( \xi_j \sigma^{\mu} \xi_j^{\dag} ) ) 
\]
The total divergence integrates out to results in zero.  

So if for $\psi_j$, a column vector of $2N$ fields (components), $\psi_j = \left( \begin{matrix} \psi_1 \\ \vdots \\ \psi_{2N} \end{matrix} \right)$ s.t. $\psi_j = \chi_j$ for $j=1, \dots , N$, $\psi_{j+N} = \xi_j$, $j=1, \dots, N$ then clearly, $\chi_j^{\dag} \overline{\sigma}^{\mu} \partial_{\mu} \chi_j + \xi_j^{\dag} \overline{\sigma}^{\mu} \partial_{\mu} \xi_j = \psi_j^{\dag} \overline{\sigma}^{\mu} \partial_{\mu} \psi_j$.  

Then as explicitly shown in the beginning of the problem, $\psi_j^{\dag} \overline{\sigma}^{\mu} \partial_{\mu} \psi_j$ invariant under $\psi_j \to U_{jk} \psi_k$.  
\[
\Longrightarrow \boxed{ U(2N) \text{ symmetry }}
\]
\item[(e)] $N$ Dirac fields with a common mass $m$, $\mathcal{L} = i \overline{\Psi}_j \gamma^{\mu} \partial_{\mu} \Psi_j - m \overline{\Psi}_j \Psi_j$ (36.80)

\[
\overline{\Psi}_j \Psi_j = ( \xi_j^a, \chi_{j \dot{a}}^{\dag} ) \left( \begin{matrix} \chi_{j\dot{a}} \\ \xi^{\dag \dot{a}} \end{matrix} \right) = \xi_j^a \chi_{ja} + \chi_{j \dot{a}}^{\dag} \xi_j^{\dag \dot{a}} = \xi_j \chi_j + \chi_j^{\dag} \xi_j^{\dag}
\]

For $\begin{aligned} &  \quad  \\ \chi_j & = \frac{1}{ \sqrt{2}} ( \psi_{2j-1 } + i \psi_{2j} ) \\ \xi_j & = \frac{1}{ \sqrt{2}} ( \psi_{2j-1} - i \psi_{2j} ) \end{aligned}$  

\[
\begin{aligned}
  & \chi_j^{\dag} \overline{\sigma}^{\mu} \partial_{\mu} \chi_j = \frac{1}{2} ( \chi_{2j-1}^{\dag} -i \psi_{2j}^{\dag}) \overline{\sigma}^{\mu} \partial_{\mu} ( \psi_{2j-1} + i \psi_{2j} ) = \frac{1}{2} ( \psi_{2j-1}^{\dag} \overline{\sigma}^{\mu} \partial_{\mu} \psi_{2j-1} + i \psi_{2j-1}^{\dag} \overline{\sigma}^{\mu} \partial_{\mu} \psi_{2j} -  i \psi_{2j}^{\dag} \overline{\sigma}^{\mu} \partial_{\mu} \psi_{2j-1} + \psi_{2j}^{\dag} \overline{\sigma}^{\mu} \partial_{\mu} \psi_{2j} ) \\
  & \xi_j^{\dag} \overline{\sigma}^{\mu} \partial_{\mu} \xi_j = \frac{1}{2} ( \chi_{2j-1}^{\dag} +i \psi_{2j}^{\dag}) \overline{\sigma}^{\mu} \partial_{\mu} ( \psi_{2j-1} - i \psi_{2j} ) = \frac{1}{2} ( \psi_{2j-1}^{\dag} \overline{\sigma}^{\mu} \partial_{\mu} \psi_{2j-1} - i \psi_{2j-1}^{\dag} \overline{\sigma}^{\mu} \partial_{\mu} \psi_{2j} +  i \psi_{2j}^{\dag} \overline{\sigma}^{\mu} \partial_{\mu} \psi_{2j-1} + \psi_{2j}^{\dag} \overline{\sigma}^{\mu} \partial_{\mu} \psi_{2j} )
\end{aligned}
\]
\[
\chi_j^{\dag} \overline{\sigma}^{\mu} \partial_{\mu} \chi_j + \xi_j^{\dag} \overline{\sigma}^{\mu} \partial_{\mu} \xi_j = \psi_{2j-1}^{\dag} \overline{\sigma}^{\mu} \partial_{\mu} \psi_{2j-1} + \psi_{2j}^{\dag} \overline{\sigma}^{\mu} \partial_{\mu} \psi_{2j}
\]

We treat $\psi_{2j-1}$, $\psi_{2j}$ as 2 left-handed Weyl fields.  Weyl fields describe spin-$\frac{1}{2}$ particles; by spin-statistics theorem, particles must be fermions $\Longrightarrow$ corresponding fields must anticomute.  

e.g. $\chi \psi = \chi^a \psi_a = - \psi_a \chi^a = - \epsilon^{ab} \psi_b \epsilon_{ac} \chi^c =  = \epsilon_{ca} \epsilon^{ab} \psi_b \chi^c = + \delta_c^b \psi_b \chi^c = \psi_b \chi^b $ \\
so $\psi_{2j-1} \psi_{2j} - \psi_{2j} \psi_{2j-1} = \psi_{2j-1}^{\dag} \psi_{2j}^{\dag} - \psi_{2j}^{\dag} \psi_{2j-1}^{\dag} = 0 $

\[
\xi_j \chi_j + \chi_j^{\dag} \xi_j^{\dag} = \frac{1}{2} ( \psi_{2j-1} \psi_{2j-1} +  \psi_{2j} \psi_{2j} + \psi_{2j-1}^{\dag} \psi_{2j-1}^{\dag} + \psi_{2j}^{\dag} \psi_{2j}^{\dag} )  
\]
\[
\Longrightarrow \mathcal{L} = i ( \psi_{2j-1}^{\dag} \overline{\sigma}^{\mu} \partial_{\mu} \psi_{2j-1} + \psi_{2j}^{\dag} \overline{\sigma}^{\mu} \partial_{\mu} \psi_{2j} ) - \frac{m}{2} ( \psi_{2j-1} \psi_{2j-1} + \psi_{2j} \psi_{2j} + \psi_{2j-1}^{\dag} \psi_{2j-1}^{\dag} + \psi_{2j}^{\dag} \psi_{2j}^{\dag} )  
\]

With $\psi_j$ a column vector of $2N$ left-handed Weyl fields $\psi_j = \left( \begin{matrix} \psi_1 \\ \vdots \\ \psi_{2N} \end{matrix} \right)$, then like (a), 
\[
\Longrightarrow \boxed{ O(2N) \text{ symmetry } }
\]
\end{itemize}




\section{ Canonical quantization of spinor fields I }

\solutionhead{37.1} Given
\[
\begin{aligned}
  \lbrace \psi_a(\mathbf{x},t) , \psi_c(\mathbf{y}, t) \rbrace & = 0  \quad \quad \quad \quad \, & (37.4) \\
  \lbrace \psi_a(\mathbf{x},t) , \Pi^c(\mathbf{y}, t) \rbrace & = i \delta_a^c \delta^3(\mathbf{x} - \mathbf{y} )  \quad \quad \quad \quad \, & (37.5)
\end{aligned}
\]
where $\Pi^a(x) \equiv \frac{ \partial \mathcal{L} }{ \partial ( \partial_0 \psi_a(x)) } = i \psi_{\dot{a}}^{\dag}(x) \overline{\sigma}^{ 0 \dot{a} a}$ \quad \, (37.2).  

$\chi, \xi$ each obey (37.4), (37.5).  

Now $\lbrace \Psi_{\alpha}(\mathbf{x},t), \Psi_{\beta}(\mathbf{y},t) \rbrace =0$, since for $\Psi \equiv \left( \begin{matrix} \chi_a \\ \xi^{\dag \dot{a}} \end{matrix} \right)$,  

\[
\begin{aligned}
  & \lbrace \chi_a(\mathbf{x},t) , \chi_b(\mathbf{y},t) \rbrace = 0 \\ 
  & \lbrace \chi_a(\mathbf{x},t) , \xi_b^{\dag}(\mathbf{y},t)  \rbrace =  0  \quad \quad \, (\chi_a, \xi_b \text{ are 2 different left-handed Weyl fields}) \\
  & \lbrace \xi_a^{\dag}(\mathbf{x},t) , \xi_b^{\dag}(\mathbf{y},t) \rbrace = 0 
\end{aligned}
\]

Recall $\lbrace \Psi_{\alpha}(\mathbf{X},t), \overline{\Psi_{\beta}}(\mathbf{y},t) \rbrace = ( \gamma^0)_{\alpha \beta} \delta^3(\mathbf{x} - \mathbf{y})$ (37.14).  Now $\overline{\Psi} = (\xi^c, \chi_{\dot{a}}^{\dag} )$, since recall $\gamma^0 = \left( \begin{matrix} & 1 \\ 1 & \end{matrix} \right) =  \left( \begin{matrix} & \delta_{a\dot{c}} \\ \delta^{\dot{a}c} & \end{matrix} \right)$.  

So from the following, then for the only nonzero results for the anticommutators, \\
$(\alpha, \beta) = (1,3), (3,1), (2,4), (4,2)$.  ($\alpha, \beta$ are 4-component spinor indices are ``top'' and ``bottom'' are distinguished by dotted or undotted.  \\


\noindent $\lbrace \chi_a(\mathbf{x},t), \xi^b(\mathbf{y},t) \rbrace = 0$.  So clearly $\lbrace \Psi_{\alpha}(\mathbf{x}) , \overline{\Psi}_{\beta}(\mathbf{y}) \rbrace = 0$ for $\alpha, \beta$ undotted.  \\
$\lbrace \xi^{\dag \dot{a}}(\mathbf{x},t), \xi_{\dot{a}}^{\dag}(\mathbf{y},t) \rbrace = 0$.  So clearly $\lbrace \Psi_{\alpha}(\mathbf{x}) , \overline{\Psi}_{\beta}(\mathbf{y}) \rbrace = 0$ for $\alpha, \beta$ dotted. \\
$\lbrace \chi_a(\mathbf{x},t), i \chi_{\dot{a}}^{\dag}(\mathbf{y},t) \overline{\sigma}^{ 0 \dot{a}c} \rbrace = \delta_a^c \delta^3(\mathbf{x}- \mathbf{y})$  

Multiply both sides by the number $\sigma^0_{c \dot{c}}$ to get back (37.7).  

\[
\lbrace \chi_a(\mathbf{x},t), \chi_{\dot{c}}^{\dag}(\mathbf{y},t) \rbrace = \sigma^0_{a\dot{c}} \delta^3(\mathbf{x} - \mathbf{y}) = \delta_{a\dot{c}} \delta^3(\mathbf{x}- \mathbf{y})
\]
Comparing with $(\gamma^0)_{\alpha \beta}$, clearly, I obtain the ``block'' of $( \gamma^0)_{\alpha \beta}$ containing $\delta_{a\dot{c}}$.  

\[
\lbrace \xi^{\dag \dot{c}}, \xi^a \rbrace = - \lbrace \xi^a(\mathbf{y},t), \xi^{\dag \dot{c}}(\mathbf{x},t) \rbrace = - \overline{\sigma}^{0 \dot{c} a} \delta^3(\mathbf{y}-\mathbf{x}) = - \overline{\sigma}^{ 0 \dot{c} a} \delta^3(\mathbf{y} - \mathbf{x})
\]
or
\[
\lbrace \xi^{\dag \dot{a}} , \xi^c \rbrace = - \overline{\sigma}^{ 0 \dot{a} c} \delta^3(\mathbf{x} - \mathbf{y}) = - \delta^{c\dot{a} } \delta^3(\mathbf{x} - \mathbf{y}) = \delta^{\dot{a} c} \delta^3(\mathbf{x} - \mathbf{y})
\]
where for the last equality, the apparent statement that $\delta^{c\dot{a}} = - \delta^{\dot{a} c}$ can be traced back to $\overline{\sigma}^{ 0 \dot{a} c}$.  

And so for the ``block'' containing $\delta^{\dot{a} c}$ (note the order of dotted $a$ first, then undotted $c$), then we reobtain the entries for $(\gamma^0)_{\alpha \beta}$.  
\[
\Longrightarrow \lbrace \Psi_{\alpha}(\mathbf{x},t) , \overline{\Psi}_{\beta}(\mathbf{y},t) \rbrace = (\gamma^0)_{\alpha \beta} \delta^3(\mathbf{x} - \mathbf{y})
\]



\section{Spinor technology}

\[
S_z = \frac{i}{4} \left[ \gamma^1, \gamma^2 \right] = \frac{1}{2} \gamma^1 \gamma^2 = \left( \begin{matrix} \frac{1}{2} \sigma_3 & \\ & \frac{1}{2} \sigma_3 \end{matrix} \right) \quad \quad \quad (38.2) 
\]
\[
\gamma^{\mu} = \left( \begin{matrix} & \sigma^{\mu}_{a \dot{c}} \\ \overline{\sigma}^{\dot{a} c}  & \end{matrix} \right)
\]
required 
\[
\begin{aligned}
  & S_z u_{\pm}(0) = \pm \frac{1}{2} u_{\pm}(0) \quad \quad \quad (38.3)  \\
  & S_z v_{\pm}(0) = \mp \frac{1}{2} v_{\pm}(0)
\end{aligned}
\]

\[
S^{\mu \nu} \equiv \left(  \begin{matrix} + (S_L^{\mu \nu})_a^{\, \, c} & \\ & (S_R^{\mu \nu})^{\dot{a}}_{\, \, \dot{c}} \end{matrix} \right) = \frac{i}{4} [ \gamma^{\mu}, \gamma^{\nu} ] 
\]

\[
s S_z u_{s'}(0) = ss' \frac{1}{2} u_{s'}(0)
\]

Then 
\[
\begin{aligned}
  & \frac{1}{2} ( 1 + 2 s S_z) u_{s'}(0) = \frac{1}{2} u_{s'}(0) + \frac{ss'}{2} u_{s'}(0) = \delta_{ss'} u_{s'}(0) \\ 
  &  \frac{1}{2} ( 1 - 2 s S_z) v_{s'}(0) = \frac{1}{2} v_{s'}(0) + \frac{ss'}{2} v_{s'}(0) = \delta_{ss'} v_{s'}(0) 
\end{aligned}
\]


\solutionhead{38.1}

Boost $K^i$ has a matrix representation
\[
K^i = \frac{i}{4} [ \gamma^i, \gamma^0] = \frac{i}{2} \gamma^i \gamma^0
\]
Now $p_i K^i = \frac{i}{2} p_i \gamma^i \gamma^0$, so then
\[
2i \widehat{p}_i K^i = \frac{ - p_i}{ |p_i |} \gamma^i \gamma^0 = \frac{- p_i }{ |p_i |} \left( \begin{matrix} \sigma^i_{a\dot{c}} & \\ & \overline{\sigma}^{i \dot{a}c} \end{matrix} \right) = - \left( \begin{matrix} \widehat{p}_3 & \widehat{p}_1 - i \widehat{p}_2 & & \\ \widehat{p}_1 + i \widehat{p}_2 & - \widehat{p}_3 & & \\ & & -\widehat{p}_3 & -\widehat{p}_1 + i \widehat{p}_2 \\ & & - \widehat{p}_1 - i \widehat{p}_2 & \widehat{p}_3 \end{matrix} \right)
\]

Solving the eigenvalue equation for one of the $2\times 2$ blocks, 
\[
\left| \begin{matrix} \lambda - \widehat{p}_3 & - (\widehat{p}_1 - i \widehat{p}_2 ) \\ -( \widehat{p}_1 + i \widehat{p}_2 ) & \lambda + \widehat{p}_3 \end{matrix} \right| = \lambda^2 - \widehat{p}_3^2 - (\widehat{p}_1^2 + \widehat{p}_2^2) = \lambda^2 - 1 = 0
\]
So then $\lambda = \pm 1$.  

Now for $A$ diagonalized to be $B$, 
\[
\begin{gathered}
  e^{cA} = e^{c U^{\dag} B U} = \sum_{j=0}^{\infty} \frac{ (c U^{\dag} B U )^j }{ j!} = U^{\dag} \sum_{j=0}^{\infty} \frac{ c^j B^j }{j!} U \text{ since } \\
  \underbrace{ (U^{\dag} B U ) \dots (U^{\dag} B U )}_{j \text{ times } } = U^{\dag} B U
\end{gathered}
\]
Now $\begin{aligned} B^{2n} & = \left( \begin{matrix} 1 & \\ & 1 \end{matrix} \right) = 1 \\ B^{2n+1} & = \left( \begin{matrix} 1 & \\ & -1 \end{matrix} \right) = B \end{aligned}$.  So
\[
e^{cA} = U^{\dag} \left( \sum_{n=0}^{\infty} \frac{ c^{2n} }{ (2n)!} + \sum_{n=0}^{\infty} \frac{ c^{2n+1} \left( \begin{matrix} 1 & \\ & -1 \end{matrix} \right) }{ (2n+1)!} \right) U = \cosh{c} + \sinh{c} A
\]
So 
\[
\exp{ (i \eta \widehat{p} \cdot K ) } = \exp{ \left( \frac{ \eta}{2} (2i \widehat{p} \cdot K ) \right) } = \cosh{ \left( \frac{ \eta }{2} \right) } + \sinh{ \left( \frac{ \eta}{2} \right) } (2i \widehat{p}\cdot K )
\]

Recalling that 
\[
\begin{aligned}
  u_+(0) = \sqrt{ m } \left( \begin{matrix} 1 \\ 0 \\ 1 \\ 0 \end{matrix} \right) \\ 
  u_-(0) = \sqrt{ m } \left( \begin{matrix} 0 \\ 1 \\ 0 \\ 1 \end{matrix} \right) \\ 
  v_+(0) = \sqrt{ m } \left( \begin{matrix} 0 \\ 1 \\ 0 \\ -1 \end{matrix} \right) \\ 
  v_-(0) = \sqrt{ m } \left( \begin{matrix} -1 \\ 0 \\ 1 \\ 0 \end{matrix} \right) 
\end{aligned} \text{ then }
\]


\problemhead{38.2}  Verify eq. (38.15).  

\solutionhead{38.2} 

Recall that $\overline{A} = \beta A^{\dag} \beta$ where $\beta = \left( \begin{matrix} & I \\ I & \end{matrix} \right)$, and 
\[
\gamma^{\mu } = \left( \begin{matrix} 0 & \sigma^{\mu}_{a\dot{c}} \\ \overline{\sigma}^{\mu \dot{a} c} & 0 \end{matrix} \right) \quad \quad \quad \, 
(\gamma^{\mu} )^{\dag} = \left( \begin{matrix} & \overline{\sigma}^{\mu * a\dot{c}}  \\ \sigma^{\mu *}_{\dot{a}c} & \end{matrix} \right) = \left( \begin{matrix} & \overline{\sigma}^{\mu \dot{c} a} \\ \sigma^{\mu}_{c\dot{a}} & \end{matrix} \right)
\]
\[
\begin{gathered}
  (\gamma^{\mu})^{\dag} \beta = \left( \begin{matrix} \overline{\sigma}^{\mu \dot{c} a } & \\ \sigma^{\mu}_{c\dot{a}} \end{matrix} \right)  \quad \quad \Longrightarrow
  \overline{\gamma^{\mu}} =  \beta (\gamma^{\mu})^{\dag} \beta = \left( \begin{matrix} & \sigma^{\mu}_{c\dot{a}} \\ \overline{\sigma}^{\mu \dot{c} a} & \end{matrix} \right) =  \left( \begin{matrix} & \sigma^{\mu}_{a\dot{c}} \\ \overline{\sigma}^{\mu \dot{a} c} & \end{matrix} \right) = \gamma^{\mu}
\end{gathered}
\]

Recall
\[
S^{\mu \nu} = \left( \begin{matrix} (S_L^{\mu \nu})_a^c & 0 \\ 0 & - (S^{\mu \nu}_R )^{\dot{a}}_{\dot{c}} \end{matrix} \right) \text{ where } 
\begin{aligned} (S_L^{\mu \nu})_a^c & = \frac{ + i }{4} ( \sigma^{\mu } \overline{\sigma}^{\nu} - \sigma^{\nu} \overline{\sigma}^{\mu} )_a^c \\  (S_R^{\mu \nu})^{\dot{a}}_{\dot{c}} & = \frac{ - i }{4} ( \overline{\sigma}^{\mu } \sigma^{\nu} - \overline{\sigma}^{\nu} \sigma^{\mu} )^{\dot{a}}_{\dot{c}} \end{aligned} 
\]

Now 
\[
\overline{\sigma}^{\mu}_{\dot{a} l} \sigma^{\nu l \dot{c}} = \epsilon_{\dot{a}}\epsilon_{lm} \overline{\sigma}^{\mu \dot{b} m } \epsilon^{ln} \epsilon^{\dot{c} \dot{a}} \sigma^{\nu}_{n\dot{d}} = \epsilon_{\dot{a} \dot{b} } \epsilon^{\dot{c} \dot{d}} \overline{\sigma}^{\mu \dot{b} l } \sigma^{\nu}_{l\dot{d}}
\]
Note that $\epsilon_{\dot{a} \dot{b}} \epsilon^{\dot{c} \dot{d}} = -\delta_{\dot{a}}^{\dot{c}} \delta_{\dot{b}}^{\dot{d}} -\delta_{\dot{a}}^{\dot{d}} \delta_{\dot{b}}^{\dot{c}}$, and so 
\[
\overline{\sigma}^{\mu}_{\dot{a} l} \sigma^{\nu l \dot{c}} = ( -\delta_{\dot{a}}^{\dot{c}} \delta_{\dot{b}}^{\dot{d}} -\delta_{\dot{a}}^{\dot{d}} \delta_{\dot{b}}^{\dot{c}}) \overline{\sigma}^{\mu \dot{b} l }  \sigma^{\nu}_{l \dot{d}} = (- \overline{\sigma}^{\mu} \sigma^{\nu} )^{\dot{c}}_{\dot{a}}
\]
where the first term in the second equality is zero since the trace of the Pauli matrices are zero.  

Likewise, $(\sigma^{\mu} \overline{\sigma}^{\nu})^a_c = (\sigma^{\mu} \overline{\sigma}^{\nu})^c_a$.  

\[
\begin{aligned}
  (S_L^{\mu \nu})_a^{c\dag} & = - \frac{i}{4} ( \overline{\sigma}^{\nu \dag} \sigma^{\mu \dag} - \overline{ \sigma}^{\mu \dag} \sigma^{\nu \dag} )_{\dot{a}}^{\dot{c}} = \frac{i}{4} ( \overline{\sigma}^{\mu} \sigma^{\nu} - \overline{\sigma}^{\nu} \sigma^{\mu} )_{\dot{a}}^{\dot{c}} = - (S_R^{\mu \nu})^{\dot{a}}_{\dot{c}} \\ 
  (S_R^{\mu \nu})^{\dot{a}\dag}_{\dot{c}} & = \frac{i}{4} ( \sigma^{\nu \dag} \overline{\sigma}^{\mu \dag} - \sigma^{\mu \dag} \overline{\sigma}^{\nu})^a_c = \frac{i}{4} ( \sigma^{\nu} \overline{\sigma}^{\mu} - \sigma^{\mu} \overline{\sigma}^{\nu} )^a_c = - \frac{i}{4} ( \sigma^{\mu } \overline{\sigma}^{\nu} - \sigma^{\nu } \overline{\sigma}^{\mu} )_a^c = - (S_L^{\mu \nu})_a^c
\end{aligned}
\]
So
\[
\begin{aligned}
  (S^{\mu \nu})^{\dag} \beta = \left( \begin{matrix} 0 & - (S_R^{\mu \nu})^{\dot{a}}_{\dot{c} } \\ +(S_L)^{\mu \nu c}_a & \end{matrix} \right) \\ 
  \beta (S^{\mu \nu})^{\dag} \beta = \left( \begin{matrix} +(S_L)^{\mu \nu c}_a &  \\  &  - (S_R^{\mu \nu})^{\dot{a}}_{\dot{c} } \\ \end{matrix} \right)
\end{aligned}
\]

Recall $\gamma_5 \equiv \left( \begin{matrix} - \delta_a^c & \\ & \delta^{\dot{a}}_{\dot{c}} \end{matrix} \right)$ \\ 
  \[
\overline{i\gamma_5} = \beta (i \gamma_5)^{\dag} \beta = \left( \begin{matrix} & 1 \\ 1  & \end{matrix} \right) (-i) \left( \begin{matrix} -\delta_{\dot{a}}^{\dot{c}} & \\ & \delta^a_{\dot{c}} \end{matrix} \right) \left( \begin{matrix} & 1 \\ 1 & \end{matrix} \right)= (-i)\left( \begin{matrix} \delta^a_c & \\ & - \delta_{\dot{a}}^{\dot{c}} \end{matrix} \right) = i \left( \begin{matrix} - \delta_a^c & \\ & \delta^{\dot{a}}_{\dot{c}} \end{matrix} \right)
\]


\[
\begin{gathered}
  \begin{aligned}
    \gamma^{\mu} \gamma_5 = \left( \begin{matrix} & \sigma^{\mu}_{a\dot{c}} \\ \overline{ \sigma}^{\mu \dot{a} c} & \end{matrix} \right) \left( \begin{matrix} - \delta_a^c & \\  & \delta^{\dot{a}}_{\dot{c}} \end{matrix} \right) = \left( \begin{matrix} & \sigma^{\mu}_{a\dot{a}} \\ - \overline{\sigma}^{\mu \dot{a} a} & \end{matrix} \right) \\ 
 (   \gamma^{\mu} \gamma_5 )^{\dag} = \left( \begin{matrix} & - \overline{\sigma}^{\mu \dag a\dot{a}} \\ \sigma^{\mu \dag}_{\dot{a} a} & \end{matrix} \right)  = \left( \begin{matrix} & - \overline{\sigma}^{\mu * a\dot{a}} \\ \sigma^{\mu *}_{\dot{a} a} & \end{matrix} \right) = \left( \begin{matrix} & - \overline{\sigma}^{\mu \dot{a}a} \\ \sigma^{\mu }_{a\dot{a} } & \end{matrix} \right) 
\end{aligned}
 \quad \quad \, \\ \begin{aligned}
   \beta( \gamma^{\mu} \gamma_5 )^{\dag} = \left( \begin{matrix} & 1 \\ 1 & \end{matrix} \right) \left( \begin{matrix} & - \overline{\sigma}^{\mu \dot{a} a} \\ \sigma^{\mu}_{ a \dot{a}} & \end{matrix} \right) = \left( \begin{matrix} \sigma^{\mu}_{a\dot{a}} & \\ & - \overline{\sigma}^{\mu \dot{a} a} \end{matrix} \right) \\ 
   \beta( \gamma^{\mu} \gamma_5 )^{\dag}\beta =  \left( \begin{matrix} \sigma^{\mu}_{a\dot{a}} & \\ & - \overline{\sigma}^{\mu \dot{a} a}  \end{matrix} \right) \left( \begin{matrix} & 1 \\ 1 & \end{matrix} \right)= \left( \begin{matrix} & \sigma^{\mu}_{a\dot{a}} \\  - \overline{\sigma}^{\mu \dot{a} a} &  \end{matrix} \right) 
\end{aligned}
\end{gathered}
\]

\[
\begin{aligned}
  i \gamma_5 S^{\mu \nu} = i \left( \begin{matrix} - (S_L^{\mu \nu})_a^c & \\ & - (S_R^{\mu \nu})^{\dot{a}}_{\dot{c}} \end{matrix} \right)  = -i \left( \begin{matrix} ( S_L^{\mu \nu})_a^c & \\ & ( S_R^{\mu \nu})^{\dot{a}}_{\dot{c}} \end{matrix} \right) \\ 
  (i \gamma_5 S^{\mu \nu})^{\dag} = i \left( \begin{matrix}  (S_L^{\mu \nu})_{\dot{a}}^{\dag \dot{c}} & \\ & - (S_R^{\mu \nu})^{\dag a}_{c} \end{matrix} \right)  = i \left( \begin{matrix} ( -S_R^{\mu \nu})^{\dot{a}}_{\dot{c}} & \\ & ( -S_L^{\mu \nu})_{a}^{c} \end{matrix} \right)  = -i \left( \begin{matrix} (S_R^{\mu \nu})^{\dot{a}}_{\dot{c}} & \\ & (S_L^{\mu \nu} )_a^c \end{matrix} \right)\\ 
(i \gamma_5 S^{\mu \nu})^{\dag} \left( \begin{matrix} & 1 \\ 1 & \end{matrix} \right) = -i \left( \begin{matrix} & (S_R^{\mu \nu }) \\ S_L^{\mu \nu} & \end{matrix} \right) \\ 
  \beta ( i \gamma_5 S^{\mu \nu})^{\dag} \beta = -i \left( \begin{matrix} S_L^{\mu \nu} & \\ & S_R^{\mu \nu} \end{matrix} \right) = \overline{(i \gamma_5 S^{\mu \nu} )} = i \gamma_5 S^{\mu \nu }
\end{aligned}
\]

\problemhead{38.3}  Verify eq. (38.22).  

\solutionhead{38.3} 
\[
\begin{gathered}
  \gamma^{\mu}\not{p} = \frac{1}{2} \lbrace \gamma^{\mu}, \not{p} \rbrace + \frac{1}{2} \left[ \gamma^{\mu}, \not{p} \right] = - p^{\mu} - 2i S^{\mu \nu} p_{\nu} \text{ since } \\
  \lbrace \gamma^{\mu}, p_{\nu} \gamma^{\nu} \rbrace = p_{\nu} \lbrace \gamma^{\mu}, \gamma^{\nu} \rbrace =  -2 g^{\mu \nu} p_{\nu} = - 2p^{\mu} \\
\frac{i}{4} [ \gamma^{\mu}, \gamma^{\nu} ] = S^{\mu \nu} \quad \, \text{(by definition)}
\end{gathered}
\]
\[
\begin{gathered}
  \not{p}'\gamma^{\mu} = \frac{1}{2} \lbrace \gamma^{\mu}, \not{p}' \rbrace - \frac{1}{2} [ \gamma^{\mu}, \not{p}' ] = - p^{'\mu} + 2i S^{\mu \nu} p_{\nu}' \\
  \lbrace \gamma^{\mu}, p_{\nu}'\gamma^{\nu} \rbrace = p_{\nu}' \lbrace \gamma^{\mu}, \gamma^{\nu} \rbrace = -2p_{\nu}' g^{\mu \nu} = -2 p^{'\mu} \\
[\gamma^{\mu}, \not{p}' ] = [\gamma^{\mu}, p_{\nu}' \gamma^{\nu} ] = p_{\nu}' (-4i S^{\mu \nu}) = -4i S^{\mu \nu} p_{\nu}'
\end{gathered}
\]

Add the 2 equations above together.  
\[
\gamma^{\mu} \not{p} + \not{p}' \gamma^{\mu} = - (p + p')^{\mu} - 2i (S^{\mu \nu} p_{\nu} - S^{\mu \nu} p_{\nu}' ) 
\]
Sandwich the above equation with $\overline{u}$,$u$ and $\overline{v}$, $v$.

\[
\overline{u}_s'( \mathbf{p})(\gamma^{\mu} \not{p} + \not{p}' \gamma^{\mu})u_s(\mathbf{p}) = \overline{u}_{s'}'(\mathbf{p}') ( - (p + p')^{\mu} - 2i S^{\mu \nu} (p_{\nu} - p_{\nu}') )u_s(\mathbf{p})
\]

Now use 
\[
\begin{aligned}
  (\not{p} + m)u_s(\mathbf{p}) & = 0 \\ 
  (-\not{p} + m)v_s(\mathbf{p}) & = 0 \\ 
\end{aligned}
\quad \quad \, \begin{aligned}
  \overline{u}_s(\mathbf{p})(\not{p} + m) & = 0 \\ 
  \overline{v}_s(-\not{p} + m) & = 0 \\ 
\end{aligned} 
\]

\[
\begin{gathered}
  \overline{u}_s'(\mathbf{p}')(\gamma^{\mu}\not{p} + \not{p}' \gamma^{\mu})u_s(\mathbf{p}) = \overline{u}_s'(\mathbf{p}') \gamma^{\mu}(-m u_s(\mathbf{p})) + - m \overline{u}_s'(\mathbf{p}') \gamma^{\mu} u_s(\mathbf{p}) = -2m \overline{u}_s'(\mathbf{p}') \gamma^{\mu} u_s(\mathbf{p}) \\ 
  \Longrightarrow \boxed{ 2m \overline{u}_s'(\mathbf{p}') \gamma^{\mu} u_s(\mathbf{p}) = \overline{u}_{s'}'(\mathbf{p}') ((p + p')^{\mu} - 2i S^{\mu \nu}( p' -p)_{\nu}) u_s(\mathbf{p}) }
\end{gathered}
\]

Likewise,
\[
\begin{gathered}
  \overline{v}_{s'}'(\mathbf{p}')(\gamma^{\mu} \not{p} + \not{p}' \gamma^{\mu}) v_s(\mathbf{p}) = \overline{v}_{s'}'(\mathbf{p}') \gamma^{\mu} m v_s(\mathbf{p}) + m \overline{v}_{s'}(\mathbf{p}') \gamma^{\mu} v_s(\mathbf{p}) = \\
  = \boxed{ 2m \overline{v}'_{s'}(\mathbf{p}') \gamma^{\mu} v_s(\mathbf{p}) = - \overline{v}_{s'}'(\mathbf{p}') ((p+p')^{\mu} - 2i S^{\mu \nu} (p'-p)_{\nu} )v_s(\mathbf{p}) }
\end{gathered}
\]

\problemhead{38.4}  Derive the Gordon identities
\[
\begin{aligned}
  \overline{u}_{s'}(\mathbf{p}') [ (p' + p)^{\mu} - 2i S^{\mu \nu} (p' - p)_{\nu}] \gamma_5 u_s(\mathbf{p}) & = 0 \\ 
  \overline{v}_{s'}(\mathbf{p}') [ (p' + p)^{\mu} - 2i S^{\mu \nu} (p' - p)_{\nu}] \gamma_5 v_s(\mathbf{p}) & = 0. \quad \quad \, (38.41) \\ 
\end{aligned}
\]

\solutionhead{38.4}  Recall that 
\[
(\gamma^{\mu} \not{p} + \not{p}' \gamma^{\mu}) = -(p+p')^{\mu} + 2i S^{\mu \nu}(p' - p)_{\nu} 
\]
Consider $\overline{u}_{s'} (\mathbf{p}') (\gamma^{\mu} \not{p} + \not{p}' \gamma^{\mu}) \gamma_5 u_s(\mathbf{p})$.  Now
\[
\overline{u}_{s'}(\mathbf{p}') \not{p}' \gamma^{\mu} \gamma_5 u_s(\mathbf{p}) = - m \overline{u}_{s'}(\mathbf{p}') \gamma^{\mu} \gamma_5 u_s(\mathbf{p})
\]
from using the Dirac equation (equation of motion) as in the previous problem.  

\textbf{An important property} is $\boxed{ \lbrace \gamma_5, \gamma^{\mu} \rbrace =0 }$ (Zee, 2003).  

\[
\begin{aligned}
  \overline{u}_{s'}(\mathbf{p}')(\gamma^{\mu} \not{p} \gamma_5 u_s(\mathbf{p}) & = \overline{u}_{s'}(\mathbf{p}')(\gamma^{\mu} p_{\nu}\gamma^{\nu} \gamma_5 u_s(\mathbf{p}) ) =  -\overline{u}_{s'}(\mathbf{p}') \gamma^{\mu} \gamma_5(p_{\nu} \gamma^{\nu} u_s(\mathbf{p}) ) = - \overline{u}_{s'}(\mathbf{p}')\gamma^{\mu} \gamma_5 \not{p}u_s(\mathbf{p}) = \\ 
  & = + m \overline{u}_{s'}(\mathbf{p}')\gamma^{\mu} \gamma_5 u_s(\mathbf{p})
\end{aligned}
\]
Thus,
\[
\boxed{ \overline{u}_{s'}(\mathbf{p}') \left[ (p' + p)^{\mu} - 2i S^{\mu \nu}(p' - p)_{\nu} \right] \gamma_5 u_s(\mathbf{p}) = 0 }
\]
Likewise,
\[
\begin{aligned}
  \overline{v}_{s'}(\mathbf{p}') (\gamma^{\mu} \not{p} + \not{p}'  \gamma^{\mu})\gamma_5 v_s(\mathbf{p}) & = \overline{v}_{s'}(\mathbf{p}')\gamma^{\mu} \not{p} \gamma_5 v_s(\mathbf{p}) + m \overline{v}_{s'}(\mathbf{p}') \gamma^{\mu} \gamma_5 v_s(\mathbf{p}) = \\
  & = - \overline{v}_{s'}(\mathbf{p}') \gamma^{\mu} \gamma_5 (m v_s(\mathbf{p})) + m \overline{v}_{s'}(\mathbf{p}') \gamma^{\mu} \gamma_5 v_s(\mathbf{p})  =0 
\end{aligned}
\]
So 
\[
\boxed{ \overline{v}_{s'}(\mathbf{p}') \left[ (p' + p)^{\mu} - 2i S^{\mu \nu}(p' - p)_{\nu} \right] \gamma_5 v_s(\mathbf{p}) = 0 }
\]

\section{Canonical quantization of spinor fields II}

\problemhead{39.1} Verify eq. (39.24).  

\solutionhead{39.1} 
Recall 
\[
\begin{aligned}
  & \Psi(x) = \sum_{s=\pm} \int \widetilde{dp} [b_s(\mathbf{p}) u_s(\mathbf{p}) e^{ipx} + d_s^{\dag}(\mathbf{p}) v_s(\mathbf{p}) e^{-ipx} ] \\
     & \overline{\Psi}(x) = \sum_{r=\pm} \int \widetilde{dq} [b_r^{\dag}(\mathbf{q}) \overline{u}_r(\mathbf{q}) e^{-iqx} + d_r^{\dag}(\mathbf{q}) \overline{v}_r(\mathbf{q}) e^{-iqx} ] 
\end{aligned}
\]
So
\[
\begin{aligned}
  \overline{\Psi}(x) \gamma^0 \Psi(x) & = \sum_{r,s = \pm} \int \widetilde{dq} \int \widetilde{dp} (b_r^{\dag} \overline{u}_r e^{-iqx} + d_r \overline{v}_r e^{iqx})\gamma^0 (b_s u_s e^{ipx} + d_s^{\dag}v_s e^{-ipx} ) = \\ 
  & = \sum_{r,s = \pm} \int \widetilde{dq} \widetilde{dp} \lbrace  b_r^{\dag} \overline{u}_r \gamma^0 b_s u_s e^{-i(q-p)x} + b_r^{\dag} \overline{u}_r \gamma^0 d_s^{\dag} v_s e^{-i(q+p)x} + \\ & \phantom{ = \sum \int dp dq b_r u_r} + d_r \overline{v}_r \gamma^0 b_s u_s e^{i(q+p) x} + d_r \overline{v}_r \gamma^0 d_s^{\dag} v_s e^{i(q-p)x}  \rbrace 
\end{aligned}
\]

Note the $x^0 = t$ part of the exponentials.  Without loss of generality, since $H$ doesn't depend on time, choose a single time $t$.  
\[
\begin{aligned}
  & i(p-q)x \to i(- \omega t + \omega t) = 0 \\ 
    & i(p-q)x \to i(- \omega t - \omega t) = -2i\omega t \text{ and likewise}  
\end{aligned}
\]

Integrate to obtain and use $\delta$ functions.  

\[
\xrightarrow{\int \widetilde{dq} \int d^3x } \begin{gathered} \sum_{r,s = \pm} \int \frac{\widetilde{dp} }{2\omega} \lbrace b_r^{\dag}(\mathbf{p} )b_s(\mathbf{p}) \overline{u}_r(\mathbf{p}) \gamma^0 u_s(\mathbf{p}) + b_r^{\dag}(-\mathbf{p})d_s^{\dag}(\mathbf{p}) \overline{u}_r(-\mathbf{p}) \gamma^0 v_s(\mathbf{p}) e^{2i \omega t} + \\
  d_r(-\mathbf{p})b_s(\mathbf{p}) \overline{v}_r \gamma^0 u_se^{-2i\omega t} + d_r(\mathbf{p}) d_s^{\dag}(\mathbf{p}) \overline{v}_r \gamma^0 v_s \rbrace \end{gathered}
\]

Recall the Gordon identities:
\[
\begin{aligned} \overline{u}_s'(\mathbf{p}) \gamma^{\mu} u_s(\mathbf{p}) & = 2 p^{\mu} \delta_{s' s} \\ 
 \overline{v}_s'(\mathbf{p}) \gamma^{\mu} v_s(\mathbf{p}) & = 2 p^{\mu} \delta_{s' s} \quad \, (38.21) \end{aligned} \quad \quad \, 
\begin{aligned}
  \overline{u}_{s'}(\mathbf{p})\gamma^0 v_s(-\mathbf{p}) & = 0 \\ 
  \overline{v}_{s'}(\mathbf{p})\gamma^0 u_s(-\mathbf{p}) & = 0 \quad \, (38.22) \\ 
\end{aligned}
\]

\[
\Longrightarrow H = \sum_{s = \pm} \int \widetilde{dp} \lbrace b_s^{\dag}(\mathbf{p}) b_s(\mathbf{p}) + d_s(\mathbf{p})d^{\dag}_s(\mathbf{p}) \rbrace
\]

\[
\Longrightarrow H = \sum_{s=\pm} \int \widetilde{dp} \lbrace b_s^{\dag}b_s - d_s^{\dag} d_s + (2\pi)^3 2 \omega \rbrace = \sum_{s=\pm} \int \widetilde{dp} \lbrace b_s^{\dag}b_s - d_s^{\dag} d_s \rbrace + 2 V_0
\]
where $V_0 = \int d^3p$ and 2 is for 2 spin factors.  

\problemhead{39.2} Use $[ \Psi(x), M^{\mu \nu} ] = -i (x^{\mu} \partial^{\nu} - x^{\nu} \partial^{\mu} )\Psi(x) + S^{\mu \nu} \Psi(x)$, plus whatever spinor identities you need, to show that 
\[
\begin{aligned}
  J_z b_s^{\dag}(p\hat{\mathbf{z}} )|0\rangle & = \frac{1}{2} s b_s^{\dag}(p \hat{\mathbf{z}})|0 \rangle \\ 
  J_z d_s^{\dag}(p\hat{\mathbf{z}} )|0\rangle & = \frac{1}{2} s d_s^{\dag}(p \hat{\mathbf{z}})|0 \rangle \\ 
\end{aligned}
\]
where $\mathbf{p} = p \hat{\mathbf{z}}$ is the three-momentum, and $\hat{\mathbf{z}}$ is a unit vector in the $z$ direction.  

\solutionhead{39.2}  Recall that 
\[
\begin{aligned}
  J_i & = \frac{1}{2} \epsilon_{ijk} M^{jk} \\ 
  J_3 & = \frac{1}{2} \epsilon_{3jk} M^{jk}
\end{aligned}
\]
$M^{jk}$'s are antisymmetric, the generators of the Lorentz group; $M^{\mu \nu} = - M^{\nu \mu}$.  \\
$J_3 = M^{12}$.  

Also, recall,

\[
\begin{aligned}
  b_s^{\dag}(\mathbf{p}) & = \int d^3x e^{ipx} \overline{\Psi}(x) \gamma^0 u_s(\mathbf{p}) \\ 
  d_s^{\dag}(\mathbf{p}) & = \int d^3x e^{ipx} \overline{v}_s(\mathbf{p}) \gamma^0 \Psi(x)
\end{aligned} \quad \quad \, \text{both obtained by Fourier transforming the mode expansion of $\Psi$}
\]

\emph{The trick} is to use this and the mode expansion to make delta functions (like using the completeness relation).  

Recall the hint 
\[
[ \Psi(x), M^{\mu \nu} ] = -i (x^{\mu} \partial^{\nu} - x^{\nu} \partial^{\mu})\Psi(x) + S^{\mu \nu} \Psi(x)
\]

Now $J_3$ is an operator on states.  
\[
\begin{gathered}
  J_z d_s^{\dag}(p\hat{z})|0\rangle = \int d^3x e^{ipx} \overline{v}_s(p\hat{z}) \gamma^0 J_3 \Psi(x)| 0\rangle = \\
  = \int d^3x e^{ipx} \overline{v}_s(p\hat{z})\gamma^0 \lbrace \Psi(x) M^{12} + i (x^1 \partial^2 - x^2 \partial^1) \Psi(x) - S^{12} \Psi(x) \rbrace |0 \rangle = * \\
\end{gathered}
\]

We were able to move $M^{12}$ because $M^{12}$ is an operator on $|0\rangle$.  $\gamma^0$ is a matrix.  $u_s(p\hat{z})$ is a ``number'' of a 4-component spinor.  They each are on different spaces, so they ``commute.'' \\

$J_3 |0 \rangle =0$ (vacuum state is spinless; it's zero).  

Since $M_{12}|0\rangle = J_3 |0 \rangle =0$, the first term in $*$ is 0.  

\[
\Longrightarrow * = \int d^3 x e^{ipx} \overline{v}_s(p\hat{z}) \gamma^0 \left( i (x^1 \partial^2 - x^2 \partial^1 ) - S^{12} \right) \sum_{s'=\pm} \int \widetilde{dp}'[ b_{s'}(p') u_{s'}(p')e^{ip'x} + d_{s'}^{\dag}(p') v_{s'}(p') e^{-ip'x} ]|0 \rangle
\]
Now
\[
\begin{gathered}
  b_{s'}(p')u_{s'}(p') e^{ip'x'}|0\rangle = u_{s'}(p')e^{ip'x'} b_{s'}(p')|0\rangle = 0  \\ 
\end{gathered}
\]
How does $S_z$ act on a 4-component spinor in a non-rest frame?  Let's be careful.  
\[
S_z v_{s'}(\mathbf{p'}) = \frac{1}{2m} \gamma_5\not{z}\not{p} v_{s'}(\mathbf{p}') = \frac{1}{2} \gamma_5 \not{z} v_{s'}(\mathbf{p'})
\]
where $(\not{p} - m) v_s(\mathbf{p}) = 0$ was used, and 
\[
\text{ rest frame $S_z$ } = \frac{-1}{2} \gamma_5 \gamma^3 \gamma^0 \xrightarrow{\text{boosted}} S_z = \frac{1}{2m} \gamma_5 \not{z}\not{p}
\]
Now $z = (0, \hat{z})$ and $z$ is all in the 3 direction.  
\[
\not{z} = z_{\mu}\gamma^{\mu} = \gamma^3
\]

So with this explicit formula for $S_z v_{s'}(\mathbf{p}')$ in boosted frame, we can conclude $S_z v_{s'}(\mathbf{p'}) = \frac{-1}{2} s' v_{s'}(\mathbf{p'})$, either using the explicit formula
\[
\begin{gathered}
  v_s(p) = \exp{ (i \eta \hat{p}\cdot \mathbf{K})} v_s(0) \text{ and } \\
  v_+(0) = \sqrt{ m} \left( \begin{matrix} 0 \\ 1 \\ 0 \\ -1 \end{matrix} \right), \quad \quad \quad v_-(0) = \sqrt{m} = \left( \begin{matrix} -1 \\ 0 \\ 1 \\ 0 \end{matrix} \right) 
\end{gathered}
\]
or use 
\[
\frac{1}{2} (1 - s \gamma_5 \not{z}) v_{s'}(\mathbf{p}) = \delta_{ss'} v_{s'}(\mathbf{p} )
\]
which is mathematically equivalent to $S_z v_{s'}(\mathbf{p}') = \frac{-1}{2} s' v_{s'}(\mathbf{p'})$.  

Likewise for $S_z u_{s'}(\mathbf{p}')  = \frac{-s'}{2} u_{s'}(\mathbf{p}')$.  Why $(-1)$ factor for $u_{s'}(\mathbf{p})$?  It comes from boosting which gives (38.27) (Srednicki, 2007):
\[
\frac{1}{2}(1 - s\gamma_5 z)u_{s'}(\mathbf{p}) = \delta_{ss'} u_{s'}(\mathbf{p})
\]



\[
\begin{gathered}
  \int d^3x e^{ipx} e^{-ip'x} = \int d^3 x e^{i(p-p')x} = (2\pi)^3 \delta^3(p-p')
  \xrightarrow{ \int \widetilde{dp}' } \frac{1}{2\omega} \sum_{s'= \pm} \overline{v}_s(p\hat{z}) \gamma^0 v_{s'}(p\hat{z}) \frac{1}{2} s' d^{\dag}_{s'} |0\rangle
\end{gathered}
\]
where $S_z$ acted on $v_{s'}(p')$, and then $\int \widetilde{dp}' \delta^3(p\hat{z} - p')$ was applied, yielding $\frac{1}{2} s' v_{s'}(p\hat{z})$.  

Now recall
\[
\begin{gathered}
  \overline{v}_{s'}(\mathbf{p})\gamma^{\mu} v_s(\mathbf{p}) = 2p^{\mu} \delta_{s's} \\ 
  \Longrightarrow \frac{1}{2} s d_s^{\dag}|0\rangle
\end{gathered}
\]

What about $(x^1 \partial^2 - x^2 \partial^1)$?  

\emph{Use this trick} of moving the partial around, 
\[
\int d^3 x e^{ipx} \overline{v}_s(p\hat{z})\gamma^0 (i x^1 \partial^2 - x^2 \partial^1)\Psi(x) |0 \rangle = i v_s(p\hat{z}) \gamma^0 \int d^3 x e^{ipx} ( (\partial_y (x \Psi(x)) ) - \partial_x(y \Psi(x)) )| 0 \rangle 
\]
Now $\mathbf{p} = p\hat{z}$.  
\[
\int dy \partial_y(x \Psi(x)) = \lim_{y\to \infty} \left. x \Psi(x) \right|_{-y}^y = 0 \text{ since } \Psi(x) \xrightarrow{x\to \infty} 0 
\]
Likewise $\int dx \partial_x(y \Psi(x)) = \lim_{x\to \infty} \left. y \Psi(x) \right|_{-x}^x = 0$.  

For the other identity, 
\[
J_z b_s^{\dag}(p\hat{z}) | 0 \rangle = J_z \int d^3x e^{ipx} \overline{\Psi}(x) \gamma^0 u_s(p\hat{z}) |0\rangle = \int d^3x e^{ipx} M^{12} \overline{\Psi}(x) \gamma^0 u_s(p\hat{z}) |0 \rangle = ** 
\]

Now
\[
\begin{aligned}
& [\overline{\Psi}(x), M^{12}] = ? \\
& [\Psi,M^{12}] = \Psi M^{12} - M^{12}\Psi \\
& [\Psi,M^{12}]^{\dag} = M^{12\dag} \Psi^{\dag} - \Psi^{\dag}M^{12\dag} = M^{21}\Psi^{\dag} - \Psi^{\dag} M^{21} = \Psi^{\dag} M^{12} - M^{12} \Psi^{\dag} = [\Psi^{ \dag}, M^{12} ] \\ 
& [\Psi^{\dag}, M^{12} ]\beta = \Psi^{\dag}M^{12} \beta - M^{12} \Psi^{\dag} \beta = \Psi^{\dag} \beta M^{12} - M^{12} \overline{\Psi} = \overline{\Psi}M^{12} - M^{12} \overline{\Psi} = [ \overline{\Psi}, M^{12} ] \\
& [\Psi, M^{12}]^{\dag} = (-i (x^1 \partial^2 - x^2 \partial^1) \Psi)^{\dag} + (S^{\mu \nu}\Psi)^{\dag} = i (x^1 \partial^2 - x^2 \partial^1 )\Psi^{\dag} + \Psi^{\dag} S^{\mu \nu \dag} \\ 
& [\Psi,M^{12}]^{\dag} \beta = i (x^1 \partial^2 - x^2 \partial^1)\Psi + \overline{\Psi} S^{12 \dag}
\end{aligned}
\]

\[
\begin{aligned}
  **  = \int d^3x e^{ipx} (\overline{\Psi}(x) M^{12} - i(x^1 \partial^2 - x^2 \partial^1) \overline{\Psi} - \overline{\Psi} S^{12} ) \gamma^0 u_s(p\hat{z}) |0 \rangle = 
\end{aligned}
\]
Consider the orbital part alone.
\[
\int d^3x e^{ipx} (-i(x\partial_y - y \partial_x )\overline{\Psi} \gamma^0 u_s(p\hat{z}) |0 \rangle = - i \lbrace \int d^3x e^{ipz}(x\partial_y - y \partial_x)\overline{\Psi}(x) \rbrace \gamma^0 u_s(p\hat{z}) |0\rangle 
\]
\[
\begin{aligned}
  & \int dy (x \partial_y \overline{\Psi} ) = \int dy (\partial_y (x\overline{\Psi}) ) = \left. \lim_{y\to \infty} (x \overline{\Psi}) \right|_{-y}^{y} = 0 \\ 
  & \text{ Likewise } \\ 
  & \int dx (y \partial_x \overline{\Psi} ) = 0 
\end{aligned}
\]

So we're left with $\int d^3x e^{ipx} \overline{\Psi} S_z \gamma^0 u_s(p\hat{z}) |0 \rangle$.   \\ 
$S_z u_s(p\hat{z}) = \frac{-s}{2} u_s(p\hat{z})$ (from (38.27) and our previous discussion).  
\[
\Longrightarrow \int d^3x e^{ipx} [ \sum_{s'= \pm} \int \widetilde{dp}' ( b_{s'}^{\dag}(p')\overline{u}_{s'}(\mathbf{p}') e^{-ip'x} + d_{s'}(\mathbf{p}') \overline{v}_{s'}(\mathbf{p}') e^{ip'x} )]\gamma^0 \left( \frac{s}{2} u_s(p\hat{z}) \right) |0 \rangle
\]
with $d_{s'}|0\rangle =0$, \\
$\Longrightarrow \sum_{s' = \pm} \int \widetilde{dp}' \int d^3x e^{i(p - p') x} \overline{u}_{s'}(p') \gamma^0 u_s(p\hat{z}) \frac{s}{2} b_{s'}^{\dag}(\mathbf{p}') | 0 \rangle$.  

$\int d^3x e^{i(p-p')x} = (2\pi)^3 \delta^3(p-p')$,  \, applying $\int \widetilde{dp}'$ and using $\overline{u}_{s'}(\mathbf{p}) \gamma^{\mu} u_s(\mathbf{p}) = 2p^{\mu} \delta_{s's}$, 
\[
\Longrightarrow \boxed{ \frac{s}{2} b_s^{\dag}(p\hat{z})|0\rangle }
\]

\section{Parity, time reversal, and charge conjugation}

\subsection*{Problems}

\problemhead{40.1} Find the transformation properties of $\overline{ \Psi} S^{\mu \nu} \Psi$ and $\overline{\Psi} i S^{\mu \nu} \gamma_5 \Psi$ under $P$, $T$, $C$.  Verify that they are both even under $CPT$, as claimed.  Do either or both vanish if $\Psi$ is a Majorana field?

\solutionhead{40.1} Recall $S^{\mu \nu} = \frac{i}{4} [ \gamma^{\mu}, \gamma^{\nu} ] = \frac{i}{2} \gamma^{\mu} \gamma^{\nu}$ for $\mu \neq \nu$.  

\[
\begin{aligned}
  & C^{-1} ( \overline{ \Psi} S^{\mu \nu } \Psi ) C   = \overline{\Psi} \mathcal{C}^T \left( \frac{i}{2} \gamma^{\mu} \gamma^{\nu} \right)^T \mathcal{C} \Psi = \frac{i}{2} \overline{\Psi} \mathcal{C}^T \gamma^{\nu T} \gamma^{\mu T} \mathcal{C} \Psi = \frac{i}{2} \overline{\Psi} (- \gamma^{\nu}) ( - \gamma^{\mu}) \Psi = - \overline{\Psi} S^{\mu \nu} \Psi \\
  & P^{-1} ( \overline{ \Psi} S^{\mu \nu } \Psi ) P   = \overline{\Psi} (\beta S^{\mu \nu} \beta) \Psi = \overline{ \Psi} \beta \frac{i}{2} \gamma^{\mu} \gamma^{\nu} \Psi = \frac{i}{2} \overline{\Psi} \mathcal{P}^{\mu}_a \gamma^a \mathcal{P}^{\nu}_b \gamma^b \Psi = \overline{\Psi} S^{\mu \nu} \Psi \\ 
  & T^{-1} ( \overline{ \Psi} S^{\mu \nu } \Psi ) T   =  \overline{\Psi} ( \gamma_5 \mathcal{C}^{-1} \left( \frac{i}{2} \gamma^{\mu} \gamma^{\nu} \right)^* \mathcal{C} \gamma_5 ) \Psi = \frac{-i}{2} \overline{\Psi} (\gamma_5 \mathcal{C}^{-1} \gamma^{\mu *} \mathcal{C} \gamma_5 \gamma_5 \mathcal{C}^{-1} \gamma^{\nu *} \mathcal{C} \gamma_5 \Psi = \\
  & = \frac{-i}{2} \overline{\Psi} ( \mathcal{T}^{\mu}_a \gamma^a \mathcal{T}^{\nu}_b \gamma_b ) \Psi = - \frac{i}{2} \overline{\Psi} \gamma^{\mu} \gamma^{\nu} \Psi = - \overline{\Psi} S^{\mu \nu} \Psi 
\end{aligned}
\]

\[
\begin{aligned}
  & C^{-1} ( \overline{ \Psi} i S^{\mu \nu } \gamma_5 \Psi ) C   = i \overline{\Psi} \mathcal{C}^T \left( S^{\mu \nu} \gamma_5 \right)^T \mathcal{C} \Psi = i \overline{\Psi} \mathcal{C}^T \left( \frac{i}{2} \gamma_5 \gamma^{\nu T} \gamma^{\mu T} \right) \mathcal{C} \Psi = \\
  & = \frac{-1}{2} \overline{\Psi} ( \mathcal{C}^T \gamma_5^T \gamma^{\nu T} \gamma^{\mu T} ) \mathcal{C} \Psi = \frac{+i}{2} \overline{\Psi} (i \gamma_5 (-\gamma^{\nu } )(- \gamma^{\mu} ) ) \Psi  = - \overline{\Psi} (i \gamma_5 \frac{i}{2} \gamma^{\mu} \gamma^{\nu} ) \Psi = - \overline{\Psi} ( i  S^{\mu \nu} \gamma_5) \Psi \\
  & P^{-1} ( \overline{ \Psi} i S^{\mu \nu } \gamma_5 \Psi ) P   = \overline{\Psi} (i \beta \frac{i}{2} \gamma^{\mu} \gamma^{\nu} \gamma_5  \beta) \Psi = \overline{ \Psi} (-i) \frac{i}{2} \mathcal{P}^{\mu}_a \gamma^a \mathcal{P}^{\nu}_b \gamma^b \gamma_5 \Psi = -\overline{\Psi} i S^{\mu \nu} \gamma_5 \Psi \\ 
  & T^{-1} ( \overline{ \Psi} i S^{\mu \nu } \gamma_5 \Psi ) T   =  \overline{\Psi}  \gamma_5 \mathcal{C}^{-1} (-i) S^{\mu \nu *} \gamma_5^*  \mathcal{C} \gamma_5  \Psi =  \overline{\Psi} \gamma_5 \mathcal{C}^{-1} (-i) \left( \frac{-i}{2} \right) \gamma^{\mu *} \mathcal{C} \gamma_5 \gamma_5 \mathcal{C}^{-1} \gamma^{\nu *} \gamma_5^* \mathcal{C} \gamma_5 \Psi = \\
  & = \frac{(-i)^2 }{2} \overline{\Psi} ( -\mathcal{T}^{\mu}_a \gamma^a)( -\mathcal{T}^{\nu}_b \gamma_b ) \gamma_5 \Psi =  \overline{\Psi} i  S^{\mu \nu} \gamma_5 \Psi 
\end{aligned}
\]

Recall, for a Majorana field, $C^{-1} \Psi C = \Psi$ and $C^{-1} \overline{\Psi} C = \overline{\Psi}$, implying $C^{-1} (\overline{\Psi} A \Psi ) C = \overline{\Psi} A \Psi$ for any combination of gamma matrices $A$.  Clearly, 
\[
\begin{aligned}
  & C^{-1} ( \overline{\Psi} S^{\mu \nu} \Psi)C = \overline{\Psi} S^{ \mu \nu} \Psi = - \overline{\Psi} S^{\mu \nu } \Psi \\ 
  & C^{-1} ( \overline{\Psi} i S^{\mu \nu} \gamma_5 \Psi ) C = \overline{\Psi} i S^{\mu \nu} \gamma_5 \Psi = - \overline{\Psi} ( i S^{\mu \nu } \gamma_5 ) \Psi 
\end{aligned}
\]
so $\overline{\Psi} S^{\mu \nu} \Psi$, $\overline{\Psi} i S^{\mu \nu} \gamma_5 \Psi$ vanish under charge conjugation for Majorana fields.  

\section{LSZ reduction for spin-one-half particles}

\section{The free fermion propagator}

\solutionhead{42.1}

Recall 
\[
S(x-y)_{\alpha \beta} = \int \frac{d^4p }{ (2\pi)^4} e^{i p(x-y) } \frac{ (- \not p + m)_{\alpha \beta} }{ p^2 + m^2 - i \epsilon } \quad \quad \quad (42.12)
\]
obtained from 
\[
S(x-y)_{\alpha \beta} \equiv ( \langle 0 | T \Phi_{\alpha}(x) \overline{\Phi}_{\beta}(y) | 0 \rangle \quad \quad \quad (42.6) 
\]

Note $[ S(x-y) \mathcal{C}^{-1}]_{\alpha \beta} = - [ S(y-x) \mathcal{C}^{-1} ]_{\beta \alpha}$ is the same as \\
\phantom{Note } $[S(x-y) \mathcal{C}^{-1}]^T = - [ S(y-x) \mathcal{C}^{-1} ]$.  \\

Note also that $\begin{aligned} & \quad \\ 
  \mathcal{C}^T & = \mathcal{C}^{-1} \\ 
  \mathcal{C}^{-1} & =  - \mathcal{C} \end{aligned}$ 
\[
[ S(x-y) \mathcal{C}^{-1} ]^T_{\alpha \beta} = ( S(x-y)_{\alpha a} \mathcal{C}^{-1}_{\alpha \beta} )^T = \mathcal{C}^{-1}_{\beta a } S(x-y)_{a \alpha} 
\]

Now $[(-\not p +m)\mathcal{C}]^T = (-\not p - m ) \mathcal{C}$.  This is because of the following:
\[
\begin{gathered}
  \mathcal{C}^T = \mathcal{C}^{\dag} = \mathcal{C}^{-1} = - \mathcal{C}  \quad \quad \quad \, (38.34)
  \mathcal{C}^{-1} \gamma^{\mu} \mathcal{C} = - (\gamma^{\mu})^T \quad \quad \quad \, (38.36)
\end{gathered}
\]  


On the other hand, Srednicki recalls Prob. 39.4(d).  

\[
((-\not p + m)_{ab} \mathcal{C}_{bc} )^T = \mathcal{C}_{cb} ( - \not p + m )_{ba} = (-\not p - m)_{ab} \mathcal{C}_{bc}
\]
So 
\[
\begin{gathered}
  \mathcal{C}^{-1}_{\beta a}(- \not p + m )_{a\alpha} = - \mathcal{C}_{\beta a }(- \not p + m)_{a \alpha } = ( \not p + m)_{\alpha a } \mathcal{C}_{a \beta} \\
   \mathcal{C}^{-1}_{\beta a} S(x-y)_{a\alpha} = \int \frac{d^4p }{(2\pi)^4} e^{i p(x-y)} \frac{ (\not p + m)_{\alpha a} \mathcal{C}_{a \beta } }{ p^2 + m^2 - i \epsilon }
\end{gathered}
\]
Let $p \to - p$. Note $d^4p \to d^4p$ (4 $p$'s, even number of $p$'s)

\[
\begin{gathered}
  \mathcal{C}^{-1}_{\beta a} S(x-y)_{a\alpha} = \int \frac{d^4p }{ (2\pi)^4} e^{-ip (x-y) } \frac{ ( - \not p + m)_{\alpha a} \mathcal{C}_{a \beta } }{ p^2 + m^2 - i \epsilon } = S(y-x)_{\alpha a} \mathcal{C}_{a \beta} =\boxed{  - S(y-x)_{\alpha a } \mathcal{C}^{-1}_{a\beta} }
\end{gathered}
\]


\section{The path integral for fermion fields}

\section{Formal development of fermionic path integrals}

\section{The Feynman rules for Dirac fields}

\solutionhead{45.1}
\begin{itemize}
  \item[(a)] For the theory $\mathcal{L}_1 = g \varphi \overline{\Psi} \Psi$, 
\[
 P^{-1} \varphi \overline{\Psi} \Psi P  = (P^{-1} \varphi P )( - i \overline{\Psi} \beta ) ( i \beta \Psi ) = ( P^{-1} \varphi P) ( \overline{\Psi} \Psi )
\]
\[
\Longrightarrow P^{-1} \varphi P = \varphi 
\]

\[
(T^{-1} \varphi T) ( \overline{\Psi} \gamma_5 \mathcal{C}^{-1} )( \mathcal{C} \gamma_5 \Psi ) = T^{-1} \varphi T \overline{\Psi} \Psi
\]
\[
T^{-1} \varphi T = \varphi
\]
where $\gamma_5^2 = 1$
\[
\mathcal{C}^{-1} \varphi \mathcal{C} ( \Psi^T \mathcal{C}) ( \mathcal{C} \overline{\Psi}^T )  = \mathcal{C}^{-1} \varphi \mathcal{C} ( \Psi^T \overline{\Psi}^T ) = \mathcal{C}^{-1} \varphi \mathcal{C} (\overline{\Psi} \Psi  ) 
\]
\[
\mathcal{C}^{-1} \varphi \mathcal{C} = \varphi
\]
Then $\varphi$ is invariant under CPT.  
  \item[(b)]
\[
P^{-1} \varphi P P^{-1} i \overline{ \Psi} \gamma_5 \Psi P = (P^{-1} \varphi P)( - \overline{\Psi} i \gamma_5 \Psi)
\]
so $P^{-1} \varphi P = - \varphi$.  

\[
T^{-1} \varphi T T^{-1} \overline{\Psi} i \gamma_5 \Psi T = T^{-1} \varphi T(- \overline{\Psi} i \gamma_5 \Psi ) 
\]
so $T^{-1} \varphi T = - \varphi$.  

\[
\mathcal{C}^{-1} \varphi \mathcal{C} \mathcal{C}^{-1} \overline{\Psi} i \gamma_5 \Psi \mathcal{C} = \mathcal{C}^{-1} \varphi \mathcal{C} \overline{\Psi} i \gamma_5 \Psi 
\]
so $\mathcal{C}^{-1} \varphi \mathcal{C} = \varphi$.  

Since $\varphi$ is odd under $T, P$ and even under $C$, $\varphi$ is even under CPT.  
\end{itemize}

\solutionhead{45.2}


\[
\boxed{ i\mathcal{T}_{ e^+ e^+ \to e^+ e^+ }  = \frac{1}{i} (ig)^2 \left[ \frac{ (\overline{v}_1 v_1' )(\overline{v}_2 v_2' ) }{ -t + M^2 } - \frac{ ( \overline{v}_2 v_2' )( \overline{v}_2 v_1' ) }{ -u + M^2 } \right] }
\]

\[
\boxed{ i\mathcal{T}_{\varphi \varphi  \to e^+ e^- } = \frac{1}{i} (ig)^2 \overline{u}'_2 \left[ \frac{ \not p'_1  - \not k_1 + m }{ -t + M^2 } + - \frac{ - \not p_2' + \not k_1 + m }{ -u + M^2 } \right] v_1' }
\]

\section{Spin Sums}

\section{ Gamma matrix technology}

\section{Spin-averaged cross sections}

\section{The Feynman rules for Majorana fields}

\section{Massless particles and spinor helicity}

% 20121016

\[
u_s(\mathbf{p}) \overline{u}_s(\mathbf{p}) = \frac{1}{2} ( 1 + 2 \gamma_5) ( - \not p) \quad \quad \quad (50.1) 
\]

\[
 \not p = - \gamma^{\mu} p_{\mu} = \left( \begin{matrix} & p_{\mu} \sigma^{\mu} \\ p_{\mu} \overline{\sigma}^{\mu} & \end{matrix} \right)
\]

\[
u_-(\mathbf{p}) \overline{u}_-(\mathbf{p}) = \frac{1}{2} ( 1 - \gamma_5) ( - \not p ) \quad \quad \quad (50.3)
\]

\[
\begin{gathered}
  \frac{1}{2} ( 1 - \gamma_5) = \frac{1}{2} ( 1 -  \left( \begin{matrix} -1 &  & & \\ & -1 & & \\ & & 1 & \\ & & & 1 \end{matrix} \right) ) = \left( \begin{matrix} 1 & & & \\ & 1 & & \\ & & &  \\ & & & \end{matrix} \right) \quad \quad \quad (50.6) \\ 
 u_-(\mathbf{p}) \overline{u}_-(\mathbf{p}) = \left( \begin{matrix} & - p _{\mu} \sigma^{\mu} \\ & \end{matrix} \right) = \left( \begin{matrix} & - p_{a \dot{a}} \\ & \end{matrix} \right) \quad \quad (50.7)
\end{gathered}
\]
We know the lower 2 components of $u_-(\mathbf{p})$ vanish (it has to, by (30.7), so $u_-(\mathbf{p}) = \left( \begin{matrix} \phi_a \\ 0 \end{matrix} \right)$


\begin{equation}
  [pk] \equiv \phi^a \kappa_a \quad \quad \quad (50.16) \text{ twistor product }
\end{equation}

\[
\phi^a \kappa_a = \epsilon^{ac} \phi_c \kappa_a = + \epsilon^{ac} \kappa_a \phi_c = - \epsilon^{ca} \kappa_a \phi_c = - \kappa^c \phi_c
\]

\begin{equation}
 [kp ] = -[pk] \quad \quad \quad (50.17)
\end{equation}

Recall (50.8) $u_-(\mathbf{p}) = \left( \begin{matrix} \phi_a \\ \end{matrix} \right)$

\begin{equation}
  \overline{u}_+(\mathbf{p}) u_-( \mathbf{k}) = [pk ] \quad \quad (50.18)
\end{equation}

Now  (50.14) is

\[
\overline{u}_+(\mathbf{p}) = (  \begin{matrix} \phi^a & \end{matrix} )
\]

Define 

\begin{equation}
\langle pk \rangle \equiv \phi^*_{\dot{a}} \kappa^{* \dot{a}} \quad \quad \quad (50.19)
\end{equation}

\[
\phi^*_{\dot{a}} = (\phi_a)^*
\]

\begin{equation}
\langle pk \rangle =  \phi^*_{\dot{a}} \kappa^{*\dot{a}} = (\kappa^a \phi_a)^* = [kp]^* \quad \quad \quad (50.20)
\end{equation}


Now (50.18)

\[
\langle kp \rangle = [pk]^* = - [kp]^* = -\langle pk \rangle 
\]

\[
\text{ Antisymmetric } \langle kp \rangle = - \langle pk \rangle \quad \quad \quad (50.21)
\]

\[
\overline{u}_-(\mathbf{p}) u_+(\mathbf{k}) = ( \begin{matrix} & \phi_{\dot{a}}^* \end{matrix} ) \left( \begin{matrix}  \\ \kappa^{* \dot{a}} \end{matrix} \right) = \phi^*_{\dot{a}} \kappa^{* \dot{a}} = \langle pk \rangle
\]

\[
\overline{u}^+(\mathbf{p}')(- \not k)u_+(\mathbf{p}) = \phi^{'a} \kappa_a \kappa_{\dot{a}}^* \phi^{* \dot{a}} = [p' k] \langle kp \rangle \quad \quad \quad (50.28)
\]
\[
\begin{aligned}
  & | p ] = u_-(\mathbf{p}) = v_+(\mathbf{p}) \\ 
  & | p \rangle  = u_+(\mathbf{p}) = v_-(\mathbf{p}) \\ 
  & [ p |  = \overline{u}_+(\mathbf{p}) = \overline{v}_-(\mathbf{p}) \\ 
  & \langle p |  = \overline{u}_-(\mathbf{p}) = \overline{v}_+(\mathbf{p})
\end{aligned} \quad \quad \quad \begin{gathered}
  [ p' | ( | k \rangle [ k | + | k ] \langle k | ) | p \rangle = [ p' k ] \langle kp \rangle \\ 
    \overline{u}_-(\mathbf{p}')(-\not k) u_-(\mathbf{p}) = \phi^{'*}_{\dot{a}} \kappa^{ * \dot{a}} \kappa^a \phi_a = \langle p' k \rangle [ kp ] \quad \quad \quad (50.29) \\
\langle p' | ( | k \rangle [ k | + | k ] < k | ) | p ] = \langle p' k \rangle [ kp ]
\end{gathered}
\]



\problemhead{50.1}

\begin{enumerate}
\item[(a)]
\[
\not p = \gamma^{\mu} p_{\mu} = \left( \begin{matrix} & p_{\mu} \sigma^{\mu} \\ p_{\mu} \overline{\sigma}^{\mu} & \end{matrix} \right) = \left( \begin{matrix} & p_{a \dot{a}} \\ p^{\dot{a} a } & \end{matrix} \right) = \left( \begin{matrix} & - \phi_a \phi_{\dot{a}}^* \\ - \phi^a (\phi^*)^{\dot{a}} & \end{matrix} \right)
\]

\[
p^{\dot{a} a} = \epsilon^{ac} \epsilon^{\dot{a} \dot{c}} p_{c \dot{c}} = \epsilon^{ac} \epsilon^{\dot{a} \dot{c}} (- \phi_c \phi^*_{\dot{c}} ) = - \phi^a (\phi^*)^{\dot{a}}
\]
Note that $\phi^{ *\dot{a}} = \epsilon^{\dot{a} \dot{c}} \phi^*_{\dot{c}}$

\[
\begin{aligned}
   & | p \rangle [ p | = u_+(\mathbf{p}) \overline{u}_+(\mathbf{p}) = \left( \begin{matrix} \quad \\ \phi^{* \dot{a}} \end{matrix} \right) ( \begin{matrix} \phi^a &  \end{matrix} ) = \left( \begin{matrix} & \\ \phi^{ * \dot{a}} \phi^a & \end{matrix} \right) = \left( \begin{matrix} & \\ \phi^a \phi^{* \dot{a}} & \end{matrix} \right) \\ 
&  | p ] \langle p | = u_-(\mathbf{p}) \overline{u}_-(\mathbf{p}) =\left( \begin{matrix} & \phi_a \phi_{\dot{a}}^* \\ & \end{matrix} \right)
\end{aligned}
\]

where the following useful facts were used:
\[
\begin{aligned}
  & u_-(\mathbf{p}) = \left( \begin{matrix} \phi_a \\ \, \end{matrix} \right)  \quad \quad \, (50.8) \\ 
 & \overline{u}_-(\mathbf{p}) = ( \begin{matrix} & \phi_{\dot{a}}^* \end{matrix}) \quad \quad \, (50.10) \\ 
  & u_+(\mathbf{p}) = \left( \begin{matrix}  \,  \\ \phi^{ * \dot{a}} \end{matrix} \right) \quad \quad \, (50.13) \\ 
  & \overline{u}_+(\mathbf{p}) = ( \begin{matrix} \phi^a & \end{matrix} ) \quad \quad \, (50.14 )
\end{aligned}
\]
\[
\Longrightarrow - \not p = + \left( |p \rangle [ p | + | p ] \langle p | \right)
\]
\item[(b)] 
\end{enumerate}

\section{Loop corrections in Yukawa theory}

\section{Beta functions in Yukawa theory}

\section{Functional determinants}

\section*{Part III - Spin One}

\section{Maxwell's equations}

Recall convention: 

\begin{align}
  &  \partial_{\mu} \equiv \frac{ \partial }{ \partial x^{\mu} }   = ( \partial_t, \nabla ) \quad \quad \, & (1.14)  \\
  & \partial^{\mu} \equiv \frac{ \partial }{ \partial x_{\mu}  }  = ( -\partial_t , \nabla ) \quad \quad \, & (1.15) 
\end{align}


Also, recall 
\[
\begin{gathered}
  \partial^{\mu} x^{\nu} = g^{\mu \nu } \quad \quad \, (1.16) \\
  x^{\mu} = (t,x)
\end{gathered}
\]


\begin{align}
  \nabla \cdot \mathbf{E} & = \rho  \quad \quad \, & (54.1) \\ 
  \nabla \times \mathbf{B} - \dot{\mathbf{E}} & = \mathbf{J} \quad \quad \, & (54.2) \\ 
  \nabla \times \mathbf{E} + \dot{ \mathbf{ B} } & = 0 \quad \quad \, & (54.3) \\ 
  \nabla \cdot \mathbf{B} & = 0 \quad \quad \, & (54.4)
\end{align}

\[
\begin{aligned}
  & \mathbf{E} = - \nabla \varphi - \dot{ \mathbf{ A } }   \quad \quad \, & (54.5)\\ 
  & \mathbf{B} = \nabla \times \mathbf{ A }   \quad \quad \, & (54.6) \\ 
  & \varphi' = \varphi + \dot{ \Gamma } \quad \quad \, & (54.7) \\ 
 &  \mathbf{A}' = \mathbf{A} - \nabla \Gamma   \quad \quad \, & (54.8) \\ 
  & A^{\mu} \equiv ( \varphi , \mathbf{A} )  \quad \quad \,  & (54.9)
\end{aligned}
\]
\[
F^{\mu \nu } \equiv \partial^{\mu} A^{\nu} - \partial^{\nu} A^{\mu} \quad \quad \, (54.10)
\]
\[
E^i = - \partial^i \varphi - \partial_t A^i = \partial^0 A^i - \partial^i A^0 = F^{0i }
\]
\[
F^{ 0i} = E^i \quad \quad \, (54.11)
\]
\[
\epsilon^{ijk} B_k = \epsilon^{ijk} \epsilon_{klm} \partial^l A^m = \partial^i A^j - \partial^j A^i = F^{ij} 
\]
\[
F^{ij } = \epsilon^{ijk} B_k \quad \quad \, (54.12)
\]
\[
\partial_{\nu} F^{\mu \nu} = J^{\mu} \quad \quad \, (54.13)
\]
indeed 
\[
J^{\mu} \equiv ( \rho , \mathbf{J}) \quad \quad \, (54.14) 
\]

\[
\begin{gathered}
  \partial_i F^{0i} = \partial_i E^i = J^0 \Longrightarrow \nabla \cdot E = \rho \\ 
  J^i = \partial_{\nu} F^{i \nu} = - \dot{E}^i + \partial_j F^{0j} = -\dot{E} + \epsilon^{ijk} \partial_j B_k = -\dot{E}^i + (\nabla \times B)^i 
\end{gathered}
\]

\[
(54.13) \Longrightarrow \partial_{\mu} \partial_{\nu} F^{\mu \nu} = \partial_{\mu} J^{\mu}
\]
\[
\partial_{\mu} \partial_{\nu} F^{\mu \nu} = - \partial_{\mu} \partial_{\nu} F^{\nu \mu} = - \partial_{\nu } \partial_{\mu} F^{\mu \nu} = - \partial_{\mu } \partial_{\nu} F^{\mu \nu } \text{ so } \partial_{\mu} \partial_{\nu} F^{\mu \nu} = 0 
\]
\begin{align}
  \partial_{\mu} J^{\mu} & = 0 \quad \quad \, (54.15) \\ 
 \dot{\partial} + \nabla \cdot J & = 0 \quad \quad \, (54.16) 
\end{align}

\[
\epsilon_{ \mu \nu \rho \sigma } \partial^{\rho} F^{\mu \nu } = 0 \quad \quad \, (54.17)
\]

If $\sigma= 0$, 
\[
\epsilon_{ijk0} \partial^k F^{ij} = \epsilon_{ijk} \partial^k F^{ij} = \epsilon_{ijk} \partial^k \epsilon^{ikl} B_l = \partial^k B_k = 0 
\]
\[
\begin{gathered}
  \epsilon_{ 0jki } \partial^k F^{0j} = \epsilon_{ 0jki} \partial^k ( \partial^0 A^j - \partial^j A^0) = \epsilon_{ 0jki } \partial^0 \partial^k  A^j - \epsilon_{ 0 jki } \partial^j \partial^k A^0  \\
  (\nabla \times E)_i + \dot{B}_i = \epsilon_{ijk} \partial^j E^k + \partial_0 B_i = \epsilon_{ijk} \partial^j E^k + \partial_0 \epsilon{jki} F^{jk} = \epsilon_{ijk} ( \partial^j F^{0k} + \partial_0 F^{jk} ) = \epsilon_{ijk} (\partial^j F^{0k} - \partial^0 F^{jk} )  
\end{gathered}
\]
Since 
\[
\epsilon_{ijl} F^{ij} = \epsilon_{ijl} \epsilon^{ijk} B_k = B_l
\]

\[
(A')^{\mu} = A^{\mu} - \partial^{\mu} \Gamma \quad \quad \, (54.18)
\]
\[
(F')^{\mu \nu} = \partial^{\mu}  (A')^{\nu} - \partial^{\nu} (A')^{\mu}
\]
\[
\begin{gathered}
  (F')^{\mu \nu} = F^{\mu \nu} - (\partial^{\mu } \partial^{\nu} - \partial^{\nu} \partial^{\mu} ) \Gamma \quad \quad \, (54.19) \\
 (F')^{\mu \nu} = F^{\mu \nu} \quad \quad \, (54.20 )
\end{gathered}
\]
treat current as external source. \\
Lorentz invariant, gauge invariant, parity and time-reversal invariant.   

\[
J^{\mu} ( A_{\mu}' - A_{\mu} ) = - J^{\mu} \partial_{\mu} \Gamma = (\partial_{\mu} J^{\mu} ) \Gamma - \partial_{\mu} (J^{\mu} \Gamma ) \quad \quad \, (54.22)
\]
$\partial_{\mu} J^{\mu} =0$ (current conservation) \\
up to total derivative, $J^{\mu } A_{\mu}$ gauge invariant. 





\section{Electrodynamics in Coulomb gauge}

Recall 
\[
\begin{aligned}
  & \varphi(x) = \int \frac{ d^4 k}{ (2\pi)^4} e^{ikx} \widetilde{ \varphi}(k) \quad \quad \, (8.6) \\ 
  & \widetilde{\varphi}(k) = \int d^4x e^{ -ikx} \varphi(x)
\end{aligned}
\]

\[
\begin{gathered}
  A_i(x) \to \int d^4x e^{-ikx} A_i(x) = \widetilde{ A}_i(k) \to (\delta_{ij} - \frac{ k_i k_j }{\mathbf{k}^2 } )\widetilde{A}_i(k) = \widetilde{A}_j - \frac{ k_j k_i }{ \mathbf{k}^2 } \widetilde{A}_i \\ 
  \int \frac{d^4k}{ (2\pi)^4} ( \widetilde{A}_j -  \frac{ \partial^i \partial^j \widetilde{A}_i }{ \mathbf{ \nabla}^2 } ) e^{ikx} 
\end{gathered}
\]

\subsection*{Problems}

\problemhead{55.1} Use eqs. (55.13)-(55.20) and $[A_i, A_j] = [\Pi_i, \Pi_j]=0$ (at equal times to verify eqs. (55.21)-(55.23).

\solutionhead{55.1} Recall
\begin{align*}
  a_{\lambda}(\mathbf{k}) & = + i \epsilon_{\lambda}(\mathbf{k}) \cdot \int d^3x e^{-ikx} \overleftrightarrow{\partial}_0 \mathbf{A}(x) \quad \quad \quad \, (55.16)  \\
  a_{\lambda}^{\dag}(\mathbf{k}) & = - i \epsilon_{\lambda}^*(\mathbf{k}) \cdot \int d^3x e^{+ikx} \overleftrightarrow{\partial}_0 \mathbf{A}(x) \quad \quad \quad \, (55.17)  
\end{align*}
where $f\overleftrightarrow{\partial}_{\mu} g = f(\partial_{\mu} g) - (\partial_{\mu} f) g$.

For $x^0 = y^0$,  
\[
\begin{aligned}
  \left[a_{\lambda}(\mathbf{k}), a^{\dag}_{\lambda'}(\mathbf{k}') \right] & = [ + i \epsilon_{\lambda}(\mathbf{k}) \cdot \int d^3x e^{-ikx} \overleftrightarrow{\partial}_0 \mathbf{A}(x), - i\epsilon^*_{\lambda'}(\mathbf{k}') \cdot \int d^3y e^{+ik'y} \overleftrightarrow{\partial}_0 \mathbf{A}(y) ] = \\
  & = \int d^3x d^3y e^{-ikx} e^{ik'y} [ \epsilon_{\lambda}(\mathbf{k}) \cdot ( \partial_0 \mathbf{A}(x) - (ik^0) \mathbf{A}(x)), \epsilon^*_{\lambda'}(\mathbf{k}') \cdot ( \partial_0 \mathbf{A}(y) + (i k^{'0})\mathbf{A}(y) ) ] = * 
\end{aligned}
\]
Recall $\Pi_i = \frac{ \partial \mathcal{L} }{ \partial \dot{A}_i} = \dot{A}_i \quad \quad \, (55.18)$.  Then
\[
* = \int d^3x d^3y e^{-ikx} e^{ik'y}[ \epsilon_{\lambda}(\mathbf{k}) \cdot (\mathbf{\Pi}(x) - i k^0 \mathbf{A}(x) ), \epsilon^*_{\lambda'}(\mathbf{k}') \cdot( \mathbf{\Pi}(y) + i k^{'0} \mathbf{A}(y)) ]  
\]

Given $[A_i, A_j] = [\Pi_i, \Pi_j] =0 \quad \, \forall \, i,j$, then
\[
\begin{aligned}
  * & = \int d^3x d^3y e^{-ikx} e^{ik'y}\lbrace ik^{'0} [ \epsilon_{\lambda}(\mathbf{k}) \cdot \mathbf{\Pi}(x), \epsilon^*_{\lambda'}(\mathbf{k}') \cdot \mathbf{A}(y) ] -ik^0 [\epsilon_{\lambda}(\mathbf{k}) \cdot \mathbf{A}(x), \epsilon^*_{\lambda'}(\mathbf{k}') \cdot \mathbf{\Pi}(y) ] \rbrace = \\
  & = i \int d^3x d^3y e^{-ikx} e^{ik'y}\lbrace ik^{'0} [ \epsilon_{i\lambda}(\mathbf{k})  \Pi^i(x), \epsilon^*_{j\lambda'}(\mathbf{k}')  A^j(y) ] -k^0 [\epsilon_{i\lambda}(\mathbf{k})  A^i(x), \epsilon^*_{j\lambda'}(\mathbf{k}')  \Pi^j(y) ] \rbrace  \quad \quad \, \text{ for indices $i,j = 1,2,3$} \\
  & = i \int d^3x d^3y e^{-ikx} e^{ik'y}\lbrace k^{'0} \epsilon_{i\lambda}(\mathbf{k}) \epsilon^*_{j\lambda'}(\mathbf{k}')(-) [ A^j(y), \Pi^i(x) ] -k^0 \epsilon_{i\lambda}(\mathbf{k}) \epsilon^*_{j\lambda'}(\mathbf{k}')  [A^i(x), \Pi^j(y)  ] \rbrace = 
\end{aligned}
\]

Recall $[A_i(\mathbf{x},t), \Pi_j(\mathbf{y},t)] = i \int \frac{ d^3k}{(2\pi)^3} e^{i\mathbf{k}\cdot (\mathbf{x} - \mathbf{y}) } \left( \delta_{ij} - \frac{k_i k_j}{ \mathbf{k}^2} \right) \quad \quad \, (55.20)$.  

\[
\begin{aligned}
  * & = \int d^3x d^y \lbrace k^{'0} \epsilon_{j \lambda}(\mathbf{k}) \epsilon^*_{i\lambda'}(\mathbf{k}') \int \frac{ d^3q}{(2\pi)^3} e^{i \mathbf{q}\cdot (\mathbf{y} - \mathbf{x}) }\left( \delta_{ij} - \frac{ q_i q_j}{\mathbf{q}^2} \right)  + k^0 \epsilon_{i\lambda}(\mathbf{k}) \epsilon^*_{j\lambda'}(\mathbf{k}') \int \frac{d^3q}{(2\pi)^3} e^{i\mathbf{q}\cdot(\mathbf{x} - \mathbf{y}) } \left( \delta_{ij} - \frac{ q_i q_j}{q^2} \right) \rbrace 
\end{aligned}
\]

Now
\[
\begin{gathered}
  e^{ik'y} = e^{i\mathbf{k}'\cdot \mathbf{y} - i k^{'0} t} \\ 
  \int d^3y e^{i \mathbf{k}'\cdot \mathbf{y} } e^{i \mathbf{q} \cdot \mathbf{y} } = (2\pi)^3 \delta^3(\mathbf{k}' + \mathbf{q})
\end{gathered}
\]

\[
\begin{aligned}
  * = \int d^3x d^3q e^{-ikx} e^{-ik^{'0} t } \lbrace & k^{'0} \epsilon_{j\lambda}(\mathbf{k}) \epsilon^*_{i\lambda'}(\mathbf{k}') e^{-i\mathbf{q}\cdot \mathbf{x}} \left( \delta_{ij} - \frac{q_i q_j}{q^2} \right) \delta^3(\mathbf{k}' + \mathbf{q}) + \\ & + k^0 \epsilon_{i\lambda}(\mathbf{k}) \epsilon^*_{j\lambda'}(\mathbf{k}') e^{i\mathbf{q}\cdot \mathbf{x}} \left( \delta_{ij} - \frac{q_i q_j}{q^2} \right) \delta^3(\mathbf{k}-\mathbf{q}) \rbrace
\end{aligned}
\]

\[
\begin{aligned}
   = \int d^3q  e^{i(k^0 - k^{'0}) t } \lbrace & k^{'0} \epsilon_{j\lambda}(\mathbf{k}) \epsilon^*_{i\lambda'}(\mathbf{k}')  \left( \delta_{ij} - \frac{q_i q_j}{q^2} \right) \delta^3(\mathbf{k}' + \mathbf{q})\delta^3(\mathbf{k} + \mathbf{q}) + k^0 \epsilon_{i\lambda}(\mathbf{k}) \epsilon^*_{j\lambda'}(\mathbf{k}')  \left( \delta_{ij} - \frac{q_i q_j}{q^2} \right) \delta^3(\mathbf{k}-\mathbf{q}) \delta^3(\mathbf{k}'-\mathbf{q})  \rbrace
\end{aligned}
\]

\[
\begin{aligned}
   =  e^{i(k^0 - k^{'0}) t } \lbrace & k^{'0} \epsilon_{j\lambda}(\mathbf{k}) \epsilon^*_{i\lambda'}(\mathbf{k}')  \left( \delta_{ij} - \frac{q_i q_j}{q^2} \right) \delta^3(\mathbf{k}' - \mathbf{k}) +  k^0 \epsilon_{i\lambda}(\mathbf{k}) \epsilon^*_{j\lambda'}(\mathbf{k}')  \left( \delta_{ij} - \frac{q_i q_j}{q^2} \right) \delta^3(\mathbf{k}'-\mathbf{k})   \rbrace
\end{aligned}
\]

$\epsilon_{j\lambda}(\mathbf{k})\epsilon^*_{i\lambda'}(\mathbf{k}') \delta_{ij} = \epsilon_{i\lambda}(\mathbf{k} \epsilon^*_{i\lambda'}(\mathbf{k}') = \delta_{\lambda'\lambda} \quad \quad \, \text{ from (55.14) }$  \\
Likewise for $\epsilon_{i\lambda}(\mathbf{k}) \epsilon^*_{j\lambda'}(\mathbf{k}') \delta_{ij} = \delta_{\lambda \lambda'} = \delta_{\lambda' \lambda}$.   \\

$\epsilon_{i\lambda}(\mathbf{k}) \epsilon^*_{j\lambda'}(\mathbf{k}') \frac{k_i k_j}{\mathbf{k}^2} = \frac{k_i \epsilon_{i\lambda}(\mathbf{k}) k_j \epsilon^*_{j\lambda'}(\mathbf{k}') }{k^2} = \frac{ (k \cdot \epsilon_{\lambda}(\mathbf{k}) )( \mathbf{k}\cdot \epsilon^*_{\lambda'}(\mathbf{k}') ) }{k^2} = 0 $ \\
\quad \, since $\mathbf{k}\cdot \epsilon_{\lambda}(\mathbf{k}) =0$ \, (55.13).   \\
Likewise for $\epsilon_{j\lambda}(\mathbf{k}) \epsilon^*_{i\lambda'}(\mathbf{k}') \frac{k_i k_j}{\mathbf{k}^2} =0$.  

Note that $k^0 = \omega = \omega(\mathbf{k})$, $k^{'0} = \omega' = \omega'(\mathbf{k})$ and if $\mathbf{k} = \mathbf{k}'$  (there's only 1 dispersion relationship for 1 physical system), $\omega = \omega'$.  

\[
\Longrightarrow * = (2\pi)^3 \omega \delta_{\lambda \lambda'} (1 + 1) \delta^3(\mathbf{k}-\mathbf{k}') = \boxed{ 2 \omega(2\pi)^3 \delta_{\lambda \lambda'}\delta^3(\mathbf{k} - \mathbf{k}') }
\]

\hrulefill

\[
\begin{aligned}
  \left[ a_{\lambda}(\mathbf{k}), a_{\lambda'}(\mathbf{k'}) \right] & = - \int d^3x d^3y e^{-ikx} e^{-ik'y} \left[ \epsilon_{\lambda}(\mathbf{k}) \cdot (\partial_0 \mathbf{A}(x) - (ik^0) \mathbf{A}(x) ), \epsilon_{\lambda'}(\mathbf{k}') \cdot(\partial_0 \mathbf{A}(y) - (i k^{'0})\mathbf{A}(y)) \right] = \\
  & = - \int d^3x d^3y e^{-ikx} e^{-ik'y} [ \epsilon_{\lambda}(\mathbf{k})\cdot(\mathbf{\Pi}(x) - ik^0 \mathbf{A}(x)), \epsilon_{\lambda'}(\mathbf{k}') \cdot(\mathbf{\Pi}(y) - ik^{'0} \mathbf{A}(y) ) ] = \\
  & = -\int d^3x d^3y e^{-ikx} e^{-ik'y} \lbrace [ \epsilon_{\lambda}(\mathbf{k}) \cdot \mathbf{\Pi}(x), \epsilon_{\lambda'}(\mathbf{k}') \cdot (-ik^{'0} \mathbf{A}(y)) ] + [ \epsilon_{\lambda}(\mathbf{k})\cdot(-ik^0 \mathbf{A}(x) ), \epsilon_{\lambda'}(\mathbf{k}') \cdot \mathbf{\Pi}(y) ] \rbrace = \\
  & = -\int d^3x d^3y e^{-ikx} e^{-ik'y} \lbrace [ \epsilon_{i\lambda}(\mathbf{k})  \Pi^i(x), -ik^{'0} \epsilon_{j\lambda'}(\mathbf{k}') A^j(y) ] + [ -ik^0 \epsilon_{i\lambda}(\mathbf{k}) A^i(x) , \epsilon_{j\lambda'}(\mathbf{k}') \mathbf{\Pi}^j(y) ] \rbrace = \\
  & = - \int d^3x d^3y e^{-ikx} e^{-ik'y} \lbrace +i k^{'0} \epsilon_{i\lambda}(\mathbf{k}) \epsilon_{j\lambda'}(\mathbf{k}') [A^j(y), \Pi^i(x)] - i k^0 [ \epsilon_{i\lambda}(\mathbf{k}) \epsilon_{j\lambda'}(\mathbf{k}') [A^i(x), \Pi^j(y)] \rbrace = \\
    & = \int d^3x d^3y e^{-ikx} e^{-ik'y} \lbrace k^{'0} \epsilon_{i\lambda}(\mathbf{k}) \epsilon_{j\lambda'}(\mathbf{k}') \int \frac{d^3 q}{ (2\pi)^3} e^{i\mathbf{q}\cdot (\mathbf{y}-\mathbf{x}) } \left( \delta_{ij} - \frac{q_i q_j}{q^2} \right) - \\ 
    & \phantom{ d^3x d^3y e^{-ikx} e^{-ik'y}= } - k^0 \epsilon_{i\lambda}(\mathbf{k})\epsilon_{j\lambda'}(\mathbf{k}') \int \frac{d^3q}{(2\pi)^3} e^{i \mathbf{q}\cdot (\mathbf{x}-\mathbf{y}) } \left( \delta_{ij} - \frac{q_i q_j}{q^2} \right) \rbrace = **
\end{aligned}
\]

\[
\begin{aligned}
 \int d^3y e^{-ik'y} e^{i \mathbf{q}\cdot \mathbf{y} }  & = e^{ik^{0'} t} \int d^3y e^{-i(\mathbf{k}' - \mathbf{q} )\cdot \mathbf{y} } = e^{ik^{0'} t} (2\pi)^3 \delta^3(\mathbf{k}' - \mathbf{q}) \\ 
 \int d^3x e^{-ikx} e^{-i\mathbf{q}\cdot \mathbf{x} } & = e^{ik^0 t} \int d^3x e^{-i (\mathbf{k} + \mathbf{q}) \cdot \mathbf{x} } = e^{ik^0 t} (2\pi)^3 \delta^3(\mathbf{k} + \mathbf{q})
\end{aligned}
\]

\[
\begin{aligned}
  ** & = (2\pi)^3 \int d^3q e^{i(k^0 + k^{'0}) t} \lbrace k^{'0} \epsilon_{j\lambda}(\mathbf{k}) \epsilon_{i\lambda'}(\mathbf{k}') \left( \delta_{ij}  - \frac{ q_i q_j}{q^2} \right) \delta^3(\mathbf{k}' - \mathbf{q}) \delta^3(\mathbf{k} + \mathbf{q}) - \\
  & \phantom{ (2\pi)^3 \int d^3q e^{i(k^0 + k^{'0}) t} } - k^0 \epsilon_{i\lambda}(\mathbf{k}) \epsilon_{j\lambda'}(\mathbf{k}') \left( \delta_{ij} - \frac{q_i q_j}{q^2} \right) \delta^3(\mathbf{k}' + \mathbf{q} )\delta^3(\mathbf{k}-\mathbf{q}) \rbrace =  \\
  & = (2\pi)^3 e^{i\omega t} \omega^0 \left( \epsilon_{i\lambda}(\mathbf{k}) \epsilon_{i\lambda'}( - \mathbf{k}) - \epsilon_{i\lambda}(\mathbf{k}) \epsilon_{i\lambda'}(-\mathbf{k}) \right) \delta^3(\mathbf{k} + \mathbf{k}') = 0  
\end{aligned}
\]

\hrulefill

\[
\begin{aligned}
  \left[ a_{\lambda}^{\dag}(\mathbf{k}), a^{\dag}_{\lambda'}(\mathbf{k}') \right] & = - [ \epsilon^*_{\lambda}(\mathbf{k}) \cdot \int d^3x e^{ikx} (\mathbf{\Pi}(x) + i k^0 \mathbf{A}(x)), \epsilon^*_{\lambda'}(\mathbf{k}') \cdot \int d^3y e^{ik'y} ( \mathbf{\Pi}(y) + i k^{'0} \mathbf{A}(y) ) ] = \\
  & = -\int d^3x d^3y e^{ikx} e^{ik'y} \lbrace [ e^*_{\lambda}(\mathbf{k}) \cdot \mathbf{\Pi}(x), i k^{'0} \epsilon^*_{\lambda'}(\mathbf{k}') \cdot \mathbf{A}(y) ] + [ i k^0 \epsilon_{\lambda}^*(\mathbf{k}) \cdot \mathbf{A}(x), \epsilon^*_{\lambda'}(\mathbf{k}')\cdot \mathbf{\Pi}(y) ] \rbrace = \\
  & = -\int d^3x d^3y e^{ikx} e^{ik'y} \lbrace i k^{'0} \epsilon^*_{i\lambda}(\mathbf{k}) \epsilon^*_{j\lambda'}(\mathbf{k}') [\Pi^i(x), A^j(y) ] + i k^0\epsilon^*_{i\lambda}(\mathbf{k}) \epsilon^*_{j\lambda'}(\mathbf{k}') [A^i(x), \Pi^j(y) ] \rbrace = \\
  & = - \int d^3x d^3y e^{ikx} e^{ik'y} \lbrace k^{'0} \epsilon^*_{i\lambda}(\mathbf{k}) \epsilon^*_{j\lambda'}(\mathbf{k}') \int \frac{d^3q}{(2\pi)^3} e^{i \mathbf{q}\cdot (\mathbf{y} - \mathbf{x}) }\left( \delta_{ij} - \frac{q_i q_j}{q^2} \right) - \\ 
  & \phantom{ - \int d^3x d^3y e^{ikx} e^{ik'y} \lbrace } -k^0 \epsilon^*_{i\lambda}(\mathbf{k}) \epsilon^*_{j\lambda'}(\mathbf{k}') \int \frac{d^3q}{(2\pi)^3} e^{i\mathbf{q}\cdot (\mathbf{x} - \mathbf{y}) } \left( \delta_{ij} - \frac{ q_i q_j}{q^2} \right) \rbrace = \\
  & = (2\pi)^3 \int d^3q e^{-i(k^0 +k^{'0} )t } \lbrace k^{'0} \epsilon^*_{i\lambda}(\mathbf{k}) \epsilon^*_{j\lambda'}(\mathbf{k}') \left( \delta_{ij} - \frac{q_i q_j}{q^2} \right) \delta^3(\mathbf{k}' + \mathbf{q})\delta^3(\mathbf{k} - \mathbf{q}) - \\
   & \phantom{ = (2\pi)^3 \int d^3q e^{-i(k^0 +k^{'0} )t } } - k^0 \epsilon^*_{i\lambda}(\mathbf{k}) \epsilon^*_{j\lambda'}(\mathbf{k}') \left( \delta_{ij} - \frac{q_i q_j}{q^2} \right) \delta^3(\mathbf{k}'-\mathbf{q}) \delta^3(\mathbf{k}+\mathbf{q}) \rbrace = \\
  & = (2\pi)^3 e^{-i2 \omega t} \lbrace k^0 \epsilon^*_{i\lambda}(\mathbf{k}) \epsilon^*_{j\lambda'}(-\mathbf{k}) \left( \delta_{ij} - \frac{k_i k_j}{k^2} \right) \delta^3(\mathbf{k}' + \mathbf{k}) - \\
&  \phantom{ = (2\pi)^3 e^{-i2 \omega t} \lbrace } -k^0 \epsilon^*_{i\lambda}(\mathbf{k}) \epsilon^*_{j\lambda'}(-\mathbf{k}) \left( \delta_{ij}  - \frac{k_i k_j}{k^2} \right) \delta^3(\mathbf{k}' + \mathbf{k}) \rbrace = 0
\end{aligned}
\]

\problemhead{55.2} Use eqs. (55.11), (55.14), (55.19), and (55.21)-(55.23) to verify eq. (55.24).  

\solutionhead{55.2} Recall that 
\[
\mathbf{A}(x) = \sum_{\lambda = \pm} \int \widetilde{dk} [ \epsilon_{\lambda}^*(\mathbf{k}) a_{\lambda}(\mathbf{k}) e^{ikx} + \epsilon_{\lambda}(\mathbf{k}) a_{\lambda}^{\dag}(\mathbf{k}) e^{-ikx} ] \quad \quad \, (55.11)
\]  

\[
\begin{aligned}
  & \epsilon_{\lambda'}(\mathbf{k}) \cdot \epsilon_{\lambda}^*(\mathbf{k}) = \delta_{\lambda' \lambda} \quad \quad \quad \, (55.14) \\ 
  & \mathcal{H} = \Pi_i \dot{A}_i - \mathcal{L} = \frac{1}{2} \Pi_i \Pi_i + \frac{1}{2} \nabla_j A_i \nabla_j A_i - J_i A_i + \mathcal{H}_{coul} \quad \quad \, (55.19) \, \\
& \mathcal{H}_{coul} = - \mathcal{L}_{coul}
\end{aligned} \quad \quad \quad \, \begin{aligned}
  & [a_{\lambda}(\mathbf{k}), a_{\lambda'}(\mathbf{k}') ] = 0  \\
  & [a_{\lambda}^{\dag}(\mathbf{k}), a_{\lambda'}^{\dag}(\mathbf{k}') ] = 0  \\
  & [a_{\lambda}(\mathbf{k}), a_{\lambda'}^{\dag}(\mathbf{k}') ] = (2\pi)^3 2\omega \delta^3(\mathbf{k}- \mathbf{k}') \delta_{\lambda \lambda'}
\end{aligned}
\]

\[
\Pi_i = \dot{A}_i = \sum_{\lambda = \pm} \int \widetilde{dk} \left[ - i k^0 \epsilon_{i\lambda}^*(\mathbf{k}) a_{\lambda}(\mathbf{k})e^{ikx} + ik^0 \epsilon_{\lambda'}(\mathbf{k}) a_{\lambda}^{\dag}(\mathbf{k}) e^{-ikx} \right] = -i \sum_{\lambda = \pm} \int \widetilde{dk} k^0 \left[ \epsilon_{\lambda}^*(\mathbf{k}) a_{\lambda}(\mathbf{k}) e^{ikx} - \epsilon_{ \lambda}(\mathbf{k}) a_{\lambda}^{\dag}(\mathbf{k}) e^{-ikx} \right] 
\]

\[
\begin{aligned}
  \Pi_i \Pi_i & = - \sum_{\lambda, \lambda' = \pm} \int \widetilde{dk} \widetilde{dk}' k^0 k^{'0}  [ \epsilon_{i\lambda}^*(\mathbf{k}) a_{\lambda}(k) e^{ikx} - \epsilon_{i\lambda} a^{\dag}_{\lambda}(\mathbf{k}) e^{-ikx} ] [ \epsilon_{i\lambda}^*(\mathbf{k}') a_{\lambda'}(\mathbf{k}') e^{ik'x} - \epsilon_{i\lambda} a^{\dag}_{\lambda'}(\mathbf{k}') e^{-ikx} ] = \\
  & = -\sum_{\lambda, \lambda' = \pm} \int \widetilde{dk} \widetilde{dk}' k^0 k^{'0}  \times  \\
& \phantom{ = -\sum_{\lambda, \lambda' = \pm} \int \widetilde{dk} \widetilde{dk}' k^0 k^{'0}  }  \begin{aligned}
    & \times \left( \right. \epsilon^*_{i\lambda}(\mathbf{k}) \epsilon_{i\lambda'}^*(\mathbf{k}') a_{\lambda}(\mathbf{k}) a_{\lambda'}(\mathbf{k}') e^{i(k +k') x} + \epsilon_{i\lambda}(\mathbf{k}) \epsilon_{i\lambda'}(\mathbf{k}') a_{\lambda}^{\dag}(\mathbf{k})a_{\lambda'}^{\dag}(\mathbf{k}') e^{-i(k+k') x} - \\ 
    & - \epsilon_{i\lambda}^*(\mathbf{k}) \epsilon_{i\lambda'}(\mathbf{k}')e^{i(k-k') x} (a_{\lambda'}^{\dag}(\mathbf{k}') a_{\lambda}(\mathbf{k}) + (2\pi)^3 2\omega \delta^3(\mathbf{k}- \mathbf{k}') \delta_{\lambda \lambda'} ) - \\
    & - \epsilon_{i\lambda}(\mathbf{k}) \epsilon_{i\lambda'}^*(\mathbf{k'}) e^{-i(k-k')x} (a_{\lambda'}^{\dag} a_{\lambda})  \left. \right)
\end{aligned}  
\end{aligned}
\]

From this point on, I'll deal with the term $ \sum_{\lambda, \lambda' = \pm} \int \widetilde{dk} \widetilde{dk}' k^0 k^{'0}(2\pi)^3 2\omega \delta^3(\mathbf{k}- \mathbf{k}') \delta_{\lambda \lambda'}$ separately from $\Pi_i \Pi_i$.  Then the rest of the expression is labeled $*$.  

Recall that $\widetilde{dk}' = \frac{ d^3k'}{ (2\pi)^3 2k^{'0}}$ and $\int d^3x e^{i\mathbf{k}\cdot \mathbf{x}} = (2\pi)^3 \delta^3(\mathbf{k})$

\[
\begin{aligned}
  &  \xrightarrow{\int d^3x} \\
  & * = - \sum_{\lambda, \lambda' =\pm} \int \widetilde{dk} \widetilde{dk}' k^0 k^{'0} (2\pi)^3 \left( \right. \epsilon_{i\lambda}^*(\mathbf{k}) \epsilon_{i\lambda'}^*(\mathbf{k}') a_{\lambda}(\mathbf{k}) a_{\lambda'}(\mathbf{k}') e^{-i(k^0 +k^{'0})t } \delta^3(\mathbf{k} + \mathbf{k}') + \\ 
  & \phantom{ * = - \sum_{\lambda, \lambda' =\pm} \int \widetilde{dk} \widetilde{dk}' k^0 k^{'0} (2\pi)^3 } \begin{aligned} & + \epsilon_{i\lambda}(\mathbf{k}) \epsilon_{i\lambda'}(\mathbf{k}') a_{\lambda}^{\dag}(\mathbf{k}) a_{\lambda'}^{\dag}(\mathbf{k}') e^{i (k^0 + k^{'0} )t } \delta^3(\mathbf{k} + \mathbf{k}') + \\
    & + - \epsilon_{i\lambda}^*(\mathbf{k}) \epsilon_{i\lambda'}^*(\mathbf{k}') \delta^3(\mathbf{k} - \mathbf{k}') (a_{\lambda'}^{\dag}(\mathbf{k}') a_{\lambda}(\mathbf{k}) ) - \\
    & \phantom{+} - \epsilon_{i\lambda}^*(\mathbf{k}) \epsilon_{i\lambda'}^*(\mathbf{k}')\delta^3(\mathbf{k}- \mathbf{k}') (a_{\lambda'}^{\dag}(\mathbf{k}) a_{\lambda}(\mathbf{k}') ) \left. \right) =
\end{aligned} \end{aligned}
\]

\[
\begin{aligned}
  & = - \sum_{\lambda, \lambda' = \pm} \int \frac{ \widetilde{dk}}{2} k^0 (\epsilon_{i\lambda}^*(\mathbf{k})  \epsilon_{i\lambda'}^*(-\mathbf{k}) a_{\lambda}(\mathbf{k}) a_{\lambda'}( - \mathbf{k}) e^{-i(2\omega t)} + \epsilon_{i\lambda}(\mathbf{k}) \epsilon_{i\lambda'}(-\mathbf{k}) a_{\lambda}^{\dag}(\mathbf{k}) a_{\lambda'}^{\dag}(-\mathbf{k}) e^{2 i \omega t} ) + \\ 
  & + - \sum_{\lambda = \pm} \int \frac{\widetilde{dk}}{2} k^0 ( - (a_{\lambda}^{\dag} a_{\lambda}) - (a_{\lambda}^{\dag} a_{\lambda} ) )
\end{aligned}
\]

\[
\nabla_j A_i = \sum_{ \lambda = \pm} \int \widetilde{dk} [ \epsilon_{i\lambda}^*(\mathbf{k}) a_{\lambda}(\mathbf{k}) ik_j e^{ikx} + \epsilon_{i\lambda}(\mathbf{k}) a^{\dag}_{\lambda}(\mathbf{k}) (-ik_j) e^{-ikx} ]
\]

\[
\begin{aligned}
  \nabla_j A_i \nabla_j A_i & = - \sum_{\lambda, \lambda'= \pm} \int \widetilde{dk} \widetilde{dk}' \left[ \right. \epsilon^*_{i\lambda}(\mathbf{k}) \epsilon_{i\lambda'}^*(\mathbf{k}') k_j k_j' a_{\lambda}(\mathbf{k}) a_{\lambda'}(\mathbf{k}') e^{i(k +k') x} + \\
    & \phantom{- \sum_{\lambda, \lambda'= \pm} \int \widetilde{dk} \widetilde{dk}' } \begin{aligned} 
& + \epsilon_{i\lambda}(\mathbf{k}) \epsilon_{i\lambda'}(\mathbf{k}') k_j k_j' a^{\dag}_{\lambda}(\mathbf{k}) a^{\dag}_{\lambda'}(\mathbf{k}') e^{-i(k +k') x} - \\
   & -    \epsilon^*_{i\lambda}(\mathbf{k}) \epsilon_{i\lambda'}(\mathbf{k}') k_j k_j' a_{\lambda}(\mathbf{k}) a^{\dag}_{\lambda'}(\mathbf{k}') e^{i(k -k') x} -    \\
      & - \epsilon_{i\lambda}(\mathbf{k}) \epsilon_{i\lambda'}^*(\mathbf{k}') k_j' k_j' a_{\lambda}^{\dag}(\mathbf{k}) a_{\lambda'}(\mathbf{k}') e^{-i(k +k') x} \left. \right] 
\end{aligned}
\end{aligned}
\]

Using the commutation relation to put the term containing $a_{\lambda}(\mathbf{k})a_{\lambda}^{\dag}(\mathbf{k}')$ into ``normal ordering'', and then removing the term $ \sum_{\lambda, \lambda'= \pm} \int \widetilde{dk} \widetilde{dk}'  \epsilon^*_{i\lambda}(\mathbf{k}) \epsilon_{i\lambda'}(\mathbf{k}') k_j k_j' (2\pi)^3 2\omega \delta^3(\mathbf{k}- \mathbf{k}') \delta_{\lambda \lambda'}  e^{i(k -k') x}$ for later, and labeling the remaining terms $**$, then

\[
\begin{gathered}
  ** \xrightarrow{ \int d^3x} \\ 
  \begin{aligned}
    & = - \sum_{\lambda, \lambda'= \pm} \int \widetilde{dk} \widetilde{dk}' (2\pi)^3 \left[ \right. \epsilon^*_{i\lambda}(\mathbf{k}) \epsilon_{i\lambda'}^*(\mathbf{k}') k_j k_j' a_{\lambda}(\mathbf{k}) a_{\lambda'}(\mathbf{k}') e^{-2i\omega t} \delta^3(\mathbf{k} + \mathbf{k}') + \\
    & \phantom{ = - \sum_{\lambda, \lambda'= \pm} \int \widetilde{dk} \widetilde{dk}' (2\pi)^3 ( ) } 
 + \epsilon_{i\lambda}(\mathbf{k}) \epsilon_{i\lambda'}(\mathbf{k}') k_j k_j' a^{\dag}_{\lambda}(\mathbf{k}) a^{\dag}_{\lambda'}(\mathbf{k}') e^{2i\omega t} \delta^3(\mathbf{k} + \mathbf{k}') \left. \right] + \\
    & + - \sum_{\lambda, \lambda' = \pm} \int \widetilde{dk} \widetilde{dk}' (2\pi)^3 \left[ \right. - \epsilon_{i\lambda}^*(\mathbf{k}) \epsilon_{i\lambda}(\mathbf{k}) k_j k_j' a_{\lambda'}^{\dag}(\mathbf{k}') a_{\lambda}(\mathbf{k}) \delta^3(\mathbf{k}- \mathbf{k}') +  \\
      & \phantom{ + - \sum_{\lambda, \lambda' = \pm} \int \widetilde{dk} \widetilde{dk}' (2\pi)^3 ( ) } + \epsilon_{i\lambda}(\mathbf{k})(\epsilon_{i\lambda'}^*(\mathbf{k}') k_j' k_j a_{\lambda}^{\dag}(\mathbf{k})a_{\lambda'}(\mathbf{k}') \delta^3(\mathbf{k}-  \mathbf{k}') \left. \right] = 
  \end{aligned}
\end{gathered}
\]

\[ 
\begin{aligned}
  & = - \sum_{\lambda, \lambda'= \pm} \int \frac{ \widetilde{dk} }{2k^0 }  \left[ \right. \epsilon^*_{i\lambda}(\mathbf{k}) \epsilon_{i\lambda'}^*(-\mathbf{k}) (-k_j k_j) a_{\lambda}(\mathbf{k}) a_{\lambda'}(-\mathbf{k}) e^{-2i\omega t}  + \\
    & \phantom{ = - \sum_{\lambda, \lambda'= \pm} \int \widetilde{dk} \widetilde{dk}' } 
 + \epsilon_{i\lambda}(\mathbf{k}) \epsilon_{i\lambda'}(-\mathbf{k}) (-k_j k_j) a^{\dag}_{\lambda}(\mathbf{k}) a^{\dag}_{\lambda'}(-\mathbf{k}) e^{2i\omega t}  \left. \right] + \\
    & +  \sum_{\lambda = \pm} \int \frac{\widetilde{dk}}{2k^0}  \left[   k_j k_j' a_{\lambda'}^{\dag}(\mathbf{k}') a_{\lambda}(\mathbf{k}) + k_j k_j a_{\lambda}^{\dag}(\mathbf{k})a_{\lambda'}(\mathbf{k}') \right] 
  \end{aligned}
\]

\emph{Photons are massless}.  So $k^0k^0 - k_j k_j =0$.  

So then adding all the terms above for $*$ and $**$,

\[
* + ** = -\sum_{\lambda = \pm} \int \widetilde{dk} \left( \frac{ \omega^2}{2\omega} (-2a_{\lambda}^{\dag} a_{\lambda} ) - \frac{k_j k_j}{2\omega}(2a_{\lambda}^{\dag} a_{\lambda} ) \right) = \sum_{\lambda = \pm} \int \widetilde{dk} \left( \frac{ \omega}{2} ( 2 a_{\lambda}^{\dag} a_{\lambda} - 2a_{\lambda}^{\dag} a_{\lambda} ) \right) = \sum_{\lambda = \pm} \int \widetilde{dk} \left( \omega ( 2a_{\lambda}^{\dag} a_{\lambda} ) \right)
\]

Consider now the 2 terms I had neglected above.  

\[
\begin{aligned}
  & + \sum_{\lambda, \lambda' = \pm} \int \widetilde{dk} \widetilde{dk}' k^0 k^{'0} \epsilon_{i\lambda}^*(\mathbf{k}) \epsilon_{i\lambda'}(\mathbf{k}') e^{i(k-k')x} (2\pi)^3 2\omega \delta^3(\mathbf{k}' - \mathbf{k})\delta_{\lambda \lambda'} - \\
  & - \sum_{\lambda, \lambda' = \pm} \int \widetilde{dk} \widetilde{dk}' (- \epsilon_{i\lambda}^*(\mathbf{k}) \epsilon_{i\lambda'}(\mathbf{k}') k_j k_j' e^{i(k-k')x} (2\pi)^3 2\omega \delta^3(\mathbf{k} - \mathbf{k}')\delta_{\lambda \lambda'} 
\end{aligned}
\]

Now $\widetilde{dk}' = \frac{ d^3k'}{(2\pi)^3 2k^{'0} }$.  Then the two terms above is equal to 
\[
\begin{aligned}
  & = \sum_{\lambda = \pm} \int \frac{\widetilde{dk}}{2k^{'0}} \left( k^{02} \epsilon_{i\lambda}^*(\mathbf{k}\epsilon_{i\lambda}(\mathbf{k}) (2\omega) + \epsilon_{i\lambda}^*(\mathbf{k}) \epsilon_{i\lambda}(\mathbf{k}) k_j k_j (2\omega) \right) = \\
  & =  \sum_{\lambda = \pm } \int \widetilde{dk} (\omega^2 \delta_{\lambda \lambda'} + \omega^2 \delta_{\lambda \lambda'} ) = 2\int \frac{d^3k}{2\omega  (2\pi)^3} \omega^2(2)  = \\
  & = \frac{2}{(2\pi)^3 }\int d^3k \omega
\end{aligned}
\]
\[
\xrightarrow{\times \frac{1}{2}} \frac{2 }{2(2\pi)^3} \int d^3k \omega = 2\mathcal{\epsilon}_0 \xrightarrow{ \int d^3x} 2 \mathcal{\epsilon}_0 V
\]

Note that $ - \int d^3x J_i A_i = - \int d^3x \mathbf{J}(x) \cdot \mathbf{A}(x)$ and that \\
$H_{coul}  =\int d^3x \mathcal{H}_{coul} = \int d^3x - \mathcal{L}_{coul} = \int d^3x \frac{1}{2} \int d^3y \frac{ \rho(\mathbf{x},t) \rho(\mathbf{y},t) }{ 4\pi |\mathbf{x} - \mathbf{y} |}$.  

So then
\[
\boxed{ H = \sum_{\lambda = \pm} \int \widetilde{dk} \omega a_{\lambda}^{\dag}(\mathbf{k}) a_{\lambda}(\mathbf{k}) + 2 \epsilon_0 V - \int d^3 x \mathbf{J}(x) \cdot \mathbf{A}(x) + H_{coul}, \quad \, (55.24)  }
\]


\section{LSZ reduction for photons} 


\problemhead{56.1} Use eqs. (55.11) and (55.21)-(55.23) to verify eqs. (56.9) and (56.10).  

\solutionhead{56.1} Recall that 

\[
\mathbf{A}(x) = \sum_{\lambda = \pm} \int \widetilde{dk} [ \epsilon_{\lambda}^*(\mathbf{k}) a_{\lambda}(\mathbf{k}) e^{ikx} + \epsilon_{\lambda}(\mathbf{k}) a^{\dag}_{\lambda}(\mathbf{k}) e^{-ikx} ] \quad \quad \, (55.11)
\]

\[
\begin{aligned}
  & \left[ a_{\lambda}(\mathbf{k}), a_{\lambda'}(\mathbf{k}') \right] = 0  \quad \, (55.21)   \\
  &  [a_{\lambda}^{\dag}(\mathbf{k}), a_{\lambda'}^{\dag}(\mathbf{k}') ] = 0 \quad \, (55.22)   \\
  &  [a_{\lambda}(\mathbf{k}), a_{\lambda'}^{\dag}(\mathbf{k}') ] = (2\pi)^3 2\omega \delta^3(\mathbf{k}' - \mathbf{k})\delta_{\lambda \lambda'} \quad \, (55.23)  
\end{aligned}
\]

We want to show (56.9), (56.10).  

\[
\begin{gathered}
  \langle 0 | T A^i(x) A^j(y) | 0 \rangle = \\
  = \sum_{\lambda, \lambda' = \pm } \int \widetilde{dk} \widetilde{dk}' \langle 0 | T [ \epsilon^{i*}_{\lambda}(\mathbf{k}) e^{ikx}a_{\lambda}(\mathbf{k}) + \epsilon^i_{\lambda}(\mathbf{k}) a_{\lambda}^{\dag}(\mathbf{k}) e^{-ikx} ] [ \epsilon^{j*}_{\lambda'}(\mathbf{k}') e^{ik'y}a_{\lambda'}(\mathbf{k}) + \epsilon^j_{\lambda'}(\mathbf{k}') a_{\lambda'}^{\dag}(\mathbf{k}') e^{-ik'y} ] |0 \rangle
\end{gathered}
\]

Note that $a |0 \rangle = 0$ and $\langle 0 | a^{\dag} =0$.   Then we obtain

\[
\begin{gathered}
  \langle 0 | T A^i(x) A^j(y) | 0 \rangle =  \\
 = \sum_{ \lambda, \lambda' = \pm} \int \widetilde{dk}\widetilde{dk'} \langle 0 | T \left( \epsilon_{\lambda}^{i*}(\mathbf{k}) \epsilon_{\lambda'}^j(\mathbf{k}') a_{\lambda}(\mathbf{k}) a_{\lambda'}^{\dag}(\mathbf{k}') e^{ikx} e^{-ik'y}  + \epsilon_{\lambda}^{i}(\mathbf{k}) \epsilon_{\lambda'}^{*j}(\mathbf{k}') a_{\lambda}^{\dag}(\mathbf{k}) a_{\lambda'}(\mathbf{k}') e^{-ikx} e^{ik'y} \right) |0\rangle 
\end{gathered}
\]

Using the commutator relation,
\[
\begin{gathered}
  \langle 0 | T A^i(x) A^j(y) | 0 \rangle =  \\  
  \begin{aligned}
    & = \sum_{\lambda, \lambda' = \pm} \int \widetilde{dk} \widetilde{dk}' \langle 0 | T(\epsilon^{i*}_{\lambda}(\mathbf{k}) \epsilon_{\lambda'}^j(\mathbf{k}') (2\pi)^3 2\omega \delta^3(\mathbf{k}' - \mathbf{k}) \delta_{\lambda \lambda'} e^{ikx} e^{-ik'y} ) | 0 \rangle  = \\
    & = \sum_{\lambda = \pm} \int \widetilde{dk} \langle 0 | T(\epsilon^{i*}_{\lambda} (\mathbf{k}) \epsilon_{\lambda}^j(\mathbf{k}) e^{i \mathbf{k}\cdot (\mathbf{x} - \mathbf{y}) } e^{-ik^0 x^0 + i k^0 y^0} ) |0 \rangle  =  \int \widetilde{dk} \langle 0 | T e^{i\mathbf{k}\cdot(\mathbf{x} - \mathbf{y})} e^{-ik^0 x^0 + i k^0 y^0} |0\rangle \sum_{\lambda = \pm} \epsilon^{i*}_{\lambda}(\mathbf{k}) \epsilon_{\lambda}^j (\mathbf{k})
\end{aligned}
\end{gathered}
\]

Note $\sum_{\lambda = \pm} \epsilon_{i \lambda}^*(\mathbf{k}) \epsilon_{j\lambda}(\mathbf{k}) = \delta_{ij} - \frac{k_i k_j}{\mathbf{k}^2}$ \quad \, (55.15)

Recall the trick of doing the contour integral over $k^0$ (this was utilized in Sec. 8 of Srednicki, and Problem 8.4).  \\

$k^2 = -(k^{0})^{2} + \mathbf{k}^2$.  

\[
\int \frac{d^4 k}{ (2\pi)^4} \frac{e^{ik(x-y)} }{k^2 - i\epsilon} = \int \frac{d^3k}{(2\pi)^4} \int_{-\infty}^{\infty} dk^0 \frac{ e^{i\mathbf{k} \cdot (\mathbf{x}- \mathbf{y} ) } e^{-ik^0 (x^0 - y^0) } }{ \mathbf{k}^2 - i\epsilon - (k^0)^2 }
\]

Consider
\[
\int_{-\infty}^{\infty} dz \frac{e^{-iz \tau}}{ A^2 - z^2} = \int_{-\infty}^{\infty} dz \frac{ e^{-iz\tau} }{ (A-z)(A+z) } = -\int_{\infty}^{\infty} \frac{ e^{-iz \tau} }{ (z+ A)(z-A) }
\]

For $z = Re^{i\theta }$, $-iz \tau = -i Re^{i\theta} \tau = - i \tau R\cos{\theta} + \tau R\sin{\theta}$.  

If $\tau >0$, consider enclosing in lower half plane of $z$ for a finite integral: $Res{z} = A$, where \\
$A^2 = \mathbf{k}^2 - i \epsilon = Re^{i\theta + i2\pi n }$ \\ 
$A = \pm \sqrt{ \mathbf{k}^2 -  i\epsilon} = \sqrt{R} e^{ i \frac{\theta}{2} + i \pi n }$.  

\[
-\int_{ -\infty}^{\infty} dz \frac{e^{-iz \tau}}{ (z+A)(z-A) }  = - \left( - 2\pi i \frac{e^{-i A\tau} }{2A} \right) = \frac{ + \pi i e^{-iA \tau}}{A}
\]

Note the \emph{direction of the contour (clockwise) and winding number}, $-1$, so that in a sense, the contour integral went around by $-2\pi$.

If $\tau < 0$, enclose in upper half plane: 
\[
-\int_{-\infty}^{\infty} dz \frac{e^{-iz\tau} }{ z^2 - A^2 } = - \left( 2\pi i \frac{ e^{ iA \tau}}{ -2A } \right) = \frac{ \pi i e^{iA \tau} }{ A }
\]

For $\epsilon \to 0$, $A = \sqrt{ \mathbf{k}^2} = \omega$, ($\omega^2 = \mathbf{k}^2$ for photons),

\[
\int_{-\infty}^{\infty} dk^0 \frac{e^{-ik^0(x^0- y^0)}}{ \mathbf{k}^2 - i\epsilon - (k^0)^2 } = \theta(x^0 - y^0) \left( \frac{ + \pi i e^{-i \omega (x^0 - y^0)} }{\omega} \right) + \theta(y^0 - x^0) \frac{ \pi i e^{i \omega(x^0 - y^0 ) } }{ \omega} 
\]

\[
\int \frac{ d^4k}{ (2\pi)^4} \frac{ e^{ik(x-y) } }{ k^2 - i \epsilon } = \int \frac{d^3k}{(2\pi)^4} e^{i \mathbf{k} \cdot( \mathbf{x} - \mathbf{y} ) } \frac{ \pi i }{\omega} e^{-i\omega |x^0 - y^0 | } = i \int \widetilde{dk} e^{ i \mathbf{k}\cdot (\mathbf{x} - \mathbf{y} ) } e^{-i \omega |x^0 - y^0 | }
\]

Then 
\[
\boxed{ \langle 0 | T A^i(x) A^j(y) |0\rangle = \frac{1}{i} \Delta^{ij}(x-y)  }\quad \, (56.9) 
\]
where $\boxed{ \Delta^{ij}(x-y) = \int \frac{d^4k}{(2\pi)^4} \frac{e^{ik(x-y)}}{ k^2 - i \epsilon} \sum_{\lambda =\pm } \epsilon_{\lambda}^{i*}(\mathbf{k}) \epsilon_{\lambda}^j(\mathbf{k})}$  

\section{The path integral for photons}  

\[
\begin{gathered}
  (k^{\mu} A^{\nu} - k^{\nu} A^{\mu})( k_{\mu} A_{\nu} - k_{\nu} A_{\mu} ) = k^2 A^{\nu} A_{\nu} - k^{\mu} A_{\mu} k^{\nu} A^{\nu} - k^{\nu} A_{\nu} k^{\mu} A_{\mu} + k^2 A^{\mu} A_{\mu}
\end{gathered}
\] 

$P$  defined :

\begin{align}
  & k^2 g^{\mu \nu} - k^{\mu} k^{\nu} = k^2 P^{\mu \nu} \quad \quad \quad \, (57.4) \\ 
  & P^{\mu \nu}(k) \equiv g^{\mu \nu} - \frac{ k^{\mu} k^{\nu} }{ k^2} \quad \quad \quad \, (57.5) 
\end{align}

This is a projection matrix because

\begin{equation}
  P^{\mu \nu}(k) P_{\nu}^{\lambda}(k) = P^{\mu \nu}(k) \quad \quad \quad \, (57.6)
\end{equation}


\[
\begin{gathered}
  (g^{\mu \nu} - \frac{k^{\mu } k^{\nu} }{k^2 } ) ( g_{\nu}^{ \, \, \lambda} - \frac{k_{\nu} k^{\lambda} }{ k^2 } ) = g^{\mu \nu } - \frac{ 2 k^{\mu } k^{\nu }}{ k^2 } + \frac{k^{\mu} k^{\lambda} }{k^2 } = P^{\mu \nu}
\end{gathered}
\]

$P^2 = 1$ \\
$P^2 x = P ( Px) = Px$ \\
$P\psi = \psi$  \\
so eigenvalues are $0$ or $1$ \\

$P^{\mu \nu}(k)k^{\nu} = 0$, so at least 1 $0$ eigenvalue.  

\begin{equation}
g_{\mu \nu} P^{\mu \nu} = 3 = d-1 \quad \quad \quad \, (57.8)
\end{equation}

$g_{\mu \nu} P^{\mu \nu} = P^{\mu}_{ \, \, \mu} =\text{Tr}{P} = d- 1$

With

\begin{equation}
  Z_0(J) = \int \mathcal{D}A \exp{ [ i S_0 ] } \quad \quad \quad \, (57.1)
\end{equation}

\begin{equation}
S_0 = \frac{1}{2} \int \frac{d^dk}{ (2\pi)^d} [ -A_{\mu}(k) (k^2 g^{\mu \nu} - k^{\mu } k^{\nu} ) A_{\nu}(-k) + J^{\mu}(k) A_{\mu}(-k) + J^{\mu}(-k) A_{\mu}(k) ] \quad \quad \quad \, (57.3)
\end{equation}

i.e. $A_{\mu} P^{\mu \nu} A_{\nu}$ \\
$Pk = P^{\mu \nu} k_{\nu} = 0$ so $A_{\mu}$ component (lies along $k_{\mu}$) doesn't contribute to quadratic term. \\

$\partial^{\mu} J_{\mu} =0 \Longrightarrow k^{\mu} J_{\mu} =0$.  $kJ = 0$.  $J$ component along $k$ doesn't contribute. 

Within subspace orthogonal to $k^{\mu}$, $P=1$ inverse of $k^2 P$ is $\frac{1}{k^2} P$

Sec. 8 (change of variables, $\chi$, integrate out a (co)factor, in $\int \mathcal{D}\chi \propto \int \mathcal{D}A$ integral)

\begin{gather}
Z_0(J) = \exp{  \left[  \frac{i}{2}  \right]} \\
\Delta^{\mu \nu}(x-y)    \quad \quad \quad \, (57.10)
\end{gather}

Lorenz gauge (also known as Landau gauge)





\section{Spinor electrodynamics }

\solutionhead{58.1} Recall the Lagrangian $\mathcal{L} = - \frac{1}{4} F^{\mu \nu} F_{\mu \nu} + i \overline{\Psi} {\not}{D} \Psi - m \overline{\Psi} \Psi$ where $D = \partial_{\mu} - i e A_{\mu}$.  

\[
i \overline{\Psi} \gamma^{\mu} ( \partial_{\mu} - i e A_{\mu} ) \Psi = i \overline{\Psi} \gamma^{\mu} \partial_{\mu} \Psi + e \overline{\Psi} \gamma^{\mu} A_{\mu} \Psi
\]

\[
P^{-1} \overline{\Psi} \gamma^{\mu} \Psi A_{\mu} P = \mathcal{P}^{\mu}_{ \nu} \overline{\Psi} \gamma^{\nu} \Psi P^{-1} A_{\mu} \Psi = \overline{\Psi} \gamma^{\nu} \Psi P^{-1} \mathcal{P}^{\mu}_{\nu} A_{\mu} P
\]
If $A_{\mu} \to \mathcal{P}^{\nu}_{\mu} A_{\nu}$ then $\mathcal{P}^{\mu}_{\nu} \mathcal{P}^{\nu}_{\mu} = 1$, resulting in $\overline{\Psi} \gamma^{\nu} \Psi A_{\nu}$
\[
\Longrightarrow \boxed{ P^{-1} A_{\mu} P = \mathcal{P}^{\nu}_{\, \mu} A_{\nu} }
\]

\[
T^{-1} \overline{\Psi} \gamma^{\mu} A_{\mu} \Psi T = T^{-1} \overline{\Psi} \gamma^{\mu} \Psi TT^{-1} A_{\mu} T = - \mathcal{T}^{\mu}_{\, \nu} \overline{\Psi} \gamma^{\nu} \Psi T^{-1} A_{\mu} T
\]
If $\boxed{ A_{\mu} \to - \mathcal{T}^{\nu}_{\, \mu} A_{\nu} }$, then $\mathcal{L}$ invariant.  

\[
\Longrightarrow T^{-1} A_{\mu} T = - \mathcal{T}^{\nu}_{ \, \mu } A_{\nu}
\]

$C^{-1} \overline{\Psi} \gamma^{\mu} A_{\mu} \Psi C = - \overline{\Psi} \gamma^{\mu} \Psi C^{-1} A_{\mu} C$.  
If $C^{-1} A_{\mu} C = - A_{\mu}$, then $\mathcal{L}$ invariant.  $\boxed{ C^{-1} A_{\mu} C= - A_{\mu} }$.  

\solutionhead{58.2} \emph{Furry's theorem}  

Recall the LSZ formula for photons,

\[
\begin{aligned}
  & a^{\dag}_{\lambda}(\mathbf{k})_{in} \to  i \epsilon_{\lambda}^{\mu *}(\mathbf{k}) \int d^4x e^{ +ikx} ( - \partial^2) A_{\mu}(x) \\
  & a_{\lambda}(\mathbf{k})_{out} \to i \epsilon_{\lambda}^{\mu}(\mathbf{k}) \int d^4x e^{-ikx} (- \partial^2) A_{\mu}(x)
\end{aligned}
\]

So scattering amplitudes involve field $A_{\mu}$ anytime with external photons.    

Charge conjugation is a symmetry in QED, so $C |0 \rangle =0$.  

\[
\langle 0  | A_{\mu} | 0 \rangle = \langle 0 | C^{-1} C A_{\mu} C^{-1} C | 0 \rangle = \langle 0 | C A_{\mu} A^{-1} | 0 \rangle = - \langle 0 | A_{\mu} | 0 \rangle \Longrightarrow \langle 0 | A_{\mu} | 0 \rangle 
\]

\[
\langle 0 | A_{\mu} A_{\nu} A_{\rho} | 0 \rangle = \langle 0 | C^{-1} C A_{\mu} C^{-1} C A_{\nu} C^{-1} C A_{\rho} C^{-1} C | 0 \rangle = (-1)^3 \langle 0 | A_{\mu} A_{\nu} A_{\rho} | 0 \rangle  \Longrightarrow \langle 0 | A_{\mu} A_{\nu} A_{\rho} | 0 \rangle  = 0
\]

By induction, scattering amplitude of an odd number of external photons is zero.  

\section{Scattering in spinor electrodynamics }

\solutionhead{59.2} Bhabha scattering, $e^+ e^- \to e^+ e^-$

\[
\begin{aligned}
  & \mathcal{T} = e^2 ( (\overline{u}' \gamma^{\mu} v' )( \overline{v} \gamma_{\mu} u ) \frac{1}{ -s - i \epsilon } - ( \overline{v} \gamma^{\mu} v')( \overline{u}' \gamma_{\mu} u ) \frac{1}{ - t - i \epsilon } )  \\ 
  & \overline{\mathcal{T}} = e^2 ( (\overline{v} \gamma^{\nu} u') ( \overline{u} \gamma_{\nu} v) \frac{1}{ -s - i \epsilon } - (\overline{v}' \gamma^{\nu} v)( \overline{u} \gamma_{\nu} u' ) \frac{1}{ -t -i \epsilon } ) \\
  & \begin{aligned} |\mathcal{T}|^2 & = e^4 \left( (\overline{u}' \gamma^{\mu} v')( \overline{v} \gamma_{\mu} u ) ( \overline{v}' \gamma^{\nu} u') (\overline{v}' \gamma^{\nu} u')( \overline{u} \gamma_{\nu} v) \frac{1}{s^2} - \frac{ (\overline{u}' \gamma^{\mu} v' )( \overline{v} \gamma_{\mu} u)( \overline{v}' \gamma^{\nu} v)( \overline{u} \gamma_{\nu} u') }{ st} - \right. \\
      & = \left.  - \frac{ (\overline{v} \gamma^{\mu} v')( \overline{u}' \gamma_{\mu} u) ( \overline{v}' \gamma^{\nu} u')( \overline{u} \gamma_{\nu} v) }{ ts} + \frac{ (\overline{v} \gamma^{\mu} v')( \overline{u}' \gamma_{\mu} u)( \overline{v}' \gamma^{\nu} v)( \overline{u} \gamma_{\nu} u' ) }{ t^2} \right) \\
      & = e^4 \left( tr( u'\overline{u}' \gamma^{\mu} v' \overline{v}' \gamma^{\nu} ) tr( u \overline{u} \gamma_{\nu} v\overline{v} \gamma_{\mu} ) \frac{1}{s^2} - \frac{ tr( u\overline{u} \gamma_{\nu} u' \overline{u}' \gamma^{\mu} v' \overline{v}' \gamma^{\nu} v\overline{v} \gamma_{\mu} )}{st} - \right. \\
      & \left. \phantom{=} - \frac{ tr( u \overline{u} \gamma_{\nu} v\overline{v} \gamma^{\mu} v'\overline{v}' \gamma^{\nu} u' \overline{u}' \gamma_{\mu} ) }{ ts} + \frac{ tr( u\overline{u} \gamma_{\nu} u' \overline{u}' \gamma_{\mu} ) tr(v\overline{v} \gamma^{\mu} v'\overline{v}' \gamma^{\nu} ) }{t^2} \right)
\end{aligned}
\end{aligned}
\]
Average over initial spins of electron and positron, sum over final spin states of electron-positron pair.  




\section{Spinor helicity for spinor electrodynamics}



\section{Scalar electrodynamics}

\begin{equation}
  D_{\mu} \equiv \partial_{\mu} - ie A_{\mu} \quad \quad \quad \, (61.5)
\end{equation}

\begin{equation}
  A^{\mu}(x) \to A^{\mu}(x) - \partial^{\mu} \Gamma(x) \quad \quad \quad \, (61.6)
\end{equation}

\begin{equation}
\mathcal{L} = -(D^{\mu} \varphi)^{\dag} D_{\mu} \varphi - m^2 \varphi^{\dag} \varphi - \frac{\lambda}{4} (\varphi^{\dag} \varphi)^2 - \frac{1}{4} F^{\mu \nu} F_{\mu \nu} \quad \quad \quad \,   (61.8) 
\end{equation}



\section{Loop corrections in spinor electrodynamics}

\solutionhead{62.1} 
\[
\begin{gathered}
  \frac{-1}{4} F^{\mu \nu} F_{\mu \nu} - \frac{1}{2} \frac{ ( \partial^{\mu} A_{\mu} )^2 }{ \xi } \\
  -\frac{1}{2} (\partial^{\mu} A^{\nu})( \partial_{\mu} A_{\nu} ) + \frac{1}{2} ( \partial^{\mu} A^{\nu} )(\partial_{\nu} A_{\mu} ) - \frac{1}{2} \frac{ (\partial^{\mu} A_{\mu})( \partial^{\nu} A_{\nu} ) }{ \xi }
\end{gathered}
\]

Apply Fourier transform to momentum space.  
\[
\Longrightarrow \frac{-1}{2} ( i k^{\mu} )( i k_{\mu}') \widetilde{A}^{\nu}(k) \widetilde{A}_{\nu}(k') + \frac{1}{2} i k^{\nu} i k_{\mu}' \widetilde{A}^{\mu}(k) \widetilde{A}_{\nu}(k') - \frac{1}{2 \xi} i k^{\mu} i k^{\nu '} \widetilde{A}_{\mu}(k) \widetilde{A}_{\nu}(k') 
\]

Applying $\int d^4x$ and noting that $\int \frac{d^4x}{ (2\pi)^4} e^{i (k+k') x} = \delta^4(k+k')$ and then applying $\int d^4k'$,
\[
\begin{gathered}
  \Longrightarrow \frac{-1}{2} \left( \widetilde{A}_{\mu}(k) (g^{ \mu \nu} k^2 - k^{\mu} k^{\nu} )\widetilde{A}_{\nu}(-k) + \widetilde{A}_{\mu}(k) \frac{k^{\mu} k^{\nu} }{\xi} \widetilde{A}_{\nu}(-k) \right)   = \frac{-1}{2} \left( \widetilde{A}_{\mu}(k) k^2 \left( g^{\mu \nu} - \frac{k^{\mu} k^{\nu} }{ k^2 } + \frac{1}{ \xi} \frac{k^{\mu } k^{\nu}}{ k^2} \right) \widetilde{A}_{\nu}(-k) \right)
\end{gathered}
\]



For $P^{\mu \nu}(k) = g^{\mu \nu} - \frac{k^{\mu} k^{\nu}}{k^2}$, $P^2=1$ (definition of a projection operator).  

Note $\frac{k^{\mu} k^{\nu}}{k^2} \left( \frac{k_{\nu} k_{\mu}}{k^2} \right) =1$ and $\left( g_{\mu \sigma}  - \frac{k_{\mu} k_{\sigma}}{k^2} \right) \frac{k^{\sigma} k^{\nu}}{ k^2} = \frac{k_{\mu} k^{\nu}}{k^2} - \frac{k_{\mu} k^{\nu}}{k^2 } = 0 $.  

So then $k^2 \left( g^{\mu \nu}  -\frac{k^{\mu } k^{\nu}}{k^2} + \frac{1}{ \xi} \frac{k^{\mu} k^{\nu}}{ k^2} \right)$ can be inverted to yield the photon propagator in $R_{\xi}$ gauge.  

So
\[
\widetilde{\Delta}_{\mu \nu} = \frac{1}{k^2 - i \epsilon } \left( g_{\mu \nu} - (1- \xi) \frac{k_{\mu} k_{\nu}}{k^2 } \right) - \frac{1}{ k^2 - i\epsilon } \left( P_{\mu \nu}(k) + \xi \frac{k_{\mu} k_{\nu}}{k^2 } \right) 
\]

When $\xi = 0$, $\widetilde{\Delta}_{\mu \nu}(k) = \frac{1}{ k^2 - i\epsilon } \left( g_{\mu \nu} - \frac{k_{\mu} k_{\nu}}{k^2 } \right) = \frac{P_{\mu \nu}(k) }{k^2 - i \epsilon }$.  

For $\xi =0$, $\partial^{\mu} A_{\mu} = 0$ for finite $\mathcal{L}$, in $\frac{ (\partial^{\mu} A_{\mu} )(\partial^{\nu} A_{\nu}) }{ \xi}$.  

\solutionhead{62.2}

\[
\Longrightarrow (\xi -1) e^2 \int dF_3 \int \frac{d^4 q}{ (2\pi)^4} \frac{q^2}{d} \left( -dm - {\not}{p} (d- 2 + x_1(d+2)) \right) \frac{1}{ (q^2 + D)^3}
\]

\[
\begin{gathered}
  \int \frac{d^4q}{(2\pi)^4} \frac{(q^2)^1}{ (q^2 + D)^3} \Longrightarrow  i\frac{\Gamma(3-1 -d/2) \Gamma(1 + d/2) }{ (4\pi)^{d/2} \Gamma(3) \Gamma(d/2) } D^{-(3-1 - d/2) } = 
    \frac{ i \left( \frac{2}{\epsilon} - \frac{1}{2} + \ln{ \frac{\mu}{D}} \right) }{ (4\pi)^2} 
\end{gathered}
\]
\[
\begin{gathered}
\Longrightarrow   \left( \xi - 1 \right) e^2 \left( \left( - m - {\not}{p} \left( 1  - \frac{2}{d} + x_1 \left( 1  + \frac{2}{d} \right) \right) \right) \frac{ i \left( \frac{2}{\epsilon }- \frac{1}{2} + \ln{ \frac{\mu }{ D} } \right) }{ (4\pi)^2 } \right) =  \\
  \xrightarrow{ d = 4 - \epsilon , \, \epsilon \to 0 } \left( \xi - 1 \right) e^2 \int dF_3 \left( - m - {\not}{p} \left( \frac{1}{2} + x_1 \left( \frac{3}{2} \right) \right) \right) \frac{ i \left( \frac{2}{ \epsilon } + \frac{1}{2} + \ln{ \frac{\mu}{D} } \right) }{ (4\pi)^2 } \\   = \left( \xi -1 \right) e^2 \left(  \left( - m - {\not}{p} \frac{1}{2} - {\not}{p} \frac{3}{2} \frac{1}{3} \right) i \frac{2}{ \epsilon (4\pi)^2 } + \mathcal{O}(1) \right) = \frac{ - i \left( \xi -1 \right) e^2}{(4\pi)^2 } \left( \frac{ 2 ( {\not}{p} + m ) }{\epsilon } + \mathcal{O}(1) \right)
\end{gathered}
\]

So then 
\[
\begin{aligned}
  Z_2 & = 1  -\frac{e^2}{8\pi^2} \frac{1}{ \epsilon} + \frac{ (1 - \xi ) e^2 }{ 8\pi^2 \epsilon } \\ 
  Z_m & = 1 - \frac{e^2}{ 2\pi^2 } \left( \frac{1}{ \epsilon } \right) + \frac{ (1 - \xi ) e^2}{ 8\pi^2 \epsilon }
\end{aligned}
\]
or
\[
\Longrightarrow \boxed{ \begin{aligned} 
    Z_2 & = 1 - \frac{e^2}{ 8 \pi^2 } \frac{1}{ \epsilon } \xi \\
    Z_m & = 1 - \frac{e^2}{ 2\pi^2} \frac{1}{ \epsilon } \left( \frac{3}{4} + \frac{\xi }{4} \right) \end{aligned} }
\]

\section{The vertex function in spinor electrodynamics}

\solutionhead{63.1}
\begin{itemize}
\item[(a)] 


So $q_{\mu} \overline{u}' V^{\mu}(k,p) u = 0$.  
\[
\Longrightarrow q_{\mu} \overline{u}' V^{\mu} u = e\overline{u}' [A {\not}{q} + B q_{\mu}( k +p)^{\mu} + C q^2 ] u 
\]
Note $\overline{u}' {\not}{q} u = \overline{u}' ({\not}{k} - {\not}{p} )u = - m \overline{u}' u + m \overline{u}'  u = 0$ and $(k-p) (k+p) = k^2 - p^2 = - m^2 + m^2 = 0$

\[
\Longrightarrow q_{\mu} \overline{u}' V^{\mu} u = e[ 0 + B \overline{u}' (k-p)(k+p) u + C q^2] u = 0 
\]
So $C=0$.  
\item[(b)] Now
\[
\overline{u}' V^{\mu}(k,p) u = e\overline{u}' [ F_1 \gamma^{\mu} - \frac{i}{m} F_2 S^{\mu \nu } q_{\nu} ]u = F_1 e \overline{u}' \gamma^{\mu} u - \frac{i}{m} eF_2 \overline{u}' S^{\mu \nu} q_{\nu} u 
\]
\emph{The trick is to \textbf{use the Gordon identity}!}
\[
\Longrightarrow \overline{u}' (p' + p)^{\mu} u = \overline{u}' [2m \gamma^{\mu} + 2i S^{\mu \nu} q_{\nu} ] u \text{ or }  \overline{u}' (p'+p)^{\mu} u  - 2m \overline{u}' \gamma^{\mu} u = 2i \overline{u}' S^{\mu \nu} q_{\nu} u 
\]

So then

\[
\overline{u}' V^{\mu}(k,p) u = F_1 e\overline{u}' \gamma^{\mu} u - \frac{1}{m} eF_2 \left( \frac{ \overline{u}' (p + p') u }{2} - m \overline{u}' \gamma^{\mu} u \right) = \overline{u}' (F_1 e + eF_2) \gamma^{\mu} u + \overline{u}' \left( \frac{ - eF_2}{ 2m } \right) ( p + p')^{\mu} u 
\]

\[
\Longrightarrow \boxed{ \begin{aligned} A & = F_1 + F_2  \\ B & = \frac{ -F_2}{ 2m } \end{aligned} }
\]


\end{itemize}


\section{The magnetic moment of the electron }




\section{Loop corrections in scalar electrodynamics}




\section{Beta functions in quantum electrodynamics}  

\solutionhead{66.1}

Note that 
\[
\begin{aligned}
  & Z_3^{1/2} A = A_0 \\ 
  & Z_2^{1/2} \Psi = \Psi_0 
\end{aligned}
\]
\[
\begin{aligned}
  &  \frac{1}{2} \ln{ Z_3 } + \ln{A} = \ln{ A_0 } \\ 
  &  \frac{1}{2} \ln{ Z_2} + \ln{ \Psi} = \ln{ \Psi_0 } 
\end{aligned} \quad \quad \quad \, 
\begin{aligned}
  &  \frac{1}{2} \ln{ \left( 1  - \frac{ 2\alpha}{ 3\pi } \left( \frac{1}{ \epsilon } - \ln{ \frac{m}{\mu } \ } \right) \right) } = \frac{1}{2} \ln{ Z_3} \\ 
  &  \frac{1}{2} \ln{ Z_2} = \frac{1}{2} \ln{ \left( 1 - \frac{ \alpha}{ 2\pi \epsilon } \xi \right) } \simeq \frac{ -\alpha}{4\pi \epsilon } \xi 
\end{aligned}
\]
\[
\frac{ d\ln{ \Psi_0 } }{ d\ln{ \mu }} = 0 = \frac{ d\ln{ \Psi_0 } }{ d \ln{ \mu }} = \frac{-1}{4\pi \epsilon} \xi \left( \frac{d\alpha}{ d\ln{ \mu} } \right) = \frac{ - \xi }{ 4\pi \epsilon } \left( -\epsilon \alpha  +\frac{2\alpha^2}{3\pi } \right)
\]

For $\ln{ Z_3}$, if the $\frac{1}{\epsilon}$ can be treated as the dominant term, even over $\ln{ \frac{m}{\mu} }$, 

\[
\frac{1}{2} \ln{ Z_3} \simeq \frac{1}{2} \left( \frac{ -2\alpha }{ 3\pi \epsilon } \right)  = \frac{ -\alpha }{ 3\pi \epsilon }
\]
\[
\frac{  - \alpha}{ 3\pi \epsilon }  + \ln{ A}  = \ln{ A_0 } \Longrightarrow \frac{d\ln{ A_0 }}{ d\ln{ \mu }} = \frac{ d\ln{ A} }{ d\ln{ \mu}  } + \frac{-1 }{ 3\pi \epsilon } \left( -\alpha \epsilon + \frac{2\alpha^2}{ 3\pi } \right) = 0 \text{ or } \frac{d\ln{ A}}{ d\ln{ \mu } } = \frac{1}{3\pi \epsilon} \left( - \alpha \epsilon + \frac{2\alpha^2}{ 3 \pi} \right) = \frac{ \alpha }{3\pi } \left( -1 + \frac{2\alpha }{ 3\pi \epsilon } \right)
\]


\solutionhead{66.3}

Recall that 
\[
\mathcal{L} = i \overline{\Psi} {\not}{\partial} \Psi - m \overline{ \Psi} \Psi - \frac{1}{4} F^{\mu \nu} F_{\mu \nu} + Z_1 e \overline{\Psi} {\not}{A} \Psi + i (Z_2 - 1) \overline{\Psi} {\not}{\partial} \Psi - (Z_m -1) m \overline{ \Psi} \Psi  - \frac{1}{4} ( Z_3 - 1 ) F^{\mu \nu} F_{\mu \nu }
\]

So then
\[
\begin{aligned}
  & Z_2^{1/2} \Psi  = \Psi_ 0 \\ 
  & Z_3^{1/2} A  = A_0 \\ 
  & Z_m Z_2^{-1} m  = m_0 
\end{aligned}
\]
\[
Z_1 Z_2^{-1} Z_3^{-1/2} \overline{\mu}^{\epsilon/2} e = e_0 
\]
Also note that 
\[
\begin{gathered}
  \alpha = \frac{e^2}{ 4 \pi } \quad \quad \quad \, \alpha_0 = \frac{e_0^2}{ 4\pi} \\ 
  \alpha_0 = Z_3^{-1} Z_2^{-2} Z_1^2 \overline{\mu}^{\epsilon} \alpha 
\end{gathered}
\]

\[
\begin{aligned}
  & Z_1 = 1 - \frac{e^2}{ 8 \pi^2 \epsilon } \xi = 1-  \frac{ \alpha }{2\pi \epsilon} \xi \\ 
  & Z_2 = 1 - \frac{ e^2 }{ 8 \pi^2 \epsilon } \xi = 1 - \frac{ \alpha}{ 2\pi \epsilon} \xi \\ 
  & Z_3 = 1 - \frac{e^2}{6\pi^2} \left[ \frac{1}{ \epsilon } - \ln{ \frac{m }{\mu}} \right] = 1 - \frac{ 2\alpha}{ 3\pi} \left[ \frac{1}{ \epsilon } - \ln{ \frac{m}{\mu} } \right] \\ 
  & Z_m = 1 - \frac{ e^2}{ 2\pi^2 \epsilon } \left( \frac{3}{4} + \frac{ \epsilon }{4} \right) = 1 - \frac{ \alpha }{ 2\pi \epsilon }( 3 + \xi ) 
\end{aligned}
\]

For $\ln{ ( Z_3^{-1} Z_2^{-2} Z_1^2 ) } = \sum_{n=1}^{\infty} \frac{ E_n(\alpha )}{ \epsilon^n }$, 
\[
\ln{ (Z_3^{-1} Z_2^{-2} Z_1^2) } = 2 \left( \frac{ -\alpha}{ 2\pi \epsilon } \xi + \dots \right) - 2  \left( \frac{-\alpha }{2\pi \epsilon} \xi  + \dots \right)  - \left( \frac{ -2 \alpha}{ 3\pi \epsilon } + \dots \right) = \frac{2\alpha }{3\pi \epsilon } + \dots 
\]

\[
\ln{ \alpha_0 } = \epsilon \ln{ \widetilde{\mu} } + \ln{ \alpha} + \frac{ 2\alpha }{ 3\pi \epsilon } \text{ so } E_1(\alpha) = \frac{ 2\alpha}{ 3\pi}
\]

\[
\begin{aligned}
  & \beta(\alpha) = \alpha^2  E_1'(\alpha) + \mathcal{O}(\alpha^3) \\ 
  & \boxed{ \beta(\alpha) = \frac{2 \alpha^2}{3\pi } + \mathcal{O}(\alpha^3) } \text{ or } \beta(e) = \frac{e^4}{ 24 \pi^3 } + \mathcal{O}(e^5) 
\end{aligned}
\]

\[
\begin{gathered}
  Z_m Z_2^{-1} m = m_0 \\ 
  \begin{aligned}
    & \ln{ Z_m Z_2^{-1} } + \ln{ m }  = \ln{  m_0 } \\ 
    & \ln{ Z_m Z_2^{-1}  } =  \ln{ \left( 1 - \frac{ \alpha}{ 2\pi \epsilon} ( 3 + \xi) \right) } - \ln{ \left( 1 - \frac{ \alpha}{2\pi \epsilon} \xi \right) } \simeq \frac{ - \alpha}{ 2\pi \epsilon } (3+ \xi) - \left( \frac{ -\alpha}{ 2\pi \epsilon } \xi \right)  = \frac{-3\alpha }{ 2\pi \epsilon }
  \end{aligned}
\end{gathered}
\]

Require $m_0$ to be unchanged.  

\[
 0 = \frac{ d\ln{ m_0} }{ d\ln{ \mu}} = \frac{1}{m } \frac{ dm}{ d\ln{ \mu }} +  \frac{ -3}{2\pi \epsilon } \frac{ d\alpha}{ d\ln{ \mu} } \text{ or } \frac{1}{m } \frac{ dm}{ d\ln{ \mu }} = \frac{ 3 }{2\pi \epsilon } \left( - \alpha \epsilon + \frac{2\alpha^2}{ 3\pi } \right)  = \frac{ 3\alpha }{ 2\pi } \left( -1 + \frac{ 2\alpha^2}{ 3\pi \epsilon } \right) 
\]

where $\frac{d\ln{ \alpha_0}}{ d\ln{ \mu} } = \epsilon + \frac{ d\ln{ \alpha} }{ d\ln{ \mu} } + \frac{ 2 }{ 3\pi \epsilon } \frac{d\alpha}{ d\ln{ \mu }} = \epsilon + \left( \frac{ 1}{ \alpha} + \frac{ 2 }{ 3\pi \epsilon } \right) \frac{ d\alpha}{ d\ln{ \mu} } = 0 \text{ or } \frac{d\alpha}{ d\ln{\mu }} = \frac{ - \epsilon}{ \frac{1}{ \alpha} + \frac{2}{ 3\pi \epsilon } } = -\alpha \epsilon + \frac{2 \alpha^2}{ 3\pi } $ (use long division for the last statement).  

\[
\boxed{ \gamma_m(\alpha) = \frac{ -3\alpha}{ 2\pi} + \mathcal{O}(\alpha^2) }
\]


\solutionhead{66.4} $M_W = 80.4 \, GeV$.  Note that the top quark mass is $174 \, GeV$, greater than the mass of the $W$ boson; the top quark is not included in the sum.  

For each quark, there is a factor of 3 for the 3 color factors.  

Neutrinos have charge 0 and so do not contribute to the sum.  

\[
\begin{gathered}
\frac{1}{ \alpha (M_W) } = \frac{1}{ \alpha( \mu) } - \frac{2}{ 3\pi } \left( (2/3)^2 \ln{ \left( \frac{ 80400}{300 } \right) }\cdot 3 + (-1/3)^2 \ln{ \left( \frac{ 80400 }{ 300 } \right)} \cdot 3 \cdot 2 + (2/3)^2 \ln{ \left( \frac{ 80.4}{1.3}  \right)} \cdot 3 + \right. \\
 \left. + (-1/3)^2 (\ln{ ( 80.4/4.3 ) })\cdot 3 + (-1)^2 \ln{ 80400/105.7 } + (-1)^2 \ln{ ( 80400/1777 )} + (-1)^2 \ln{ ( 80400/0.511) }  \right)= 128.5
\end{gathered}
\]
so that 
\[
\alpha(M_W) = 0.007782
\]




\section{Ward identities in quantum electrodynamics I}




\solutionhead{67.2}

Consider $\overline{v}' \left[ {\not}{\epsilon}' \left( \frac{ - {\not}{p} + {\not}{k} + m }{ m^2  - t } \right)  {\not}{\epsilon } + {\not}{\epsilon }  \left(  \frac{ - {\not}{p } + {\not}{l } + m }{ m^2 - u } \right) {\not}{\epsilon}' \right] u$.  

Now replace ${\not}{\epsilon}$ with ${\not}{k}$.  

\[
\overline{v}' \left[ {\not}{\epsilon}' \left( \frac{ - {\not}{p} + {\not}{k} + m }{ m^2  - t } \right)  {\not}{k} + {\not}{k }  \left(  \frac{ - {\not}{p } + {\not}{l } + m }{ m^2 - u } \right) {\not}{\epsilon}' \right] u
\]

Consider
\[
  \overline{v}' {\not}{\epsilon}' \left(  - {\not}{p} + {\not}{k} + m  \right)  {\not}{k} u = \overline{v}' ( - {\not}{\epsilon}' {\not}{p} {\not}{k} + {\not}{\epsilon}' {\not}{k} {\not}{k} + m {\not}{\epsilon}' {\not}{k} ) u 
\]

Now (\emph{trick is})
\[
{\not}{p} {\not}{k} = p_{\mu} k_{\mu} \gamma^{\mu} \gamma^{\nu} = p_{\mu} k_{\nu} (- \gamma^{\nu} \gamma^{\mu} - 2 g^{\mu \nu} ) = - {\not}{k} {\not}{p} - 2(pk)
\]

\[
\text{ so } - {\not}{\epsilon}' {\not}{p} {\not}{k} = {\not}{\epsilon}' ( {\not}{k} {\not}{p} +  2(pk) )
\]

Also, (\emph{the trick is})

\[
{\not}{k}  {\not}{k} = k_{\mu} k_{\nu} \gamma^{\mu} \gamma^{\nu} = k_{\mu} k_{\nu} ( -\gamma^{\nu} \gamma^{\mu} - 2g^{\mu \nu } ) = - {\not}{k}  {\not}{k}  - 2k^2 = - {\not}{k}  {\not}{k}  \quad \quad \, (\text{ since } k^2 = 0)
\]
\[
\Longrightarrow {\not}{k}  {\not}{k}  = 0
\]

Also $ {\not}{k}  {\not}{p} u = - m {\not}{k} u$ since $({\not}{p} + m  ) u = 0$

\[
\Longrightarrow (- {\not}{p} + {\not}{k} + m ) {\not}{k} u = ( {\not}{k} {\not}{p} + 2(pk) + m {\not}{k} ) u = 2(pk) u 
\]

For the ${\not}{k}(- {\not}{p} + {\not}{l} + m ) {\not}{\epsilon}' u$ term, 

By momentum conservation ($p + q = k + l$), 

\[
{\not}{k} ( - {\not}{p}  + {\not}{l } + m ) = {\not}{k} ( \not{q} - \not{k} + m)  = - {\not}{q} {\not}{k} - 2(qk) - {\not}{k} {\not}{k} + m {\not}{k} = - {\not}{q} {\not}{k} - 2(qk) + m {\not}{k} 
\]

Since $ \overline{v} {\not}{q} = m \overline{v}{\not}{q}$,  so applying $\overline{v}$ from the left, $- m {\not}{k} - 2(qk) + m {\not}{k} = -2 (qk)$.  

\[
\Longrightarrow \overline{v}' \left[ {\not}{\epsilon}' \frac{ 2(pk) }{ m^2 - t} + \frac{ -2(qk) }{ m^2 -u } {\not}{\epsilon} ' \right] u 
\]

Note that 
\[
\begin{aligned}
  & (pk) ( m^2 - u)  = (pk) ( m^2 + (p-l)^2 ) = (pk) ( m^2 + -m^2  - 2pl ) = -2(pk)( pl) \\ 
  & (qk)( m^2  -t) = (qk) (m^2 + (p-k)^2 ) = (qk) ( m^2 + -m^2 - 2pk) = -2(qk) ( pk) 
\end{aligned}
\]
\[
(p-l)^2 = -m^2 - 2pl = (q-k)^2 = - m^2 - 2qk \text{ so } pl = qk 
\]

\[
\Longrightarrow \boxed{ \overline{v}' \left[ {\not}{\epsilon}' \left( \frac{ - {\not}{p} + {\not}{k} + m }{ m^2  - t } \right)  {\not}{k} + {\not}{k }  \left(  \frac{ - {\not}{p } + {\not}{l } + m }{ m^2 - u } \right) {\not}{\epsilon}' \right] u =0  }
\]



\section{Ward identities in quantum electrodynamics II}


\section{Nonabelian gauge theory }\label{Sec:NonabelianGaugeTheory}

\subsection*{Notes on the material in section \ref{Sec:NonabelianGaugeTheory}}

20090417

Prof. Gukov says:

\subsection*{ \underline{Proof of $Z_1 = Z_2$ } }

\[
\mathcal{L}_{QED} = \dots  + \underbrace{ i Z_{2}  \overline{ \Psi} \not \partial \Psi  + Z_1 e \overline{ \Psi} \not A \Psi }_{ i Z_2 \overline{ \Psi } \not D \Psi  \text{ iff } Z_1 = Z_2 } + \dots 
\]

(sic the proof is) in book.  

\subsection*{ \underline{ Yang-Mills theory} }

\[
\begin{aligned}
G = U(1) \to & G = \overbrace{ U(N), SU(N), SO(N), Sp(2N) }^{ \text{ classical groups (continuous ones) } }
& \underbrace{ G_2, F_4, E_6, E_7, E_8 }_{ \text{ rank 2 group } \quad \, \text{ exceptional - don't belong to infinite families } }
\end{aligned}
\]


$\mathfrak{g} = $ Lie $(G)$ generated by $T_a$, $a=1, \dots, \text{dim}{(G)}$ \, \, set of matrices or in general, elements

$U \in G$, $U = \exp{ [ -i \Gamma^a T^a ] }$
\[
[T^a, T^b] = i f^{abc} T^c
\]
$f^{abc}$ real for Lagrangian to be real, real structure constants.  

\underline{normalization} $\text{Tr}{ (T^a T^b) } = \frac{1}{2} \delta^{ab}$ metric or Killing form.  

\underline{Ex}: $G = SU(2) \simeq S^2$  (isomorphic to $S$ sphere).  $\text{dim}{ (G) } = 3$.  3 independent generators.  \\
$T^a = \frac{1}{2} \sigma^a$.  $f^{abc} = \epsilon^{abc}$ totally antisymmetric.  




\solutionhead{69.1} Recall the expansion in terms of traceless, Hermitian generators 
\[
A_{\mu }(x) = A^a_{\mu}(x) T^a
\]
and the gauge transformation
\[
A_{\mu}(x) \to U(x) A_{\mu}(x) U^{\dag}(x) + \frac{i}{g} U(x) \partial_{\mu} U^{\dag}(x)
\]
Recall the infinitesimal expansion of the unitary transformation in this adjoint representation, to first order:
\[
\begin{aligned}
  & U_{jk} = \delta_{jk} - ig \theta^a(x) T^a_{jk}  \\ 
  & U_{jk}^{\dag} = \delta_{jk} + ig \theta^a T^a_{jk}
\end{aligned}
\]
Then
\[
\begin{gathered}
  UA_{\mu}U^{\dag} = (\delta_{jk} - ig\theta^a T^a_{ik} )A_{\mu}^b T^b_{kl} ( \delta_{lj} + ig \theta^c T^c_{lj} ) + \frac{i}{g} (\delta_{ik} - ig \theta^a T^a_{ik} )( ig (\partial_{\mu} \theta^b) T^b_{kj} ) = \\ 
  = (A_{\mu}^b T^b_{il} - ig \theta^a A^b_{\mu} (T^a T^b)_{il } )( \delta_{lj} + ig \theta^c T^c_{lj} ) + \frac{i}{g} ( ig (\partial_{\mu} \theta^b ) T^b_{ij} )  = A_{\mu}^b T^b_{ij} - ig \theta^a A_{\mu}^b (T^a T^b)_{ij} + ig \theta^c A_{\mu}^b T^b_{il} T^c_{lj} - (\partial_{\mu} \theta^b) T^b_{ij} = \\
  = (A_{\mu}^b - \partial_{\mu} \theta^b ) T^b_{ij} + i g \theta^a A^b_{\mu} (T^bT^a - T^a T^b)_{ij} = (A_{\mu} - \partial_{\mu} \theta^a ) T^a_{ij} - ig \theta^a A^b_{\mu} i f^{abc} T_{ij}^c = \\
  = (A_{\mu}^a - \partial_{\mu} \theta^a) T^a_{ij} + g \theta^b A_{\mu}^c f^{bca} T^a_{ij} = (A_{\mu}^a - \partial_{\mu} \theta^a + gf^{abc} \theta^b A_{\mu}^c ) T^a_{ij} = 
\end{gathered}
\]
Now for $UA_{\mu}U^{\dag} + \frac{i}{g} U \partial_{\mu} U^{dag}$, 
\[
\begin{gathered}
  (1 - ig \theta^a T^a) A_{\mu} ( 1 + ig \theta^b T^b ) + \frac{i}{g} ( 1 - ig \theta^a T^a ) ( ig (\partial_{\mu} \theta^b) T^b) = A_{\mu} - ig \theta^a T^a A_{\mu} + A_{\mu} i g \theta^b T^b + - (\partial_{\mu} \theta^b )T^b = \\
   = A_{\mu} +  ig \theta^a ( A_{\mu} T^a - T^a A_{\mu} ) - ( \partial_{\mu} \theta^a ) T^a 
\end{gathered}
\]




\solutionhead{69.2} 
\[
[T^a T^a, T^b ] = T^a [T^a,T^b] + [T^a,T^b] T^a = T^a i f^{abc} T^c + i f^{abc} T^c T^a = i f^{abc} T^a T^c + i f^{cba} T^a T^c  = i f^{abc} T^a T^c (1-1) = 0 
\]




\section{Group representations}

\solutionhead{70.1} 

Recall the index $T(R)$, such that $tr(T^a T^b) = T(R) \delta^{ab}$.  

Recall that matrix $T^a T^a$ commutes with every generator, and so must be a number times the identity.  

\[
T^a T^a = C(R) 1
\]

\[
\Longrightarrow tr(T^a T^a) = C(R) D(R)
\]
where $D(R) = $ dimension of representation $R$.  

The fundamental representation of a group is spanned by $D(A)$ generator matrices.  

$D(A) = $ number of generators of the group.

$tr(T^a T^a) = T(R) \delta^{aa} = D(A) T(R)$.  

\[
\Longrightarrow \boxed{ D(A) T(R) = C(R) D(R) }
\]





\solutionhead{70.2}
\begin{itemize}
\item[(a)] For $SU(N)$, $N \otimes \overline{N} = 1 \oplus A$.  

\[
T(N \otimes \overline{N} )  = T(N) D(\overline{N}) + D(N) T(\overline{N}) = \frac{1}{2} N  + \frac{1}{2} N = N 
\]
\[
T(1 \oplus A) = T(1) + T(A) = \boxed{ T(A) = N }
\]
where $\boxed{ T(1) = 0 } $, the index of a singlet is 0.  
\item[(b)] From the definition of an index of a representation, $tr(T^a_R T^b_R) = T(R) \delta^{ab}$, 
\[
tr(T_A^a T_A^b) = T(A) \delta^{ab} = (T^a_A)^{cd} (T^b_A)^{dc} = -i \epsilon^{acd} (-i \epsilon^{bdc} ) = - \epsilon^{cda} \epsilon^{dcb} = \epsilon^{cda} \epsilon^{cdb} = \delta^{dd} \delta^{ab} - \delta^{db} \delta^{ad} = \delta^{dd} \delta^{ab} = 2 \delta^{ab}
\]
$T(A)=2$ for $SU(2)$.  

Also, another way to do this is to consider the adjoint representation explicitly.  
\[
\begin{aligned}
  T^{1} & = - i \left( \begin{matrix} & &  \\ & & 1 \\ & -1 & \end{matrix} \right) = i \left( \begin{matrix} & & \\ & & -1 \\ & 1 & \end{matrix} \right) \\
  T^2 & = -i \left( \begin{matrix} & & -1 \\ & & \\ 1 & & \end{matrix} \right) = i \left( \begin{matrix}  & & 1 \\ & & \\  -1 & & \end{matrix} \right) \\ 
  T^3 & = -i \left( \begin{matrix} & 1 & \\ -1 & & \\ & & \end{matrix} \right) = i \left( \begin{matrix} & -1 & \\ 1 & & \\ & & \end{matrix} \right) 
\end{aligned}
\]

$T^1 T^1 = -1 \left( \begin{matrix} & & \\ & & -1 \\ & 1 & \end{matrix} \right) \left( \begin{matrix} & & \\ & & -1 \\ & 1 & \end{matrix} \right) = -1 \left( \begin{matrix} & & \\ & -1 & \\ & & -1 \end{matrix} \right) = 2$
\item[(c)] Given that $N$ of $SU(N)$ transforms as $2 \oplus (N-2)1s$ that means that 
\[
\Longrightarrow \boxed{ N = 2 \oplus \underbrace{1 \oplus \dots \oplus 1 }_{N-2} }
\]
Recall $N \otimes \overline{N} = 1 \oplus A$ (70.26).  

\[
\overline{N}  = 2 \oplus \underbrace{1 \oplus \dots \oplus 1 }_{N-2} \text{ since } \overline{2} = 2
\]

\[
N \otimes \overline{N} = ( 2 \oplus  \underbrace{1 \oplus \dots \oplus 1 }_{N-2} ) \otimes ( 2 \oplus  \underbrace{1 \oplus \dots \oplus 1 }_{N-2} ) = 2 \otimes 2 \oplus \underbrace{ 2 \oplus \dots \oplus 2 }_{2(N-2)} \oplus \underbrace{ 1 \oplus \dots \oplus 1 }_{ (N-2)^2 }
\]

Now 
\[
2 \otimes 2 = 3_S \oplus 1_A
\]
\[
\Longrightarrow N \otimes \overline{N} = 3 \oplus \underbrace{ 2 \oplus \dots \oplus 2 }_{ 2 (N-2)} \oplus \underbrace{ 1 \oplus \dots \oplus 1 }_{ (N-2)^2 + 1 }
\]

So 
\[
\boxed{ A = 3 \oplus \underbrace{ 2 \oplus \dots \oplus 2}_{2 (N-2) } \oplus \underbrace{ 1 \oplus \dots \oplus 1}_{(N-2)^2 } }
\]

\item[(d)] Recalling $T(R_1 \otimes R_2) = T(R_1) D(R_2) + D(R_1) T(R_2)$,
\[
T(2 \otimes 2) = T(2) D(2) + D(2) T(2) = \frac{1}{2} 2 + \frac{1}{2} 2 =2 
\]
Now
\[
\begin{aligned}
  & A = 3 \oplus \underbrace{2 \oplus \dots \oplus 2}_{2 (N-2) } \oplus \underbrace{ 1 \oplus \dots \oplus 1 }_{(N-2)^2 } \\ 
  & A = 2 \otimes 2 \oplus \underbrace{ 2 \oplus  \dots \oplus 2}_{2 (N - 2) } \oplus \underbrace{ 1 \oplus \dots \oplus 1 }_{ (N-2)^2 -1 }
\end{aligned}
\]
Recall that $T(2) = \frac{1}{2}$ for the fundamental representation of $SU(N)$.  

\[
T(A) = T(2\otimes 2) + T(2)  2 (N-2) + T(1) ((N-2)^2 - 1) = 2 + \frac{1}{2} 2 (N-2) = \boxed{ N = T(A) }
\]

\end{itemize}



\solutionhead{70.6} 

Recalling that 
\[
(D_{\rho} F_{\mu \nu})^a = \partial_{\rho} F_{\mu \nu}^a - ig A_{\rho}^c (T^c_A )^{ab} F_{\mu \nu}^b
\]
and
\[
F^c_{\mu \nu} = \partial_{\mu } A_{\nu}^c - \partial_{\nu} A_{\mu}^c + gf^{abc} A_{\mu}^a A_{\nu}^b
\]


\[
\begin{aligned}
  & (D_{\mu} F_{\nu \rho})^a +   (D_{\nu} F_{\rho \mu})^a +   (D_{\rho} F_{\mu \nu})^a = \\
    & = \partial_{\mu} F_{\nu \rho}^a - ig A_{\mu}^c (T^c_A)^{ab} F^b_{\nu \rho } +  \\
    & \phantom{ ,}  + \partial_{\nu} F_{\rho \mu}^a - ig A_{\nu}^c (T^c_A)^{ab} F^b_{\rho \mu } + \\
    & \phantom{ ,}  + \partial_{\rho} F_{\mu \nu}^a - ig A_{\rho}^c (T^c_A)^{ab} F^b_{\mu \nu } =     \\
    & = \partial_{\mu} ( \partial_{\nu} A_{\rho}^a - \partial_{\rho } A_{\nu}^a + gf_{bca} A_{\nu}^b A_{\rho}^c  ) - igA_{\mu}^c (T_A^c)^{ab} ( \partial_{\nu} A_{\rho}^b - \partial_{\rho } A_{\nu}^b + gf^{deb} A_{\nu}^d A_{\rho}^e ) + \\
    & \phantom{,} +  \partial_{\nu} ( \partial_{\rho} A_{\mu}^a - \partial_{\mu } A_{\rho}^a + gf_{bca} A_{\rho}^b A_{\mu}^c  ) - igA_{\nu}^c (T_A^c)^{ab} ( \partial_{\rho} A_{\mu}^b - \partial_{\rho } A_{\rho}^b + gf^{deb} A_{\nu}^d A_{\rho}^e ) + \\
    & \phantom{,} +  \partial_{\rho} ( \partial_{\mu} A_{\nu}^a - \partial_{\nu } A_{\mu}^a + gf_{bca} A_{\mu}^b A_{\nu}^c  ) - igA_{\rho}^c (T_A^c)^{ab} ( \partial_{\mu} A_{\nu}^b - \partial_{\nu } A_{\mu}^b + gf^{deb} A_{\mu}^d A_{\nu}^e ) = \\
  & = gf^{bca} \partial_{\mu} ( A^b_{\nu} A^c_{\rho} ) - gA_{\mu}^c f^{cab} (\partial_{\nu} A_{\rho}^b - \partial_{\rho} A_{\nu}^b ) - g^2 A_{\mu}^d f^{cab} f^{deb} A_{\nu}^d A_{\rho}^e  + \\ 
  & = gf^{bca} \partial_{\nu} ( A^b_{\rho} A^c_{\mu} ) - gA_{\nu}^c f^{cab} (\partial_{\rho} A_{\mu}^b - \partial_{\mu} A_{\rho}^b ) - g^2 A_{\nu}^d f^{cab} f^{deb} A_{\rho}^d A_{\mu}^e  + \\ 
  & = gf^{bca} \partial_{\rho} ( A^b_{\mu} A^c_{\nu} ) - gA_{\rho}^c f^{cab} (\partial_{\mu} A_{\nu}^b - \partial_{\nu} A_{\mu}^b ) - g^2 A_{\rho}^d f^{cab} f^{deb} A_{\mu}^d A_{\nu}^e  
\end{aligned}
\]

Looking at the above equation, the $O(g)$ terms cancel.  

For $O(g^2)$ term, the trick is to use the Jacobi identity.  Making all those terms have a common factor of $A_{\mu}^c A_{\nu}^d A_{\rho}^e$ by rewriting the labels (indices).  

\[
f^{cab} f^{bde} + f^{dab} f^{bec} + f^{eab} f^{bcd} = 0 
\]



\section{The path integral for nonabelian gauge theory }




\section{The Feynman rules for nonabelian gauge theory}

\solutionhead{72.1}

For a renormalizable, Lorentz invariant theory with $N$ complex scalars, indexed by $i$, the free Lagrangian is 
\[
\mathcal{L} = - (D^{\mu} \varphi)^{\dag} ( D_{\mu} \varphi) -  m^2 \varphi^{\dag} \varphi - \frac{\lambda}{4} (\varphi^{\dag} \varphi)^2 
\]
where $\varphi^{\dag} \varphi = \varphi^{i \dag} \varphi_i$, $(\varphi^{\dag} \varphi)^2 = \varphi^{i \dag} \varphi_i \varphi^{j \dag} \varphi_j$ and $(D^{\mu} \varphi)^{\dag} (D_{\mu} \varphi) = (D^{\mu} \varphi^i)^{\dag}  (D_{\mu} \varphi_i )$.  

The covariant derivative is 
\[
D_{\mu} \varphi_i = \partial_{\mu} \varphi_i - ig A_{\mu}^a t^{aj}_i \varphi_j
\]
where the generators $t^{a}$ are for the gauge group.  

Now
\[
\begin{aligned}
  & (D^{\mu} \varphi^i)^{\dag} = \partial^{\mu} \varphi^{i \dag} + ig A^{\mu a} t_k^{ai} \varphi^{k \dag} 
\end{aligned} 
\]

so that 
\[
\begin{aligned} 
   (D^{\mu} \varphi^i)^{\dag} D_{\mu} \varphi_i & = ( \partial^{\mu} \varphi^{i \dag} + ig A^{\mu b} t_k^{bi} \varphi^{k\dag} ) ( \partial_{\mu} \varphi_i - ig A_{\mu}^a t^{aj}_i \varphi_j )  = \\
  & = \partial^{\mu} \varphi^{i \dag} \partial_{\mu} \varphi_i - ig A_{\mu}^a t_i^{aj} ( (\partial^{\mu} \varphi^{i \dag} ) \varphi_j - \varphi^{i \dag} \partial^{\mu} \varphi_j ) + g^2 A^{\mu b} A_{\mu}^a t_k^{bi} t_i^{aj} \varphi^{k \dag} \varphi_j
\end{aligned}
\]

So there is a term of order $g^2$, ($\mathcal{O}(g^2)$): $g^2 g_{\mu \nu} A^{\mu a} A^{\nu b} t_j^{ai} t_{i}^{bk} \varphi^{j \dag} \varphi_k = g^2 g_{\mu \nu} A^{\mu a} A^{\nu b} \varphi^{\dag} t^a t^b \varphi$.  

So for 
\[
\mathcal{L}_0 = -\partial^{\mu} \varphi^{ i \dag} \partial_{\mu} \varphi_i - m^2 \varphi^{\dag} \varphi 
\]
and $\mathcal{L}_1$, which includes counterterms, 
\[
\begin{aligned}
  \mathcal{L}_1 & = - ( Z_{\varphi} - 1) \partial^{\mu} \varphi^{ i \dag} \partial_{\mu} \varphi_i  - (Z_m - 1)  m^2 \varphi^{\dag} \varphi + Z_1 ig A_{\mu}^a t_i^{aj} ( (\partial^{\mu} \varphi^{i \dag} ) \varphi_j - \varphi^{i \dag} \partial^{\mu} \varphi_j ) - Z_2 g^2 g_{\mu \nu} A^{\mu a} A^{\nu b} t_k^{ai} t_i^{bj} \varphi^{k \dag} \varphi_j - \\
   & \phantom{ = } - Z_4 \frac{ \lambda}{4} ( \varphi^{i \dag} \varphi_i \varphi^{j \dag} \varphi_j )
\end{aligned}
\]

Consider doing a Fourier transform to transform the fields into momentum space (to deal with derivatives):

\[
\begin{aligned}
   (\partial^{\mu} \varphi^{ i \dag}) \varphi_j  & \to - i p \varphi^{ i \dag} \varphi_j \\ 
   \partial^{ i \dag} \partial_{\mu}  \varphi_j & \to iq \varphi^{ i \dag} \varphi_j 
\end{aligned}
\]

Then doing the Fourier transform on the following term,
\[
Z_1 ig A_{\mu}^a t_i^{aj} (( \partial^{\mu} \varphi^{i \dag} ) \varphi_j - \varphi^{i \dag} \partial_{\mu} \varphi_j ) \to Z_1 g A_{\mu}^a t_i^{aj} ( p + q)_{\mu} \varphi^{i \dag} \varphi_j
\]

\section{The beta function in nonabelian gauge theory }




\section{ BRST symmetry }



\section{Chiral gauge theories and anomalies }

\section{ Anomalies in global symmetries}

\section{ Anomalies and the path integral for fermions}

\section{ Background field gauge }

\section{ Gervais-Neveu gauge}

\section{ The Feynman rules for $N \times N$ matrix fields}

\begin{equation}
  \text{Tr}{ T^a T^b} = \delta^{ab} \quad \quad \quad (80.2) \quad a = 1 \dots N 
\end{equation}

\[
\text{Tr}{ T^a T^b} = (T^a T^b)_{ii} = T_{ij}^a T_{ji}^b
\]
\begin{equation}
B(x) = B^a T^a \quad \quad \quad (80.3) 
\end{equation}

\[
B_{ij} = B^a T_{ij}^a
\]

\begin{equation}
  \mathcal{L} = \text{Tr}{ ( \frac{-1}{2} \partial^{\mu} B \partial_{\mu} B + \frac{1}{3} g B^3  - \frac{1}{4} \lambda B^4) } \quad \quad \quad (80.5)
\end{equation}


20131119

LSZ formula in momentum space.  

Recall the Fourier transform from momentum space to position space for a field, in Srednicki's (physicist) convention

\[
\varphi(x_i) = \int \frac{d^d k}{ (2\pi)^d } e^{ik_i x_i } \widetilde{\varphi}(k_i)
\]

Recall the LSZ formula, cf. Srednicki (5.24), 5. The LSZ reduction formula, that's in position space

\[
\langle f | i \rangle = i^{n+n'} \int d^dx_1 e^{ik_1 x_1} (- \partial_1^2 + m^2 ) \dots dx_{1'}^d e^{- i k_1' x_1' } (- \partial_{1'}^2 + m^2 ) \dots \langle 0 | T \varphi(x_1) \dots \varphi(x_{1'}) \dots | 0 \rangle
\]

Let's consider the Fourier transform 

\[
\langle f | i \rangle = i^{n+n'} \int d^dx_1 e^{ik_1 x_1} (- \partial_1^2 + m^2 ) \dots dx_{1'}^d e^{- i k_1' x_1' } (- \partial_{1'}^2 + m^2 ) \dots \langle 0 | T \int \frac{ d^d \widetilde{k}_1 }{ (2\pi)^d } e^{i \widetilde{k}_1 x_1 } \widetilde{\varphi}(\widetilde{k}_1  ) \dots \int \frac{d^d \widetilde{k}_{1'} }{ (2\pi)^d } e^{  i \widetilde{k}_{1'} x_{1' }   } \widetilde{\varphi}(\widetilde{k}_{1'}) \dots | 0 \rangle 
\]

Now we get out delta functions

\[
\int \frac{ d^d x_{1'}}{(2\pi)^d} e^{i (-k_{1'} + \widetilde{k}_{1'}  ) x_1' } = \delta_{-k_1 + \widetilde{k}_1 }
\]

Just apply for 1 of the coordinates, $x_{1'}$,

\[
\int \frac{d^d \widetilde{k}_{1'} }{ (2\pi)^d } \int d^d x_{1'} e^{ -i k_{1'} x_{1' } } (-\partial_{1'}^2 + m^2 ) e^{ i \widetilde{k}_{1'} x_{1'} }  = \int \frac{d^d \widetilde{k}_{1'} }{ (2\pi)^d } \int d^d x_{1'} e^{-i k_{1'} x_{1'} }( \widetilde{k}_{1'}^2 + m^2 ) e^{ i \widetilde{k}_{1'} x_{1'} } = \int d^d \widetilde{k}_{1'} (\widetilde{k}_{1'}^2 + m^2 ) \delta_{ - k_{1'} + \widetilde{k}_{1' } }
\]

\[
\Longrightarrow \langle f | i \rangle = i^{n+n'} (k_1^2  + m^2 ) \dots ( k_{1'}^2 + m^2 ) \dots \langle 0 | T \widetilde{\varphi}(-k_1) \dots \widetilde{\varphi}(k_{1'} ) \dots | 0 \rangle
\]

So this is the LSZ formula in momentum space.  

20131119

$\phi^4$ theory review

Consider $N=1$ in 

\[
\mathcal{L} = \text{tr}{ \left( \frac{-1}{2} \partial^{\mu} B \partial_{\mu} B + \frac{1}{3} g B^3 - \frac{1}{4} \lambda B^4 \right) } \quad \quad \quad (80.5)
\]

$B(x) = B^a(x) T^a$

\[
\begin{gathered}
  \mathcal{L} = \frac{-1}{2} \partial^{\mu} B^a \partial_{\mu} B^a + \frac{g}{3} \text{tr}{ ( T^a T^b T^c )} B^a B^b B^c  - \frac{\lambda}{4} \text{tr}{ T^a T^b T^c T^d } B^a B^b B^c B^d \quad \quad \, (80.6)
\end{gathered}
\]

$N=1$ \\
$T^a$ was a complete set of $N^2$ \, $N\times N$ hermitian matrices.  

When $N=1$, $T^a$ is just a number (namely 1) and there's only 1 matrix, $a=1$

\[
\mathcal{L} = \frac{-1}{2} \partial^{\mu} B \partial_{\mu} B + \frac{g}{3} B^3 - \frac{\lambda}{4} B^4
\]

\[
\widetilde{\Delta}{ (k^2)} = \frac{1 }{ k^2 - i \epsilon } \sim \frac{1}{k^2} \quad \quad \quad (80.7)
\]

$2ig$ 

\[
\begin{aligned}
& \exp{ [ iS ] } = \exp{ [i \int d^d x \left( \frac{-1}{2} Z_B \partial^{\mu} B \partial_{\mu} B - \frac{1}{2} Z_m m^2 B^2 + \frac{g}{3} Z_g \widetilde{\mu}^{\epsilon/2} B^3 + \frac{-\lambda}{4} Z_{\lambda} B^4  \right) ] } \\ 
& \exp{ \left[ i \int d^d x \mathcal{L}_3  \right] } \exp{ \left[ i \int d^dx \mathcal{L}_4 \right] } = \exp{ \left[ i \int d^d x \frac{g}{3} Z_g \widetilde{\mu}^{\epsilon/2} B^3 \right] } \exp{ \left[ i \int d^dx \frac{ -\lambda}{4} Z_{\lambda} B^4 \right] }
\end{aligned}
\]

\[
Z_0(J) = \langle 0 | 0 \rangle_J = \int \mathcal{D} \varphi \exp{ \left[ i \int d^dx [ \mathcal{L}_0 + J \varphi ] \right] } = \exp{ \left[ \frac{i}{2} \int d^dx d^dx' J(x) \Delta(x-x') J(x') \right] }
\]
in this case we want the 4th order term from $Z_0(J)$.  There'll be a $\frac{1}{4!}$ factor from the Taylor expansion.

\[
\exp{ \left[ i \int d^dx \mathcal{L}_4 \right] } = \exp{ \left[ i \int d^dx \frac{-\lambda}{4} Z_{\lambda} B^4 \right] } \Longrightarrow \exp{ \left[ i \int d^dx \frac{-\lambda Z_{\lambda} }{4} \left( \frac{1}{i} \delta_{J(x)} \right)^4 \right]}Z_0
\]

To see explicitly the Taylor expansion
\[
(1 + i \int d^dx \frac{-\lambda Z_{\lambda}}{4} \frac{1}{i^4} \delta_{J(x)}^4 + \dots ) Z_0
\]

$V=1$ vertices.  $P=4$ propagators.  

\[
i \int d^dx \frac{-\lambda Z_{\lambda}}{ 4} \frac{1}{ i^{4V}} \delta_{J(x)}^4 \frac{1}{4!} \left( \frac{i}{2} \int d^dy d^dy' J_y \Delta_{y-y'} J_{y'} \right)^4 = i \frac{-\lambda Z_{\lambda}}{ 4} i^{P - 4V} \int d^dx \left( \int d^dy J_y \Delta_{y-x} \right)^4
\]
this is a term in $Z$ (not $Z_0$)


Recall the LSZ formula (5.24)

\[
\langle f | i \rangle = i^{n+n'} \int d^4x_1 e^{ik_1x_1}(-\partial_1^2 + m^2 ) \dots d^4x_{1'} e^{ik_{1'}x_{1'}}(-\partial_{1'}^2 + m^2 ) \dots \langle 0 | T\varphi(x_1) \dots \varphi(x_{1'} ) \dots | 0 \rangle
\]

$1\dots n$ incoming \\
$1 \dots n'$ outgoing \\

For the term  
\[
\langle 0 | T\varphi(x_1) \dots \varphi(x_{1'}) \dots | 0 \rangle
\]

so $\frac{1}{i} \delta_{J_{x_1}} \frac{1}{i} \delta_{J_{x_2}} \frac{1}{i} \delta_{J_{x_3}} \frac{1}{i} \delta_{J_{x_4} }$ must be applied.  

Note that factors of $e^{i(k_1 + k_2- k_3 - k_4) x}$ may appear.  

Recall

\[
\delta^{(d)}(x-y) = \int \frac{d^d k}{(2\pi)^d} e^{ik(x-y)}
\]

Apply.  

\[
\begin{gathered}
  i^{-4} i \frac{-\lambda Z_{\lambda}}{4} i^{P - 4V} \int d^dx 4! \Delta_{x_1 - x} \Delta_{x_2 -x} \Delta_{x_3 -x} \Delta_{x_4 - x} =  \\
  = i^{ P - 4V -4 } \frac{ -i \lambda Z_{\lambda }}{ 4} \int d^dx(4!) \int \frac{ d^d \widetilde{k}_1 d^d \widetilde{k}_2 d^d\widetilde{k}_3 d^d\widetilde{k}_4 }{ ( (2\pi)^{4d} ) } \frac{ e^{i \widetilde{k}_1(x_1 -x) } e^{i \widetilde{k}_2(x_2 -x) }e^{i \widetilde{k}_3(x_3 -x) }e^{i \widetilde{k}_4(x_4 -x) } }{ ( \widetilde{k}_1^2 + m^2 ) \dots (\widetilde{k}_4^2 + m^2 ) } = \\ 
 = i^{P-4V -4} \frac{- i \lambda Z_{\lambda }}{ 4 } (4!) \int \frac{ d^dk_1 \dots d^dk_4}{ (2\pi)^{4d} } \frac{1}{ ( \widetilde{k}_1^2 + m^2 ) \dots ( \widetilde{k}_4^2 + m^2 ) } e^{ i ( \widetilde{k}_1 x_1 + \widetilde{k}_2 x_2 + \widetilde{k}_3 x_3 + \widetilde{k}_4 x_4 ) } \delta_{k_1 + k_2 +k_3 + k_4 } (2\pi)^d 
\end{gathered}
\]

apply 

\[
\int d^dx_1 e^{ik_1x_1} (-\partial_1^2 + m^2 ) d^dx_2 e^{ik_2x_2} (-\partial_2^2 + m^2 ) d^dx_3 e^{ik_3x_3} (-\partial_3^2 + m^2 ) d^dx_4 e^{ik_4x_4} (-\partial_4^2 + m^2 ) 
\]

and remember that 

\[
\int d^dx_4 e^{ -i ( k_4 - \widetilde{k}_4 ) x_4 } = \delta_{k_4 - \widetilde{k}_4 } (2\pi)^d
\]

so think about 
\[
\int d^dx_1 d^dx_2 d^dx_3 d^dx_4
\]

\[
i^{ P -4V -4 } \left( \frac{ -i \lambda Z_{\lambda }}{ 4} \right) 4! \int d^dx e^{i (k_1 + k_2 + k_3 +k_4)x } (2\pi)^d \delta_{k_1 + k_2 + k_3 + k_4 } = i^{ P - 4V - 4} \left( \frac{ -i \lambda Z_{\lambda }}{ 4} \right) 4! (2\pi)^d 
\]

\[
\Longrightarrow = -6 i\lambda Z_{\lambda}
\]

From (10.14) Srednicki,

\[
\langle f | i \rangle = (2\pi)^d \delta^d(k_{\text{in}}- k_{\text{out}} ) i \mathcal{T}
\]


$\phi^4$ theory, Feynman diagrams, Feynman rules in momentum space 



\[
\begin{gathered}
  ( \partial^{\mu} B \partial_{\mu} B )_{ii} = \partial^{\mu} B_{ij} \partial_{\mu} B_{ji} \\ 
  \begin{aligned}
    & B^3_{ii} = B_{ij} B_{jk} B_{ki} \\ 
    & B^4_{ii} = B_{ij} B_{jk} B_{kl} B_{li}
\end{aligned}
\end{gathered}
\]

\begin{equation}
\mathcal{L} = \frac{-1}{2} \partial^{\mu} B^a \partial_{\mu} B^a + \frac{1}{3} g \text{Tr}{ ( T^a T^b T^c ) } B^a B^b B^c - \frac{1}{4} \lambda \text{Tr}{ (T^a T^b T^c T^d )} B^a B^b B^c B^d \quad \quad (80.6)
\end{equation}

\[
\begin{gathered}
  \partial^{\mu} B_{ij} \partial_{\mu} B_{ji} = \partial^{\mu} B^a(x) T^a_{ij} \partial_{\mu} B^b(x) T^b_{ji} = ( \partial^{\mu} B^a) ( \partial_{\mu} B^b) T^a_{ij} T_{ji}^b = (\partial^{\mu} B^a)( \partial_{\mu} B^b) \text{Tr}{T^a T^b} = \partial^{\mu} B^a \partial_{\mu} B^b  \\
  B^3_{ii} = B^a B^b B^c T^a_{ij} T^b_{jk} T^c_{ki} = B^a B^b B^c \text{Tr}{ (T^a T^b T^c)}
\end{gathered}
\]

follow arrows backward for matrix index contraction.  

\[
\begin{gathered}
  B(x)_i^{ \, \, j } = B^a(x) (T^a)_i^{ \, \, j } 
\end{gathered}
\]


\begin{equation}
 \widetilde{\Delta}_{i \, \,  k }^{ \, \, j \, \, l }(k^2) = \frac{ (T^a)_i^{\, \, j } (T^a)_k^{ \, \, l } }{ k^2 - i \epsilon } \quad \quad (8.08)
\end{equation}

arrow points from up index to down index.  






% ========================================================================================================
% ======================================== The Feynman rules for N \times N fields 80 ====================
%====================================================================================================
% ======================================== Problems ========================================


\subsection*{ Problems }

\problemhead{1} $N=1$ 





\section{ Scattering in quantum chromodynamics}

\section{ Wilson loops, lattice theory, and confinement}


theory is confining if all finite-energy states invariant under a global gauge transformation. \\
$U(1)$ is not confining, because $\exists \, $ finite-energy states (such as state of a single electron) that have nonzero electric charge, and hence change by phase under global gauge transformation.  

confinement is a nonperturbative phenomenon; it can't be seen at any finite order in the kind of weak-coupling perturbation theory that we've been doing (this is why we had no trouble calculating quark and gluon scattering amplitudes) \\
the inverse lattice spacing $\frac{1}{a}$ then acts as an ultraviolet cutoff (section 29) 

Consider 2 spacetime pts. $x^{\mu}, x^{\mu} + \epsilon^{\mu}$ 

Wilson link 
\begin{equation}
  W(x+\epsilon, x) \equiv \exp{  [ ig \epsilon^{\mu} A_{\mu}(x) ] } \quad \quad \quad \, (82.1)
\end{equation}

path $P$ starts at $x$, ends at $y= x + \epsilon_1 + \dots + \epsilon_n$

\begin{equation}
W_P(y,x) \equiv W(y, -\epsilon_n) \dots W(x+\epsilon_1 +  \epsilon_2, x+ \epsilon_1) W(x+\epsilon_1,x) \quad \quad \quad \, (82.8)
\end{equation}


\begin{equation}
  W^{\dag}_P(y,x) = W_{-P}(x,y) \quad \quad \quad \, (82.10)
\end{equation}

consider path that returns to its starting pt., forming a closed, oriented curve $C$ in spacetime.  

Wilson loop

\begin{equation}
  W_C \equiv \text{tr}{ W_C(x,x) } \quad \quad \quad \, (82.11)
\end{equation}

$W_C$ gauge invariant by (82.9)


curve $C$ in spacetime.  

\[
x^{\mu} = x^{\mu}(t) = F^{\mu}(t) = F(t)
\]

$F: \mathbb{R}^1 \to M$, $M = \mathbb{R}^4$ (Euclidean) or $M$ Lorentzian 4-dim. spacetime or $M$ 4-dim. spacetime.  

$A$ 1-form on $M$, $A = A_{\mu} dx^{\mu} = A_{\mu}(x) dx^{\mu}$ \\
oriented region $(U,o)$ in $\mathbb{R}^p = \mathbb{R}^1$ 

\[
\begin{gathered}
  \int_{ (U,o;F)} A = \int_{ (U,o)} F^*A = \int_{ (U,o) } (F^* A) \left[ \frac{  \partial }{ \partial t } \right] dt = o(t) \int_U A\left( F_* \frac{ \partial }{ \partial t} \right) dt = o(t) \int_U A \left( \dot{x}^{\mu} \frac{ \partial }{ \partial x^{\mu} } \right) dt = o(t) \int_U A_{\mu} \dot{x}^{\mu} dt = \\
  = \int_a^b A_{\mu}(x^{\nu}(t) ) \dot{x}^{\mu} dt
\end{gathered}
\]
with
\[
F_* \frac{ \partial }{ \partial t } = \frac{ \partial F^{\mu} }{ \partial t} \frac{ \partial }{ \partial x^{\mu} } = \dot{x}^{\mu} \frac{ \partial }{ \partial x^{\mu }}
\]

\[
\begin{gathered}
  \oint_C A_{\mu} dx^{\mu} = \int_a^b A_{\mu}{ (x^{\nu}(t) )} \dot{x}^{\mu} dt = \int_a^b A_{\mu}(x) \dot{x}^{\mu} dt = \int_a^b A_{\mu} \dot{x}^{\mu} dt \\ 
  g \oint_C dx^{\mu} = g \int_a^b \dot{x}^{\mu}(t) dt 
\end{gathered}
\]






\[
dV = \int_0^{ 2 \pi } d\phi_3 \int_0^{\pi} d\phi_2 \int_0^{\pi} d\phi_1 \int_0^{\infty} \sin^2{ ( \phi_1)} \sin{ ( \phi_2) } d\phi_1 d\phi_2 d\phi_3 
\]
now 
\[
\int_0^{\pi} d\phi_2 \sin{ (\phi_2) } = - \int_{1}^{-1} d\cos{ ( \phi_2) } = \int_{-1}^1 d\cos{ (\phi_2) } = 2
\]

cf. \url{http://en.wikipedia.org/wiki/N-sphere} \\
wikipedia : $n$-sphere

\subsection*{ On Bessel functions }

In Mathematica, you would input:  \\

\begin{verbatim}
 Integrate[ ( Sin[t] )^2 * Exp[ -I * l * Cos[t]], {t, 0, Pi } ] 
\end{verbatim}

to get 

\begin{verbatim}  
ConditionalExpression[(\[Pi] BesselJ[1, Abs[l]])/Abs[l],  l \[Element] Reals] }
\end{verbatim}


\[
\int_0^{\pi} d\theta e^{i z \cos{\theta} } e^{ i n \left(  \theta - \frac{\pi}{2} \right) } + \int_0^{\pi} d\theta e^{iz \cos{\theta} } e^{}
\]

\[
\begin{gathered}
  i^{-n} \int_0^{\pi} e^{ i z \cos{\theta} } \cos{ (n\theta ) } d\theta = \\
  = \frac{ e^{ - i n \pi } }{2} \int_0^{\pi} d\theta e^{ i z \cos{\theta }} ( e^{ i n \theta } + e^{- i n \theta } ) = \underbrace{ \frac{1}{2} \int_0^{\pi} d\theta e^{ i z \cos{\theta } } e^{ in ( \theta - \pi/2 ) }  }_{ \text{ do $\theta = t + \frac{ \pi} {2} $ } }  + \underbrace{  \frac{1}{2} \int_0^{\pi} d\theta e^{ i z \cos{\theta } } e^{ - i n ( \theta + \pi/2 )  }  }_{ \text{ do $\theta = -t - \frac{\pi}{2}$ }}  = \\ 
  = \frac{1}{2} \int^{  \frac{ \pi}{2} }_{ \frac{- \pi}{2} } dt e^{ -iz \sin{t } } e^{ int } + \frac{1}{2} \int_{  \frac{ -\pi}{2} }^{ - \frac{ 3 \pi}{2} } (-dt) e^{-iz \sin{t} } e^{int } = \frac{1}{2} \int_{-\pi}^{ \pi/2} dt e^{- iz \sin{t}} e^{ int } + \frac{1}{2} \int_{ -3\pi /2}^{ -\pi} dt e^{ -iz \sin{t} } e^{int} = \\
   = \frac{1}{2} \int_{-\pi}^{ \pi} dt e^{int} e^{-iz \sin{t}}
\end{gathered}
\]

So 

\[
\boxed{ J_n(z) = \frac{i^{-n} }{ \pi } \int_0^{\pi} e^{iz \cos{t}  } \cos{ (n \theta ) } d\theta = \frac{1}{2 \pi} \int_{-\pi}^{ \pi} d\tau e^{in\tau } e^{-ix \sin{\tau } } }
\]

cf. \url{ http://dlmf.nist.gov/10.9 } \\
dlmf - Digital Library of Mathematical Functions \\

cf. \url{http://en.wikipedia.org/wiki/Bessel_function} \\
wikipedia - Bessel function \\

How do we do integrals on Bessel functions?

Consider using recursion relations.  Now the Bessel function can be defined via a Taylor series (cf. wikipedia - Bessel function)

\[
J_{\alpha}(x)  = \sum_{m=0}^{\infty} \frac{ (-1)^m }{ m! \Gamma{ (m+ \alpha + 1 ) } } \left( \frac{x}{2} \right)^{ 2m + \alpha }
\]

\[
\begin{gathered}
  J'_{\alpha}(x) = \sum_{m=0}^{\infty} \frac{ (-1)^m  ( 2m+ \alpha ) }{ m! \Gamma{ (m+\alpha + 1) } } \frac{ x^{2m + \alpha - 1 } }{ 2^{2m + \alpha } } = \sum_{m=0}^{\infty} \frac{ (-1)^m }{  (m-1)! (m+\alpha)! } \frac{x^{ 2m + \alpha - 1 }  }{ 2^{2m + \alpha - 1 } } + \frac{\alpha}{x} J_{\alpha}{(x)}  = -J_{\alpha + 1}(x) + \frac{\alpha}{x} J_{\alpha}(x)
\end{gathered}
\]

So that 

\[
\boxed{ J_{\alpha}'(x) = -J_{\alpha + 1}(x) + \frac{\alpha}{x} J_{\alpha}(x) }
\]

cf. \url{http://educ.jmu.edu/~lucassk/Papers/BESS.pdf}.  Lukas, Stone.  \emph{Evaluating infinite integrals involving Bessel functions of arbitrary order}. 1994. pp. 5, Eqn. 10 \\

So for $\alpha =0$, 
\[
\begin{gathered}
  J_1(x) = J_0'(x) \\ 
 \int J_1(x) dx = \left. J_0(x)  \right|_0^{\infty}
\end{gathered}
\]
Use again from Lukas, Stone, Eqn. (9) on pp. 5, for large $x$:

\begin{equation}
J_n{(x)} \sim  \sqrt{ \frac{2}{ \pi x } } \cos{  x - \frac{n \pi }{2} - \frac{\pi}{4} } \quad \quad \quad (10)
\end{equation}

and so $J_n(x) \xrightarrow{ x\to \infty} 0$ and $J_0(0)=1$.  




%%%%%% KONEC On Bessel Functions

Applying this to the case at hand,

\[
\begin{gathered}
  \int_0^{\pi} d\theta \sin^2{ \theta } e^{ ik |x-y| \cos{\theta } } = \int_0^{\pi} d\theta \sin^{\theta} e^{ i l \cos{\theta } } =  \left. \frac{ e^{ - i l \cos{\theta  }} }{ il } \sin{\theta} \right|_0^{\pi} - \int \frac{ d\theta ( - e^{ il \cos{\theta} }  )\cos{\theta } }{ il } = + \int_0^{\pi}  \frac{ d \theta \cos{\theta } e^{ il \cos{\theta } } }{ il } = \\
  = \frac{ \pi J_1(l) }{ l }
\end{gathered}
\]


\[
\begin{gathered}
\int_0^{\infty} du J_1(u) = \int_0^{\infty} du ( -J_0'(u) + \frac{0}{u} J_0(u) ) = - \left. J_0(u) \right|_0^{\infty} = - ( 0 - J_0(0) ) = 1
\end{gathered}
\]

Consider this statement from pp.510, Srednicki: \\

Since $\Delta_{\mu \nu}(x-y)$ depends only on $x-y$, \\

What does that mean?  Note that $C$ is a curve in spacetime.  Doing the double integral over the same curve $C$

\[
\begin{gathered}
  \langle 0 | W_C | 0 \rangle = \exp{ \left[ \frac{-g^2}{2} \oint_C dx_{\mu} \oint_C dy_{\nu} \Delta_{\mu \nu}(x-y) \right] } = \quad \quad \quad \, (82.16) \\ 
  = \exp{ \left[ \frac{-g^2}{2} \int_a^b \dot{x}^{\mu}(t) dt \int_{a'}^{b'} \dot{y}^{\nu}(t') dt' \Delta_{\mu \nu}(x-y) \right] } = \exp{ \left[ \frac{-g^2}{2} \int_a^b \dot{x}^{\mu} dt \int_{a'}^{b'} \dot{y}^{\nu} dt' \frac{ \delta_{\mu  \nu } }{ 4\pi^2 ( x-y)^2 } \right] }  = \\
   = \exp{ \left[ \frac{-g^2}{2} \oint_C dx^{\mu} \oint_C dy^{\nu} \frac{ \delta_{\mu \nu } }{ 4\pi^2 (x-y)^2 }  \right] } = \exp{ \left[ \frac{-g^2}{2} \oint_C \oint_C dx \cdot dy \frac{1}{ 4\pi^2 ( x- y)^2 } \right] }
\end{gathered}
\]

since $\Delta_{\mu \nu}(x-y)$ depends only on $x-y$, \\
\quad the double line integral in (82.16) will yield a factor of perimeter $P$ of the curve $C$.  \\
(cf. EY imagine doing $\oint_C dy^{\nu} = \int_{a'}^{b'} \dot{y}^{\nu} dt'$ first, going around $C$, after choosing some pt. on $C$ to be $x^{\mu}$.  \\
it doesn't matter which pt. on $C$ to be $x^{\mu}$, since for $(x-y)^2$, all the distances from $x^{\mu}$ to $y^{\nu}$ get touched.  \\
Then the 2nd. integral, $\oint_C dx^{\mu} = \int_a^b \dot{x}^{\mu} dt$ gives $P$, perimeter, but in spacetime, as it's an integral of the 1-form $dx^{\mu}$) \\

There's also an ultraviolent divergence as $x\to y$, we will cut this off at length scale $a$ 

\begin{equation}
  \langle 0 | W_C | 0 \rangle = \exp{ \left[ - (\widetilde{c} g^2 / a) P \right] } \quad \quad \quad \, (82.18)
\end{equation}

if $x,y$ are on the same side, $dx \cdot dy = dx dy$ \\

keeping in mind ``reasonable'' integration up to the lattice space cutoff (spacetime becomes discretized) 

\[
\begin{gathered}
  \int_0^L \int_0^L \frac{dx dy}{ (x-y)^2 } = \int_a^{L-a} dx \left( \int_0^{x-a} dy + \int_{x+a}^L dy \right) \frac{1}{ (x-y)^2 } = \int_a^{L-a} dx \left( \left. \frac{1}{ x-y} \right|_0^{x-a} + \left. \frac{1}{x-y} \right|_{x+a}^L \right) = \\
  = \int_a^{L-a} dx \left( \frac{1}{a} - \frac{1}{x} + \frac{1}{x-L} - \frac{1}{ -a} \right) = \\
  = \int_a^{L-a} dx \left( \frac{2}{a} + - \frac{1}{x} + \frac{1}{x-L} \right) = \frac{2}{a} (L-a-a ) - \left. \ln{x} \right|_a^{L-a} + \left. \ln{(x-L)} \right|_a^{L-a} = \\
  = \frac{2L}{a} - 4 - \ln{ (L-a) } + \ln{a} + \ln{(-a)} - \ln{(a-L) } = \frac{2L}{a} -4 + \ln{ \left( \frac{a^2}{ (L-a)^2 } \right) } = \frac{2L}{a} - 4 + 2 \ln{ \left( \frac{a/L}{ 1 - a/L} \right) } = \\
  \frac{2L}{a} + 2\ln{ \frac{a}{L} }-  4 + - 2\ln{ \left( 1-  \frac{a}{L} \right) } = \frac{2L}{a} + 2\ln{ \frac{a}{L} } + O(1)
\end{gathered}
\]


If $x,y$ on opposite short sides, $\sim \frac{R^2}{ T^2} \ll 1$.  ignore. \\
If $x,y$ on opposite long sides, 

\[
\begin{gathered}
  \int_0^T dx \int_0^T dy \frac{1}{ (x-y)^2 + R^2 } = \int_0^T dy   \left. \arctan{ \left( \frac{x-y}{R} \right) }  \frac{1}{R} \right|_0^T = \frac{1}{R} \int_0^T dy \left(  \arctan{ \left( \frac{T-y}{R} \right) } - \arctan{ \left( \frac{-y}{R} \right)} \right) = \\
  = \frac{1}{R} \int_0^T dy \left( \arctan{ \left( \frac{T-y}{R} \right) } + \arctan{ \frac{y}{R} } \right) = \\
  = \frac{1}{R} \left. \left( - \arctan{ \left( \frac{T-y}{R} \right) } ( T-y) + \frac{R}{2} \ln{ ((T-y)^2 + R^2 ) } + \arctan{ \left( \frac{y}{R} \right) } y + \frac{-R}{2} \ln{ (y^2 + R^2) } \right) \right|_0^T  = \\
  = \frac{1}{R} \left( \arctan{ \left( \frac{T}{R} \right) } T  + \frac{R}{2} \ln{ R^2 } - \frac{R}{2} \ln{ (T^2 + R^2 ) } + \arctan{ \left( \frac{T}{R} \right) } T + \frac{-R}{2} \ln{ (T^2 + R^2 ) } + \frac{R}{2} \ln{ (R^2 ) }    \right) \cong \\
  \cong \frac{1}{R} \left( \pi T + R \ln{ \left( \frac{R^2}{ T^2 + R^2 } \right) } \right) = \\
  = \frac{\pi T}{R} + \ln{ \left( \frac{ R^2 / T^2 }{ 1 + R^2 /T^2 } \right) } = \frac{\pi T}{R}  - 2\ln{ \frac{T}{R} } + O\left( \frac{R^2}{ T^2} \right) 
\end{gathered}
\]



\section{ Chiral symmetry breaking}

\section{ Spontaneous breaking of gauge symmetries}

\solutionhead{84.1}
\begin{enumerate}
\item[(a)] Given $V(\Phi) = \frac{1}{2} m^2 Tr{ \Phi^2} + \frac{1}{4} \lambda_1 Tr{ \Phi^4} + \frac{1}{4} \lambda_2 (Tr{ \Phi^2 })^2 $, 
\[
\begin{aligned}
  & \Phi^2 = v^2 \left( \begin{matrix} \alpha_1^2 & & \\ & \ddots & \\ &  & \alpha_N^2 \end{matrix} \right) \\ 
  &  \Phi^4 = v^4 \left( \begin{matrix} \alpha_1^4 & & \\ & \ddots & \\ &  & \alpha_N^4 \end{matrix} \right) 
\end{aligned} \quad \quad \quad 
\begin{aligned}
&  Tr{ \Phi^2 } = v^2 \sum_{i=1}^N \alpha_i^2 = v^2 \\ 
&  Tr{ \Phi^4 } = v^4 \sum_{i=1}^N \alpha_i^4
\end{aligned}
\]
So then
\[
V = \frac{m^2}{2} v^2 + \frac{ \lambda_1}{4} v^4 \sum_{i=1}^N \alpha_i^4 + \frac{\lambda_2}{4} v^4
\]
\[
\frac{ \partial V}{\partial v} = m^2 v + ( \lambda_1 \sum_{i=1}^N \alpha_i^4 + \lambda_2) v^3 = 0 \Longrightarrow v= 0 \text{ or } v^2 = \frac{ -m^2 }{ \lambda_1 \sum_{i=1}^N \alpha_i^4 + \lambda_2 } = \frac{ -m^2}{ \lambda_1 C + \lambda_2 } 
\]
\[
\Longrightarrow V = \frac{m^2}{2} \left( \frac{ -m^2 }{ \lambda_1 C + \lambda_2 } \right) + \frac{ \lambda_1}{4} \left( \frac{m^4}{ (\lambda_1 C + \lambda_2)^2 } \right) C + \frac{ \lambda_2}{4} \left( \frac{ m^4}{ (\lambda_1 C + \lambda_2)^2 } \right) = \frac{ -m^4}{ 2 (\lambda_1 C + \lambda_2 ) } + \frac{ m^4}{ 4 (\lambda_1 C + \lambda_2) } = \boxed{ \frac{ -m^4}{ 4 (\lambda_1 C + \lambda_2 ) } }
\]
\item[(b)] Consider 
\[
\frac{ \partial V}{ \partial v} = m^2 v \left( 1 + \frac{ (\lambda_1 C + \lambda_2) v^2 }{ m^2 } \right)
\]
$m^2 <0$.  If $\lambda_1 C + \lambda_2 \leq 0$, then $v^2 = \frac{-m^2}{ \lambda_1 C + \lambda_2 }$ isn't an allowed critical point.  

Consider large $v$.  $\frac{ \partial V}{ \partial v} \simeq m^2 v \left( \frac{ (\lambda_1 C + \lambda_2) v^2}{ m^2} \right) \to -\infty$ as $v\to -\infty$.  So $V$ is not bounded below.  
\item[(c)] $v = \frac{-m^4}{ 4(\lambda_1 C + \lambda_2 )} \leq 0$ as established above (note, for $v=0$, $V=0$).  

\[
\begin{aligned}
  \min{ (\lambda_1 C + \lambda) } & \leq (\lambda_1 C + \lambda_2) \\ 
  \frac{1}{ \lambda_1 C + \lambda_2 } & \leq \frac{1}{ \min{ ( \lambda_1 C + \lambda_2) } }
\end{aligned}  \quad \quad \Longrightarrow \frac{-1}{ \lambda_1 C + \lambda_2 } \geq \frac{-1}{ \min{ ( \lambda_1 C + \lambda_2 ) } }
\]
\[
\Longrightarrow V \geq \frac{-m^4}{ 4 \min{ (\lambda_1 C + \lambda_2) } }
\]
\item[(d)] From previous part, extremize $f(\alpha) = \lambda_1 \sum_{i=1}^N \alpha_i^4 + \lambda_2$.  
\[
\begin{aligned}
  & \sum_{i=1} \alpha_i = 0 \\ 
  & \sum_{i=1} \alpha_i^2 = 1 
\end{aligned} \quad \quad  \, \begin{aligned}
  & g(\alpha) = \sum_{i=1} \alpha_i = 0 \\ 
  & h(\alpha) = \sum_{i=1} \alpha_i^2 - 1 
\end{aligned}
\]
Let 
\[
\Lambda = f - ag -bh
\]
with $a,b$ as Lagrange multipliers.  

\[
\partial_j \Lambda = 4 \lambda_1 \alpha_j^3 - a - 2b \alpha_j = 0 \text{ or } \alpha_j^3 - \frac{b}{2\lambda_1} \alpha_j - \frac{a}{4 \lambda_1 } = 0 
\]
Now $ \frac{4\lambda_1 \sum_j \alpha_j^3 }{ N} = a$ using the constraints and if $\alpha_j \neq 0$ \, $\forall \, \alpha_j$, 
\[
2\lambda_1 \sum_j \alpha_j^4 = b
\]
then $\alpha_j^3 - \sum_{i} \alpha_i^4 \alpha_j - \frac{ \sum_i \alpha_i^3 }{ N } = 0$.

In any case, the solution to this cubic will be 3 possibly complex numbers that sum to 0.  

\item[(e)] $\Phi = v diag{(\alpha_1, \dots, \alpha_N)}$.  

Since $\sum_i \alpha_i =0$, $\sum_i \alpha_i^2 =1$, and $\alpha_i$ take on at most 3 different values, $u_1$, $u_2$, $u_3$, and $u_1 + u_2 + u_3 =0$.  But for solutions to the cubic, they are complex conjugates; so there are only 2 possible values for $\alpha_i$.  

Since $\sum_i \alpha_i = 0$, $N_+ = N_- = \frac{1}{2} N$ if $N$ even, $N_{\pm} = \frac{1}{2} (N \pm 1)$ if $N$ odd (so that $N_+ + N_- = N$).  

\[
\Longrightarrow SU(N_+) \times SU(N_-) \times U(1) 
\]
\end{enumerate}

\section{ Spontaneously broken abelian gauge theory}

\section{ Spontaneously broken nonabelian gauge theory}

\section{The Standard Model: gauge and Higgs sector}

\section{The Standard Model: lepton sector}

\section{The Standard Model: quark sector}

\section{Electroweak interactions of hadrons}

\section{Neutrino masses}

\section{Solitons and monopoles }

\section{Instantons and theta vacua}

\section{Quarks and theta vacua}

\section{Supersymmetry}

cf. Bilal 96

$\begin{aligned} & \quad \\ 
  & A = A(y) \\
  & F = F(y) \end{aligned}$ \quad \, complex scalar field, \, $\psi_a(y)$ left-handed Weyl field

\begin{equation}
  D^*_{\dot{a}}\Phi(x,\theta, \theta^*) = 0 \quad \quad \quad \, (95.21)
\end{equation}
left-handed chiral superfield. 

\begin{equation}
  y^{\mu} = x^{\mu} - i\theta^c\sigma^{\mu}_{ c\dot{c}} \theta^{ * \dot{c}} \quad \quad \quad \, (95.23)
\end{equation}

\begin{equation}
  \Phi(y,\theta) = A(y) + \sqrt{2} \theta \psi(y) + \theta \theta F(y) \quad \quad \quad , (95.25)
\end{equation}

conventions

\[
\begin{aligned}
  & \theta \psi \equiv \theta^a \psi_a \\ 
  &  \theta^2 \equiv \theta \theta \equiv \theta^a \theta_a = -2 \theta^1 \theta^2
\end{aligned}
\]



\begin{equation}
  \theta_a \theta_b = \frac{1}{2} \theta \theta \epsilon_{ab} \quad \quad \, \theta^a \theta^b = -\frac{1}{2} \theta \theta \epsilon^{ab} \quad \quad \quad \, (95.26)
\end{equation}

\begin{equation}
  (\theta \sigma^{\mu} \theta^*) ( \theta \sigma^{\nu} \theta^*) = \frac{-1}{2} \theta^2 \theta^* \theta^* g^{\mu \nu} \quad \quad \quad \, (95.28)
\end{equation}

\[
\begin{aligned}
  & A(y) = A(x) + - i \theta \sigma^{\mu} \theta^* \partial_{\mu} A + \frac{ - (\theta \sigma^{\mu} \theta^*)( \theta \sigma^{\nu} \theta^*) }{ 2 } \partial^2_{\mu \nu} A \\ 
  & \psi(y) = \psi(x) + i \theta\sigma^{\mu} \overline{\theta} \partial_{\mu} \psi \\ 
  & \theta^a (\theta^c \sigma^{\mu}_{c\dot{c} } \theta^{ * \dot{c}} ) \partial_{\mu} \psi_{\dot{a}} = \frac{-1}{2} \theta \theta \theta^* \overline{\sigma}^{\mu} \partial_{\mu} \psi 
\end{aligned}
\]

\begin{equation}
  \Phi(x,\theta,\theta^*) = A(x) + \sqrt{2} \theta \psi(x) + \theta^2 F(x) - i (\theta \sigma^{\mu} \theta^* ) \partial_{\mu} A(x) + \frac{-1}{\sqrt{2}} i \theta^2 \theta^* \overline{\sigma}^{\mu} \partial_{\mu} \psi(x) + \frac{1}{4} \theta^2 (\theta^*)^2 \partial^2 A(x) \quad \quad \quad \, (95.29)
\end{equation}




\begin{equation}
  \left. W^a W_a \right|_F = 2 i \lambda^a \sigma^{\mu}_{ a \dot{a}} \partial_{\mu} \lambda^{\dag \dot{a}} - \frac{1}{2} F^{\mu \nu } F_{\mu \nu } - \frac{i}{2} \widetilde{F}^{\mu \nu } F_{\mu \nu } + D^2 \quad \quad \quad \, (95.76)
\end{equation}



\problemhead{95.1} Recall $\lbrace Q_{aA}, Q^{\dag}_{\dot{a} B} \rbrace = -2\delta_{AB} \sigma^{\mu}_{a\dot{a}} P_{\mu}$ (95.6).  $P_0$ is the Hamiltonian.   
\[
-2\delta_{AB} ( \sigma_{1\dot{1}}^0 P_0 + \sigma_{1\dot{1}}^3 P_3 ) = - 2\delta_{AB} (-H + P_3) = \lbrace Q_{1A}, Q^{\dag}_{\dot{1}B} \rbrace
\]

If massive, $P_3 =0$, 

\[
2 \langle \varphi | H  | \varphi \rangle = \langle \varphi | Q_{1A} Q^{\dag}_{\dot{1} A} | \varphi \rangle + \langle \varphi | Q^{\dag}_{ \dot{1} A } Q_{1A} | \varphi \rangle = \| Q^{\dag}_{\dot{1} A} | \varphi \rangle \|^2 + \| Q_{1 A} | \varphi \rangle \|^2 
\]

$H$ positive definite.  \\

If massless, then consider $a, \dot{a} = 2, \dot{2}$ for $\sigma^3{2\dot{2}}P_3 = -H$. Then similar to above, $H$ positive definite.     \\

If $H | \varphi \rangle = 0$, $\| Q_{1A} | \varphi \rangle \|^2 = 0 $ or $Q_{1A} | \varphi \rangle = 0$.  Likewise for $Q_{2A} | \varphi \rangle =0$  \\

If massless, $P_3 = H$ so consider $\sigma^{\mu}_{2\dot{2}}$ for $-2\delta_{AB} (-2H) = 4 \delta_{AB} H = \lbrace Q_{2A}, Q^{\dag}_{ \dot{2} B} \rbrace$ and for $a = \dot{a} = 1$, $\lbrace Q_{1A}, Q^{\dag}_{\dot{1} B} \rbrace =0$ so reobtain $\| Q^{\dag}_{\dot{1} A} | \varphi \rangle \|^2 + \| Q_{1A} | \varphi \rangle \|^2 =0$

\problemhead{95.2}
\begin{enumerate}
\item[(a)] If $\langle 0 | F | 0 \rangle \neq 0$, using first of eqs. (95.34), $\lbrace \psi_c, Q_a \rbrace = -i \sqrt{2} \epsilon_{ac} F$, 
\[
\langle 0 | -i \sqrt{2} \epsilon_{12} F | 0 \rangle = i \sqrt{2} \langle 0 | F | 0 \rangle = \langle 0 | \lbrace \psi_2, Q_1 \rbrace | 0 \rangle = \langle 0 | \psi_2 Q_1 | 0 \rangle  + \langle 0 | Q_1 \psi_2 | 0 \rangle
\]
If $Q_1 | 0 \rangle =0$, $\langle 0 | Q_1 \psi_2 |  0 \rangle \neq 0$ implying $Q_1^{\dag} | 0 \rangle \neq 0$ and vice versa.  
\item[(b)] $[V(x,\theta,\overline{\theta}), Q_a] = -i \mathcal{Q}_aV(x,\theta, \overline{\theta})$   \\

Recall that 
\[
V(x,\theta, \overline{\theta} ) = C(x) + \theta \chi(x) + \overline{\theta} \chi^{\dag}(x) + \theta^2 M(x) + \overline{\theta}^2 \overline{M}(x) + \theta \sigma^{\mu} \overline{\theta} v_{\mu}(x) + \theta^2 \overline{\theta} \lambda^{\dag}(x) + \overline{\theta}^2 \theta \lambda(x) + \frac{1}{2} \theta^2 \overline{\theta}^2 D(x)
\]

Now \\
$[\overline{\theta}^2 \theta^a \lambda_a, Q_b ]= \overline{\theta}^2 \theta^a \lambda_a Q_b - Q_b \overline{\theta}^2 \theta^a \lambda_a = \overline{\theta}^2 \theta^a \lbrace \lambda_a, Q_b \rbrace$

Consider $\overline{\theta}^2 \theta_a$ components.  

\[
\begin{gathered}
  \partial_{\theta^b} \frac{1}{2} \theta^2 \overline{\theta}^2 D = \frac{1}{2} \theta_b \overline{\theta}^2 D = \frac{1}{2} \theta^a \epsilon_{ ab} \overline{\theta}^2 D
\end{gathered}
\]

\[
\begin{gathered}
  -i \sigma^m_{b\dot{b}} \overline{\theta}^{\dot{b}} \partial_m \theta^{\alpha} \sigma^n_{\alpha \dot{\alpha}} \overline{\theta}^{\dot{\alpha}} v_n = \frac{i}{2} \sigma^m_{b\dot{b}} \sigma^n_{\alpha \dot{\alpha} } \epsilon^{\dot{b} \dot{\alpha} } \overline{\theta}^2 \theta^{\alpha} \partial_m v_n = \frac{i}{2} \sigma^m_{b\dot{b}} \epsilon_{\alpha c} \epsilon_{\dot{ \alpha} \dot{c}} \overline{\sigma}^{n \dot{c } c } \epsilon^{\dot{b} \dot{\alpha}} \overline{\theta}^2 \theta^{\alpha} \partial_m v_n  = \frac{i}{2} \sigma^m_{b\dot{b}} \overline{\sigma}^{ n \dot{b} c} \epsilon_{\alpha c} \overline{\theta}^2 \theta^{\alpha} \partial_m v_n = \\
   = \frac{ -i }{2} (\sigma^m \overline{\sigma}^n )_{ba} \overline{\theta}^2 \theta^a \partial_m v_n
\end{gathered}
\]

So then 
\[
\lbrace \lambda_a, Q_b \rbrace = \frac{1}{2} \epsilon_{ab} D - \frac{i}{2} (\sigma^m \overline{\sigma}^n)_{ba} \partial_m v_n
\]

Knowing that $\langle 0 | D| 0 \rangle \neq 0$, 

\[
\langle 0 | \frac{1}{2} \epsilon_{21} D | 0 \rangle = \langle 0 | \lambda_2 Q_1 | 0 \rangle + \langle 0 | Q_1 \lambda_2 | 0 \rangle + \frac{i}{2} (\sigma^m \overline{\sigma}^n)_{12} \langle 0 | \partial_m v_n | 0 \rangle
\]

\end{enumerate}

\problemhead{95.3}
\begin{enumerate}
\item[(a)] Given the O'Raifeartaigh model, 
\[
W = mBC + \kappa A (C^2-v^2)
\]
we must also consider the $D$-terms due to these hermitian terms:
\[
\begin{aligned}
  & \left. A^{\dag} A \right|_D = - \partial^{\mu} A_A^{\dag} \partial_{\mu} A_A + i \psi_A^{\dag} \overline{\sigma}^{\mu} \partial_{\mu} \psi_A + F_A^{\dag} F_A \\ 
  & \left. B^{\dag} B \right|_D = - \partial^{\mu} A_B^{\dag} \partial_{\mu} A_B + i \psi_B^{\dag} \overline{\sigma}^{\mu} \partial_{\mu} \psi_B + F_B^{\dag} F_B \\
  & \left. C^{\dag} C \right|_D = - \partial^{\mu} A_A^{\dag} \partial_{\mu} A_A + i \psi_C^{\dag} \overline{\sigma}^{\mu} \partial_{\mu} \psi_C + F_C^{\dag} F_C \\
\end{aligned}
\]

Now to compute the $F$-terms using
\[
W(\Phi)_F = \frac{ \partial W(A) }{ \partial \Phi_i } F_i - \frac{1}{2} \frac{ \partial^2 W(A)}{ \partial \Phi_i \partial\Phi_j } \psi_i \psi_j 
\]

\[
\begin{aligned}
& \partial_A W = \kappa (C^2-  v^2) \\ 
& \partial_B W = m C \\ 
&  \partial_C W = mB + 2 \kappa AC
\end{aligned} \quad \quad \quad \, 
\begin{aligned}
  & \partial_{CA} W  = 2 \kappa C \\ 
 &  \partial_{CB} W = m \\ 
  & \partial_{AC} W = 2 \kappa C \\ 
  & \partial_{BC} W = m \\
  & \partial_{CC} W = 2 \kappa A
\end{aligned}
\]

\[
\begin{gathered}
  W(\Phi)_F = \kappa (A_C^2 - v^2) F_A + m A_C F_B + (m A_B + 2 \kappa A_A A_C ) F_C - \\
  - \frac{1}{2} \left( 2 \kappa A_C \psi_C \psi_A + m \psi_C \psi_B + 2 \kappa A_C \psi_A \psi_C + m \psi_B \psi_C + 2\kappa A_A\psi^2_C \right) = \\
  = \kappa (A_C^2 - v^2)F_A + mA_C F_B + (m A_B + 2\kappa A_A A_C) F_C - 2\kappa A_C \psi_C \psi_A - m \psi_C \psi_B - \kappa A_A\psi_C^2
\end{gathered}
\]
\[
\mathcal{L} = \left. \Phi_i^{\dag} \Phi_i \right|_D + (  \left. W(\Phi) \right|_F + \text{ h.c.})
\]

So
\[
\frac{  \partial \mathcal{L}}{ \partial F_A } = F_A^{\dag} + \kappa (A_C^2 - v^2) = 0 \text{ or } F_A^{\dag} = - \kappa (A_C^2 - v^2)
\]
Similarly for the other fields.   \\

So recalling 
\[
\mathcal{L} = - \partial^{\mu} A_i^{\dag} \partial_{\mu} A_i + i \psi_i^{\dag} \overline{\sigma}^{\mu} \partial_{\mu} \psi_i - \left| \frac{ \partial W(A) }{ \partial \Phi_i } \right|^2 - \frac{1}{2} \left[ \frac{ \partial^2 W(A) }{ \partial \Phi_i \partial \Phi_j } \psi_i \psi_j + \text{ h.c. } \right]
\]

So clearly, 
\[
V = \left| \frac{ \partial W(A) }{ \partial \Phi_i } \right|^2 + \frac{1}{2} \left[ \frac{ \partial^2 W(A) }{ \partial \Phi_i \partial \Phi_j } \psi_i \psi_j + \text{ h.c. } \right]
\]

\[
V = | \kappa |^2 | A_C^2 - v^2 |^2 + m^2 |A_C|^2 + | m A_B + 2\kappa A_A A_C |^2 + 2\kappa A_C \psi_C \psi_A + m\psi_C \psi_B + \kappa \psi_C^2 A_A
\]

\[
\begin{aligned}
  & F_A = - \kappa ((A_C^{\dag})^2 - v^2) \\ 
  & F_B = - m A_C^{\dag} \\ 
  & F_C = - ( mA_B^{\dag} + 2\kappa A^{\dag}_C A^{\dag}_A )
\end{aligned}
\]


If $A_C^2 = v^2$, then $F_B \neq 0$.  If $|A_C|^2 =0$, then $F_A = |\kappa|^2| 0 - v^2|^2 = |\kappa |^2 v^4 \neq 0$
\item[(b)] From $V$ above, the scalar potential is 
\[
V = \kappa^2( v^4 - v^2 A_C^2 - v^2 \overline{A}_C^2  + |A_C|^4 ) + m^2 |A_C|^2 + m^2 |A_B|^2 + 2\kappa m A_B \overline{A}_A \overline{A}_C + 2 \kappa m \overline{A}_B A_A A_C + 4 \kappa^2 |A_A|^2 |A_C|^2
\]

Identify quadratic terms: $- \kappa^2 v^2 A_C^2 - \kappa^2 v^2 \overline{A}_C^2 + m^2 |A_C|^2 + m^2 |A_B|^2$
\[
\begin{aligned}
  m_{A_C}^2 = m^2 \\ 
  m_{A_B}^2 = m^2 \\ 
  m_{A_A}^2 = 0 
\end{aligned}
\]

\textbf{Big trick}: If $A_C = c+id$, 
\[
\begin{aligned}
  A_C^2 = c^2 + 2i cd + - d^2 \\ 
  \overline{A}_C^2 = c^2 + - 2icd - d^2 \\ 
  |A_C|^2 = c^2 + d^2 
\end{aligned} \quad \quad \quad \Longrightarrow \quad \begin{aligned} m^2_C = m^2 + - 2\kappa^2 v^2 \\ m^2_d = m^2 + 2\kappa^2 v^2 \end{aligned}
\]

If $m^2 > -2\kappa^2 v^2$, absolute minimum of scalar potential when $A_C = A_B = 0$.  So $V = \kappa^2 v^4$ at minimum.   \\

Masses for 6 real scalars (bosons): $0,0, m^2, m^2, m^2 \pm 2\kappa^2 v$  



\end{enumerate}

\problemhead{95.4} \emph{Supersymmetric quantum electrodynamics}
\begin{enumerate}
\item[(a)] 
\[
\begin{gathered}
   \left. \Phi^{\dag} e^{-2gV } \Phi \right|_D +  \left. B^{\dag} e^{+2gV} B \right|_D  +  \left. \frac{1}{4} W^a W_a \right|_F + \text{ h.c. } + \left. W \right|_F = \\
   = - ( D^{\mu} A)^{\dag} D_{\mu} A + i \psi^{\dag} \overline{\sigma}^{\mu} D_{\mu} \psi + F^{\dag} F + \sqrt{2} g \psi^{\dag} \lambda^{\dag} A + \sqrt{2} g A^{\dag} \lambda \psi - g A^{\dag} DA  + \\ 
   + - ( D^{\mu} A_B)^{\dag} D_{\mu} A_B + i \psi_B^{\dag} \overline{\sigma}^{\mu} D_{\mu} \psi_B + F^{\dag}_B F_B - \sqrt{2} g \psi^{\dag}_B \lambda^{\dag} A_B  - \sqrt{2} g A_B^{\dag} \lambda \psi_B + gA^{\dag}_B D A_B  + \\
    + i \lambda^{\dag} \overline{\sigma}^{\mu} \partial_{\mu} \lambda - \frac{1}{4} F^{\mu \nu} F_{\mu \nu} + \frac{1}{2} D^2 + m(A_B F + AF_B - \psi_B \psi)
\end{gathered}
\]
\item[(b)]
\[
\begin{aligned}
  \partial_F \mathcal{L} & = F^* + mA_B = 0 \\ 
  \Longrightarrow F & = - mA^*_B \\ 
  \partial_{F_B} \mathcal{L} & = F_B^* + mA = 0 \\ 
  \Longrightarrow F_B & = - mA^* 
\end{aligned} \quad \quad \quad \quad \begin{gathered}
  \partial_D \mathcal{L} = D -g(|A|^2 - |A_B|^2) = 0 \text{ or } D = g(|A|^2 - |A_B|^2)
\end{gathered}
\]
so then
\[
\begin{gathered}
     - ( D^{\mu} A)^{\dag} D_{\mu} A + i \psi^{\dag} \overline{\sigma}^{\mu} D_{\mu} \psi +   \sqrt{2} g \psi^{\dag} \lambda^{\dag} A + \sqrt{2} g A^{\dag} \lambda \psi  - \\
 - ( D^{\mu} A_B)^{\dag} D_{\mu} A_B + i \psi_B^{\dag} \overline{\sigma}^{\mu} D_{\mu} \psi_B  - \sqrt{2} g \psi^{\dag}_B \lambda^{\dag} A_B  - \sqrt{2} g A_B^{\dag} \lambda \psi_B      + i \lambda^{\dag} \overline{\sigma}^{\mu} \partial_{\mu} \lambda - \frac{1}{4} F^{\mu \nu} F_{\mu \nu}  - m(\psi_B \psi)
\end{gathered}
\]
\end{enumerate}

\problemhead{95.5} 
\begin{enumerate}
\item[(a)] Consider supersymmetric gauge theory with $U(1)$ factor.  So supersymmetry transformation of the $D$ component of a vector field is a total derivative.  Here $D$ is the auxiliary field for the $U(1)$ gauge field or i.e. the vector field $V$, which transforms in this way: $V \to V + i (\Xi^{\dag} - \Xi)$ 

So $e \xi D$ would transform as a total derivative under supersymmetry in the action, leaving the action invariant.  

Recall, for any vector field, 
\[
\begin{aligned}
  & [D, Q_a] = - \sigma_{a\dot{c}}^{\mu} \partial_{\mu} \lambda^{\dag \dot{c}} \\ 
  &  [D, Q_{\dot{a}}^{\dag} ] = \partial_{\mu} \lambda^c \sigma^{\mu}_{c\dot{a}}
\end{aligned}
\]

or i.e. 
\[
\delta D^a = \frac{i}{\sqrt{2}} ( - \epsilon^{\dag} \overline{\sigma}^{\mu} D_{\mu } \lambda^a + D_{\mu} \lambda^{\dag a} \overline{\sigma}^{\mu} \epsilon )
\]

If nonabelian, $[D,Q_a]$, $[D,Q^{\dag}_{\dot{a}}]$ not a total derivative.  
\item[(b)] Recall that 
\[
\begin{aligned}
  & \left. \Phi^{\dag} e^{-2gV } \Phi \right|_D = - ( D^{\mu} A)^{\dag} D_{\mu} A + i \psi^{\dag} \overline{\sigma}^{\mu} D_{\mu} \psi + F^{\dag} F + \sqrt{2} g \psi^{\dag} \lambda^{\dag} A + \sqrt{2} g A^{\dag} \lambda \psi - g A^{\dag} DA \\  
  & \left. B^{\dag} e^{+2gV} B \right|_D = - ( D^{\mu} A_B)^{\dag} D_{\mu} A_B + i \psi_B^{\dag} \overline{\sigma}^{\mu} D_{\mu} \psi_B + F^{\dag}_B F_B - \sqrt{2} g \psi^{\dag}_B \lambda^{\dag} A_B  - \sqrt{2} g A_B^{\dag} \lambda \psi_B + gA^{\dag}_B D A_B  \\
  & \left. \frac{1}{4} W^a W_a \right|_F + \text{ h.c. } = i \lambda^{\dag} \overline{\sigma}^{\mu} \partial_{\mu} \lambda - \frac{1}{4} F^{\mu \nu} F_{\mu \nu} + \frac{1}{2} D^2
\end{aligned}
\]


For SQED, we're told to include gauge-invariant superpotential $W = m B \Phi$.  Recall, 

\[
\left. W \right|_F = m ( A_B F + AF_B - \psi_B \psi)
\]

Solving for the equation of motions for $F^{\dag}$ and $F_B^{\dag}$ (i.e. just take the partial derivatives) then 
\[
\begin{aligned}
  & F = - mA^*_B \\ 
  & F_B = - mA^*
\end{aligned}
\]

Then the scalar potential is immediately obtained

\[
\left. W \right|_F = m^2 |A_B|^2 + m^2 |A|^2 - m \psi_B \psi
\]


Having added 
\[
\mathcal{L}_{FI} = e\xi D
\]

then solving for the equation of motion for $D$,
\[
\partial_D \mathcal{L} = e\xi + D + - g|A|^2 + g |A_B|^2 = 0 \text{ or } D = g(|A|^2 - |A_B|^2) - e\xi
\]
\item[(c)] The potential is 
\[
\begin{gathered}
  V = m^2 | A|^2 + m^2 |A_B|^2 + \frac{1}{2} D^2 = m^2 (|A|^2 + |A_B|^2) + \frac{1}{2} ( g(|A|^2  - |A_B|^2) - e\xi )^2 = \\
   = m^2 ( |A|^2 + |A_B|^2) + \frac{1}{2} (g^2 ( |A|^2 - |A_B|^2 )^2  - 2 e \xi g(|A|^2 - |A_B|^2 ) + e^2 \xi^2 ) = \\
   = (m^2 - e \xi g) |A|^2 + (m^2 + e\xi g) |A_B|^2 + \frac{g^2}{2} ( |A|^2 - |A_B|^2 )^2 + \frac{ e^2 \xi^2 }{2} 
\end{gathered}
\]

Note that 
\[
\begin{aligned}
  & \partial_B W = m \Phi \\ 
  &  \partial_{\Phi}W = mB
\end{aligned}
\]
Suppose nonzero vacuum expectation values.  Then for 
\[
m^2 > e \xi g \text{ or } \xi < \frac{m^2}{eg}
\]
supersymmetry is spontaneously broken ($V$ nonzero).  
\end{enumerate}

\problemhead{95.6} \emph{$R$ symmetry}
\begin{enumerate}
\item[(a)] Upon inspection of $\Phi^{\dag} e^{-2gV } \left. \Phi \right|_D$, if $\lambda_a \to e^{ - i \alpha} \lambda_a$, and noting that 
\[
\Phi(\theta) = A + \sqrt{2} \theta \psi + \theta^2 F
\]
and letting $\theta_a \to e^{i \alpha} \theta_a$, 
so that $\Phi(\theta') = e^{-ir \alpha} \Phi(\theta)$, then
\[
\Phi(\theta') = e^{-ir\alpha} A + \sqrt{2} e^{-ir \alpha} \theta e^{-i (r-1) \alpha} \psi + e^{-2i\alpha} \theta^2 e^{-i (r-2) \alpha} F
\]
then $A^{\dag} \lambda \psi$ term leads to $e^{ir\alpha} A^{\dag} e^{-i\alpha} \lambda e^{-i(r-1)\alpha} \psi$

So for the scalar component of the chiral superfield, it has $R$-charge $r$ and the fermionic component has $R$-charge $r-1$.  
\item[(b)] $W(\Phi)_F$ is $\theta^2$ coefficient.  Now $d^2\theta \to e^{2i\alpha} d^2\theta$.  $W(\Phi)_F \to e^{-2i\alpha} W(\Phi)_F$.  $+2$ $R$-charge.  
\item[(c)] $\left. \Phi^{\dag}e^{-2gV} \Phi\right|_D + \left. B^{\dag} e^{+2gV} B \right|_D$ has $R$ symmetry.  \\
$\frac{1}{4} \left. W^a W_a\right|_F + \text{ h.c. } \left. W\right|_F$ will have $+2$ $R$-charge.
\end{enumerate}






\end{document} 
